\medterm{A \& E Medicine} Accident and emergency (A\&E) medicine provides rapid assessment, stabilization, diagnosis, and initial treatment for sudden illness and injury.

\medterm{A Fib} Alternate name; see \textit{Atrial fibrillation}.

\medterm{A-T} The adenine--thymine base pair in DNA. The two bases join by hydrogen bonds and help form the structure of DNA.

\medterm{Aagenaes Syndrome} A rare inherited disorder that causes reduced flow of bile from the liver beginning in infancy and long-lasting swelling from lymph-fluid buildup, usually in the legs. Liver problems may come and go with age, but some people develop vitamin deficiency, poor growth, or liver scarring.

\medterm{Aarskog Syndrome} Alternate name; see \textit{Aarskog-Scott Syndrome}.

\medterm{Aarskog-Scott Syndrome} A rare inherited developmental disorder, often related to FGD1 variants, characterized by distinctive facial features, short stature, skeletal differences, and genital anomalies.

\synonyms
Aarskog Syndrome

\medterm{Aase Syndrome} Alternate name; see \textit{Aase-Smith Syndrome II}.

\medterm{Aase-Smith Syndrome I} A rare condition present at birth that may cause excess fluid in the brain, a split in the roof of the mouth, joints that are fixed in bent positions, and unusual hand and facial features. Some people also have heart defects.

\medterm{Aase-Smith Syndrome II} A rare inherited disorder that causes reduced production of red blood cells and thumbs with three bones instead of the usual two.

\medterm{Abadie Sign} Reduced pain sensation when firm pressure is applied to the Achilles tendon, classically described in tabes dorsalis.

\medterm{Abandoned Lead} A wire from a pacemaker, defibrillator, or other implanted electrical device that is no longer connected or used but remains inside the body. A working lead carries electrical signals to the heart and senses its rhythm; its metal core conducts electricity and its silicone or polyurethane covering protects nearby tissue. These materials are chosen to be biocompatible, corrosion-resistant, strong, and flexible, but no implanted material is completely risk-free. A lead may be left in place when removal would be more dangerous because scar tissue can firmly attach it to a vein or the heart. Retained leads can later complicate infection treatment, blood-vessel access, or placement of another device.

\synonyms
Retained Cardiac Device Lead; Abandoned Pacemaker Lead; Abandoned ICD Lead

\medterm{Abasia} Inability to walk even though the legs may have enough strength to move in other situations. It may result from a problem with coordination, the nervous system, or the way movement is planned, rather than paralysis.

\medterm{Abate} To become less intense or to reduce in severity, frequency, or amount; symptoms, inflammation, pain, or fever may abate.

\medterm{Abatement} A reduction in the severity, intensity, frequency, or amount of a symptom, disease process, exposure, or other condition.

\medterm{ABCDE Assessment} A systematic emergency assessment of airway, breathing, circulation, disability, and exposure. Life-threatening
abnormalities are treated as they are identified, with repeated reassessment.

\medterm{ABCDE Rule For Melanoma} A screening mnemonic for suspicious pigmented lesions: asymmetry, border irregularity, color variation,
diameter greater than about 6 mm, and evolution. It does not diagnose melanoma; suspicious lesions require clinical
assessment.

\medterm{Abdomen} The body region between the diaphragm and pelvis. It contains most digestive organs, the kidneys, spleen, pancreas, and
major blood vessels and is described clinically by four quadrants or nine regions.

\medterm{Abdominal} Relating to the abdomen, the body region between the chest and pelvis that contains organs of the digestive, urinary, and reproductive systems.

\medterm{Abdominal Abscess} A localized collection of pus within the abdomen, usually caused by bacterial infection, inflammation, surgery, or
organ perforation. It may cause fever and abdominal pain and often requires antibiotics and drainage.

\medterm{Abdominal Adhesion} Alternate index label; see \textit{Abdominal Adhesions}.

\medterm{Abdominal Adhesions} Bands of scar-like tissue that form inside the abdomen and make two or more organs stick to each other or to the abdominal wall. They most often develop after surgery but may also follow infection, inflammation, endometriosis, or radiation; some are present from birth. They often cause no symptoms, but may cause long-lasting abdominal pain, infertility when they affect pelvic reproductive organs, or a partial or complete intestinal blockage. A blockage may cause pain, bloating, constipation, vomiting, or inability to pass gas.

\medterm{Abdominal Angina} Pain in the abdomen after eating, caused by reduced blood flow through the arteries supplying the intestines. It is usually due to hardening and narrowing of these arteries.

\medterm{Abdominal Aorta} The portion of the aorta below the diaphragm that extends to its bifurcation into the common iliac arteries. It supplies
the abdominal organs, pelvic organs, and lower limbs through its branches.

\medterm{Abdominal Aortic Aneurysm} A weakened, widened area of the abdominal aorta, usually defined as a diameter of at least 3 cm or more than 50 percent above the expected size. Many cause no symptoms, but rapid growth or pain can signal danger. Rupture may cause sudden severe abdominal or back pain, dizziness, fainting, or collapse and requires immediate emergency care.

\synonyms
AAA; Abdominal Aneurysm

\medterm{Abdominal Aortic Aneurysm Repair} Surgical or endovascular treatment of an abdominal aortic aneurysm to prevent rupture or treat a rupture.

\synonyms
AAA Repair

\medterm{Abdominal Aortic Aneurysm Screening} Ultrasound examination used to detect an abdominal aortic aneurysm before symptoms occur. Screening
recommendations depend on age, sex, smoking history, and clinical risk; follow-up intervals depend on aneurysm size.

\medterm{Abdominal Aortic Ultrasound} An ultrasound examination of the abdominal aorta used to look for an aneurysm, measure its size, or monitor it over time. It is painless and does not use radiation.

\medterm{Abdominal Binder} A wide supportive garment placed around the abdomen to provide gentle compression after surgery, injury, or selected abdominal-wall conditions.

\medterm{Abdominal Bloating} A feeling of fullness, pressure, or tightness in the abdomen. It may occur with or without visible enlargement and is commonly caused by intestinal gas, constipation, food intolerance, or functional gastrointestinal disorders. Persistent or severe bloating may require evaluation.

\medterm{Abdominal Bruit} An abnormal blowing or whooshing sound heard over the abdomen with a stethoscope,
usually caused by turbulent blood flow through a narrowed or enlarged artery. It may occur with
renal-artery stenosis or an abdominal aortic aneurysm. Further assessment depends on symptoms and
vascular risk factors.

\medterm{Abdominal Cavity} The cavity between the diaphragm and pelvis containing the stomach, liver, gallbladder,
spleen, pancreas, kidneys, and intestines.

\medterm{Abdominal Circumference} The distance around the abdomen. In pregnancy, fetal abdominal circumference is measured by ultrasound at a standardized level that includes the stomach, portal sinus, and spine; it is used with other measurements to estimate fetal size and assess growth.

\medterm{Abdominal Colic} Cramping abdominal pain caused by contractions or obstruction of a hollow organ,
such as the intestine, ureter, or bile duct. Pain often comes in waves and may be accompanied by
nausea, vomiting, sweating, or altered bowel or urinary function. Severe or persistent pain requires
prompt medical assessment.

\medterm{Abdominal Compartment Syndrome} A dangerous rise in pressure inside the abdomen, usually above 20 millimeters of mercury, that causes new problems with organ function. Causes include severe fluid buildup, an overly swollen or blocked intestine, internal bleeding, and some conditions after abdominal surgery.

\medterm{Abdominal CT} Computed tomography of the abdomen used to evaluate selected causes of abdominal pain, trauma,
bowel obstruction, inflammation, bleeding, and masses. The protocol and use of contrast
depend on the suspected condition and the patient's clinical status.

\medterm{Abdominal CT Scan} Alternate name; see \textit{Abdominal CT}.

\medterm{Abdominal Distension} Visible or measured enlargement of the abdomen, often with a feeling of fullness or tightness. Causes include gas, constipation, fluid, pregnancy, enlarged organs, bowel obstruction, or bleeding; severe pain, vomiting, fever, or inability to pass stool or gas requires urgent assessment.

\synonyms
Abdominal Swelling; Swollen Abdomen; Abdominal Distention

\medterm{Abdominal Distention} Alternate name; see \textit{Abdominal Distension}.

\medterm{Abdominal Exploration} Surgical inspection of the abdominal cavity to identify or treat a suspected source of pain, bleeding, infection, obstruction, or other disease.

\medterm{Abdominal Fat} Adipose tissue stored beneath the skin or around abdominal organs; excess visceral fat is metabolically active and associated with insulin resistance, fatty liver, and cardiovascular risk.

\medterm{Abdominal Fluid} Fluid in the space around the abdominal organs. A small amount may be normal, but excess fluid, called ascites, can cause swelling, discomfort, breathing difficulty, or infection and needs evaluation for its cause.

\synonyms
Peritoneal Fluid

\medterm{Abdominal Girth} The distance around the abdomen, measured at a specified level. It may be used to monitor abdominal distention, fetal growth, obesity, ascites, or changes in intra-abdominal volume.

\medterm{Abdominal Guarding} Tensing of the abdominal muscles when the abdomen is touched. It may occur when the lining around the abdominal organs is irritated or inflamed.

\medterm{Abdominal Hernia} Protrusion of an organ or tissue through a weakness or opening in the abdominal wall. Common types include incisional, umbilical, epigastric, and ventral hernias.

\medterm{Abdominal Hysterectomy} Surgical removal of the uterus through an incision in the abdominal wall; the cervix, ovaries, and fallopian tubes may or may not also be removed.

\medterm{Abdominal Mass} A palpable or imaging-detected abnormal collection, enlargement, or lesion within the
abdomen. Causes include distended bowel or bladder, organ enlargement, hernia,
inflammatory collection, benign tumour, and malignancy; assessment depends on location,
mobility, tenderness, systemic features, and appropriate imaging.

\medterm{Abdominal Migraine} A migraine-related disorder characterized mainly by recurrent episodes of moderate to severe midline abdominal pain, often with nausea or vomiting and little or no headache.

\medterm{Abdominal MRI Scan} A magnetic resonance imaging study that produces detailed images of the abdomen using magnetic fields and radiofrequency energy.

\medterm{Abdominal Muscles} The muscles forming the front and sides of the abdominal wall. They support the trunk, help with posture and breathing, and increase intra-abdominal pressure during coughing, lifting, and defecation.

\synonyms
Abdominal Wall Muscles

\medterm{Abdominal Neoplasm} An abnormal growth in an organ or tissue of the abdomen. It may be noncancerous or cancerous and can cause pain, swelling, bleeding, changes in bowel function, or no symptoms.

\synonyms
Abdominal Tumor

\medterm{Abdominal Pain} Pain or discomfort in the area between the chest and pelvis. It may be felt deep inside, come from the abdominal wall, or be felt in the abdomen even when the cause is elsewhere. Sudden severe pain, rigid muscles, fainting, persistent vomiting, fever, blood in stool, or pregnancy requires prompt medical assessment.

\synonyms
Stomach Pain

\medterm{Abdominal Point Tenderness} Pain felt in one specific area when the abdomen is gently pressed. It may occur with appendicitis, gallbladder inflammation, or irritation of the abdominal lining; it does not diagnose a condition by itself.

\synonyms
Focal Abdominal Tenderness; Localized Peritoneal Tenderness; Point Pain in Abdomen

\medterm{Abdominal Pregnancy} An ectopic pregnancy implanted in the peritoneal cavity rather than the uterus, fallopian tubes, or ovaries.
It is rare and can cause severe internal bleeding. Diagnosis requires specialist imaging and management.

\medterm{Abdominal Quadrant} One of four topographical anatomical divisions of the abdomen created by perpendicular vertical and horizontal imaginary planes intersecting at the umbilicus (right upper, left upper, right lower, and left lower quadrants), utilized in clinical triage to localize abdominal pain and organ pathology.

\medterm{Abdominal Radiography} A regular X-ray of the abdomen. It may be taken while a person is lying down, standing, or lying on their side.

\medterm{Abdominal Reflex} Contraction of the abdominal muscles after stimulation of the abdominal skin. An absent or asymmetric reflex may
indicate neurologic disease, although it can be absent normally in some people.

\medterm{Abdominal Regions} The nine named areas used to describe locations on the abdominal wall: the right and left hypochondriac, epigastric, right and left lumbar, umbilical, right and left iliac, and hypogastric regions.

\medterm{Abdominal Retractor} A surgical instrument used to hold back the abdominal wall or organs to improve access and visibility during an operation.

\synonyms
Balfour Retractor; Bookwalter Retractor

\medterm{Abdominal Rigidity} Involuntary tensing of the abdominal wall, usually caused by irritation of the peritoneum. It may occur with
peritonitis, perforation, severe infection, or intra-abdominal bleeding and requires clinical assessment.

\medterm{Abdominal Sounds} Sounds produced by movement of gas and fluid through the intestines; changes in frequency or character may support, but do not establish, a diagnosis.

\medterm{Abdominal Swelling} Alternate name; see \textit{Abdominal Distension}.

\medterm{Abdominal Tap} Alternate name; see \textit{Abdominocentesis}.

\medterm{Abdominal Thrusts} A first-aid maneuver for severe choking in a conscious adult or child older than 1 year. Give 5 back blows followed by 5 quick inward and upward pushes above the navel, repeating until the object comes out or the person becomes unconscious. Use chest thrusts for a pregnant or very large person; infants and unconscious people need different first aid.

\medterm{Abdominal Ultrasound} An ultrasound examination that uses sound waves to produce images of the abdomen and assess its structure and blood flow when applicable.

\medterm{Abdominal Wall} The layered muscular, fascial, and skin-covered boundary surrounding the abdominal cavity.
It supports and protects abdominal organs and helps with posture, breathing, coughing, and straining.
Weakness or defects may produce abdominal or incisional hernias; trauma, infection, and tumors may
also affect the wall.

\medterm{Abdominal Wall Fat Pad Biopsy} Removal of a small sample of abdominal fat to look for amyloid, usually with a special Congo-red stain. It is less invasive than biopsy of an affected organ, but a negative result does not rule out amyloidosis.

\medterm{Abdominal Wall Hernia} Alternate name; see \textit{Abdominal Hernia}.

\medterm{Abdominal Wall Surgery} Surgery to repair a defect or disease of the abdominal wall, often a hernia. It may be open or laparoscopic, and mesh may or may not be used depending on the defect, location, and patient factors.

\medterm{Abdominal X-Ray} Alternate name; see \textit{Abdominal Radiography}.

\medterm{Abdominocentesis} Inserting a needle through the abdominal wall into the space around the abdominal organs to remove or test extra fluid, called ascites. It is also called paracentesis.

\medterm{Abdominopelvic Cavity} The connected body space below the diaphragm containing abdominal and pelvic organs, including digestive, urinary, and reproductive structures, protected by surrounding muscles, bones, and lining tissues.

\medterm{Abdominoperineal Resection} An operation that removes the anus, rectum, and usually part of the sigmoid colon. Because the normal anal outlet is removed, the surgeon creates a permanent colostomy, bringing the colon through the abdominal wall to drain stool into a bag. It may be used for selected low rectal or anal cancers.

\synonyms
APR

\medterm{Abdominoplasty} A surgical procedure that removes excess abdominal skin and may tighten the abdominal wall. It is commonly called a tummy tuck.

\synonyms
Tummy Tuck

\medterm{Abducens Nerve} The sixth cranial nerve. It controls the muscle that moves the eye outward. Damage to this nerve can prevent outward eye movement and cause double vision.

\medterm{Abducens Nerve Palsy} Dysfunction of the sixth cranial nerve resulting in paralysis or paresis of the lateral rectus muscle. Clinically characterized by horizontal diplopia that worsens on gaze toward the affected side and convergent strabismus (esotropia) in primary gaze. Common etiologies include microvascular ischemia (diabetes mellitus, hypertension), elevated intracranial pressure (false localizing sign), trauma, and cavernous sinus lesions.

\synonyms
Sixth Nerve Palsy; CN VI Palsy

\medterm{Abducent Nerve} Alternate name; see \textit{Abducens Nerve}.

\medterm{Abduct} To move an arm, leg, finger, or other body part away from the center line of the body.

\medterm{Abduction} Movement of a limb or other body part away from the median plane of the body. In ocular motility, it is movement of an eye outward, away from the nose.

\medterm{Abductor Muscle} A muscle that moves a body part away from the midline or, for a digit, away from the axis of the hand or foot.

\medterm{Aberrant Conduction} Transient abnormal intraventricular conduction of a supraventricular electrical impulse through the ventricular conduction system, typically occurring when an impulse reaches bundle branch tissue during its relative refractory period. Most commonly presents as a wide QRS complex with right bundle branch block morphology following a sudden change in cycle length (Ashman phenomenon).

\synonyms
Aberrancy; Ventricular Aberration

\medterm{Aberration} A deviation from normal structure, function, or appearance; in optics, it is an imperfection that prevents light from forming a sharply focused image.

\medterm{Abetalipoproteinemia} A rare inherited disorder in which the body cannot make normal particles that carry fat and cholesterol in the blood. It causes poor absorption of dietary fat, deficiencies of vitamins carried with fat, abnormal red blood cells, and possible nerve or eye disease.

\medterm{Abfraction} Loss of tooth material near the gum line, possibly related to repeated mechanical stress or biting forces. It may cause a notch and tooth sensitivity; a dentist can assess the cause and need for treatment.

\medterm{Abiotrophy} Progressive loss of function as cells in a tissue or organ gradually degenerate, often because of an inherited or developmental problem. It may cause worsening weakness, loss of sensation, poor coordination, or other problems depending on the affected part of the body.

\medterm{Ablate} To remove or destroy tissue using surgery, heat, cold, chemicals, or another energy source. Ablation may treat abnormal heart rhythms, tumors, pain pathways, or other selected conditions.

\medterm{Ablation} A procedure that removes or destroys tissue or stops an unwanted body function using surgery, heat, cold, electricity, medicines, or other methods.

\medterm{Ablative Surgery} Surgery that destroys or removes abnormal tissue, often using excision, heat, cold, electrical energy, or other tissue-destructive methods.

\medterm{Ablepharon} A rare congenital condition in which one or both eyelids are absent or severely underdeveloped. It can leave the eye exposed and may require protective treatment and reconstructive surgery.

\medterm{Abnormal} Different from the usual healthy structure, function, or appearance. An abnormal finding may be harmless or may indicate a disease, depending on its type and context.

\medterm{Abnormal (Dysfunctional) Uterine Bleeding} Bleeding from the uterus that is abnormal in timing, regularity, duration, or amount and is not explained by pregnancy. Causes include structural lesions, ovulatory or hormonal disorders, medicines, bleeding disorders, and endometrial disease; evaluation depends on age and risk factors.

\medterm{Abnormal Color} A change in the usual color of skin, moist tissues, body fluids, or another structure. Its meaning depends on the color, location, duration, and associated symptoms.

\synonyms
Discoloration

\medterm{Abnormal Fissure} An atypical crack, groove, or cleft in a tissue or organ. It may be congenital or develop from injury, inflammation, scarring, pressure, or disease, and its importance depends on its location and depth.

\medterm{Abnormal Hemoglobins Testing} Blood testing that looks for unusual forms of hemoglobin, the oxygen-carrying protein in red blood cells. It helps diagnose inherited disorders such as sickle-cell disease and thalassemia; results may require specialized laboratory analysis.

\medterm{Abnormal Labor} Labor that is unusually slow, prolonged, arrested, or associated with abnormal contractions, fetal position, or fetal intolerance. Evaluation considers cervical change, contraction pattern, pelvic factors, fetal status, infection, and maternal condition before treatment or delivery decisions.

\medterm{Abnormal Tooth Shape} Abnormal tooth shape may result from developmental defects, genetic conditions, trauma, wear, or tumors and can affect function, appearance, and eruption.

\medterm{Abnormal Uterine Bleeding} Bleeding from the uterus that is unusually heavy, prolonged, irregular, or occurs between periods, after sex, or after menopause. Causes include pregnancy-related problems, hormonal changes, fibroids, infection, medicines, and cancer. Persistent or heavy bleeding may indicate an underlying condition.

\medterm{Abnormal Virilization} Development of masculinizing physical features in a person whose expected sex development does not include them, often from excess androgens.

\synonyms
Virilization

\medterm{Abnormality} A difference from the usual structure, function, development, or test result. An abnormality may be harmless, temporary, or a sign of disease; its importance depends on the finding, symptoms, and clinical evaluation.

\synonyms
Anomaly

\medterm{Abnormally Dark Or Light Skin} Skin that has become darker or lighter than a person's usual skin tone. Darkening is called
hyperpigmentation and lightening is called hypopigmentation. Causes include inflammation or injury, sun exposure,
medicines, infections, hormonal disorders, and genetic conditions.

\medterm{Abnormally Excessive Sweating} See \textit{Hyperhidrosis}.

\medterm{ABO Blood Group System} The major red-cell blood-group system based on whether A and/or B antigens are present on the red-cell surface and the corresponding antibodies are present in plasma. The four principal phenotypes are A, B, AB, and O; correct ABO matching is critical for safe transfusion.

\synonyms
ABO Blood Group
ABO Blood Groups
ABO Blood Types
ABO System

\medterm{ABO Blood Typing} A laboratory test that determines a person's ABO and Rh blood groups. It is performed before transfusion,
transplantation, and some obstetric procedures.

\medterm{ABO Incompatibility} An immune reaction that can occur when blood with an incompatible A, B, or O type is transfused, or when maternal and fetal blood types are incompatible. It can destroy red blood cells, cause a severe transfusion reaction, or affect a newborn.

\medterm{Abortion} The ending of a pregnancy, either spontaneously or intentionally. An intentional abortion uses medicines or a medical procedure; spontaneous abortion is also called miscarriage.

\medterm{Abortion Complications} Potential complications after spontaneous or induced abortion, including hemorrhage, retained tissue, infection, uterine injury, cervical trauma, thromboembolism, and emotional or psychological distress.

\medterm{Abortive} Incomplete, unsuccessful, or stopped before full development; an abortive infection produces limited disease without the usual complete clinical course.

\medterm{Abortive Polio} A mild, nonparalytic poliovirus infection causing nonspecific fever, sore throat, headache, nausea, or malaise without the neurologic findings of paralytic polio.

\medterm{Above The Knee Amputation} Amputation of the lower limb through the thigh above the knee joint, usually performed for severe vascular disease, infection, trauma, or cancer.

\medterm{ABR Test} Auditory brainstem response testing records electrical activity from the auditory nerve and brainstem through scalp electrodes to assess hearing and auditory-pathway function, especially in infants or people unable to provide reliable behavioral responses.

\medterm{Abrade} To wear away tissue by friction or erosion; an abrasion is the resulting scrape or superficial loss of skin.

\medterm{Abrasion} A superficial injury in which the skin is scraped or rubbed away, usually by friction. In dentistry, abrasion means mechanical wearing away of tooth structure by an external force.

\medterm{Abrin} A highly poisonous protein found in rosary pea seeds. It stops cells from making proteins, causing them to die. Swallowing chewed seeds or exposure to abrin powder can cause vomiting, diarrhea, dehydration, breathing problems, organ failure, and death. Suspected exposure requires urgent medical care; there is no routine antidote.

\medterm{Abruptio Placenta} Alternate name; see \textit{Abruptio Placentae}.

\medterm{Abruptio Placentae} Premature separation of the placenta from the uterine wall before delivery. It can cause vaginal bleeding,
abdominal pain, uterine tenderness, fetal distress, and maternal shock and requires urgent obstetric assessment.

\medterm{Abs} An informal name for the muscles at the front and sides of the abdomen. They help bend and rotate the trunk, support the spine, protect abdominal organs, and increase pressure for coughing, lifting, breathing, or bowel movements.

\medterm{Abscess} A pocket of pus that forms in tissue, an organ, or a body cavity, usually because of an infection. It may cause pain, swelling, warmth, and fever; treatment often requires drainage and sometimes antibiotics.

\medterm{Abscess, Bartholin's} A pocket of pus caused by infection in a Bartholin gland near the vaginal opening. It usually causes a painful, swollen lump and may need drainage; a fever or severe pain needs medical care.

\medterm{Abscess, Cutaneous} Alternate index label; see \textit{Boils}.

\medterm{Abscessed Tooth} A collection of pus caused by bacterial infection of the tooth or surrounding tissues, usually arising from
untreated dental decay, a cracked tooth, or periodontal disease. Pain, facial swelling, fever, and sensitivity may
occur. Dental treatment and drainage are usually required; swelling that affects swallowing or breathing is an
emergency.

\medterm{Abscission} The separation or removal of tissue, either naturally or by cutting. In surgery it may describe removing a growth, damaged tissue, or part of an organ; the exact meaning depends on the medical context.

\medterm{Absence Epilepsy} A generalized epilepsy characterized by brief episodes of impaired awareness, often with staring and subtle movements.

\synonyms
Childhood Absence Epilepsy; Pyknolepsy; Petit Mal Epilepsy

\medterm{Absence Of Sweating} Anhidrosis is diminished or absent sweating, which may be focal or generalized. It impairs thermoregulation and can result from autonomic dysfunction, peripheral neuropathy, skin disease, medications, or congenital conditions. Severe generalized anhidrosis can cause heat intolerance and exertional hyperthermia.

\medterm{Absence Of The Breast} Congenital absence of breast development, called amastia; it may affect one or both sides and can occur with other developmental anomalies.

\medterm{Absence Of The Nipple} Athelia: a rare congenital absence of one or both nipples, which may occur with hereditary disorders, chromosomal abnormalities, or developmental syndromes.

\medterm{Absence Seizure} A brief seizure causing sudden staring and reduced awareness, usually lasting only seconds with rapid recovery. The person may stop speaking or responding and may make small repetitive movements; an EEG helps confirm the diagnosis.

\synonyms
Petit Mal Seizure; Typical Absence Seizure

\medterm{Absence Seizure Cross-Reference} Alternate name; see \textit{Absence seizure}.

\medterm{Absent Breath Sounds} Lack of audible air movement over part or all of a lung. Causes include pneumothorax, pleural effusion, or
airway obstruction; sudden onset requires urgent assessment.

\medterm{Absent Eye} Anophthalmia, a congenital absence or severe underdevelopment of the eye globe and ocular tissue; it may affect one or both eyes and may occur with other birth defects.

\medterm{Absent Menstrual Periods} Alternate index label; see \textit{Amenorrhea}.

\medterm{Absent Menstrual Periods, Primary} Alternate index label; see \textit{Amenorrhea, Primary}.

\medterm{Absent Menstrual Periods, Secondary} Alternate index label; see \textit{Amenorrhea, Secondary}.

\medterm{Absent Pulmonary Valve} A birth defect in which the valve between the heart and the lung arteries is missing or severely underdeveloped. It can enlarge those arteries and cause breathing problems, poor blood flow to the lungs, or other heart defects.

\medterm{Absent Reflex} Failure of an expected automatic movement to occur when a tendon or other body area is stimulated. It may suggest a problem in a nerve, nerve root, spinal cord, or muscle, although reflexes can normally differ between people and sides.

\medterm{Absent Thirst} Reduced or absent desire to drink despite a potential need for fluids, which can occur with hypothalamic disease, aging, altered consciousness, medications, or severe illness.

\medterm{Absolute Eosinophil Count} A laboratory blood test measuring the exact number of eosinophils per microliter of blood to evaluate allergic disorders, asthma, parasitic infections, drug reactions, or hypereosinophilic syndromes.

\synonyms
AEC; Absolute Eosinophils; Blood Eosinophil Count

\medterm{Absolute Neutrophil Count} The number of neutrophils, a type of white blood cell that fights many bacterial and fungal infections, in a volume of blood. A low count increases infection risk, especially when it is severely reduced.

\synonyms
ANC

\medterm{Absorb} To take in a substance, such as through the skin or intestine, or to reduce the intensity of radiation or light by taking up its energy.

\medterm{Absorption} The process by which a drug enters the bloodstream from its site of administration, such
as the digestive tract, skin, or muscle.

\medterm{Abstinence} Deliberate avoidance of a substance or behavior, such as alcohol, drugs, tobacco, or sexual activity. Stopping some substances, especially alcohol, sedatives, or opioids, can cause dangerous withdrawal and may require medical supervision.

\medterm{Abulia} Marked reduction in motivation, initiative, and spontaneous activity, usually caused by dysfunction of frontal-subcortical brain circuits or related neurologic disease.

\medterm{Abusive Head Trauma} Injury to an infant or young child caused by inflicted blunt impact, violent shaking, or
both. Findings may include subdural hemorrhage, retinal hemorrhage, encephalopathy,
fractures, or other injuries; diagnosis requires careful multidisciplinary assessment
because no single finding is diagnostic.

\synonyms
Shaken Baby Syndrome; Non-Accidental Head Injury

\medterm{Abutment} A natural tooth, dental implant, or component that supports a dental prosthesis. Abutments may hold a bridge, denture, or implant crown in position and transfer chewing forces to the supporting tooth, implant, or bone.

\medterm{AC Joint} The acromioclavicular joint, a gliding joint between the acromion of the scapula and the clavicle at the top of the shoulder.

\medterm{Acalculia} Loss of the ability to do or understand numbers and calculations that a person could previously manage. It may result from injury or disease affecting areas of the brain involved in calculation.

\medterm{Acalculous Cholecystitis} Inflammation of the gallbladder without gallstones. It usually occurs suddenly in people who are critically ill, injured, recovering from surgery, unable to eat, or in shock. Reduced bile flow and poor blood supply can contribute, and severe cases may cause tissue death or a hole in the gallbladder. Fever and pain in the upper-right abdomen may occur; ultrasound is commonly used for evaluation.

\medterm{Acanthamoeba} A free-living amoeba found in soil and water that can cause corneal infection, skin or sinus infection, or rarely granulomatous amoebic encephalitis. Contact-lens exposure to contaminated water increases keratitis risk, while invasive disease is more likely with impaired immunity.

\medterm{Acanthamoeba Keratitis} A rare but potentially sight-threatening infection of the cornea caused by the free-living amoeba
Acanthamoeba. It is strongly associated with contaminated contact-lens water or poor lens hygiene and may cause
severe eye pain, redness, photophobia, blurred vision, and a ring-shaped corneal infiltrate. Prompt ophthalmic
assessment and prolonged antiamoebic treatment are required.

\medterm{Acanthocyte} An abnormal red blood cell with irregular, thorny or spike-like projections of varying length and width on its outer membrane. They are seen on a blood smear in advanced liver disease (such as cirrhosis), lipid metabolism disorders, abetalipoproteinemia, or neuroacanthocytosis.

\synonyms
Spur Cell

\medterm{Acanthocytosis} An abnormal hematologic condition characterized by the presence of circulating red blood cells with irregular, thorny or spiked projections (acanthocytes or spur cells). It is classically observed in abetalipoproteinemia, advanced liver cirrhosis (spur cell anemia), severe malnutrition, and neuroacanthocytosis syndromes.

\synonyms
Spur Cell Formation

\medterm{Acantholysis} Loss of attachment between cells in the outer layer of the skin, causing them to separate. This can produce fragile blisters or erosions and occurs in some autoimmune, inherited, infectious, or drug-related skin disorders.

\medterm{Acanthoma} A usually benign skin growth made of thickened squamous epidermal cells. Its appearance can resemble other lesions, so a changing or uncertain lesion may need examination or biopsy.

\medterm{Acanthosis Nigricans} Dark, thickened, velvety skin, usually in folds such as the neck, armpits, or groin. It is often associated with insulin resistance, obesity, type 2 diabetes, hormonal disorders, or medicines; rarely, sudden extensive disease can signal an underlying cancer. Treating the cause may improve the skin changes.

\medterm{Acapnia} An unusually low carbon-dioxide level in arterial blood, usually caused by breathing too quickly or deeply. It may cause light-headedness, tingling, or a change in the blood's acid level.

\medterm{Acardia} Absence or severe malformation of the fetal heart, usually occurring in the setting of a
rare twin perfusion sequence. The affected fetus cannot sustain independent circulation.

\medterm{Acariasis} Illness or skin infestation caused by mites or ticks. It may cause itching, a rash, burrows, skin nodules, or deeper tissue inflammation; some mites and ticks can also transmit infection.

\medterm{Acaricide} A pesticide used to kill mites and related arachnids such as ticks. It may be used on animals, plants, or environments; exposure can irritate skin and may cause poisoning if misused.

\medterm{Acatalasemia} A rare inherited condition in which catalase, an enzyme that breaks down hydrogen peroxide, is absent or greatly reduced. Many people have no symptoms, but some develop mouth ulcers, tissue damage, or increased sensitivity to oxidative stress.

\medterm{Acatalasia} A rare inherited condition in which catalase, an enzyme breaking down hydrogen peroxide, is absent or greatly reduced. Some people have mouth ulcers, infections, or tissue damage, while others have no symptoms.

\synonyms
Catalase Deficiency

\medterm{Accelerated Growth} An increase in a child's growth velocity above the expected rate for age and developmental stage, determined from serial height or length measurements. It may occur during normal puberty or catch-up growth; abnormal acceleration, especially with early pubertal signs or advanced bone age, may suggest precocious puberty or another endocrine disorder.

\medterm{Accelerated Idioventricular Rhythm} A heart rhythm that starts in the lower chambers and is usually slower than most dangerous ventricular rhythms. It has wide beats and may occur during recovery of blood flow to the heart. It is often temporary, but the person and the cause still need assessment.

\synonyms
AIVR

\medterm{Accelerometry} Measurement of movement or acceleration using a sensor. It can assess physical activity, walking, tremor, falls, sleep, or rehabilitation progress, depending on the device and the recording method.

\medterm{Acceptable Daily Intake} The estimated amount of a substance that a person can consume every day throughout life without a meaningful health risk, based on available evidence and safety margins. It is usually expressed per kilogram of body weight.

\synonyms
ADI

\medterm{Acceptance-Based Therapies} Psychotherapeutic approaches that help a person notice and accept thoughts, feelings, or bodily sensations while reducing unhelpful attempts to control or avoid them.

\synonyms
Acceptance-Based Psychotherapy; Acceptance-Focused Therapy

\medterm{Access Needle} A hollow needle used to enter a blood vessel, body cavity, organ, or fluid collection. It may be used to take a sample, drain fluid, give treatment, or help place a catheter.

\synonyms
Introducer Needle; Percutaneous Access Needle; Vascular Access Needle

\medterm{Accessory} Additional or supplementary; in anatomy, an accessory structure is an extra part or a structure that assists the usual organ or system.

\medterm{Accessory Digestive Organ} An organ that supports digestion but is not part of the alimentary canal, such as the liver, pancreas, gallbladder, or salivary glands.

\medterm{Accessory Dwelling Unit} A secondary, independent residential unit on the same property as a primary home; it is a housing and community-health term rather than a disease or body structure.

\medterm{Accessory Muscles} Muscles recruited to assist the primary muscles of respiration or another movement when the normal workload is increased; use of respiratory accessory muscles may indicate respiratory distress.

\medterm{Accessory Nerve} The eleventh cranial nerve (CN XI), a primarily motor nerve whose spinal component supplies the sternocleidomastoid and trapezius muscles. It enables head rotation and shoulder elevation; injury may cause weak head rotation, impaired shoulder shrug, and shoulder droop.

\synonyms
Cranial Nerve XI; Spinal Accessory Nerve

\medterm{Accessory Pathway} An abnormal electrical connection between the atria and ventricles that can bypass the atrioventricular node and cause pre-excitation or tachyarrhythmias, including atrioventricular reentrant tachycardia.

\medterm{Accessory Placenta} An extra placental lobe separate from the main placenta but connected to it by fetal blood vessels; it can increase the risk of retained placental tissue or vessel injury during delivery.

\medterm{Accident And Emergency Medicine} Accident and emergency medicine is the specialty concerned with immediate evaluation and treatment of acute illness, injury, poisoning, and life-threatening emergencies.

\medterm{Accidental Death} Death caused by an unintentional injury or event rather than disease, suicide, or homicide. The cause is determined through medical, scene, and sometimes legal investigation.

\synonyms
Unintentional Death

\medterm{Accidental Fall} An unintentional movement to the ground or a lower level. Falls can cause bruises, fractures, head injury, or internal bleeding, especially in older adults; loss of consciousness, severe pain, or confusion requires urgent care.

\synonyms
Fall Injury

\medterm{Accommodation} The ability of the eye to change the shape of its lens so that near objects can be focused clearly.

\medterm{Accommodation Reflex} The coordinated adjustment of the lens, ciliary muscle, and pupils that shifts the eye from distance to near focus; it includes lens thickening, convergence, and pupillary constriction.

\medterm{Accommodative Insufficiency} Inability of the eye to increase lens power adequately for near focusing, causing eyestrain, blurred near vision, headache, or difficulty reading.

\medterm{Accouchement} Accouchement is the process of childbirth, including labor, delivery of the infant, and expulsion of the placenta.

\medterm{Accretion} Gradual buildup or growth of material that makes an object or body tissue larger. Examples include crystals forming in fluid, extra bone developing after injury, or plaque collecting on teeth.

\medterm{ACE Blood Test} A blood test measuring angiotensin-converting enzyme. It may support evaluation of sarcoidosis and some other inflammatory or liver disorders, but an abnormal result is not specific and cannot diagnose a disease by itself.

\medterm{ACE Inhibitors} Medicines that reduce formation of angiotensin II, lowering blood pressure and pressure within the kidney's filtering units. They are used for hypertension, heart failure, and selected kidney or cardiovascular conditions; cough, high potassium, kidney-function changes, and rare facial or airway swelling are important risks, and they are generally avoided in pregnancy.

\medterm{Acebutolol} Acebutolol is a beta-adrenergic blocker with intrinsic sympathomimetic activity used in some settings for hypertension and cardiac rhythm disorders; it can slow heart rate and conduction.

\medterm{Acellular} Without cells. An acellular material may consist of fluid, minerals, or the supporting material around cells and may be natural, diseased, or used in a medical product.

\medterm{Acentric Chromosome} A chromosome fragment that lacks a centromere and therefore usually cannot attach normally to the spindle during cell division.

\medterm{Acephalostomia} An extremely rare and severe birth defect involving absence or major failure of development of the head and mouth. It is usually incompatible with survival and is identified during prenatal or newborn examination.

\medterm{Acetabular Labrum} A fibrocartilaginous rim attached to the margin of the acetabulum that deepens the hip
socket, helps maintain a fluid seal, and contributes to hip-joint stability.

\medterm{Acetabular Notch} The indentation in the inferior margin of the acetabulum, bridged by the transverse
ligament of the acetabulum and forming a passage for structures entering the hip joint.

\medterm{Acetabuloplasty} A surgical procedure that reshapes or augments the acetabular rim or socket to improve
coverage of the femoral head in selected structural hip disorders.

\medterm{Acetabulum} The cup-shaped socket of the hip formed by the fused ilium, ischium, and pubis. It articulates with the femoral head
and is deepened by the acetabular labrum.

\medterm{Acetaldehyde} A chemical produced when the body metabolizes alcohol. Accumulation may contribute to flushing, nausea, and vomiting.

\medterm{Acetaminophen} A widely used medication that relieves mild-to-moderate pain and reduces fever, but lacks significant anti-inflammatory effects. While safe at recommended therapeutic doses, excessive amounts can cause severe, life-threatening liver injury, especially when taken with alcohol or other products containing acetaminophen.

\synonyms
Paracetamol; APAP

\medterm{Acetaminophen And Codeine Overdose} A medical emergency caused by ingesting toxic amounts of a combination pain medicine. It produces dual life-threatening effects: codeine suppresses the central nervous system, causing extreme drowsiness, slowed or stopped breathing, and coma, while acetaminophen depletes protective liver antioxidants, leading to progressive liver failure.

\synonyms
Co-Codamol Overdose; Tylenol with Codeine Overdose

\medterm{Acetaminophen Overdose} Taking more acetaminophen than the body can safely process, causing serious liver damage. A person may have few symptoms at first, followed later by nausea, right-sided abdominal pain, jaundice, or liver failure. Get emergency help or contact a poison center immediately; the antidote N-acetylcysteine works best when started early but may still help when given later.

\synonyms
Acetaminophen Toxicity; Acetaminophen Poisoning; Paracetamol Toxicity; Paracetamol Overdose; APAP Overdose

\medterm{Acetaminophen Poisoning} Alternate name; see \textit{Acetaminophen Overdose}.

\medterm{Acetaminophen Toxicity} Alternate name; see \textit{Acetaminophen Overdose}.

\medterm{Acetate} A form of acetic acid and a small molecule used by cells in energy and building processes. The body can convert acetate into acetyl-CoA, which helps make or break down fats and other substances.

\medterm{Acetazolamide} A medicine that changes salt and fluid handling in the kidneys. It is used for glaucoma, prevention or treatment of altitude illness, and selected fluid, metabolic, or neurologic conditions.

\medterm{Acetone} A volatile ketone body produced during increased fat metabolism. It may contribute to the characteristic breath odor of ketosis and ketoacidosis.

\medterm{Acetone Body} A ketone produced when the body breaks down fat for energy. Acetone is one of three main ketones and can contribute to the distinctive breath odor of prolonged fasting or diabetic ketoacidosis.

\medterm{Acetone Breath} A distinctive sweet, fruity odor on the breath caused by the exhalation of volatile acetone, a ketone body produced during accelerated lipolysis and ketogenesis. It is a characteristic physical examination finding in diabetic ketoacidosis (DKA), as well as in severe starvation, prolonged fasting, and ketogenic metabolic states.
\synonyms{Fruity Breath; Ketotic Breath; Diabetic Breath Odor}

\medterm{Acetone Poisoning} Harm from swallowing or breathing a large amount of acetone, a solvent found in some products. It may cause irritation, dizziness, vomiting, drowsiness, low blood pressure, or breathing problems; significant exposure is a medical emergency requiring prompt treatment.

\medterm{Acetonuria} The presence of acetone or other ketone bodies in the urine, which may occur with fasting, vomiting, uncontrolled diabetes, or other states of increased fat metabolism.

\medterm{Acetyl-CoA Carboxylase} An enzyme that begins the process of making fatty acids by changing acetyl-CoA into malonyl-CoA. Its activity responds to the cell’s energy supply and to hormones, helping balance fat production with energy needs.

\medterm{Acetyl-Coenzyme A} A substance used in many energy and building processes inside cells. It helps release energy from food and provides material for making fatty acids, cholesterol, and ketones.

\medterm{Acetylation} A chemical change that adds a small acetyl group to a protein, medicine, or DNA-related molecule. This can alter how the molecule works, how long it lasts, how it is moved, or how strongly a gene is active.

\medterm{Acetylcholine} A major neurotransmitter synthesized from choline and acetyl-CoA by choline acetyltransferase. Operating throughout the central and peripheral nervous systems, it mediates neuromuscular transmission at motor endplates, autonomic ganglionic neurotransmission, parasympathetic postganglionic signaling, and central cognitive/memory processing.

\synonyms
ACh

\medterm{Acetylcholine Receptor Antibody} Acetylcholine receptor antibody testing detects antibodies associated with autoimmune myasthenia gravis and is interpreted with the neurologic examination and other studies.

\medterm{Acetylcholinesterase} A rapid-acting serine hydrolase enzyme localized at synaptic clefts and neuromuscular junctions that hydrolyzes the neurotransmitter acetylcholine into choline and acetate, terminating signal transmission and allowing tissue repolarization.

\synonyms
AChE; Acetylcholine Acetylhydrolase

\medterm{Acetylcholinesterase Inhibitors} A group of medicines that slow the breakdown of acetylcholine, a chemical messenger used by nerves. This increases acetylcholine activity and can help manage conditions such as Alzheimer disease and myasthenia gravis or reverse some medicine-related muscle paralysis.

\medterm{Acetylcysteine} A medicine that loosens thick mucus so it is easier to clear. It is also the main antidote for acetaminophen overdose because it restores glutathione, a substance the liver uses to neutralize the overdose’s toxic breakdown product.

\medterm{Acetylsalicylic Acid} Alternate index label; see \textit{Aspirin}.

\medterm{ACh (Acetylcholine)} Alternate name; see \textit{Acetylcholine}.

\medterm{Achalasia} A disorder in which the muscle ring between the esophagus and stomach does not relax properly, and the esophagus cannot move food normally into the stomach. It causes difficulty swallowing and regurgitation.

\synonyms
Cardiospasm; Esophageal Achalasia

\medterm{Achalasia Of The Cardia} A disorder in which the valve between the esophagus and stomach does not relax properly and the esophagus loses its normal squeezing movement. Food and liquid can collect in the esophagus, causing progressive difficulty swallowing, regurgitation, chest discomfort, and weight loss.

\medterm{AChE (Acetylcholinesterase)} Alternate name; see \textit{Acetylcholinesterase}.

\medterm{Achenbach's Syndrome} A benign condition of recurrent, painful, spontaneous hematoma formation in the pulp of
a finger or thumb, typically affecting middle-aged women. It presents with sudden onset
of throbbing pain and blue-black discoloration of the fingertip, resolving spontaneously
over one to two weeks without treatment.

\medterm{Achilles Tendinitis} Painful irritation or degeneration of the tendon joining the calf muscles to the heel. It usually follows repeated loading and can cause stiffness, swelling, tenderness, and pain during walking or exercise.

\synonyms
Achilles Tendinopathy; Calcaneal Tendinitis

\medterm{Achilles Tendinopathy} Pain and structural disorder of the Achilles tendon, usually caused by repeated loading. It causes posterior
heel pain and stiffness and is managed with activity modification, rehabilitation, and treatment of contributing factors.

\medterm{Achilles Tendon} The thickest and strongest tendon in the human body, connecting the calf muscles (gastrocnemius and soleus) to the heel bone (calcaneus). It transmits muscle force to the foot, enabling the ankle to point downward (plantarflexion) during walking, running, jumping, and standing on tiptoe.

\synonyms
Calcaneal Tendon; Tendo Calcaneus

\medterm{Achilles Tendon Repair} Surgery that reconnects a completely or severely torn Achilles tendon. The ankle is protected in a cast or brace afterward, followed by gradual exercises to restore movement and strength; recovery takes several months.

\medterm{Achilles Tendon Rupture} A partial or complete tear of the Achilles tendon, usually caused by sudden force. It causes posterior ankle
pain, weakness in plantar flexion, and difficulty walking or standing on tiptoe.

\medterm{Achilles Tendonitis} A commonly used term for pain and inflammation or degeneration of the Achilles tendon. It may affect the tendon
at its heel insertion or its mid-portion and is often related to repetitive loading.

\medterm{Achillobursitis} Inflammation of a bursa near the Achilles tendon, causing pain, tenderness, swelling, or friction-related discomfort at the back of the heel.

\medterm{Achlorhydria} A condition in which the stomach makes little or no hydrochloric acid. It may result from autoimmune gastritis, long-term acid-suppressing medicines, or loss of acid-producing cells and can impair absorption of iron and vitamin B12.

\medterm{Achondrogenesis} A group of severe inherited disorders in which cartilage and bone development is profoundly impaired. Affected fetuses have very short limbs and poorly formed bones, and most cases are fatal before or shortly after birth.

\medterm{Achondroplasia} An autosomal dominant bone-growth disorder caused by a disease-causing change in the \textit{FGFR3} gene. It causes short arms and legs, short stature, a larger head, and a prominent forehead. One altered gene copy usually causes the condition; two copies cause a severe, often fatal form.

\synonyms
Chondrodystrophy; Achondroplastic Dwarfism

\medterm{Achoo Syndrome} The photic sneeze reflex: repeated sneezing triggered by sudden exposure to bright light, especially after entering sunlight from darkness.

\medterm{Achromatopsia} An inherited disorder causing absent or greatly reduced color vision, often with light sensitivity and poor visual acuity.

\synonyms
Complete Color Blindness

\medterm{Achromia} Achromia is loss or marked reduction of pigment in skin, hair, eye, or another tissue; causes include genetic disorders, autoimmune depigmentation, injury, and chemical exposure.

\medterm{Aciclovir} Aciclovir, also called acyclovir, is an antiviral drug that inhibits herpesvirus DNA polymerase and is used for infections such as herpes simplex and varicella-zoster.

\medterm{Acid} A substance that increases acidity by releasing hydrogen ions or accepting electrons. Strong acids can damage tissue, while acids also take part in normal digestion, metabolism, and body-fluid regulation.

\medterm{Acid Base Balance} The body’s control of acids and bases in its fluids, especially the blood. The lungs remove carbon dioxide and the kidneys adjust acids and bicarbonate to keep blood pH within a narrow safe range.

\medterm{Acid Indigestion} Dyspepsia or upper-abdominal discomfort that may include burning, fullness, belching, or nausea; it can accompany reflux, gastritis, ulcer disease, or other conditions.

\medterm{Acid Loading Test} A specialized test of how well the kidneys remove acid from the body. After a carefully controlled acid dose, urine is measured for acidity and other substances to investigate suspected kidney-tubule problems.

\medterm{Acid Mucopolysaccharides} An older name for long sugar chains found in supporting tissues such as cartilage, skin, and joint fluid. Abnormal breakdown or storage of these substances causes a group of inherited storage disorders.

\medterm{Acid Phosphatase} An enzyme that works best in acidic conditions and is found in several tissues, including prostate and blood cells.

\synonyms
ACP

\medterm{Acid Reflux} Backward flow of acidic stomach contents into the esophagus, causing heartburn, sour regurgitation, cough, or throat irritation. Frequent reflux may be gastroesophageal reflux disease and can injure the esophagus.

\medterm{Acid Soldering Flux Poisoning} Injury from swallowing or inhaling acidic chemicals used in soldering. It may burn the mouth, throat, skin, or eyes and can cause vomiting, breathing difficulty, or dangerous changes in body chemistry; significant exposure requires urgent medical treatment.

\medterm{Acid Sphingomyelinase Deficiency} An inherited lysosomal storage disorder caused by deficient acid sphingomyelinase, usually because of disease-causing variants in the \textit{SMPD1} gene. Sphingomyelin accumulates in cells and may cause enlarged liver or spleen, lung disease, bone changes, or neurologic impairment; it is also called Niemann--Pick disease types A and B.

\medterm{Acid-Base Analysis} Evaluation of blood pH, bicarbonate, and carbon dioxide to identify metabolic or respiratory disturbances,
including those caused by poisoning.

\medterm{Acid-Base Disorder} A problem that makes the blood too acidic or too alkaline. It may result from changes in breathing, kidney function, metabolism, or severe illness. Blood tests help identify the cause and whether the body is compensating.

\medterm{Acid-Base, Fluid, and Electrolyte Disorders} Disorders that disturb blood acidity, body-water balance, or concentrations of minerals such as sodium, potassium, calcium, magnesium, and phosphate. They may arise from kidney, heart, endocrine, gastrointestinal, medication, or critical-illness problems and can affect the brain, muscles, and heart.

\medterm{Acid-Fast Bacilli} Rod-shaped bacteria that resist decolorization by acid-alcohol solutions during laboratory staining procedures due to the high lipid and mycolic acid content of their cell walls. The group classically includes \textit{Mycobacterium tuberculosis} and \textit{Mycobacterium leprae}.

\synonyms
AFB; Acid-Fast Bacteria

\medterm{Acid-Fast Bacillus Smear} Microscopic examination of a specimen, often sputum, for acid-fast organisms using a
special stain. A positive result supports mycobacterial infection but does not by itself
identify the species or establish drug susceptibility; culture or molecular testing is
also required.

\medterm{Acid-Fast Stain} A special dye test used to find mycobacteria and a few related organisms in samples such as sputum. A positive result does not identify the species or show which medicines will work.

\medterm{Acidemia} An abnormally low arterial blood pH (below 7.35), indicating an increased hydrogen ion concentration in circulating blood. It is the measurable laboratory state resulting from an underlying respiratory or metabolic acidosis.

\medterm{Acidophil} A cell or tissue part that absorbs acidic laboratory dyes strongly, especially the pink dye eosin. The pattern helps pathologists recognize certain cells in tissue samples, including some pituitary cells.

\medterm{Acidophilus} Lactobacillus acidophilus, a lactic-acid-producing bacterium found in the intestinal and genital microbiota and used in some fermented foods and probiotic products.

\medterm{Acidosis} A condition in which the body’s fluids become too acidic. It can result from breathing problems, kidney disease, severe infection, diabetes, poisoning, or other illness and may impair organ function when severe.

\medterm{Aciduria} Aciduria means an abnormally acidic urine or the presence of excess organic acids in urine; it may reflect diet, dehydration, medication, or an inherited metabolic disorder.

\medterm{Acinar Cell} A cell in a small cluster that makes and releases a fluid or substance. In the pancreas, acinar cells release digestive enzymes into small ducts.

\medterm{Acinetobacter Baumannii Resistance} Resistance means that Acinetobacter baumannii is not killed by one or more antibiotics. It can cause difficult hospital infections, especially in intensive-care units, so culture and susceptibility testing guide treatment.

\medterm{Acinus} A small working unit of a gland made of cells that release their product into a duct. Examples include the secretory units of the pancreas and salivary glands.

\medterm{ACL Injury} A sprain or tear of the strong ligament inside the knee that helps prevent the shinbone from sliding too far forward. It commonly follows twisting or sudden stopping and may cause a pop, swelling, pain, or a feeling that the knee gives way.

\medterm{ACL Reconstruction} Surgery that replaces a torn anterior cruciate ligament inside the knee with a tendon graft. It may restore stability after injury, especially when the knee gives way or the person plans activities requiring pivoting; rehabilitation is required afterward.

\medterm{Acne} A common inflammatory disorder of hair follicles and oil glands. Blocked follicles form blackheads or whiteheads, while inflammation can cause papules, pustules, nodules, scarring, and dark marks, especially on the face, chest, and back.

\synonyms
Acne Vulgaris; Pimples

\medterm{Acne Conglobata} A rare, severe nodulocystic form of acne with groups of inflammatory nodules, cysts, interconnected abscesses, and draining sinuses, usually on the trunk, buttocks, or limbs. It can be persistent and leave atrophic or hypertrophic scars, so early dermatologic treatment is important.

\medterm{Acne Fulminans} A sudden, severe ulcerating form of acne associated with fever, painful joints, malaise, and sometimes weight loss or abnormal blood counts. It requires urgent specialist care; isotretinoin can precipitate or worsen it if introduced too quickly.

\medterm{Acne Inversa} See \textit{Hidradenitis suppurativa}.

\medterm{Acne Keloid} A chronic inflammatory condition causing firm scar-like bumps around hair follicles, usually on the back of the neck.

\synonyms
Acne Keloidalis Nuchae

\medterm{Acne Keloidalis Nuchae (AKN)} Alternate name; see \textit{Acne Keloid}.

\medterm{Acne Rosacea} Acne rosacea is a chronic inflammatory facial skin disorder causing flushing, persistent redness, visible vessels, papules, pustules, or eye symptoms; it is distinct from acne vulgaris.

\medterm{Acne Vulgaris} The common acne disorder caused by follicular plugging, increased sebum production, Cutibacterium acnes activity, and inflammation. It produces comedones, papules, pustules, nodules, or cysts on the face, chest, back, and sometimes upper arms; deeper disease can scar.

\medterm{Acneiform Eruptions} Monomorphic follicular papules and pustules that resemble acne but usually lack comedones. They may be caused by medicines, occupational or chemical exposures, infections, or other disorders; timing, distribution, and trigger help distinguish them from acne vulgaris.

\medterm{Acoustic} Relating to sound, hearing, or how sound waves are produced, transmitted, measured, or received. Acoustic testing can assess hearing, rooms, medical devices, tissues, or signals used during diagnosis.

\medterm{Acoustic Nerve} The acoustic nerve, properly called the vestibulocochlear nerve, carries signals for hearing and balance from the inner ear to the brain. It is the eighth cranial nerve.

\medterm{Acoustic Neuroma} A benign, slow-growing neoplasm arising from Schwann cells of the vestibular division of the vestibulocochlear nerve (cranial nerve VIII) within the internal auditory canal or cerebellopontine angle. It classically presents with insidious, progressive unilateral sensorineural hearing loss, asymmetric high-pitched tinnitus, and dysequilibrium.

\synonyms
Vestibular Schwannoma; Acoustic Neurinoma; Acoustic Neurilemmoma

\medterm{Acoustic Reflex} An automatic tightening of a small middle-ear muscle in response to a loud sound. Testing it helps assess the middle ear, hearing pathways, and facial nerve.

\synonyms
Stapedial Reflex; Stapedius Reflex; Middle-Ear Muscle Reflex

\medterm{Acoustic Shock} Auditory and other symptoms after exposure to a sudden or intense sound, potentially including ear pain, tinnitus, hyperacusis, dizziness, or temporary hearing change.

\medterm{Acoustic Trauma} Acute damage to the auditory system from a brief, intense sound such as an explosion or
firearm. It may cause tinnitus, temporary or permanent sensorineural hearing loss, and
tympanic or ossicular injury; urgent audiologic assessment is indicated when hearing
loss is sudden.

\medterm{Acquired} Developed after birth rather than being present at birth or inherited; an acquired condition may result from infection, exposure, injury, aging, or another later event.

\medterm{Acquired C1-Inhibitor Deficiency Angioedema} A bradykinin-mediated disorder caused by acquired deficiency or dysfunction of C1 inhibitor, often associated with lymphoproliferative or autoimmune disease. Recurrent swelling of the skin, bowel wall, or upper airway usually occurs without hives or itching; low C1q and C1-inhibitor testing help distinguish it from hereditary or histamine-mediated angioedema.

\medterm{Acquired Cystic Kidney Disease} Development of multiple fluid-filled cysts in kidneys with long-standing chronic kidney disease, especially in people receiving dialysis. It can cause blood in the urine, pain, infection, or rarely kidney tumors and is monitored with clinical assessment and imaging.

\medterm{Acquired Feminization} Development of typically female secondary sexual characteristics after birth in a person not expected to develop them, often from hormonal change.

\medterm{Acquired Immunity} Immune protection that develops after infection, vaccination, or transfer of antibodies, involving adaptive immune responses directed toward specific antigens.

\medterm{Acquired Immunodeficiency Syndrome} See \textit{HIV/AIDS}.

\medterm{Acquired Lipodystrophy} A rare disorder characterized by selective, progressive loss of adipose tissue (lipoatrophy) occurring after birth without an inherited genetic defect. Frequently associated with autoimmune diseases, panniculitis, or antiretroviral therapy (e.g., Lawrence syndrome, Barraquer-Simons syndrome), it leads to ectopic lipid accumulation, severe insulin resistance, hypertriglyceridemia, and hepatic steatosis.

\synonyms
Acquired Generalized Lipodystrophy; AGL; Lawrence Syndrome; Acquired Partial Lipodystrophy

\medterm{Acquired Methemoglobinemia} An acquired hematologic disorder in which hemoglobin iron is oxidized from the ferrous ($\text{Fe}^{2+}$) to the ferric ($\text{Fe}^{3+}$) state by exogenous oxidizing agents (such as dapsone, benzocaine, lidocaine, nitrites, or sulfonamides), producing clinical cyanosis, chocolate-brown arterial blood, and cellular hypoxia refractory to supplemental oxygen.

\synonyms
Toxic Methemoglobinemia; Drug-Induced Methemoglobinemia; Chemically Induced Methemoglobinemia

\medterm{Acquired Nasolacrimal Duct Obstruction} Blockage of the tear-drainage system that develops after birth. It may cause excessive tearing, discharge, or repeated inflammation of the lacrimal sac and can result from aging, infection, injury, inflammation, scarring, or a mass.

\medterm{Acquired Platelet Function Defect} Impaired platelet adhesion or aggregation that develops after birth. Causes include antiplatelet medicines, kidney failure, hematologic disease, and other systemic disorders; platelet counts may remain normal, but mucocutaneous bleeding can occur.

\medterm{Acquired Tracheomalacia} Acquired tracheomalacia is weakening of the tracheal wall after injury, inflammation, compression, or prolonged ventilation, causing airway collapse and breathing symptoms.

\medterm{Acquisition} The process by which new information or a skill is learned and initially encoded, especially in discussions of memory and learning.

\synonyms
Learning Acquisition; Memory Acquisition

\medterm{Acquisition Duration} The time required to collect a measurement, image, or body signal. Longer acquisition may improve detail or accuracy but can increase motion problems, discomfort, radiation exposure, or equipment time.

\medterm{Acquisition Time} The time during which a medical imaging system collects data. The required time depends on the imaging method,
body region, protocol, and patient movement.

\medterm{Acral Lentiginous Melanoma} A subtype of melanoma arising on acral skin (palms, soles, or under nails), appearing as
a dark, irregularly pigmented patch or plaque. It is more common in people with darker
skin tones and often presents at a later stage due to less frequent surveillance of
these sites. Prognosis depends on thickness and stage at diagnosis.

\medterm{Acro-Osteolysis} Resorption or destruction of the terminal bones of the fingers or toes. It can occur with inherited disorders, autoimmune disease, severe vascular disease, trauma, or repeated occupational exposure and may cause shortened, painful, or deformed digits.

\medterm{Acrocentric Chromosome} A chromosome whose central joining point lies close to one end, giving it a very short section and a longer section. In humans, chromosomes 13, 14, 15, 21, and 22 have this shape.

\medterm{Acrochordon} Alternate name; see \textit{Skin Tag}.

\medterm{Acrocyanosis} Persistent bluish or purplish discoloration and coolness of the hands, feet, nose, or ears, usually without pain, caused by narrowing of small skin blood vessels. It is often worse in cold conditions and may be primary or due to another disorder.

\medterm{Acrodermatitis} Inflammation affecting skin on the hands, feet, or other body extremities. Causes vary and include infections, nutritional deficiency, inherited disorders, medicines, and immune reactions; appearance and treatment depend on cause.

\medterm{Acrodermatitis Chronica Atrophicans} A late cutaneous manifestation of Lyme borreliosis, usually caused by persistent \textit{Borrelia afzelii} infection in Europe. It begins with unilateral reddish-violet swelling on an extremity and can progress over months or years to thin, atrophic, fibrotic skin; diagnosis uses clinical findings, serology, and sometimes biopsy.

\medterm{Acrodermatitis Enteropathica} A zinc-deficiency disorder classically causing periorificial and acral dermatitis, diarrhea, and alopecia. The primary inherited form results from impaired intestinal zinc absorption; acquired deficiency may follow malnutrition, malabsorption, or other illness, and untreated disease can impair growth, healing, and immunity.

\medterm{Acrodynia} A painful condition with red, swollen, sweaty hands and feet, classically associated with mercury poisoning and historically called pink disease in children.

\medterm{Acrodysostosis} A rare inherited disorder of bone development causing short, broad fingers and toes, a small nose and middle face, and short stature. Some people have learning difficulties or reduced response to parathyroid and other hormones.

\medterm{Acrokeratosis Paraneoplastica} A rare skin condition causing thick, wart-like growths on the hands and feet, often around the nails, as a reaction to an internal cancer. It is most often linked with cancers of the digestive or respiratory tract and needs medical evaluation.

\medterm{Acromegaly} A disorder caused by excess growth hormone after height growth has ended, usually from a noncancerous pituitary adenoma. It enlarges the hands, feet, jaw, and soft tissues and may cause headaches, joint problems, diabetes, or vision changes.

\medterm{Acromioclavicular Joint} The small joint at the top of the shoulder where the collarbone meets the shoulder blade. It allows limited gliding as the arm moves and is held stable mainly by two groups of strong ligaments.

\medterm{Acromioclavicular Joint Injuries} Sprains or dislocations of the joint between the acromion and clavicle, usually caused by a fall or direct
shoulder impact. Severity ranges from ligament sprain to complete displacement and is assessed clinically and with imaging.

\synonyms
Acromioclavicular Joint Injury

\medterm{Acromion} A bony projection from the shoulder blade that forms the highest point of the shoulder. It joins the collarbone and forms a roof over the shoulder joint, with a tendon and lubricating sac passing beneath it.

\medterm{Acroparaesthesia} Acroparaesthesia is tingling, pins-and-needles, burning, or altered sensation in the hands or feet, commonly related to nerve compression, neuropathy, vascular disease, or metabolic causes.

\medterm{Acrophobia} An excessive, persistent fear of heights that causes marked anxiety, avoidance, or impairment beyond the ordinary caution associated with high places.

\medterm{Acropustulosis of Infancy} A recurrent itchy eruption of tiny vesicles and pustules, mainly on the palms and soles of infants and young children. Crops become less frequent and usually resolve by early childhood; scabies may precede or mimic it.

\medterm{Acrosin} An enzyme released from the cap-like structure at a sperm’s tip. It helps sperm pass through the egg’s surrounding layers during fertilization, although successful fertilization requires many additional steps.

\medterm{Acrosome} A cap-like structure on the head of a sperm cell containing enzymes that help it attach to and enter the egg during fertilization.

\medterm{Acrosome Reaction} The release of enzymes from the cap at the front of a sperm cell after it contacts an egg. These enzymes help the sperm pass through the egg’s outer coating; other steps are also needed for fertilization.

\medterm{ACTH Blood Test} A blood test that measures the level of ACTH, a hormone made by
the pituitary gland. It is usually checked with a cortisol test to help find problems
of the adrenal or pituitary glands.

\medterm{ACTH Stimulation Test} A test that shows whether the adrenal glands can make
cortisol. A synthetic form of ACTH is given, and cortisol levels are measured before
and after it is given. A small rise may suggest adrenal insufficiency.

\medterm{Actigraphy} Recording movement over days or weeks with a small wearable sensor, usually worn on the wrist. The pattern of movement helps estimate when a person is asleep or awake and how active they are, but it does not measure sleep directly and cannot by itself diagnose every sleep problem.

\synonyms
Actigraphic Monitoring

\medterm{Actin} A protein that helps cells keep their shape and move. In muscle, actin works with myosin to produce contraction; in other cells, it supports movement, internal transport, and changes in cell structure.

\synonyms
Actin Filament; Microfilament Protein

\medterm{Actin Myofilament} A thin muscle fiber made mainly of actin that interacts with myosin, allowing muscle cells to shorten and produce force during contraction and movement.

\synonyms
Thin Filament

\medterm{Actin-Binding Protein} A protein that attaches to actin filaments and helps build, organize, move, or anchor the cell's internal framework. Examples include myosin, gelsolin, and alpha-actinin.

\medterm{Acting Out} Expressing emotional conflict or distress through behavior rather than words.

\medterm{Actinic} Relating to ultraviolet radiation from sunlight or UV lamps; chronic actinic exposure can cause sunburn, photoaging, and actinic keratoses.

\medterm{Actinic Cheilitis} A precancerous sun-damage condition of the lips, usually the lower lip, caused by long-term ultraviolet exposure. It may cause persistent dryness, scaling, whitening, thickening, or blurring of the lip border and can progress to squamous cell carcinoma.

\medterm{Actinic Granuloma} A usually harmless inflammatory skin condition that occurs on sun-damaged skin. It causes ring-shaped patches with a thinner center and a raised red border, most often on the neck, arms, or backs of the hands.

\medterm{Actinic Keratosis} A rough, scaly, sandpaper-like lesion caused by cumulative ultraviolet damage, usually on sun-exposed skin such as the face, scalp, ears, neck, forearms, or hands. It is a keratinocytic precancer and may progress to cutaneous squamous-cell carcinoma, so changing, thick, tender, or bleeding lesions need assessment and treatment.

\synonyms
Solar Keratosis; AK; Senile Keratosis

\medterm{Actinic Keratosis Pathology} Atypical keratinocytes with disordered maturation in the lower epidermis, often with parakeratosis, alternating orthokeratosis, and solar elastosis. It is a premalignant squamous lesion that can progress to cutaneous squamous cell carcinoma.

\medterm{Actinic Purpura} Bruising caused by fragile, sun-damaged skin and connective tissue, usually on the backs of the hands and
forearms in older adults. The purple patches may appear after minor trauma and are generally harmless, although
medicines such as anticoagulants can worsen them. New, widespread, or unexplained bruising warrants assessment.

\medterm{Actinomycetoma} A slowly worsening bacterial infection of the skin and deeper tissue, usually affecting the foot after soil contamination. It causes swelling, firm lumps, draining channels, and grains in the discharge; untreated disease can damage bone.

\medterm{Actinomycin D} Actinomycin D, or dactinomycin, is an antineoplastic antibiotic that blocks DNA-dependent RNA synthesis and is used for selected cancers; myelosuppression and tissue injury are important risks.

\medterm{Actinomycosis} A long-lasting infection caused by a type of bacteria called Actinomyces that normally lives in the mouth, intestines, and female genital area. It starts when these bacteria get past a damaged surface, such as after a dental problem, dental work, abdominal surgery, or use of an intrauterine device. It most often affects the face and neck, chest, or belly and pelvis, where it forms hard, growing lumps and draining tunnels that contain small yellow grains called sulfur granules. It can spread to nearby tissue and look like cancer. Diagnosis uses cultures or tissue samples, and treatment is a long course of antibiotics, usually high-dose penicillin.

\medterm{Action Potential} A brief electrical signal that travels along a nerve or muscle cell. It begins when the cell’s electrical charge changes enough to trigger the signal and allows nerves to communicate and muscles to contract.

\medterm{Action Tremor} An involuntary rhythmic oscillation of a body part occurring during active voluntary muscle contraction. Subtypes include postural tremors (maintaining a position against gravity), kinetic tremors (moving through an arc), and intention tremors (worsening near the target).

\medterm{Activated Charcoal} A porous form of carbon that adsorbs some substances in the gastrointestinal tract and can reduce
their absorption after selected poisonings. It is not effective for every toxin and can
cause aspiration or other complications.

\medterm{Activated Partial Thromboplastin Time} A blood test that measures clotting through the intrinsic and common coagulation pathways.
It may be prolonged by unfractionated heparin, clotting-factor deficiencies, inhibitors,
or some liver and systemic disorders.

\medterm{Activated Protein C Resistance} An inherited thrombophilic disorder characterized by an impaired anticoagulant response of plasma to activated protein C (APC). Approximately 90 to 95 percent of cases are caused by the factor V Leiden mutation (G1691A), which alters the APC cleavage site on factor V, predisposing affected individuals to recurrent deep venous thrombosis and pulmonary embolism.

\synonyms
APC Resistance; Factor V Leiden Thrombophilia

\medterm{Active Caries} Tooth decay that is continuing to progress. The affected surface is often soft, rough, or chalky and requires dental
assessment and management.

\medterm{Active Management Of Third Stage} A set of measures used after birth to reduce postpartum hemorrhage, usually including a uterotonic
medicine, controlled cord traction when appropriate, and assessment of uterine tone and bleeding.

\medterm{Active Phase Of Labor} The part of the first stage of labor in which the cervix dilates progressively, generally beginning at about 6 cm
and continuing to full dilation. Progress varies between individuals.

\medterm{Active Site} The part of an enzyme where a target molecule attaches and a chemical reaction takes place. Its shape and chemical properties help the enzyme recognize particular molecules and change them into products.

\medterm{Active Surveillance} Close, structured monitoring of a cancer or precancerous condition without immediate
treatment. It includes scheduled examinations, laboratory tests, imaging, or biopsies,
with treatment started if evidence of progression develops.

\medterm{Active Transport} The movement of a substance across a cell membrane using cellular energy. It often moves the substance against a concentration or electrical gradient and helps cells control salts, nutrients, waste products, and other substances.

\medterm{Activities Of Daily Living} Routine tasks required for personal care and independent living, including basic activities and
more complex instrumental activities. Performance is assessed to describe functional
status and support care planning.

\medterm{Acupressure} A complementary treatment in which firm pressure is applied to selected points on the body, often with the fingers. It may provide short-term relief of symptoms such as nausea or pain, but evidence of benefit varies by condition.

\medterm{Acupuncture} A treatment originating in traditional Chinese medicine in which thin needles are inserted into the skin and sometimes moved or electrically stimulated. It may help some pain conditions or treatment-related nausea, but benefits vary.

\medterm{Acute} Describing symptoms or a condition that begins or worsens quickly and usually lasts a limited time. It is commonly contrasted with chronic, which develops or continues over a longer period.

\medterm{Acute Abdomen} Sudden or severe abdominal pain with clinical features suggesting a potentially serious intra-abdominal disorder that may require urgent investigation or surgery.

\medterm{Acute Acalculous Cholecystitis} Acute gallbladder inflammation without gallstones, usually occurring in critically ill, injured, postoperative, or fasting patients. Ischaemia, bile stasis, and infection may contribute; it can progress rapidly to necrosis or perforation and requires urgent assessment.

\medterm{Acute Adrenal Crisis} Acute adrenal crisis is life-threatening cortisol deficiency causing hypotension, vomiting, abdominal pain, electrolyte disturbance, hypoglycemia, or shock; urgent glucocorticoid and fluid treatment is required.

\medterm{Acute Agitation} A state of heightened motor activity, emotional tension, irritability, and impaired
behavioral control that may threaten the patient or others. Management begins with
de-escalation and assessment for delirium, intoxication, withdrawal, medical illness,
and psychosis.

\medterm{Acute Alcohol Intoxication} Neurobehavioral and physiologic impairment caused by recent ethanol exposure. Severity
varies with concentration, tolerance, co-ingestants, and comorbidity; dangerous effects
include respiratory depression, aspiration, hypoglycemia, hypothermia, trauma, and coma.

\medterm{Acute Angle-Closure Glaucoma} An ophthalmic emergency in which sudden closure of the anterior-chamber angle causes a rapid
rise in intraocular pressure. Severe eye pain, blurred vision, halos, headache, nausea,
vomiting, and a red eye may occur; urgent ophthalmic treatment is required to protect
vision.

\medterm{Acute Aortic Syndrome} A group of sudden, life-threatening problems in the wall of the aorta, the body’s largest artery. The wall may tear, bleed within itself, or develop an ulcer, causing severe chest or back pain and requiring emergency assessment.

\medterm{Acute Bacterial Prostatitis} A sudden bacterial infection of the prostate causing fever, chills, pain between the scrotum and anus or in the lower back, painful urination, and sometimes inability to pass urine. It requires prompt antibiotics and medical care; a forceful prostate examination is avoided because it can spread bacteria into the bloodstream.

\medterm{Acute Bacterial Rhinosinusitis} Bacterial infection of the paranasal sinuses suspected when symptoms persist without
improvement for at least 10 days, become substantially worse after initial improvement,
or are severe with high fever and purulent discharge. Treatment decisions depend on
severity, complications, resistance risk, and local guidance.

\medterm{Acute Bilirubin Encephalopathy} Brain dysfunction in a newborn caused by very high bilirubin levels. Poor feeding, sleepiness, weak muscle tone, irritability, a high-pitched cry, unusual movements, seizures, or coma may develop. It is an emergency requiring rapid treatment to prevent permanent brain injury.

\synonyms
Bilirubin Neurotoxicity; Kernicterus Spectrum Disorder

\medterm{Acute Blood Loss} Rapid loss of blood that reduces the amount circulating through the body. It can cause weakness, fast heartbeat, low blood pressure, confusion, organ injury, or shock and requires urgent control of bleeding and replacement of lost volume.

\medterm{Acute Bronchitis} Short-term inflammation of the large airways in the lungs, usually caused by a virus such as a cold or flu. It causes a cough that may bring up clear, yellow, or green mucus and may cause chest soreness, mild fever, tiredness, or wheezing. The cough often lasts 1 to 3 weeks. Because viruses cause more than 90\% of cases, antibiotics usually do not help; care focuses on rest, fluids, and symptom relief.

\synonyms
Chest Cold; Tracheobronchitis

\medterm{Acute Bursitis} Sudden, acute inflammation of a synovial bursa, typically triggered by direct trauma, acute crystal deposition (gout/pseudogout), bacterial infection (most commonly \textit{Staphylococcus aureus}), or unaccustomed repetitive strain. It presents with rapid-onset localized swelling, erythema, warmth, exquisite focal tenderness, and restricted active joint motion.

\synonyms
Acute Bursal Synovitis; Traumatic Bursitis; Septic Bursitis

\medterm{Acute Caries} Tooth decay that progresses rapidly, often producing a soft lesion or cavity. Risk is increased by factors such as frequent
sugar exposure, reduced saliva, or poor oral hygiene.

\medterm{Acute Cerebellar Ataxia} Sudden or rapidly developing incoordination caused by cerebellar dysfunction. In children it may follow an infection, but stroke, intoxication, inflammation, tumor, and other urgent causes must be considered.

\medterm{Acute Chest Syndrome} A potentially life-threatening pulmonary complication of sickle cell disease characterized by a new lung infiltrate plus fever or respiratory symptoms such as chest pain, cough, hypoxemia, or breathing difficulty.

\medterm{Acute Cholangitis} Alternate name; see \textit{Ascending Cholangitis}.

\medterm{Acute Cholecystitis} Acute inflammation of the gallbladder, most often caused by persistent cystic-duct
obstruction by a gallstone. It commonly causes sustained right-upper-quadrant pain,
fever, nausea, and a positive Murphy sign; ultrasound and timely antimicrobial and
surgical assessment are standard.

\medterm{Acute Colonic Pseudo-Obstruction} Severe widening of the large intestine without a physical blockage, usually in people who are seriously ill, injured, or recovering from surgery. It can cause pain, swelling, nausea, and vomiting. If it is not treated, the bowel may lose its blood supply or tear.

\synonyms
Ogilvie Syndrome; Colonic Pseudo-Obstruction

\medterm{Acute Compartment Syndrome} A rapid rise in pressure inside a closed muscle compartment that reduces blood flow and can permanently damage muscles and nerves. It usually follows a fracture, crush injury, or severe swelling and requires emergency assessment.

\synonyms
Compartment Syndrome; Intracompartmental Hypertension

\medterm{Acute Coronary Syndrome} A clinical spectrum of acute myocardial ischemia caused most often by coronary plaque disruption and thrombosis. It includes unstable angina, non-ST-elevation myocardial infarction (NSTEMI), and ST-elevation myocardial infarction (STEMI); ECG findings and serial cardiac troponin distinguish the major presentations and guide urgent treatment.

\medterm{Acute Coronary Syndromes} Another name for \textit{Acute Coronary Syndrome}.

\medterm{Acute Cough} A cough lasting less than three weeks. Common causes include viral respiratory infection, acute bronchitis, and pneumonia.

\medterm{Acute Cutaneous Lupus Erythematosus (ACLE)} A lupus-specific inflammatory skin eruption that commonly presents as a transient malar or generalized erythematous rash, often triggered by sunlight and associated with systemic lupus activity. The malar rash usually spares the nasolabial folds; evaluation should consider systemic symptoms and organ involvement.

\medterm{Acute Cystitis} Acute bacterial infection and inflammation confined to the lower urinary tract and bladder mucosa, typically caused by uropathogenic \textit{Escherichia coli} and presenting with acute dysuria, urinary frequency, urgency, suprapubic pain, and hematuria.

\synonyms
Acute Bacterial Cystitis; Bladder Infection; Acute UTI

\medterm{Acute Dacryocystitis} Acute infection and inflammation of the lacrimal sac, most commonly caused by
obstruction of the nasolacrimal duct leading to stasis and secondary bacterial
infection. It presents with painful swelling, redness, and tenderness at the medial
canthus with purulent discharge.

\medterm{Acute Decompensated Heart Failure} Sudden or rapidly worsening heart failure, causing breathlessness, tiredness, leg swelling, or fluid in the lungs. It may be triggered by infection, a heart attack, rhythm problems, uncontrolled blood pressure, or missed treatment and can require urgent hospital care.

\medterm{Acute Diarrhea} Passage of loose or watery stools, usually three or more times daily, lasting less than two weeks. Infection, food-related
illness, medicines, and other intestinal disorders can cause it. Oral rehydration is important; blood in stool,
severe dehydration, persistent vomiting, or severe pain requires prompt assessment.

\medterm{Acute Disease} An illness that begins suddenly or develops over a short period. It often lasts a limited time, although it can be mild, severe, or life-threatening and may sometimes lead to a longer-term condition.

\synonyms
Acute Illness

\medterm{Acute Disseminated Encephalomyelitis} A rare sudden inflammation of the brain and spinal cord, usually after an infection and less often after vaccination. It can cause confusion, weakness, poor coordination, vision problems, or seizures, especially in children.

\medterm{Acute Ear Infection} A sudden, rapid-onset bacterial or viral infection of the middle ear space behind the eardrum, characterized by earache, fever, hearing muffled sounds, and bulging of the tympanic membrane.

\synonyms
Acute Otitis Media; AOM; Middle Ear Infection

\medterm{Acute Endometritis} Acute inflammation of the endometrium, usually caused by ascending infection or retained tissue after pregnancy-related events or uterine instrumentation; fever, pelvic pain, and abnormal bleeding may occur.

\medterm{Acute Epididymitis} Sudden inflammation of the epididymis, usually caused by infection and sometimes involving the testis. It causes unilateral scrotal pain and swelling and must be distinguished urgently from testicular torsion.

\medterm{Acute Esophageal Necrosis} Ischemic necrosis of the esophageal mucosa, most commonly caused by low-flow states such
as shock, sepsis, or thromboembolism, presenting with hematemesis and chest pain.
Endoscopy shows circumferential black necrotic mucosa, typically in the distal
esophagus.

\medterm{Acute Exacerbation} A rapid worsening of a chronic pulmonary condition, with increased breathlessness, cough, oxygen need, or other symptoms; infection, aspiration, embolism, and other triggers may contribute.

\medterm{Acute Exacerbation Of Asthma} A progressive increase in dyspnea, cough, wheezing, or chest tightness with reduced
expiratory airflow compared with the patient’s baseline. Severe attacks may cause
exhaustion, silent chest, hypercapnia, respiratory failure, and need for urgent
treatment.

\medterm{Acute Exacerbation Of COPD} A sudden worsening of chronic obstructive pulmonary disease, with more breathlessness, cough, or sputum than usual. It is often triggered by a respiratory infection, air pollution, or another stress on the lungs.

\medterm{Acute Fatty Liver Of Pregnancy} A rare, potentially life-threatening liver disorder that usually develops late in pregnancy. Fat builds up in liver cells and may cause nausea, abdominal pain, jaundice, low blood sugar, abnormal bleeding, kidney failure, or other organ failure.

\medterm{Acute Febrile Illness} A sudden illness characterized by fever, often with headache, chills, muscle aches, gastrointestinal symptoms, or other signs of infection; the cause may be infectious or noninfectious.

\medterm{Acute Febrile Neutrophilic Dermatosis (Sweet Syndrome)} An inflammatory neutrophilic skin disease causing abrupt tender red or violaceous papules, plaques, or nodules with fever and neutrophilia. It may be triggered by infection, inflammatory disease, medicines, pregnancy, or an underlying malignancy; biopsy confirms the characteristic neutrophilic infiltrate.

\medterm{Acute Fever in Adults} A fever in an adult with recent onset, commonly caused by infection, inflammation, medication, heat illness, or other conditions. Severity, vital signs, immune status, source, and associated symptoms guide evaluation; confusion, breathing difficulty, stiff neck, shock, or severe localized pain requires urgent care.

\medterm{Acute Flaccid Myelitis} A rare serious illness, usually in children, causing sudden weakness and limpness of one or more arms or legs. Facial weakness, swallowing difficulty, or breathing failure may occur; immediate hospital assessment, imaging, and breathing monitoring are needed.

\medterm{Acute Gastritis} Sudden inflammation or irritation of the stomach lining. It may cause upper-abdominal pain, nausea, vomiting, or bleeding and can result from anti-inflammatory medicines, alcohol, severe illness, infection, or other irritants.

\synonyms
Acute Erosive Gastritis

\medterm{Acute Generalized Exanthematous Pustulosis} A sudden skin reaction, usually to a medicine, that causes many small pus-filled bumps on widespread red skin. Fever may occur. It often starts within days of exposure and improves after the triggering medicine is stopped.

\medterm{Acute Glomerulonephritis} Acute immune or infectious inflammation of the glomeruli producing a nephritic syndrome: hematuria with dysmorphic red cells or red-cell casts, variable proteinuria, edema, hypertension, and reduced kidney function. Causes include postinfectious disease, IgA nephropathy, vasculitis, and lupus; complement, serology, and sometimes kidney biopsy help identify the cause.

\medterm{Acute Granulocytic Leukemia} See \textit{Acute myelogenous leukemia}.

\medterm{Acute Hemorrhagic Edema of Infancy} A usually self-limited cutaneous small-vessel vasculitis of infants, characterized by fever, edema, and large bruise-like or targetoid purpuric lesions on the face and extremities. Internal-organ involvement is uncommon, but serious causes of purpura must be excluded; most cases resolve within one to three weeks.

\medterm{Acute Hemorrhagic Leukoencephalitis} A rare, fulminant inflammatory demyelinating disease of the brain, considered a severe form of acute disseminated encephalomyelitis. It can cause rapidly worsening headache, fever, confusion, seizures, weakness, coma, and brain swelling and requires emergency neurologic care.

\medterm{Acute Hepatitis} Sudden inflammation and injury of the liver, developing over days to months. It may cause tiredness, nausea, pain under the right ribs, dark urine, jaundice, or no symptoms. Causes include viruses, alcohol, medicines, toxins, and autoimmune disease.

\medterm{Acute HIV Syndrome} A flu-like illness that may occur about 2–6 weeks after HIV infection, when the virus is multiplying rapidly. Fever, sore throat, swollen glands, rash, muscle aches, headache, nausea, diarrhea, or mouth ulcers may occur, but some people have no symptoms.

\medterm{Acute Hypercapnic Respiratory Failure} A sudden inability to breathe out enough carbon dioxide, causing it to build up in the blood. Severe lung disease, asthma, weak breathing muscles, sedative medicines, or chest problems can cause confusion, sleepiness, or dangerous breathing failure.

\medterm{Acute Hypoxemic Respiratory Failure} A sudden inability of the lungs to put enough oxygen into the blood. Pneumonia, fluid in the lungs, severe inflammatory lung injury, collapsed lung tissue, or a blood clot in the lungs may cause severe breathlessness, bluish skin, confusion, or organ injury.

\medterm{Acute Idiopathic Polyneuritis} Guillain--Barré syndrome, an acute immune-mediated disorder of peripheral nerves that usually causes progressive symmetric weakness and reduced reflexes and may impair breathing.

\medterm{Acute Illness} A disease or health problem that begins suddenly or over a short period and generally lasts a limited time, although its severity may range from mild to life-threatening.

\medterm{Acute Inflammation} A rapid immune response to injury or infection that usually lasts hours to days. It can cause redness, warmth, swelling, pain, and reduced function as extra blood, fluid, and immune cells enter the affected area.

\medterm{Acute Inflammatory Demyelinating Polyneuropathy} See \textit{Guillain-Barre syndrome}.

\medterm{Acute Inflammatory Demyelinating Polyradiculoneuropathy} The most common form of Guillain–Barré syndrome. The immune system damages the protective covering of nerves, causing progressive weakness, tingling, and reduced reflexes that may spread to the breathing muscles.

\medterm{Acute Intermittent Porphyria} An inherited disorder of haem production caused by reduced hydroxymethylbilane synthase activity. Attacks may cause severe abdominal pain, nausea, constipation, dark urine, nerve symptoms, confusion, or seizures; medicines, hormonal changes, fasting, and illness can trigger them.

\medterm{Acute Interstitial Nephritis} Sudden inflammation of the tissue between the kidney’s filtering structures, most often caused by a reaction to a medicine. Infections and autoimmune disease can also cause it, leading to blood or protein in urine and a sudden decline in kidney function.

\medterm{Acute Ischemic Stroke} Sudden loss of brain, spinal-cord, or retinal function because a blood vessel is blocked and the tissue does not receive enough oxygen. Emergency imaging distinguishes it from bleeding and determines urgent treatment.

\medterm{Acute Kidney Failure} A sudden decline in kidney function that prevents adequate removal of waste, salt, and water. It may cause reduced urine, swelling, nausea, confusion, or dangerous electrolyte and acid changes.

\medterm{Acute Kidney Injury} An abrupt decline in kidney function, defined by a serum-creatinine rise of at least 0.3 mg/dL within 48 hours, a rise to at least 1.5 times baseline within 7 days, or urine output below 0.5 mL/kg/hour for 6 hours. Causes are commonly prerenal hypoperfusion, intrinsic kidney injury, or postrenal obstruction; complications include azotemia, fluid and electrolyte disorders, and acidosis.

\synonyms
AKI

\medterm{Acute Kidney Injury Staging} A standardized classification based on changes in serum creatinine and urine output.
KDIGO stages 1 through 3 quantify the severity of abrupt kidney-function decline and
support prognosis, monitoring, and treatment decisions.

\medterm{Acute Leukemia} A fast-growing blood cancer in which immature abnormal blood cells crowd the bone marrow and interfere with normal blood-cell production. It can cause anemia, infections, and bleeding and needs urgent treatment.

\medterm{Acute Limb Ischemia} Sudden reduction in arterial blood flow to a limb that threatens tissue viability,
usually caused by embolism, thrombosis, or graft occlusion. The six classic findings are
pain, pallor, pulselessness, paresthesia, paralysis, and poikilothermia; urgent vascular
assessment and revascularization are required.

\medterm{Acute Limb Ischemia Six Ps} Six warning findings of suddenly reduced artery blood flow to a limb: pain, paleness, absent pulse, tingling or numbness, weakness or paralysis, and a cold limb. This is an emergency requiring urgent vascular assessment.

\medterm{Acute Liver Failure} Rapid deterioration of liver function in a previously healthy individual without prior
liver disease. Causes include acetaminophen toxicity, viral hepatitis, drug reactions,
autoimmune hepatitis, and Wilson disease. Features include coagulopathy, encephalopathy,
and multiorgan failure. Presents with jaundice, confusion, and bleeding.

\synonyms
Fulminant Hepatic Failure; Acute Hepatic Necrosis

\medterm{Acute Low Back Pain} Lumbosacral pain lasting less than 4 to 6 weeks, most commonly caused by mechanical musculoligamentous strain, facet joint irritation, or acute lumbar disc herniation. Management emphasizes conservative therapy, staying active, and monitoring for neurological red flags.

\synonyms
Acute Lumbago; Mechanical Low Back Strain; Lumbosacral Strain

\medterm{Acute Lymphangitis} Acute inflammation of lymphatic vessels, typically resulting from streptococcal or
staphylococcal infection spreading from a distal wound. It presents with red, tender
streaks extending proximally from the site of infection, along with fever, malaise, and
regional lymphadenopathy. Treatment requires systemic antibiotics.

\medterm{Acute Lymphoblastic Leukemia} Acute lymphoblastic leukemia is a rapidly progressive cancer of immature lymphoid precursor cells that crowds the bone marrow and can involve blood, lymph nodes, or the central nervous system.

\medterm{Acute Lymphocytic Leukemia} A malignant proliferation of immature lymphoid precursors in the bone marrow and blood,
commonly presenting with cytopenias, fever, bone pain, or lymphadenopathy.

\synonyms
ALL; Acute Lymphoblastic Leukemia

\medterm{Acute Lymphoid Leukemia} See \textit{Acute lymphocytic leukemia}.

\medterm{Acute Megakaryoblastic Leukemia} A rare subtype of acute myeloid leukemia characterized by the proliferation of
megakaryoblasts in the bone marrow and peripheral blood.

\medterm{Acute Mesenteric Ischemia} A vascular surgical emergency caused by an abrupt reduction in intestinal blood flow (most commonly from superior mesenteric artery thromboembolism), classically presenting with severe abdominal pain out of proportion to physical examination findings, rapidly progressing to bowel necrosis and peritonitis.

\synonyms
Acute Mesenteric Infarction; AMI

\medterm{Acute Mitral Regurgitation} Sudden backward flow of blood from the left ventricle into the left atrium during systole, caused by acute leaflet, chordal, papillary-muscle, or ventricular injury. It can rapidly produce pulmonary edema, low cardiac output, or shock.

\medterm{Acute Mountain Sickness} A high-altitude illness caused by inadequate acclimatization, usually after ascent above about 2,500\,m (8,200\,ft). It is characterized by headache plus symptoms such as nausea, loss of appetite, dizziness, fatigue, or disturbed sleep. Worsening symptoms, confusion, inability to walk normally, or breathlessness at rest indicate progression to life-threatening high-altitude cerebral or pulmonary edema.

\medterm{Acute Myeloblastic Leukemia} See \textit{Acute Myeloid Leukemia}.

\medterm{Acute Myelogenous Leukemia} Alternate name; see \textit{Acute Myeloid Leukemia}.

\medterm{Acute Myeloid Leukemia} A clonal hematopoietic neoplasm characterized by the rapid proliferation and accumulation of abnormal, poorly differentiated myeloid progenitor cells (blasts) in the bone marrow and peripheral blood. This suppresses normal hematopoiesis, leading to progressive cytopenias manifested as fatigue from anemia, recurrent or opportunistic infections from neutropenia, and petechiae, ecchymoses, or mucosal bleeding from thrombocytopenia.

\synonyms
AML; Acute Myelogenous Leukemia; Acute Granulocytic Leukemia; Acute Non-Lymphocytic Leukemia

\medterm{Acute Myocardial Infarction} Sudden death of heart muscle caused by prolonged inadequate coronary blood flow, usually from an acutely obstructed coronary artery; it is a medical emergency.

\synonyms
AMI

\medterm{Acute Myocardial Infarction Pathology} Histopathologic evolution of myocardial infarction, from early ischemic injury and coagulative necrosis through inflammation, granulation tissue, and scar formation. The findings depend on the interval between infarction and examination.

\medterm{Acute Necrotizing Ulcerative Gingivitis} A painful, rapidly developing infection of the gums characterized by punched-out gingival papillae, ulceration, spontaneous bleeding, gray necrotic tissue, and foul breath. It is associated with anaerobic oral bacteria and risk factors such as poor oral hygiene, smoking, stress, malnutrition, or immunosuppression; urgent dental evaluation, cleaning, and treatment are needed.

\medterm{Acute Nephritic Syndrome} A clinical syndrome of glomerular inflammation marked by hematuria, variable proteinuria, edema, hypertension, and reduced kidney function.

\medterm{Acute Nonlymphocytic Leukemia} See \textit{Acute myelogenous leukemia}.

\medterm{Acute Obstructive Uropathy} Sudden impairment of urinary flow caused by obstruction of the urinary tract,
potentially producing hydronephrosis, pain, infection, and acute kidney injury.

\medterm{Acute Onset Pain} Pain that begins suddenly, often over seconds, minutes, or hours.

\synonyms
Sudden-Onset Pain

\medterm{Acute Otitis Media} Infection of the middle ear space with acute onset of symptoms and signs of middle ear
inflammation. Diagnosis requires middle ear effusion plus acute symptoms (ear pain,
fever, irritability) or signs (bulging tympanic membrane; erythema alone is not diagnostic). Most common
pathogen: Streptococcus pneumoniae.

\medterm{Acute Pain} Pain that begins suddenly and usually lasts for a limited period, often because of injury, illness, or a procedure.

\medterm{Acute Pancreatitis} Sudden inflammation of the pancreas, most often caused by gallstones or alcohol. It usually causes severe upper-abdominal pain that may reach the back, with nausea or vomiting; complications can include fluid loss, infection, organ failure, or pancreatic collections.

\medterm{Acute Pericarditis} Rapid-onset inflammation of the sac around the heart, often idiopathic or viral and sometimes related to autoimmune disease, heart injury, kidney disease, cancer, or medicines. It typically causes sharp chest pain that worsens with breathing or lying down and improves when sitting up and leaning forward; a friction rub, diffuse ST-segment elevation with PR-segment depression, or a pericardial effusion may occur.

\medterm{Acute Poisoning} Harmful effects that develop after a short-term exposure to a toxic substance by swallowing, inhaling, injection, skin contact, or another route. Effects depend on the substance, amount, route, and time of exposure.

\medterm{Acute Proctitis} Sudden inflammation of the rectal mucosa causing rectal pain, urgency, tenesmus, bleeding, mucus, or discharge. Infectious, inflammatory, ischemic, traumatic, and radiation-related causes are possible.

\medterm{Acute Promyelocytic Leukemia} A subtype of acute myeloid leukemia that can quickly cause severe bleeding or abnormal clotting. It is linked to a specific chromosome change and needs immediate specialist treatment.

\medterm{Acute Pulmonary Edema} Acute or rapidly worsening accumulation of fluid in the lungs; see \textit{Pulmonary Edema}.

\medterm{Acute Pulmonary Histoplasmosis} An acute lower respiratory tract infection caused by inhaling microconidia of the dimorphic fungus \textit{Histoplasma capsulatum}, commonly found in soil enriched with bird or bat guano. Manifestations range from mild flu-like symptoms to severe pneumonia with mediastinal lymphadenopathy.

\synonyms
Primary Pulmonary Histoplasmosis; Acute Histoplasmosis Pneumonia; Cave Disease

\medterm{Acute Pyelonephritis} A sudden bacterial infection of the kidney, usually spreading upward from the bladder. It can cause fever, chills, pain in the side or back, nausea, and painful or frequent urination. Prompt medical treatment is important because the infection can spread to the bloodstream and damage the kidney.

\medterm{Acute Radiation Sickness} Illness caused by receiving a high dose of penetrating radiation over a short time, potentially damaging bone marrow, intestines, skin, and other organs.

\medterm{Acute Radiation Syndrome} See \textit{Radiation sickness}.

\medterm{Acute Rehabilitation} Multidisciplinary rehabilitation provided during or soon after hospitalization for conditions such as stroke,
spinal-cord injury, traumatic brain injury, or major trauma. Treatment is tailored to functional needs.

\medterm{Acute Renal Artery Occlusion} Sudden obstruction of blood flow through one or both main renal arteries or their primary branches, most commonly caused by thromboembolism, in situ thrombosis, or trauma. It presents with severe flank or abdominal pain, hematuria, refractory hypertension, and rapid deterioration of kidney function.

\synonyms
Acute Renal Infarction; Renal Artery Thrombosis; Renal Artery Embolism

\medterm{Acute Renal Failure} See \textit{Acute kidney injury}.

\medterm{Acute Respiratory Disease} A general term for an illness affecting breathing that begins suddenly or over a short time. It may include infections, asthma attacks, or other lung problems and is not one specific diagnosis.

\synonyms
ARD

\medterm{Acute Respiratory Distress Syndrome} An acute, diffuse inflammatory lung injury causing hypoxemic respiratory failure and bilateral lung opacities, usually developing within one week of an illness or injury and not primarily explained by heart failure or fluid overload. Common triggers include sepsis, pneumonia, aspiration, trauma, pancreatitis, and major transfusion; treatment addresses the cause and supports breathing with oxygen, lung-protective ventilation, and, in severe cases, prone positioning.

\medterm{Acute Respiratory Failure} A sudden inability of the respiratory system to maintain adequate oxygenation, carbon-dioxide removal, or both. It may result from lung disease, airway obstruction, neuromuscular weakness, or impaired respiratory drive and requires urgent evaluation.

\medterm{Acute Retinal Necrosis} A rapidly progressive necrotizing viral retinitis, usually caused by herpes viruses, with vitreous inflammation, peripheral retinal necrosis, occlusive vasculitis, and risk of retinal detachment.

\medterm{Acute Retroviral Syndrome} The acute illness that may occur shortly after HIV infection, commonly involving fever, rash, sore throat, lymph-node enlargement, muscle aches, or headache.

\medterm{Acute Rhinitis} Sudden inflammation of the nasal lining, commonly caused by a viral infection or an irritant, producing congestion, nasal discharge, and sneezing.

\medterm{Acute Rhinosinusitis} See \textit{Acute sinusitis}.

\medterm{Acute Sinusitis} Acute inflammation of the mucosa of one or more paranasal sinuses, usually following a
viral upper-respiratory infection. It causes nasal obstruction, purulent or discoloured
discharge, facial pressure or pain, and reduced smell; persistent, severe, or worsening
symptoms may indicate acute bacterial rhinosinusitis or a complication.

\synonyms
Acute Rhinosinusitis; Paranasal Sinus Infection

\medterm{Acute Stress Disorder} A trauma-related mental-health condition beginning after exposure to actual or threatened death, serious injury, or sexual violence and lasting from three days to one month. Symptoms may include intrusive memories, avoidance, negative mood, dissociation, and increased arousal.

\medterm{Acute Symptomatic Hyponatremia} A recently developing low blood sodium concentration accompanied by symptoms, particularly neurologic symptoms such as headache, confusion, seizures, or reduced consciousness. Clinical significance depends on the sodium level, rate of change, underlying cause, and associated illness.

\medterm{Acute Transverse Myelitis} Inflammation across a segment of the spinal cord that can cause rapidly developing weakness, sensory changes, and bladder or bowel dysfunction.

\medterm{Acute Tubular Necrosis} Sudden injury and death of cells in the kidney’s tiny tubes, which normally help filter and balance the blood. It can follow prolonged low blood flow, severe illness, or toxic effects from medicines and other substances, causing acute kidney injury.

\synonyms
ATN

\medterm{Acute Urinary Retention} Sudden inability to pass urine despite a full bladder, often causing severe lower-abdominal pain and swelling. Causes include blockage, medicines, nerve problems, or impaired bladder contraction.

\medterm{Acute Viral Illness} A short-duration illness caused by a virus. Symptoms depend on the affected organ system and may
include fever, fatigue, cough, vomiting, diarrhea, or rash.

\medterm{Acute Vision Loss} Sudden reduction or loss of vision in one or both eyes caused by retinal vascular occlusion, retinal detachment, optic-nerve disease, stroke, acute glaucoma, inflammation, trauma, or other emergency. It requires immediate medical or ophthalmic assessment.

\medterm{Acute Weeping Dermatitis} A form of acute dermatitis with inflamed, often blistered skin that releases clear or yellowish fluid and may form crusts. It is a descriptive pattern that can occur with contact dermatitis, eczema, infection, or other inflammatory skin disorders.

\medterm{Acute-On-Chronic Liver Failure} Acute deterioration of pre-existing chronic liver disease associated with organ failure
and high short-term mortality. It is commonly precipitated by infection, alcohol-related
hepatitis, bleeding, or other acute insults and requires urgent multidisciplinary
management.

\medterm{Acute-Phase Reactants} A group of plasma proteins whose circulating concentrations increase or decrease by at least 25 percent during inflammatory states, tissue injury, or infection. Positive reactants include C-reactive protein, ferritin, fibrinogen, and hepcidin; negative reactants include albumin and transferrin.

\synonyms
Acute-Phase Proteins

\medterm{Acyclovir} An antiviral drug that inhibits herpesvirus DNA polymerase and is used for infections caused by herpes simplex virus and varicella-zoster virus.

\medterm{Ad Lib} A prescription or clinical instruction meaning as desired or as tolerated, from the Latin ad libitum.

\medterm{ADA Deficiency} A rare inherited disorder in which cells lack enough adenosine deaminase, an enzyme needed to break down certain cell chemicals. Toxic substances then build up and damage immune cells, often causing severe infections from early infancy.

\medterm{Adactyly} Absence of one or more fingers or toes from birth.

\medterm{Adam's Apple} The common name for the laryngeal prominence, the visible projection at the front of the
neck formed by the angle where the two laminae of the thyroid cartilage meet.

\medterm{Adamantine} Very hard or resembling diamond; in dentistry, relating to tooth enamel or enamel-like tissue. The term may describe an unusually hard structure or the appearance of a lesion containing enamel-forming tissue.

\medterm{Adamantinoma} A rare, low-grade malignant primary bone neoplasm characterized by a distinctive biphasic histology comprising epithelial cells embedded within an osteofibrous stroma. It shows a dramatic anatomical predilection for the anterior cortex of the mid-shaft (diaphysis) of the tibia (accounting for over 85\% to 90\% of cases), occasionally involving the adjacent fibula. Most commonly diagnosed in young adults aged 20 to 40 years, it presents as a slow-growing, painful pretibial swelling or pathological fracture. Plain radiographs characteristically show eccentric, multilocular, radiolucent cortical lesions with surrounding sclerosis producing a ``soap-bubble'' appearance. Because adamantinoma is resistant to chemotherapy and radiotherapy, standard treatment requires wide en bloc surgical resection with negative margins and bone reconstruction.

\synonyms
Adamantinoma of Long Bones; Tibial Adamantinoma; Malignant Angioblastoma of Bone

\medterm{Adams-Stokes Syndrome} A clinical syndrome characterized by sudden, transient episodes of syncope or loss of consciousness caused by abrupt bradyarrhythmias, high-grade or complete atrioventricular block, or ventricular asystole leading to acute cerebral hypoperfusion.

\synonyms
Stokes-Adams Syndrome; Adams-Stokes Attack

\medterm{Adaptation} Adaptation is the adjustment of a cell, organism, or person to a changing internal or external condition; in medicine it may involve physiologic compensation, behavior, or rehabilitation.

\medterm{Adaptive Design} A clinical-study design allowing planned changes based on accumulating data while preserving scientific validity. Changes may involve sample size, treatment groups, or enrollment and must follow a prespecified plan.

\medterm{Adaptive Equipment} Tools or devices modified to help a person perform daily activities safely despite disability, weakness, pain, or sensory loss. Examples include bathroom rails, reachers, adapted utensils, wheelchairs, and communication aids.

\medterm{Adaptive Immunity} The part of the immune system that learns to recognize specific germs or abnormal cells. B cells make targeted antibodies, T cells coordinate or destroy affected cells, and memory cells help produce a faster response to a later exposure.

\medterm{Addiction} A commonly used, non-diagnostic term for compulsive substance use or
recurrent behavior despite harm, with impaired control, craving, or difficulty stopping.
In clinical practice, the specific diagnosis is usually named as a substance-use disorder
or a recognized disorder due to addictive behavior; physical dependence alone does not
establish addiction.

\medterm{Addiction, Alcohol} Legacy label; see \textit{Alcohol Use Disorder}.

\medterm{Addiction, Gambling} Legacy label; see \textit{Gambling Disorder}.

\medterm{Addiction, Nicotine} Legacy label; see \textit{Nicotine Dependence}.

\medterm{Addiction, Sexual} Legacy label; see \textit{Compulsive Sexual Behavior Disorder}.

\medterm{Addison Anemia} An older name for pernicious anemia, a vitamin B12-deficiency megaloblastic anemia caused by inadequate intrinsic factor and impaired B12 absorption.

\synonyms
Addisonian Anemia; Addison's Anemia

\medterm{Addison Disease} Primary adrenal insufficiency caused by damage to the adrenal glands, so they produce too little cortisol and often too little aldosterone. It may cause fatigue, weight loss, low blood pressure, darkened skin, and adrenal crisis.

\synonyms
Addison's Disease
Chronic Adrenal Insufficiency
Primary Adrenal Insufficiency

\medterm{Addison's Disease} Destruction or dysfunction of the adrenal cortex causing deficiency of cortisol and
aldosterone, presenting with fatigue, hyperpigmentation, weight loss, and salt craving.
Lifelong hormone replacement is required.

\synonyms
Hypoadrenocorticism; Chronic Adrenocortical Insufficiency

\medterm{Addisonian Anemia} Alternate name; see \textit{Addison Anemia}.

\medterm{Addisonian Crisis} Life-threatening emergency resulting from acute cortisol deficiency in patients with
adrenal insufficiency. Precipitated by infection, trauma, surgery, or abrupt
discontinuation of chronic glucocorticoid therapy.

\medterm{Additive Genetic Effects} The contributions of individual alleles that combine approximately additively to influence a measurable trait, without requiring interaction between those alleles.

\medterm{Adduct} To move an arm, leg, finger, toe, or other body part toward the center line of the body.

\medterm{Adduction} Movement of a body part toward the body's midline or, for a finger or toe, toward the hand or foot's central axis. For example, bringing an arm back toward the side is shoulder adduction.

\medterm{Adductor Muscle} A muscle that moves a body part toward the body’s midline or, for a digit, toward the axis of the hand or foot.

\medterm{Adductor Muscles} A group of muscles on the inner thigh that pull the leg toward the body’s midline and help stabilize the hip. Strain of these muscles is a common cause of groin pain, especially during sports.

\medterm{Adenine} A purine nitrogenous nucleobase that forms complementary base pairs with thymine in DNA (via two hydrogen bonds) and with uracil in RNA. It also forms essential components of adenosine, cellular energy carriers (ATP, ADP, AMP), second messengers (cAMP), and key metabolic cofactors (NAD+, FAD).

\synonyms
6-Aminopurine; Ade

\medterm{Adenitis} Inflammation of a gland or lymph node, usually caused by infection or an immune disorder. It may cause swelling, tenderness, redness, or fever; treatment depends on the affected site and cause.

\medterm{Adenitis, Mesenteric} See \textit{Mesenteric lymphadenitis}.

\medterm{Adenocarcinoma} A cancer that begins in cells lining or forming glands. It can arise in organs such as the colon, breast, lung, pancreas, prostate, stomach, or uterus and may invade or spread.

\medterm{Adenocarcinoma In Situ Of Cervix} A noninvasive glandular precursor lesion of the cervix, usually associated with high-risk human papillomavirus; abnormal glandular cells remain confined to the epithelium but may progress to invasive adenocarcinoma.

\medterm{Adenocarcinoma Of The Colon} A malignant gland-forming epithelial tumor of the colon, commonly presenting with
altered bowel habits, occult bleeding, anemia, or an obstructing mass.

\medterm{Adenocarcinoma Of The Lung} The most common type of non-small-cell lung cancer. It begins in gland-forming cells, often toward the outer lung, and can occur in smokers or nonsmokers. Some tumors have gene changes, such as EGFR, ALK, or KRAS, that guide treatment choices.

\synonyms
Pulmonary Adenocarcinoma; Bronchial Adenocarcinoma

\medterm{Adenohypophysis} The front part of the pituitary gland, a small hormone-producing organ at the base of the brain. It releases hormones that regulate growth, thyroid activity, adrenal function, reproduction, milk production, and other body processes.

\medterm{Adenoid Basal Carcinoma} A rare, usually indolent cervical carcinoma composed of small basaloid nests, often associated with squamous intraepithelial lesions and requiring distinction from more aggressive tumors.

\medterm{Adenoid Cystic Carcinoma} A rare malignant cancer that arises in secretory glands, most commonly the major and minor salivary glands of the head and neck, as well as the breast, skin, and airways. It tends to grow slowly but is notable for tracking along nerve pathways (perineural invasion), causing pain or numbness, and recurring or spreading to the lungs years after initial treatment.

\synonyms
Cylindroma

\medterm{Adenoid Cystic Carcinoma Of The Vulva} A rare malignant tumor of vulvar Bartholin-type glands, often causing pain because of perineural invasion. It is usually slow-growing but may recur locally or metastasize after long intervals.

\medterm{Adenoid Hypertrophy} Enlargement of the adenoid tissue located in the nasopharynx behind the nasal cavity. Commonly occurring in children due to recurrent respiratory infections or allergies, it can obstruct nasal breathing, leading to persistent mouth breathing, loud snoring, obstructive sleep apnea, a nasal tone of voice, and recurrent middle-ear fluid accumulation (otitis media with effusion).

\synonyms
Enlarged Adenoids; Adenoidal Hyperplasia

\medterm{Adenoid Removal} Alternate name; see \textit{Adenoidectomy}.

\medterm{Adenoidectomy} Surgical removal of the adenoids (pharyngeal tonsil) from the back of the nasal passage. It is commonly performed in children to treat chronic nasal obstruction, pediatric obstructive sleep apnea, recurrent ear infections, or persistent rhinosinusitis that does not improve with medical therapy.

\synonyms
Adenoid Excision; Adenoid Removal

\medterm{Adenoiditis} Acute or chronic inflammation and infection of the adenoid tissue behind the nose. It causes persistent nasal congestion, purulent nasal discharge, mouth breathing, snoring, sore throat, and muffled speech.

\medterm{Adenoids} A collection of lymphatic tissue located in the upper airway at the back of the nasopharynx, forming part of Waldeyer's ring. In infants and young children, they help sample and trap inhaled germs to build immunity, but typically shrink in size during late childhood and adolescence.

\synonyms
Pharyngeal Tonsil; Nasopharyngeal Tonsil

\medterm{Adenolipoma} A benign tumor composed of glandular cells and mature fat cells. It may occur in the skin, breast, thyroid, parathyroid, or salivary glands.

\medterm{Adenolymphoma} A benign salivary-gland tumor, also called Warthin tumor, composed of oncocytic epithelium and lymphoid stroma and occurring most often in the parotid gland.

\medterm{Adenoma} A usually noncancerous gland-forming tumor arising from epithelial tissue. Some adenomas produce hormones or other secretions, and some—especially certain colon and rectal adenomas—contain precancerous changes and may later become cancerous.

\medterm{Adenoma Malignum} A rare, slow-looking but invasive cancer of the cervix that produces mucus. Persistent watery vaginal discharge or abnormal bleeding may occur. Diagnosis requires examination of cervical tissue by an experienced pathologist because the cells may look almost normal under the microscope.

\medterm{Adenoma-Carcinoma Sequence} The usual step-by-step pattern in which some colon polyps develop genetic changes, become increasingly abnormal, and eventually turn into colon cancer over several years. Removing suitable polyps can interrupt this progression.

\medterm{Adenomas} Plural of adenoma; see \textit{Adenoma}.

\medterm{Adenomatosis} A condition in which multiple gland-forming tumors called adenomas develop in one organ or tissue. The tumors are often benign, but their number and location may affect cancer risk.

\medterm{Adenomatous Polyp} A benign gland-forming polyp, most commonly in the colon, that may contain dysplasia and
can progress through the adenoma--carcinoma sequence. Tubular, villous, and
tubulovillous types are recognized.

\medterm{Adenomatous Polyposis} An autosomal dominant genetic disorder caused by germline mutations in the \textit{APC} tumor suppressor gene, characterized by the progressive development of hundreds to thousands of adenomatous polyps throughout the colon and rectum, carrying a nearly 100 percent lifetime risk of colorectal cancer if left untreated.

\synonyms
Familial Adenomatous Polyposis; FAP

\medterm{Adenomyoepithelioma Of The Breast} A rare biphasic breast tumor composed of proliferating epithelial and myoepithelial cells. Most are benign, but local recurrence and malignant transformation can occur, so complete excision is the usual treatment.

\medterm{Adenomyoma} A benign uterine lesion composed of endometrial glands and stroma within a proliferation of smooth muscle; it may cause abnormal uterine bleeding or mimic other uterine tumors.

\medterm{Adenomyosis} A condition in which tissue normally lining the uterus grows into its muscular wall. It can enlarge the uterus and cause heavy or painful periods, ongoing pelvic pain, or no symptoms; it is more common during the later reproductive years.

\medterm{Adenopathy} Disease or enlargement of a gland, most often used for lymph-node enlargement or abnormal lymph-node involvement.

\medterm{Adenosine} A purine nucleoside made of adenine linked to ribose that participates in energy transfer and cell signaling. Given rapidly by injection, it can briefly slow conduction through the atrioventricular node and is used for selected supraventricular tachycardias.

\medterm{Adenosine Deaminase Deficiency} An inherited purine-metabolism disorder in which toxic deoxyadenosine metabolites accumulate and damage lymphocytes, commonly causing severe combined immunodeficiency.

\medterm{Adenosine Diphosphate} A nucleotide formed when ATP loses one phosphate group. ADP can be converted back to ATP to support cellular energy use, and released ADP also activates nearby platelets and promotes platelet aggregation during clot formation.

\synonyms
ADP

\medterm{Adenosine Monophosphate} A small molecule made of adenine, ribose, and one phosphate group. It is formed when cells break down or recycle energy-carrying molecules and helps regulate energy use and cell signaling.

\synonyms
AMP; Adenylic Acid

\medterm{Adenosine Triphosphate} The main short-term energy source inside cells. Cells make it from nutrients and use it to power muscle contraction, movement of substances across cell membranes, chemical reactions, and many other essential processes; it is commonly called ATP.

\medterm{Adenosis} An increase in the number or size of glands or gland-like tissue. It is often noncancerous, but its meaning depends on the organ involved, such as the breast or lung, and the specific pattern.

\medterm{Adenovirus} A group of DNA viruses that commonly cause respiratory illness, including cold-like symptoms, sore throat, fever, bronchitis, croup, or pneumonia. Depending on the type, adenoviruses can also cause conjunctivitis, gastroenteritis, cystitis, or rarely neurologic disease. They spread through respiratory secretions, close contact, contaminated surfaces, and fecal--oral exposure. Most infections are mild, but severe disease is more likely in infants, people with heart or lung disease, and people with weakened immune systems.

\medterm{Adenovirus Infection} Illness caused by an adenovirus, commonly with sore throat, cough, fever, pink eye, diarrhea, or vomiting, especially in children. The virus spreads through droplets, close contact, contaminated surfaces, or infected stool. One pattern, pharyngoconjunctival fever, brings sore throat, red eyes, and fever together and spreads easily in daycare or swimming-pool settings. Most infections clear on their own, but severe breathing or widespread disease can occur in infants, older adults, people with heart or lung disease, and people with weak immune systems. Treatment is supportive, with fluids and fever relief; severe illness with breathing difficulty or dehydration may require hospital care.

\medterm{Adenoviruses} Plural of adenovirus; see \textit{Adenovirus}.

\medterm{Adenylate Cyclase} An enzyme that changes ATP into cyclic AMP, a chemical messenger inside cells. It is often activated by signals received at the cell surface and helps control many cell functions.

\synonyms
Adenylyl Cyclase

\medterm{Adequate Intake} A formal nutritional reference standard within the Dietary Reference Intakes (DRI) system. It represents the recommended average daily nutrient intake level based on observed or experimentally determined approximations in healthy populations, used when scientific evidence is insufficient to calculate an Estimated Average Requirement (EAR) and Recommended Dietary Allowance (RDA) (such as for vitamin K, potassium, or infant nutrition).

\synonyms
AI; Dietary Reference Intake: Adequate Intake

\medterm{Adherence} The extent to which a person's behavior corresponds with agreed healthcare recommendations, including taking medications, following diets, or executing lifestyle changes. Non-adherence may stem from adverse effects, regimen complexity, cognitive impairment, financial constraints, or lack of support.

\synonyms
Compliance; Treatment Adherence; Medication Adherence

\medterm{Adherens Junction} A connection that helps neighboring cells stick together. It links proteins on one cell to the internal support fibers of another, helping tissues keep their shape, strength, and protective barrier.

\medterm{Adhesiolysis} A procedure that cuts or removes bands of scar tissue that abnormally join organs or body surfaces. It may relieve pain or release a blocked intestine, but the bands can form again after surgery.

\medterm{Adhesion} A band of scar-like tissue that abnormally joins internal surfaces or organs, often after surgery, infection, inflammation, or injury, sometimes causing pain or blockage.

\medterm{Adhesion Molecules} Proteins on cells that help cells attach to one another or to the material supporting them. They help organize tissues, guide immune cells to inflamed areas, and can influence how cancer spreads.

\medterm{Adhesive Capsulitis} Alternate name; see \textit{Frozen Shoulder}.

\medterm{Adhesive Tape} A material used to secure dressings, devices, or tubing to the skin. It is available in paper, cloth, plastic, and other forms.

\medterm{Adie Tonic Pupil} A pupil that is unusually large and reacts slowly to light but constricts better when focusing on a near object. It usually affects one eye and results from injury to nerves controlling the pupil; the cause may be unknown or related to another disorder.

\medterm{Adipocere} A grayish-white, waxy, soap-like substance formed in decomposing bodies exposed to moisture and lack of oxygen (such as damp soil or submerged water). It is produced by the bacterial hydrolysis and hydrogenation of body fats into fatty acids (saponification). This process retards putrefaction and can preserve body contours, facial features, and traumatic injuries for months or years, assisting forensic identification.

\synonyms
Grave Wax; Saponification Wax; Mortuary Wax

\medterm{Adipocyte} A fat cell that stores energy and releases hormones and other signals affecting appetite, blood sugar, inflammation, and metabolism. White fat mainly stores energy, while brown and beige fat can produce heat; enlarged or excessive fat cells can contribute to insulin resistance.

\medterm{Adipocytes} Fat cells that store energy mainly as triglycerides and also release hormones and signaling molecules. White adipocytes primarily store energy, while brown and beige adipocytes can produce heat.

\synonyms
Fat Cell

\medterm{Adipogenesis} The process by which immature cells develop into fat cells. It helps form and maintain body-fat tissue, which stores energy and produces hormones; abnormal increases in the number or size of fat cells can contribute to obesity and metabolic disease.

\medterm{Adipokine} A chemical messenger released by body fat. Adipokines help control appetite, blood sugar, fat use, inflammation, blood pressure, and immune activity; abnormal levels may contribute to obesity, diabetes, and heart or blood-vessel disease.

\medterm{Adiponectin} A hormone produced mainly by fat cells that helps regulate fatty-acid use, glucose metabolism, insulin sensitivity, and inflammation. Levels are often lower with obesity and insulin resistance, but a blood level is not diagnostic by itself.

\medterm{Adipose} Relating to body fat or composed of adipose tissue, which stores energy and also functions as an endocrine and immune-active organ.

\medterm{Adipose Tissue} A specialized connective tissue composed mainly of adipocytes that store energy as
triglycerides. It provides thermal insulation, cushioning for organs, and endocrine
regulation through hormones and signaling molecules such as leptin and adiponectin.

\medterm{Adiposis Dolorosa} Adiposis dolorosa, or Dercum disease, is a rare disorder characterized by painful fatty deposits, often with obesity, fatigue, and other systemic symptoms.

\medterm{Adiposity} The amount and location of body fat. Too much body fat, especially around the abdomen, can increase the risk of diabetes, heart disease, high blood pressure, and joint or movement problems.

\medterm{Adipsia} Little or no feeling of thirst despite the body needing fluid. It can lead to dehydration and dangerously high blood sodium, especially when caused by illness affecting the brain, medicines, aging, or severe illness.

\medterm{Adjacent} Next to, close to, or sharing a boundary with another structure, tissue, organ, or lesion. The term helps describe location, spread of disease, or relationships seen on imaging or examination.

\medterm{Adjustment Disorder} Emotional or behavioral symptoms in response to an identifiable stressor, beginning within three months of the stressor. The symptoms cause significant distress or impairment and do not meet criteria for another mental disorder. They generally do not persist for more than six months after the stressor or its consequences have ended.

\medterm{Adjustment Disorder With Anxiety} An emotional or behavioral response to an identifiable stressor that causes distress or
impairment but does not meet criteria for another mental disorder.

\medterm{Adjustment Disorders} Alternate index label; see \textit{Adjustment Disorder}.

\medterm{Adjuvant} An added treatment or substance that enhances the effect of a primary treatment; in vaccines, an adjuvant strengthens the immune response to an antigen.

\medterm{Adjuvant Chemotherapy} Systemic antineoplastic drug therapy administered following definitive primary treatment, such as complete surgical tumor resection, designed to eradicate occult micrometastatic disease and reduce the risk of cancer recurrence.

\medterm{Adjuvant Therapy} Additional treatment given after the main treatment to reduce the chance of disease returning. Examples include chemotherapy, radiation, hormone treatment, or immune treatment after surgery for selected cancers.

\medterm{Adlerian Theory} A psychological theory developed by Alfred Adler, also called individual psychology. It emphasizes feelings of inferiority, goal-directed striving, a person’s style of life, and social interest—the capacity to belong to and contribute to the community. Early family experiences and relationships are considered important influences on development.

\synonyms
Individual Psychology

\medterm{ADME} The four processes that describe what the body does to a medicine: absorption into the bloodstream, distribution through tissues, chemical breakdown, and excretion from the body. These processes influence a medicine’s dose, timing, and effects.
\synonyms
Pharmacokinetics

\medterm{Administered Activity} The measured amount of radioactive material given to a patient for a scan or treatment. The amount is recorded with the radioactive substance, method, date, and time so the exposure and treatment can be checked safely.

\medterm{Administration Procedure} A method used to give a medicine, fluid, vaccine, nutrition, or other treatment, such as by mouth, injection, inhalation, skin application, or infusion. The method affects speed, dose, and risks.

\medterm{Adnexa} Structures next to an organ. In gynecology, the uterine adnexa are the ovaries, fallopian tubes, and supporting ligaments and tissues around the uterus.

\medterm{Adnexal Dysplasia} Abnormal cell development in tissue beside the uterus, such as an ovary or fallopian tube. The change may be related to development or may represent a precancerous process, so its importance depends on the organ and tissue findings.

\medterm{Adnexal Mass} A growth or enlargement involving structures beside the uterus, such as an ovary or fallopian tube.

\medterm{Adnexal Torsion} Twisting of an ovary, fallopian tube, or both around their supporting tissues. The twist can reduce or stop blood flow. It usually causes sudden one-sided pelvic pain, often with nausea or vomiting, and requires urgent medical assessment and usually surgery.

\medterm{Adnexal Tumors} Tumors arising in structures beside the uterus, especially the ovaries, fallopian tubes, or supporting tissues. They may be benign, borderline, or malignant and are evaluated with examination, imaging, laboratory tests, and sometimes surgery or biopsy.

\medterm{Adnexal Tumors And Masses} Abnormal growths or masses involving the ovary, fallopian tube, or nearby supporting tissues; they may be benign or malignant.

\medterm{Adnexitis} Inflammation of one or more uterine adnexa, usually the fallopian tubes and sometimes the ovaries. It commonly occurs as part of pelvic inflammatory disease and may cause lower abdominal pain, fever, abnormal discharge, or pain during sex; untreated infection can lead to an abscess, infertility, or ectopic pregnancy.

\synonyms
Salpingo-Oophoritis

\medterm{Adolescence} The developmental period between childhood and adulthood, involving physical maturation, puberty, and major psychological and social changes.

\medterm{Adolescent Development} The physical, cognitive, emotional, and social changes that occur as a person develops from childhood into adulthood, including puberty and the development of greater independence.

\medterm{Adolescent Health} The health and wellbeing of young people (10-19 years), addressing the physical,
psychological, and social transitions of puberty and adolescence.

\medterm{Adolescent Medicine} A medical specialty focused on the physical, mental, reproductive, and social health needs of adolescents and young adults, including prevention, risk behavior, chronic disease, and confidentiality-sensitive care.

\medterm{Adolescent Obesity} Alternate index label; see \textit{Pediatric Obesity}.

\medterm{Adolescent Schizophrenia} See \textit{Childhood schizophrenia}.

\medterm{Adoptive Cell Therapy} A personalized cellular immunotherapy in which a patient's own lymphocytes are expanded
or engineered ex vivo and reinfused to attack the tumor; examples include
tumor-infiltrating lymphocyte (TIL) therapy and chimeric antigen receptor (CAR) T-cell
therapy.

\medterm{Adoptive Immunotherapy} Treatment in which immune cells are collected, activated or genetically modified, and returned to the patient to attack disease, especially cancer. Fever, inflammation, and other immune reactions may occur.

\medterm{ADP-Ribosylation} A chemical change in which a small molecule is attached to a protein or other cell molecule. This can switch the molecule’s activity on or off and is used by normal cell processes and by some bacterial toxins.

\medterm{Adrenal Cancer} A rare cancer that begins in one of the adrenal glands above the kidneys. Some tumors start in the outer cortex and others in the inner medulla; some produce excess hormones, while others cause pain or other symptoms as they grow.

\medterm{Adrenal Carcinoma} A rare malignant tumor arising from the adrenal cortex. It may produce excess cortisol, aldosterone, androgens, or estrogens, or present as an incidental mass or pressure symptom; hormonal evaluation and imaging help guide diagnosis and staging.

\medterm{Adrenal Cortex} The outer part of each adrenal gland that produces steroid hormones, including cortisol, aldosterone, and adrenal androgens. These hormones help regulate stress responses, blood pressure, salt balance, and metabolism.

\medterm{Adrenal Cortex Physiology} The outer layer of each adrenal gland makes hormones that control salt and blood pressure, stress responses, metabolism, and some sexual development. Its activity is regulated by the kidneys, blood potassium, and the brain’s hormone-control system.

\medterm{Adrenal Cortical Carcinoma} Alternate name; see \textit{Adrenocortical Carcinoma}.

\medterm{Adrenal Crisis} A sudden, life-threatening shortage of adrenal hormones, especially cortisol. It may occur with severe adrenal disease or after abruptly stopping long-term steroid treatment, particularly during infection, injury, or surgery. Symptoms can include severe weakness, vomiting, abdominal pain, confusion, low blood pressure, and collapse; emergency treatment is required.

\medterm{Adrenal Failure} Inadequate production of adrenal hormones, especially cortisol and sometimes aldosterone; severe acute deficiency can cause adrenal crisis with hypotension and shock.

\medterm{Adrenal Gland} One of two hormone-producing glands above the kidneys. The outer cortex makes aldosterone, cortisol, and adrenal androgens; the inner medulla makes adrenaline and noradrenaline. Adrenal hormones help regulate blood pressure, metabolism, stress responses, and salt balance.

\medterm{Adrenal Gland Cancer} Alternate name for adrenal cancer, a malignant tumor arising in an adrenal gland.

\medterm{Adrenal Gland Removal} Surgical excision of an adrenal gland, performed for tumors, hyperplasia, or other adrenal disorders. Approaches include open, laparoscopic, or retroperitoneal; perioperative management requires assessment of hormone production and potential adrenal insufficiency.

\medterm{Adrenal Gland Tumor} A growth in one of the adrenal glands above the kidneys. It may be noncancerous or cancerous and may or may not produce hormones; a pheochromocytoma is one possible type.

\medterm{Adrenal Glands} Adrenal glands sit above the kidneys and produce cortisol, aldosterone, adrenal androgens, and catecholamines that regulate stress, blood pressure, salt balance, and metabolism.

\medterm{Adrenal Hemorrhage} Bleeding into an adrenal gland, often associated with severe infection, serious illness, a bleeding disorder, blood-thinning medicines, or injury. Bleeding into both glands can cause sudden loss of adrenal hormones, very low blood pressure, and shock.

\medterm{Adrenal Hyperplasia} Enlargement or overgrowth of adrenal tissue that may increase production of adrenal
hormones. Causes include congenital adrenal hyperplasia, chronic stimulation by ACTH, and some
adrenal tumors. Symptoms depend on the hormone affected and may include abnormal genital development,
early puberty, excess body hair, high blood pressure, or electrolyte disturbances.

\medterm{Adrenal Incidentaloma} An unexpected mass found in an adrenal gland during imaging for another reason. It may be benign or cancerous and may or may not produce hormones, so evaluation usually includes imaging review and selected hormone tests.

\medterm{Adrenal Insufficiency} A condition in which the body does not make enough cortisol and sometimes not enough aldosterone. Primary adrenal insufficiency results from damage or dysfunction of the adrenal glands; secondary and tertiary forms result from inadequate pituitary or hypothalamic signals, including suppression after prolonged steroid treatment. Features may include fatigue, low blood pressure, weight loss, salt imbalance, or adrenal crisis.

\medterm{Adrenal Mass} An abnormal growth or enlargement in or near an adrenal gland. It may be benign or cancerous and may or may not produce hormones; evaluation usually involves imaging and, when appropriate, hormone tests and specialist review.

\medterm{Adrenal Medulla} The inner part of each adrenal gland, which sits above a kidney. It releases adrenaline and noradrenaline during stress, helping increase heart rate, blood pressure, alertness, and blood flow to muscles.

\medterm{Adrenal Medulla Physiology} The inner part of each adrenal gland releases adrenaline and noradrenaline when the body’s stress-response nerves are activated. These hormones quickly increase alertness, heart activity, blood pressure, and blood flow to muscles.

\medterm{Adrenal Myelolipoma} A usually benign adrenal tumor composed of mature fat and blood-forming tissue. It is often
found incidentally on imaging and causes no symptoms, but larger lesions may produce mass
effect or bleeding. Management depends on size, symptoms, growth, and diagnostic
certainty.

\medterm{Adrenal Nodule} A localized lump in an adrenal gland. It may be a harmless nonfunctioning growth, a hormone-producing tumor, or rarely cancer; assessment depends on its size, imaging features, growth, and hormone production.

\medterm{Adrenal Tumors} Tumors arising in the adrenal cortex or medulla may be benign or malignant and may or may not produce
hormones. Evaluation uses clinical assessment, hormone testing when indicated, and
imaging to assess size and appearance; treatment depends on the diagnosis, hormone
secretion, symptoms, growth, and risk of malignancy.

\medterm{Adrenal Venous Sampling} A catheter-based test that measures hormones in blood from each adrenal gland. It helps determine whether one gland is producing too much aldosterone, which can guide treatment for difficult-to-control high blood pressure.

\medterm{Adrenalectomy} Surgery to remove one or both adrenal glands, usually for a tumor or severe hormone overproduction. Removing one gland may leave enough hormone production; removing both requires lifelong hormone replacement and an emergency steroid plan.

\medterm{Adrenaline} Also called epinephrine, a hormone and neurotransmitter released during stress. It increases heart rate and contractility, raises blood pressure in many vascular beds, relaxes airway muscle, and is the first-line emergency medicine for anaphylaxis.

\medterm{Adrenarche} An early stage of sexual development, typically occurring around age six to eight, marked by increased production and secretion of weak adrenal androgens (such as dehydroepiandrosterone) by the zona reticularis of the adrenal cortex. It contributes to early pubic hair growth, axillary odor, and oily skin.

\medterm{Adrenergic Agent} A medicine or substance that stimulates or blocks receptors responding to adrenaline (epinephrine) and noradrenaline (norepinephrine). These agents can change heart rate, blood pressure, airway size, and pupil size.

\medterm{Adrenergic Agonist} A medicine or substance that copies the effects of adrenaline or noradrenaline by activating specific cell receptors. Depending on the receptor, it may raise blood pressure, speed the heart, open the airways, or improve heart pumping.

\synonyms
Sympathomimetic

\medterm{Adrenergic Antagonists} Medicines or other substances that block adrenaline-like signaling at alpha or beta adrenergic receptors. They can lower blood pressure, slow the heart, reduce cardiac workload, or relax smooth muscle; effects and risks depend on the receptor targeted and the medicine used.

\synonyms
Adrenergic Blockers

\medterm{Adrenergic Bronchodilator Overdose} Excess exposure to a medicine that opens the airways by stimulating the body’s stress-response system. It may cause shaking, agitation, headache, fast heartbeat, low potassium, chest pain, or abnormal rhythms and needs prompt medical assessment.

\medterm{Adrenergic Receptor} A cell-surface receptor activated by adrenaline (epinephrine), noradrenaline (norepinephrine), or related medicines. Alpha and beta receptors help control heart rate, blood pressure, airway size, pupil size, and other body responses.

\medterm{Adrenergic Receptors} Cell-surface receptors activated by adrenaline and noradrenaline. Alpha and beta subtypes regulate cardiovascular, respiratory, metabolic, and other physiological responses.

\medterm{Adrenocortical Carcinoma} A rare cancer that begins in the outer layer, or cortex, of an adrenal gland, one of the glands above the kidneys. Some tumors are functioning and make too much of hormones such as cortisol, aldosterone, or adrenal androgens; others are nonfunctioning and make no excess hormones. Hormone-producing tumors may cause weight gain, high blood pressure, muscle weakness, unusual hair growth, or other changes in sexual development. A nonfunctioning tumor may cause no early symptoms, but a growing or spreading tumor may cause an abdominal lump, abdominal or back pain, or a feeling of fullness.

\medterm{Adrenocorticotropic Hormone} A hormone made by the pituitary gland that tells
the adrenal glands to make cortisol. Cortisol helps the body respond to stress and
helps control blood pressure, blood sugar, and inflammation.

\synonyms
Adrenocorticotrophic Hormone (ACTH); ACTH; Corticotropin

\medterm{Adrenogenital Syndrome} A group of disorders, usually caused by congenital adrenal hyperplasia, in which excess adrenal androgens cause atypical genital development or early puberty.

\medterm{Adrenoleukodystrophy} A group of inherited disorders, most often X-linked, in which very-long-chain fatty acids accumulate and damage myelin, the adrenal glands, or both. Childhood cerebral disease can cause behavior and learning changes, seizures, and loss of neurologic function; adrenal insufficiency may also occur.

\synonyms
ALD

\medterm{Adrenomedullin} A hormone-like substance made by many tissues that widens blood vessels and helps regulate blood pressure, circulation, fluid balance, kidney activity, and responses to injury.

\medterm{Adrenomyeloneuropathy} An adult-onset form of X-linked adrenoleukodystrophy that mainly affects the spinal cord and peripheral nerves. It may cause leg stiffness, weakness, walking difficulty, bladder or bowel problems, and sometimes adrenal insufficiency.

\medterm{Adson Test} A physical examination manoeuvre used to assess for thoracic outlet syndrome and compression of the subclavian artery by a cervical rib or scalene muscles. The examiner palpates the radial pulse while the patient abducts and extends the arm, extends the neck, and rotates the head toward the ipsilateral side with deep inspiration; diminution or obliteration of the radial pulse indicates a positive test.

\synonyms
Adson Manoeuvre; Adson's Test

\medterm{Adsorption} The sticking of molecules to the outside surface of a solid or liquid rather than entering its substance. Activated charcoal works partly by adsorbing some poisons in the digestive tract.

\medterm{Adult Day Services} Community-based daytime programs providing supervision, social activities, health-related support, or assistance with daily activities for adults who cannot safely remain alone throughout the day.

\synonyms
Adult Day Care; Adult Day Program

\medterm{Adult Hypertension} Blood pressure that stays too high over time. The exact level used for diagnosis depends on the guideline and the person's health. Untreated high blood pressure raises the risk of heart disease, stroke, and kidney disease.

\synonyms
Essential Hypertension; Primary Hypertension; Systemic Hypertension

\medterm{Adult Immunisation} The practice of administering vaccines to adults to protect against preventable infectious diseases, maintain waning immunity from childhood vaccinations, and reduce illness in high-risk groups. Essential adult vaccines include annual influenza, tetanus-diphtheria-pertussis (Td/Tdap) boosters every 10 years and during pregnancy, pneumococcal vaccines for adults aged 65 and older or those with chronic medical conditions, the recombinant herpes zoster vaccine for adults aged 50 and older to prevent shingles, and targeted vaccines for COVID-19, hepatitis B, HPV, and RSV.
\synonyms{Adult Immunization; Adult Vaccination; Preventive Adult Vaccines}

\medterm{Adult Informed Consent} The ethical and legal clinical communication process whereby a competent adult patient is provided with comprehensive, understandable information regarding the nature, indication, benefits, material risks, and alternatives of a proposed medical intervention or procedure, enabling an autonomous and voluntary decision.

\synonyms
Informed Medical Consent; Surgical Consent; Patient Autonomous Decision-Making

\medterm{Adult Nutritional Support in Illness} Dietary strategies providing increased caloric and protein intake during acute or chronic illness, convalescence, or hypermetabolic states to prevent involuntary weight loss and muscle wasting.

\synonyms
High-Calorie Diet in Illness; Nutritional Therapy in Recovery

\medterm{Adult Obstructive Sleep Apnea} Recurrent nocturnal collapse or partial obstruction of the pharyngeal airway in adults, producing repetitive episodes of hypopnea or apnea, transient arterial oxyhemoglobin desaturations, sympathetic surge, and fragmented sleep architecture leading to daytime somnolence and increased cardiovascular morbidity.

\synonyms
Adult OSA; Obstructive Sleep Apnea Syndrome; Upper Airway Sleep Apnea

\medterm{Adult Primary Brain Tumor} A neoplasm arising de novo within the brain parenchyma, cranial nerves, or meninges in adults, encompassing diffuse gliomas (astrocytoma, oligodendroglioma, glioblastoma), meningiomas, vestibular schwannomas, and primary central nervous system lymphomas.

\synonyms
Adult Primary Brain Tumor; Primary Intracranial Neoplasm; Adult Glioma

\medterm{Adult Spinal Deformity} Abnormal curvature or alignment of the spine in an adult, including scoliosis, kyphosis, or loss of
normal sagittal balance. It may cause back pain, nerve compression, difficulty walking,
or functional disability.

\medterm{Adult T-Cell Leukemia} A rare and aggressive peripheral T-cell neoplasm caused by infection with the human T-lymphotropic virus type 1 (HTLV-1). Characterized by generalized lymphadenopathy, hepatosplenomegaly, osteolytic bone lesions, severe hypercalcemia, skin infiltration, and atypical multi-lobed circulating lymphocytes known as flower cells.

\synonyms
Adult T-Cell Leukemia-Lymphoma; ATLL

\medterm{Adult-Onset Asthma} Adult-onset asthma is variable airway inflammation and narrowing that begins in adulthood, producing episodic wheezing, cough, chest tightness, or shortness of breath. Diagnosis uses symptoms and lung-function testing, and treatment reduces triggers and airway inflammation.

\synonyms
Late-Onset Asthma; Adult-Onset Reactive Airway Disease

\medterm{Adult-Onset Still Disease} A rare inflammatory disease that causes recurring high fevers, joint pain or swelling, a salmon-colored rash, and sore throat. It may also inflame the liver, heart, lungs, or other organs.

\medterm{Advance Directive} A legal document stating a person's wishes about future medical care and, in some cases, naming a health-care proxy to make decisions if the person cannot decide or communicate. It may address resuscitation, life support, artificial nutrition, and other treatments and is reviewed as circumstances change.

\synonyms
Living Will; Medical Power of Attorney; Healthcare Directive

\medterm{Advance Statements} Alternate name; see \textit{Living Will}.

\medterm{Advanced Cancer} Cancer that has grown beyond its original site. It may be locally advanced when it has
invaded nearby tissues or metastatic when it has spread to distant organs.

\medterm{Advanced Cardiovascular Life Support} A standardized set of advanced resuscitation measures for cardiac arrest and other
life-threatening cardiovascular emergencies. It includes high-quality CPR, defibrillation
when indicated, airway management, medicines, and treatment of reversible causes.

\synonyms
ACLS

\medterm{Advanced Glycation End-Products} Compounds formed when sugar attaches to proteins or fats without the usual enzyme-controlled process. They increase with aging and persistently high blood sugar and may promote inflammation and damage to blood vessels, nerves, kidneys, and other tissues.

\medterm{Advanced Prostate Cancer} Prostate cancer that has grown beyond the prostate or spread to distant parts of the body. It may involve nearby tissues, lymph nodes, bones, or other organs.

\medterm{Advanced Sleep-Wake Phase Disorder} A body-clock disorder in which a person becomes sleepy much earlier than desired and wakes unusually early. The person may sleep normally on this early schedule but cannot stay asleep until a more convenient time.

\medterm{Adventitia} The outer layer of connective tissue around a blood vessel or hollow organ. It supports and anchors the structure and contains small nerves and blood vessels that nourish the wall.

\synonyms
Tunica Adventitia

\medterm{Adventitious} Abnormal, unexpected, or arising from an unusual location or process; adventitious lung sounds include crackles, wheezes, and other added breath sounds.

\medterm{Adverse Effect} An unintended and harmful or bothersome effect of a medicine or other treatment given at normal doses or for an intended purpose.

\medterm{Adverse Event} Any new or worsening health problem that happens during or after a medicine, procedure, or other treatment, whether or not the treatment caused it. Clinicians record its timing, severity, and likely cause separately.

\medterm{Adverse Reaction} An unintended and harmful response to a medicine, treatment, or other exposure. It may be predictable from the medicine's effects, such as bleeding with an anticoagulant, or unpredictable, such as an allergic reaction. Severity ranges from mild nausea or rash to serious organ injury, breathing difficulty, or shock; suspected severe reactions need urgent medical attention.

\synonyms
Adverse Drug Reaction

\medterm{Aedes Aegypti} Aedes aegypti is a mosquito that can transmit dengue, yellow fever, chikungunya, Zika, and other viruses.

\medterm{Aegophony} An altered voice sound heard through a stethoscope when a person says “E” but it sounds like “A.” It can occur over lung tissue made denser by infection or fluid, but it does not identify the cause by itself.

\medterm{Aerobe} A microorganism that grows or produces energy using oxygen. Some aerobes cannot live without oxygen, while others can grow either with or without it. The distinction helps laboratories understand how an organism may behave in the body.

\medterm{Aerobic} Requiring or using oxygen. In microbiology, an aerobic organism grows with oxygen; in the body, aerobic energy production uses oxygen to release energy from nutrients for muscles and other cells.

\medterm{Aerobic Bacteria} Bacteria that need oxygen to grow. They may cause infection in tissues exposed to air, such as the lungs, skin, or surface of a wound; laboratory tests identify the specific species.

\medterm{Aerobic Bacterium} An aerobic bacterium grows using oxygen, although some species can tolerate or use alternative metabolic pathways; oxygen requirements help guide culture and infection assessment.

\medterm{Aerobic Exercise} Physical activity using large muscles continuously while the body produces energy with oxygen. Walking, cycling, swimming, and dancing can improve heart, lung, blood-vessel, muscle, mood, and endurance health.

\synonyms
Cardio Exercise

\medterm{Aerobiosis} Life or biological activity that occurs when oxygen is available. Many human cells use oxygen to release energy from food, and some microorganisms need oxygen to grow.

\medterm{Aerophagia} Excessive swallowing of air, which may cause frequent belching, abdominal bloating, discomfort, or increased intestinal gas. It can occur with rapid eating, chewing gum, anxiety, or swallowing disorders.

\medterm{Aerophagy} Repeated swallowing of too much air, which can cause frequent belching, abdominal fullness, bloating, or discomfort. It may occur with rapid eating, chewing gum, anxiety, or difficulty swallowing.

\synonyms
Aerophagia

\medterm{Aerophobia} An excessive fear of flying or air travel that causes marked anxiety, avoidance, or interference with daily life. It is a specific phobia and may improve with structured psychological treatment.

\medterm{Aerosol} Very small solid particles or liquid drops floating in a gas such as air. Aerosols can carry inhaled medicines, pollutants, or infectious particles and may remain in the air or settle on surfaces.

\medterm{Aerospace Medicine} The field of medicine concerned with maintaining the health, safety, and performance of people who travel or work in aircraft and spacecraft. It addresses effects such as changes in pressure, acceleration, low oxygen, radiation, isolation, and microgravity.

\medterm{Aerotitis} Ear pain or middle-ear injury caused by failure to equalize pressure during rapid changes in altitude, such as during aircraft ascent or descent.

\medterm{Aesthetic Dentistry} Dental care intended to improve the appearance of the teeth, gums, or smile. Procedures may include
tooth-colored restorations, veneers, whitening, or selected crowns.

\medterm{Aesthetic Medicine} Medical care intended to improve appearance or reduce visible signs of aging through procedures, devices, medicines, or skin care; appropriate assessment includes risks, contraindications, and realistic expectations.


\synonyms
Cosmetic Medicine; Aesthetic Dermatology

\medterm{Aetiology} The study of causation, or the specific originating cause or combination of pathogenic factors (such as genetic, environmental, infectious, or autoimmune mechanisms) that lead to the development of a disease or clinical disorder.

\synonyms
Etiology; Causation; Pathogenesis

\medterm{Affect} The observable expression of a person’s emotional state, including facial expression, voice, posture, and emotional responsiveness.

\medterm{Affective Disorder} Alternate name; see \textit{Mood Disorders}.

\synonyms
Mood Disorder

\medterm{Affective Flattening} Markedly reduced outward expression of emotion, including limited facial expression, voice modulation, and gestures.
It may occur in schizophrenia and other psychiatric or neurologic conditions.

\medterm{Afferent} Carrying signals toward a central organ or structure; afferent nerves transmit sensory information from the body toward the brain and spinal cord.

\medterm{Afferent Arteriole} A small branch of the renal arterial circulation that conveys blood directly into the glomerular capillary tuft of an individual nephron, playing a pivotal role in regulating glomerular filtration pressure and renal blood flow through autoregulatory and tubuloglomerular mechanisms.

\medterm{Afferent Loop Syndrome} A blockage or poor emptying in a section of small intestine made during some stomach operations. Digestive juices build up behind the blockage, causing upper-abdominal pain, nausea, vomiting, or poor absorption of nutrients.

\medterm{Afferent Lymphatic Vessel} A lymphatic vessel that carries lymph into a lymph node.

\medterm{Afferent Neuron} A sensory nerve cell carrying signals from receptors in the body toward the spinal cord and brain, where information is processed and appropriate responses begin.

\medterm{Afferent Pathway} A sensory nerve route carrying information from the skin, muscles, organs, or other tissues toward the spinal cord or brain, where the signals are processed and responses can begin.
\synonyms
Sensory Pathway

\medterm{Afferent Pupillary Defect} An asymmetry in pupillary light reactivity caused by a unilateral or asymmetric lesion in the anterior visual afferent pathway (optic nerve or retina). Identified during the swinging-flashlight test when swinging the light beam from the unaffected eye to the affected eye results in paradoxical bilateral pupillary dilation rather than constriction.

\synonyms
Marcus Gunn Pupil; Relative Afferent Pupillary Defect; RAPD

\medterm{Affinity} The strength of attraction or binding between two molecules, such as an antibody and antigen, receptor and ligand, or enzyme and substrate.

\medterm{Affinity Chromatography} A laboratory method for separating and purifying a substance by allowing it to stick to a partner designed to recognize it. Other substances are washed away, and the desired substance is then released.

\medterm{AFib} Alternate name; see \textit{Atrial Fibrillation}.

\medterm{Afibrinogenaemia} Alternate spelling; see \textit{Afibrinogenemia}.

\medterm{Afibrinogenemia} A rare congenital or severe acquired coagulation disorder characterized by the complete or near-complete absence of circulating fibrinogen (factor I, <0.1 g/L). Inherited in an autosomal recessive pattern (mutations in \textit{FGA}, \textit{FGB}, or \textit{FGG}), it manifests with umbilical cord bleeding in neonates, spontaneous mucosal and intracranial hemorrhages, hemarthroses, and prolonged bleeding times.

\synonyms
Congenital Afibrinogenemia; Fibrinogen Deficiency; Afibrinogenaemia

\medterm{Aflatoxin} A toxin made by certain molds that can contaminate foods such as grains, nuts, and spices. Long-term exposure can damage the liver and increase the risk of liver cancer.

\medterm{African Sleeping Sickness} An infection caused by \textit{Trypanosoma brucei} parasites and spread by infected tsetse flies in sub-Saharan Africa. It may cause fever and swollen lymph nodes, followed by sleep changes and nervous-system problems; untreated disease can be fatal.

\medterm{African Trypanosomiasis} An infection caused by \textit{Trypanosoma} parasites and spread by infected tsetse flies in sub-Saharan Africa. It may cause fever and swollen lymph nodes, followed by sleep changes and nervous-system problems; untreated disease can be fatal.

\synonyms
African Sleeping Sickness

\medterm{After-Cataract} Clouding of the thin membrane behind the artificial lens after cataract surgery. It can develop months or years later and cause blurred or hazy vision, glare around lights, or reduced sharpness of vision.

\synonyms
Posterior Capsule Opacification; Secondary Cataract

\medterm{Afterbirth} The placenta and the membranes that surrounded the baby, which leave the womb after childbirth. If any part remains inside, it can cause heavy bleeding or infection.

\medterm{Afterdepolarization} An abnormal electrical oscillation occurring during or immediately following myocardial action potential repolarization, capable of reaching threshold and triggering cardiac arrhythmias. Early afterdepolarizations occur during plateau or repolarization phases, while delayed afterdepolarizations occur after repolarization is complete.

\medterm{Afterload} The pressure or resistance the heart must overcome to push blood out of its ventricles. Higher blood pressure or narrowing of blood vessels increases the heart’s workload and can make pumping less efficient.

\medterm{Afterpains} Cramp-like pains in the lower abdomen during the first few days after childbirth, caused by the womb tightening and shrinking. They may become stronger during breastfeeding and help reduce bleeding.

\medterm{Aftershave Poisoning} Harm from swallowing aftershave, which may contain alcohol, fragrance, or other chemicals. It can cause vomiting, drowsiness, low blood sugar, breathing problems, or coma; significant ingestion is an emergency requiring prompt poison-center assessment and medical treatment.

\medterm{Agammaglobulinaemia} A serious immune deficiency in which the blood contains very few or no antibodies, usually because antibody-producing B cells do not develop normally. It causes repeated bacterial infections, often beginning in infancy or early childhood.

\medterm{Agammaglobulinemia} A serious immune deficiency in which the blood contains very few or no antibodies. It often results from abnormal B-cell development and causes repeated bacterial infections of the ears, sinuses, lungs, or bloodstream.

\medterm{Agar} A jelly-like substance made from red seaweed. Laboratories add it to nutrient mixtures to grow and identify bacteria and fungi on solid surfaces, such as culture plates. It is also used in some bulk laxatives.

\medterm{Agatston Score} A number calculated from a special CT scan that measures calcium in the heart’s arteries. A higher score usually means more coronary plaque and higher heart-disease risk, but it does not directly show blockage. Doctors interpret it with age, sex, blood pressure, cholesterol, diabetes, smoking, family history, and other health conditions.

\synonyms
Coronary Artery Calcium (CAC) Score; Agatston Calcium Score

\medterm{Age Spots} Flat, tan, brown, or black spots that develop on chronically sun-exposed areas of the skin, most commonly the face, backs of the hands, forearms, and shoulders. Caused by localized proliferation of pigment cells (melanocytes) from years of ultraviolet (UV) radiation exposure and aging, they are benign and harmless. Despite the colloquial name ``liver spots,'' they have no connection to liver function or liver disease.

\synonyms
Liver Spots; Solar Lentigines; Senile Lentigines; Sun Spots

\medterm{Age-Related Hearing Loss} Gradual hearing loss in both ears caused by aging. It often affects high-pitched sounds first and can make speech harder to understand.

\medterm{Age-Related Macular Degeneration, Dry} See \textit{Dry macular degeneration}.

\medterm{Age-Related Macular Degeneration, Wet} See \textit{Wet macular degeneration}.

\medterm{Age-Standardized Rate} A disease or death rate recalculated as if different populations had the same age distribution. This makes comparisons fairer when one population has more older or younger people than another.

\medterm{Ageing} Ageing is the natural, lifelong process of physical and biological change. Over time, body reserves may decrease and the risk of some diseases may increase, but the effects vary between people.

\synonyms
Aging

\medterm{Agenesis} Failure of an organ or body structure to form before birth, so it is absent. Agenesis may affect one organ, such as a kidney, or a specific body structure, such as the corpus callosum.

\medterm{Agenesis Of The Corpus Callosum} A birth condition in which the band of nerve fibres connecting the brain’s two halves is partly or completely absent. Effects vary widely and may include seizures, developmental delay, learning difficulties, or no noticeable symptoms.

\medterm{Ageusia} Complete loss of the ability to taste. It may result from infections, mouth or nose problems, medicines, chemotherapy, head injury, nerve damage, nutritional deficiency, or neurological disease. Loss of smell can also make food seem tasteless.

\medterm{Agger Nasi} The most anterior and superior ethmoidal air cell, located immediately anterior to the attachment of the middle nasal turbinate and serving as an important anatomical landmark in functional endoscopic sinus surgery.

\medterm{Agglutination} Visible clumping that occurs when antibodies or other binding substances join cells or particles together. Laboratory workers use this reaction to help identify blood groups, infections, and other substances in patient samples.

\medterm{Agglutination Test} A laboratory test that mixes a patient sample with a known antibody or antigen. Visible clumping suggests that the matching substance is present and may help identify blood groups or infections.

\medterm{Agglutinin} An antibody or other substance that causes cells or particles to stick together in visible clumps. Agglutinins are used in laboratory tests, including blood-group testing and some infection tests.

\medterm{Agglutinogen} An antigen on a cell or particle that can react with a matching antibody and cause visible clumping. Agglutinogens are especially important in identifying blood groups.

\medterm{Aggression} Words or actions intended to threaten, harm, or intimidate another person or damage property. Aggression may result from emotional distress, pain, mental illness, brain disorders, medicines, or substance use.

\medterm{Aggressive} Describing hostile or harmful behavior, or a disease that grows, damages nearby tissue, or spreads quickly. The meaning depends on whether it describes a person’s behavior or a medical condition.

\medterm{Aggressive Angiomyxoma} A rare, usually slow-growing soft-tissue tumor that most often develops in the vulva, perineum, or pelvis of women of reproductive age. It can grow into nearby tissue and often returns after treatment, but rarely spreads to distant organs.

\medterm{Aggressive Cancer} A cancer that is likely to grow, invade nearby tissue, or spread quickly. Doctors assess its aggressiveness using the cancer type, cell appearance, genetic features, stage, growth rate, and response to treatment.

\medterm{Agility} The ability to move, stop, or change direction quickly while staying balanced and controlled. Doctors and therapists may assess agility when evaluating movement, coordination, recovery from injury, or physical function.

\medterm{Agitation} An uncomfortable state of inner tension, restlessness, or increased activity. It may cause pacing, irritability, or inability to remain calm and can result from medical illness, pain, medicines, substance use, or mental health conditions.

\medterm{Agnathia} A rare birth condition in which the lower jaw is partly or completely absent. It may cause major facial abnormalities and serious problems with breathing, swallowing, and feeding.

\medterm{Agnosia} An inability to recognize familiar objects, people, sounds, or other sensations even though the related sense works normally. It usually results from damage to areas of the brain that interpret sensory information.

\medterm{Agonist} A drug or other ligand that binds to a receptor and activates it, producing a biological response. A full agonist can produce the maximum response in a given system, whereas a partial agonist produces a submaximal response even when occupying receptors.

\medterm{Agonist Muscle} The main muscle responsible for producing a particular movement. It is also called the prime mover.

\medterm{Agoraphobia} An anxiety disorder involving persistent, intense fear of places or situations where escape may be difficult, help may be unavailable, or panic-like symptoms may feel embarrassing. People may avoid crowds, public transportation, shops, or leaving home alone; the fear and avoidance cause significant distress or interfere with daily functioning.

\medterm{Agranulocytes} A group of white blood cells that do not have prominent granules when viewed under a microscope. The group mainly includes lymphocytes and monocytes, which help defend the body against infection.

\medterm{Agranulocytosis} A serious condition in which the blood has an extremely low number of granulocytes, especially neutrophils, white blood cells that fight infection. It may be caused by medicines, chemotherapy, infection, autoimmune disease, or bone-marrow problems and can cause fever, sore throat, and severe infections.

\medterm{Agraphesthesia} An inability to recognize a number or letter traced on the skin with the eyes closed, despite normal basic touch sensation. It may indicate damage to the opposite parietal lobe of the brain.

\medterm{Agraphia} Loss of a previously learned ability to write because of brain damage or disease. It may follow a stroke, head injury, brain tumour, or neurodegenerative disease and may occur with problems speaking or reading.

\medterm{Aicardi Syndrome} A rare disorder affecting mainly girls, usually involving abnormal or absent tissue connecting the brain’s two halves, distinctive retinal defects, and seizures beginning in infancy. Other brain and eye abnormalities may also occur.

\medterm{Aicardi-Goutieres Syndrome} A rare inherited disorder in which abnormal immune activity causes inflammation and damage mainly in the brain, skin, and immune system. It may begin in infancy with developmental regression, seizures, brain changes, and painful skin lesions.

\medterm{AIDS} The most advanced stage of HIV infection, when the immune system is severely weakened. AIDS is diagnosed when the CD4 count is below 200 cells/mm$^3$ or certain serious infections or cancers develop.

\synonyms
Acquired Immunodeficiency Syndrome; Stage 3 HIV Infection

\medterm{AIDS Dementia} A severe form of HIV-associated brain disease occurring mainly in advanced HIV infection. It may cause worsening memory and thinking, slowed movement, difficulty walking, and changes in behaviour.

\medterm{AIDS-Associated Nephropathy} Kidney disease caused by HIV, usually in untreated or poorly controlled infection. It often causes large amounts of protein in the urine and may rapidly worsen kidney function, leading to kidney failure.

\synonyms
HIV-Associated Nephropathy

\medterm{Ainhum} A rare disorder in which a painful fibrous band slowly tightens around a toe, usually the little toe. The band can damage bone and blood supply and may eventually cause loss of the affected part.

\synonyms
Dactylolysis Spontanea

\medterm{Air Abrasion} A dental technique that uses compressed air to spray fine abrasive particles onto a tooth. It can remove limited decay, stains, or old filling material and prepare the tooth for a restoration.

\medterm{Air Bag} An inflatable vehicle safety device that rapidly expands during a crash to cushion impact and reduce serious injury. It works with, but does not replace, a properly worn seat belt.

\medterm{Air Bronchogram} A branching pattern of air-filled airways seen within an area of lung that has lost its normal air, usually because the air sacs contain fluid, pus, blood, or cells. It is a sign on a chest X-ray or CT scan.

\medterm{Air Contaminants} Harmful substances in indoor or outdoor air, including smoke, dust, gases, chemicals, allergens, and germs. Health effects depend on the substance, amount breathed, duration of exposure, and a person’s health.

\medterm{Air Embolism} A blockage of blood flow caused by air or another gas bubble entering a blood vessel. It can occur during medical procedures, injury, or rapid pressure changes and may cause sudden breathing, heart, or brain problems.

\medterm{Air Passages} The connected tubes that carry air from the nose and mouth to the lungs. They include the throat, voice box, windpipe, and branching airways; swelling, mucus, or an object can narrow or block them.

\medterm{Air Polishing} A dental procedure that directs air, water, and a fine powder onto teeth to remove plaque and surface stains. Common powders include glycine, erythritol, and sodium bicarbonate.

\medterm{Air Pollution} Air contaminated with harmful gases, particles, smoke, chemicals, or germs. Breathing polluted air can irritate the eyes and lungs and increase the risk of asthma attacks, infections, heart disease, stroke, and other health problems.

\medterm{Air Quality} A measure of how clean or polluted the air is, based on the amounts of harmful substances it contains. Poor air quality can affect breathing and heart health, especially in people with existing medical conditions.

\medterm{Air Sac} A small, thin-walled space that contains air. In the human lung, an air sac is called an alveolus and is where oxygen enters the blood and carbon dioxide leaves it; birds also have larger air sacs that help breathing.

\synonyms
Air Space

\medterm{Air Trapping} A condition in which extra air remains in the lungs after breathing out because narrowed or damaged airways prevent complete exhalation. It commonly occurs in asthma, chronic obstructive pulmonary disease, and bronchiolitis.

\medterm{Air-Conditioner Lung} Lung inflammation caused by an immune reaction to mold, fungi, bacteria, or other particles released from contaminated air-conditioning or humidifying systems. Symptoms may include cough, shortness of breath, fever, and tiredness.

\medterm{Air-Sickness} Motion sickness caused by the movement of an aircraft. Mismatched signals from the eyes and inner ear can cause nausea, dizziness, sweating, headache, and sometimes vomiting.

\medterm{Airborne Disease} An infection that can spread when a person breathes in tiny particles containing germs that remain in the air. The risk depends on the germ, ventilation, distance, and length of exposure.

\medterm{Airborne Precautions} Measures used to prevent infections that remain suspended in the air from spreading. They include a special single-patient room with filtered exhaust air, limiting movement outside the room, and a properly fitted respirator such as an N95. They are used for tuberculosis, measles, and chickenpox.

\medterm{Airplane Ear} Ear pain, fullness, popping, or temporary hearing loss caused by unequal air pressure on the two sides of the eardrum during takeoff or landing. It is more likely with a blocked nose or ear infection.

\synonyms
Barotitis Media

\medterm{Airway} Any passage that carries air between the outside of the body and the lungs, including the nose, mouth, throat, voice box, windpipe, and branching tubes. Swelling, mucus, or an object can narrow or block an airway and restrict breathing.

\medterm{Airway Clearance Techniques} Methods used to loosen and remove mucus from the breathing tubes. They include controlled coughing, breathing exercises, drainage positions, devices that create resistance during exhalation, and vibrating vests. The choice depends on the illness and the person’s strength and ability.

\medterm{Airway Management} The methods and equipment used to keep a person’s airway open and support breathing when airflow is threatened. They may include positioning, suction, airway devices, or a breathing tube connected to a ventilator.

\medterm{Airway Obstruction} Partial or complete blockage of airflow through the breathing passages. Causes include an inhaled object, mucus, swelling, injury, tumour, or relaxed airway muscles. It may cause noisy breathing, wheezing, low oxygen, severe breathing difficulty, or respiratory arrest.

\medterm{Airway Resistance} The difficulty air encounters while moving through the breathing tubes. Narrowing from muscle tightening, swelling, mucus, or loss of support increases resistance and makes the breathing muscles work harder.

\medterm{Akathisia} A movement disorder causing an uncomfortable inner restlessness and a strong urge to move, often with pacing, rocking, or repeated shifting of the legs. It is commonly a side effect of antipsychotic or other medicines.

\medterm{Akinesia} An inability to initiate or perform voluntary movement, or a marked lack of spontaneous movement. It may occur in Parkinson disease and related disorders, severe depression, or other conditions affecting the brain’s movement-control systems.

\medterm{Akinetic} Describing a person or condition with little or no voluntary movement. It may result from a neurological disease, medicine effect, severe illness, or damage to the brain’s movement-control systems.

\medterm{Akinetic Mutism} A state in which a person is awake but makes little or no voluntary movement or speech. It usually results from severe damage or dysfunction in brain areas that control motivation, attention, and movement.

\medterm{Akinetopsia} A rare neuropsychological visual deficit characterized by an inability to perceive visual motion, resulting in moving objects appearing as a succession of static freeze-frames. It typically results from bilateral cortical lesions in the middle temporal visual area (area MT/V5).

\synonyms
Motion Blindness; Visual Motion Agnosia

\medterm{Alagille Syndrome} A genetic disorder in which too few or abnormal bile ducts can cause liver disease and jaundice. It may also affect the heart, kidneys, eyes, blood vessels, bones, and facial features. Most cases involve changes in the JAG1 gene.

\medterm{Alanine} An amino acid, or protein-building molecule, made by the body and obtained from food. Cells use it to make proteins and to help move nitrogen and produce glucose for energy.

\medterm{Alanine Aminotransferase} An enzyme found mainly inside liver cells. When liver cells are injured, alanine aminotransferase can leak into the blood, causing a raised test result. The result helps assess liver injury but does not identify its cause by itself.

\synonyms
ALT

\medterm{Alanine Transaminase Blood Test} A blood test that measures alanine aminotransferase, an enzyme released when liver cells are damaged. A high result may occur with hepatitis, medicine-related injury, fatty liver disease, or other conditions and must be interpreted with other information.

\medterm{ALARA Principle} A foundational radiation safety standard mandating that occupational and medical ionizing radiation exposures be kept as low as reasonably achievable, utilizing dose optimization, exposure time reduction, increased distance, and protective shielding.

\synonyms
As Low As Reasonably Achievable; ALARA; Radiation Safety Principle

\medterm{Alarmin} A signal released by injured or severely stressed cells that alerts the immune system. It helps start inflammation and attract protective cells to damaged tissue, even before the body identifies whether infection is present.

\medterm{Albendazole} An antiparasitic medicine that damages the internal structures parasites need to absorb nutrients and survive. It treats several worm infections, including some intestinal and tissue infections, when prescribed for the appropriate parasite.

\medterm{Albers-SchöNberg Disease} A usually inherited form of osteopetrosis in which bones become unusually dense because old bone is not removed normally. The bones can still break easily and may cause pain or pressure on nerves, especially in the skull.
\synonyms
Autosomal Dominant Osteopetrosis Type II (ADO II); Marble Bone Disease

\medterm{Albinism} A group of inherited genetic conditions in which the body produces little or no melanin, the pigment responsible for skin, hair, and eye coloration. Characterized by pale skin, white or light blond hair, and light-colored irises, it frequently causes visual impairments—including reduced visual sharpness, rapid involuntary eye movements (nystagmus), extreme sensitivity to bright light (photophobia), and abnormal optic nerve development. People with albinism have an elevated risk of severe sunburn and skin cancer, requiring strict sun protection.

\synonyms
Oculocutaneous Albinism; OCA; Congenital Achromia; Hypopigmentation Disorder

\medterm{Albino} An older noun for a person or animal with albinism, an inherited condition involving reduced melanin. For people, “person with albinism” is usually preferred because it describes the condition without defining the person by it.

\medterm{Albright Hereditary Osteodystrophy} A genetic condition that may cause short stature, a round face, obesity, dental problems, extra bone under the skin, and shortening of some finger bones. It is linked to changes in the \textit{GNAS} gene and may occur with hormone-resistance disorders.

\synonyms
AHO; Pseudohypoparathyroidism Type 1A Phenotype

\medterm{Albright's Hereditary Osteodystrophy} A group of inherited features that may include short stature, a round face, shortened fingers or toes, extra bone under the skin, and obesity. Some people also have resistance to parathyroid hormone or other hormones.

\medterm{Albumen} The white part of an egg, made mostly of water and proteins. Albumen is a food substance and is different from albumin, the main protein in blood plasma.

\medterm{Albumin} The main protein in blood plasma, made by the liver. It helps keep fluid inside blood vessels and carries substances such as hormones, fatty acids, and some medicines through the bloodstream.

\synonyms
Serum Albumin

\medterm{Albumin Clearance} A measure of how quickly the blood protein albumin leaves the bloodstream or another body space. It is used mainly in specialized kidney, liver, or research tests to study leakage, transport, or removal of the protein.

\medterm{Albumin-To-Creatinine Ratio} A urine test that estimates how much albumin, a blood protein, is leaking through the kidneys by comparing it with creatinine in the same sample. A persistent result above 30 mg/g may indicate kidney damage; results above 300 mg/g indicate severe leakage.

\medterm{Albuminocytologic Dissociation} A characteristic cerebrospinal fluid finding marked by an elevated protein concentration in the absence of a commensurate increase in white blood cell count. Classically seen in Guillain-Barr'{e} syndrome after the first week of illness, as well as in chronic inflammatory demyelinating polyneuropathy and spinal cord tumors.

\synonyms
Albumino-Cytological Dissociation; Protein-Cell Dissociation

\medterm{Albuminuria} An abnormally high amount of albumin in the urine, showing that the kidneys are allowing this blood protein to pass through. It may be temporary or persistent; ongoing albuminuria can indicate kidney damage and increased cardiovascular risk.

\medterm{ALCAPA Syndrome} A rare congenital cardiovascular anomaly in which the left coronary artery arises abnormally from the pulmonary artery instead of the aorta. Because the pulmonary artery carries desaturated blood at low pressure, myocardial perfusion is compromised, leading to retrograde coronary flow (coronary steal), infantile heart failure, myocardial infarction, and mitral regurgitation if uncorrected.

\synonyms
ALCAPA; Bland-White-Garland Syndrome; Anomalous Left Coronary Artery from the Pulmonary Artery

\medterm{Alcohol} Ethanol, a substance in beer, wine, and spirits that affects the central nervous system. Excessive consumption impairs motor coordination and cognitive judgment, produces acute intoxication, and can lead to physical dependence, metabolic toxicity, and chronic organ damage.

\medterm{Alcohol Dehydrogenase} An enzyme found mainly in the liver that begins breaking down alcohol. It changes alcohol into acetaldehyde, a harmful intermediate that is then processed into less toxic substances.

\synonyms
ADH; Ethanol Dehydrogenase

\medterm{Alcohol Intolerance} An unpleasant reaction after drinking alcohol, often caused by inherited difficulty breaking down alcohol or sensitivity to an ingredient in the drink. Symptoms may include flushing, headache, nausea, rapid heartbeat, or worsening asthma.

\medterm{Alcohol Intoxication} Temporary impairment caused by drinking alcohol. It can cause poor judgment, confusion, slurred speech, loss of coordination, vomiting, reduced consciousness, slowed breathing, or coma.

\synonyms
Alcoholic Intoxication

\medterm{Alcohol Poisoning} A life-threatening overdose of alcohol that severely depresses brain function and can cause vomiting, inability to wake, seizures, slow or irregular breathing, low body temperature, coma, or death.

\medterm{Alcohol Swab} A single-use pad moistened with alcohol, usually used to clean intact skin before an injection or other procedure. The skin is allowed to dry before the procedure.

\medterm{Alcohol Use Disorder} A medical condition in which repeated alcohol use causes difficulty controlling drinking, strong cravings, or problems at home, work, school, or in relationships. Symptoms may be mild, moderate, or severe.

\synonyms
Alcohol Abuse
Alcohol Abuse and Alcoholism
Alcohol Addiction
Alcohol Dependence
Alcohol Dependence (Alcohol Abuse and Alcoholism)
Alcohol Dependency
Alcoholism

\medterm{Alcohol Withdrawal} Symptoms that can occur after a person who drinks heavily for a long time suddenly stops or greatly reduces alcohol. Symptoms include shaking, sweating, anxiety, nausea, hallucinations, seizures, and delirium; severe withdrawal can be life-threatening.

\synonyms
Alcohol Withdrawal Syndrome

\medterm{Alcohol-Associated Liver Disease} Liver injury associated with alcohol use. It ranges from steatosis and inflammation to fibrosis, cirrhosis, liver failure, or liver cancer; severity depends on alcohol exposure, genetics, nutrition, and other health factors.

\medterm{Alcohol-Related Psychosis} Hallucinations, delusions, or other psychotic symptoms associated with alcohol use, intoxication, withdrawal, or a coexisting alcohol-related disorder. Other medical and psychiatric causes must be excluded.

\medterm{Alcoholic Cardiomyopathy} A disease of the heart muscle caused by long-term heavy alcohol use. The heart becomes enlarged and weak, reducing its ability to pump blood and potentially causing heart failure or abnormal heart rhythms.

\medterm{Alcoholic Foetor} A characteristic odor on the breath and body resulting from recent ethanol consumption and the pulmonary excretion of volatile alcohol and its primary metabolite, acetaldehyde. It serves as an immediate clinical sign of acute alcohol exposure or intoxication during physical assessment.
\synonyms{Alcoholic Fetor; Alcohol Breath; Fetor Alcoholicus}

\medterm{Alcoholic Hepatitis} An acute inflammation and injury of the liver caused by heavy alcohol use. It commonly causes jaundice, fatigue, fever, nausea, and abdominal pain and can become severe enough to cause liver failure.

\synonyms
Alcohol-Associated Hepatitis; Alcohol-Induced Hepatitis

\medterm{Alcoholic Ketoacidosis} A dangerous buildup of acids in the blood that can occur after prolonged heavy alcohol use, poor food intake, and repeated vomiting. It may cause abdominal pain, nausea, dehydration, rapid breathing, and a normal or low blood-sugar level.

\medterm{Alcoholic Liver Disease} Alternate name; see \textit{Alcohol-Associated Liver Disease}.

\medterm{Alcoholic Myopathy} Muscle damage caused by heavy alcohol use. It may cause muscle pain, tenderness, weakness, or wasting; severe acute injury can destroy muscle tissue and damage the kidneys.

\synonyms
Alcohol-Induced Myopathy; Ethanol-Induced Myopathy

\medterm{Alcoholic Neuropathy} Nerve damage associated with long-term heavy alcohol use and related nutritional deficiency. It commonly causes burning pain, numbness, tingling, or weakness, especially in the feet and lower legs.

\medterm{Alcuronium} A medicine that temporarily relaxes or paralyzes skeletal muscles during anesthesia. Because it can stop the muscles needed for breathing, it is given with controlled airway and breathing support.

\medterm{Aldolase Blood Test} A blood test that measures aldolase, an enzyme found mainly in muscle and liver. A high result may occur with muscle injury or disease, but it is not specific and is interpreted with symptoms and other tests.

\medterm{Aldosterone} A hormone made by the adrenal glands that helps the kidneys retain sodium and water and remove potassium. It helps control blood volume and blood pressure.

\medterm{Aldosterone Blood Test} A blood test that measures the hormone aldosterone. It helps investigate high or low blood pressure and abnormal potassium levels and is usually interpreted together with a renin test.

\medterm{Aldosterone Escape} A physiological phenomenon in primary hyperaldosteronism where persistent mineralocorticoid excess causes initial sodium and water retention, but subsequent increased atrial natriuretic peptide secretion and pressure natriuresis prevent progressive hypervolemia and severe edema.

\medterm{Aldosterone-to-Renin Ratio} A clinical screening blood test used in the evaluation of secondary hypertension to detect primary aldosteronism (Conn syndrome). A significantly elevated ratio of plasma aldosterone concentration to plasma renin activity prompts confirmatory testing.

\synonyms
ARR

\medterm{Aleukemic Leukemia} A presentation of acute or chronic leukemia in which malignant blasts or abnormal leukocytes are absent or present in extremely low numbers in the peripheral blood smear, despite extensive malignant infiltration of the bone marrow.

\medterm{Alexander Disease} A rare genetic disorder in which abnormal GFAP protein damages the brain’s supporting cells and white matter. It may begin in infancy, childhood, or adulthood and can cause enlarged head size, developmental problems, movement difficulty, seizures, or progressive weakness.

\synonyms
Fibrinoid Leukodystrophy; GFAP Leukodystrophy

\medterm{Alexia} Loss of the ability to understand written words despite adequate eyesight, usually because of brain damage from a stroke, injury, tumour, or degenerative disease. It may occur with other language or reading problems.

\medterm{Alexia Without Agraphia} A classical neurological disconnection syndrome in which a patient can write spontaneously and to dictation but is completely unable to read printed or handwritten text (including their own writing). Caused by an infarction of the left posterior cerebral artery affecting the left visual cortex and the splenium of the corpus callosum, isolating the intact right visual cortex from the language centers of the left hemisphere.

\synonyms
Pure Alexia; Dejerine Syndrome; Pure Word Blindness

\medterm{Alexithymia} Difficulty recognizing, understanding, or describing one’s own emotions. It is a trait or symptom rather than a diagnosis and may occur with autism, depression, trauma, stress, or other conditions.

\medterm{Alezzandrini Syndrome} An exceptionally rare acquired disorder of unilateral melanocyte destruction involving ipsilateral facial vitiligo and poliosis with retinal or visual abnormalities and sometimes hearing loss. The cause is uncertain; ophthalmic and audiologic follow-up is important because vision or hearing may progressively worsen.

\medterm{Alfacalcidol} A medicine related to vitamin D that the liver changes into an active form. It increases absorption of calcium and phosphate and is used for selected disorders of bones, kidneys, or mineral balance.

\medterm{Algia} A word ending that means pain. It appears in terms such as neuralgia for nerve pain, myalgia for muscle pain, and arthralgia for joint pain; it does not identify the cause.

\medterm{Alglucerase} An older enzyme-replacement medicine for Gaucher disease, a disorder in which fatty substances build up in cells. It replaces an enzyme the body lacks; newer enzyme preparations are used more often.

\medterm{Algor Mortis} The gradual cooling of the body after death toward the surrounding temperature. The amount of cooling can help estimate the time since death, but it is affected by the environment and other conditions.

\medterm{Aliasing Artifact} An ultrasound error in which blood flow moving faster than the selected display range appears to change direction or speed. It reflects a limitation of the image display, not necessarily abnormal blood flow.

\medterm{Alien Hand Syndrome} A neurological condition in which an upper limb performs autonomous, purposeful-appearing actions without conscious voluntary control, often accompanied by the patient's sensation that the limb is foreign or acting independently. It is associated with lesions of the corpus callosum, supplementary motor area, or medial frontal cortex.

\synonyms
Alien Limb Phenomenon; Anarchic Hand

\medterm{Alimentary} Relating to eating, digestion, or the movement of food through the body. In medical language, it usually refers to the digestive tract and the organs that help break down and absorb food.

\medterm{Alimentary Canal} The continuous muscular tube through which food passes from the mouth to the anus. It includes the throat, oesophagus, stomach, small intestine, and large intestine.

\medterm{Alkalemia} An abnormally elevated arterial blood pH (greater than 7.45), indicating a decreased hydrogen ion concentration in circulating blood. It represents the net systemic state caused by underlying metabolic or respiratory alkalosis.

\medterm{Alkali} A basic substance that can neutralize acids. Strong alkalis, such as some cleaning chemicals, can burn the skin, eyes, and digestive tract, while weaker alkalis have food, laboratory, or medical uses.

\medterm{Alkaline Diet} A diet based mainly on foods claimed to make the body more alkaline. The lungs and kidneys tightly control blood acidity, so food does not make the blood alkaline, although many such diets include nutritious foods.

\medterm{Alkaline Phosphatase} An enzyme found mainly in the liver, bile ducts, bones, intestine, and placenta. A high blood level may reflect liver or bile-duct disease, increased bone activity, pregnancy, or another condition and must be interpreted with other tests.

\medterm{Alkaline Phosphatase Test} A quantitative serum laboratory assay measuring alkaline phosphatase activity to evaluate suspected cholestatic liver disease, biliary obstruction, or increased bone turnover such as Paget disease or bone metastases.

\synonyms
ALP Blood Test; Serum ALP; Alkaline Phosphatase Level

\medterm{Alkaloid} A naturally occurring chemical containing nitrogen, often made by plants, that can strongly affect the body. Examples include morphine, quinine, nicotine, and atropine.

\synonyms
Plant Alkaloid; Basic Nitrogenous Organic Compound

\medterm{Alkalosis} A condition in which the blood becomes too alkaline. It may result from repeated vomiting, certain medicines, hormone disorders, or breathing too rapidly, and can cause muscle cramps, tingling, confusion, or abnormal heart rhythms.

\medterm{Alkaptonuria} A rare inherited metabolic disorder caused by reduced activity of homogentisate 1,2-dioxygenase, usually because of disease-causing variants in the \textit{HGD} gene. Homogentisic acid accumulates, darkens urine, and can deposit in cartilage, tendons, skin, and other tissues in a process called ochronosis. Over time it may cause arthritis and stiffness of the spine and large joints, urinary stones, or changes involving heart valves.

\medterm{Alkylating Agents} Drugs or chemicals that damage DNA by attaching chemical groups to it. Some are used to treat cancer by stopping cells from multiplying, but they can also injure healthy cells and cause side effects.

\medterm{Allantoin} A compound used in some skin products to soften, moisturize, and protect irritated or rough skin. Its effect and safety depend on the product and how it is used.

\medterm{Allantois} A temporary structure in the early embryo that helps with blood-vessel development and contributes to the umbilical cord and part of the urinary system.

\medterm{Allele} One of two or more versions of a gene. A person usually inherits one allele from each parent; different alleles can influence traits or the risk of inherited disease.

\medterm{Allelic} Relating to an allele, one of the different versions of a gene inherited from a parent.

\medterm{Allen Test} A bedside test that checks whether blood can reach the hand through both of its main arteries. It may be done before taking blood from or placing a catheter in the radial artery.

\medterm{Allergen} A usually harmless substance that triggers an allergic immune response in a sensitive person. Common allergens include foods, pollen, medicines, animal proteins, insect venom, and natural rubber.

\medterm{Allergen Immunotherapy} Treatment that gives gradually increasing amounts of a known allergen by injections or under the tongue. Over time, it can reduce allergic symptoms for selected allergies, but it can rarely cause a serious reaction.

\medterm{Allergic and Environmental Asthma} Asthma triggered or worsened by allergens or environmental exposures such as dust mites, molds, pollens, animals, smoke, air pollution, or workplace agents. Management combines exposure reduction, controller therapy, trigger planning, and evaluation for occupational or allergic disease when indicated.

\medterm{Allergic Bronchopulmonary Aspergillosis} An allergic reaction to the fungus Aspergillus in the airways, occurring mainly in people with asthma or cystic fibrosis. It can cause worsening breathing symptoms, mucus plugs, airway widening, and high levels of allergy-related cells or antibodies.

\medterm{Allergic Cascade} The chain of immune events that follows exposure to an allergen. Immune cells release substances such as histamine, which can cause itching, swelling, mucus, wheezing, hives, or anaphylaxis.

\medterm{Allergic Conjunctivitis} Alternate name; see \textit{Eye Allergies}.

\medterm{Allergic Contact Dermatitis} A delayed type-IV hypersensitivity reaction after skin contact with an allergen to which the immune system has been sensitized. It causes itching, redness, swelling, blisters, scaling, or chronic fissuring, often extending beyond the contact area; patch testing can help identify the trigger.

\medterm{Allergic Dermatitis} Skin inflammation caused by an immune reaction to an allergen. It commonly causes itching, redness, swelling, blisters, or scaling and may occur after contact with a triggering substance.

\medterm{Allergic Diseases} Conditions in which the immune system reacts abnormally to a usually harmless substance called an allergen, such as pollen, dust mites, mold, food proteins, medicines, insect venom, or animal proteins. Some reactions involve immunoglobulin E (IgE) antibodies, while others do not. Allergic diseases may affect the nose, eyes, skin, lungs, or digestive system and can cause sneezing, itching, hives, swelling, vomiting, wheezing, or anaphylaxis. Examples include allergic rhinitis, food allergy, drug allergy, allergic conjunctivitis, atopic dermatitis, urticaria, and allergic asthma. Food intolerance is different because it does not involve an immune reaction.

\medterm{Allergic Proctocolitis} A benign, non-IgE-mediated food protein-induced inflammatory disorder of the distal colon and rectum, seen predominantly in young infants exposed to cow's milk or soy protein in formula or breast milk, classically causing blood- and mucus-streaked stools in an otherwise thriving baby.

\synonyms
Food Protein-Induced Allergic Proctocolitis; FPIP

\medterm{Allergic Reaction} An immune response to a usually harmless substance. It may cause sneezing, itching, hives, swelling, vomiting, or breathing difficulty and can sometimes become life-threatening anaphylaxis.

\synonyms
Allergy

\synonyms
Allergic Reactions

\medterm{Allergic Rhinitis} Inflammation inside the nose caused by an allergic reaction to inhaled substances such as pollen, dust mites, mold, or animal proteins. It causes sneezing, itching, a blocked or runny nose, and sometimes itchy, watery eyes.

\synonyms
Hay Fever; Pollinosis; Allergic Rhinoconjunctivitis

\medterm{Allergic Urticaria} An acute, IgE-mediated type I hypersensitivity skin reaction triggered by specific external allergens such as foods, drugs, insect venoms, or latex. Cross-linking of allergen-specific IgE on the surface of dermal mast cells triggers rapid degranulation and release of histamine, prostaglandins, and leukotrienes, causing localized vasodilation and increased vascular permeability that manifests as intensely pruritic, transient, erythematous wheals that resolve within 24 hours.
\synonyms{Acute Allergic Urticaria; IgE-Mediated Urticaria; Allergic Hives}

\medterm{Allergy} An immune response to a usually harmless substance such as pollen, food, medicine, insect venom, or animal proteins. Symptoms range from sneezing, itching, and hives to swelling, breathing difficulty, or life-threatening anaphylaxis.

\synonyms
Allergic Hypersensitivity; Type I Hypersensitivity

\medterm{Allergy Shots} Injections used for allergen immunotherapy. They contain gradually increasing amounts of a specific allergen to reduce sensitivity over time and are given in a setting where a serious allergic reaction can be recognized and treated.

\medterm{Allergy Skin Test} An in vivo diagnostic procedure in which small amounts of suspected allergens are introduced into or onto the skin via prick, puncture, scratch, or intradermal injection to assess immediate IgE-mediated hypersensitivity through localized wheal-and-flare reactions.

\synonyms
Skin Prick Test; Epicutaneous Test; Allergy Skin Testing; Scratch Test

\medterm{Allergy Test} A skin or blood test used to look for sensitivity to a suspected allergen. A positive result shows sensitization but does not always prove that exposure causes symptoms, so the result must be considered with the person’s history.

\synonyms
Skin Prick Test; Specific IgE Test; Allergen Challenge

\medterm{Allis Clamp} A surgical instrument with small teeth used to hold firm tissue during an operation. It provides a strong grip but can injure delicate tissue if applied with too much pressure.

\medterm{Allocheiria} A neurological sign in which a person feels or identifies a touch on the opposite side of the body from where it was applied. It reflects impaired localization of sensation and may occur with disorders affecting the parietal lobe.

\synonyms
Allochiria

\medterm{Allodynia} Pain caused by something that normally should not hurt, such as light touch, clothing, or gentle pressure. It can occur with nerve injury, migraine, shingles, fibromyalgia, or other pain conditions.

\medterm{Allogeneic} Coming from another person of the same species who is genetically different, as in an allogeneic blood, tissue, stem-cell, or organ transplant.

\medterm{Allograft} An organ, tissue, or group of cells transplanted from one genetically different individual to another individual of the same species. The recipient’s immune system may reject it, so immunosuppressive medicines may be needed.

\medterm{Allosteric Regulation} Control of a protein or enzyme when another molecule attaches away from its main working site. The attachment changes the protein’s shape and can increase or decrease how well it works.

\medterm{Aloe} A plant used in some skin products, medicines, and supplements. Aloe gel applied to skin may soothe minor irritation, while products containing aloe latex taken by mouth can cause cramps, diarrhoea, and dangerous fluid or salt loss.

\medterm{Aloe Vera} A plant whose clear inner-leaf gel is used in some skin products and supplements. The gel may soothe minor irritation, while oral latex or whole-leaf products can cause diarrhoea, fluid or salt loss, medicine interactions, and toxicity.

\medterm{Alogia} Markedly reduced speech or speech that conveys little information. It may occur as a negative symptom of schizophrenia and can also be associated with depression, brain disease, or other conditions.

\medterm{Alongshan Virus} A virus spread by some ticks and reported mainly in parts of Asia and Europe. It may cause fever, headache, fatigue, nausea, or neurological symptoms, but its ability to cause human disease is still being studied.

\medterm{Alopecia} Hair loss from the scalp or another part of the body. It may be nonscarring, in which the hair follicles remain and hair may regrow, or scarring, in which inflammation or injury destroys the follicles and can cause permanent hair loss. Causes include inherited pattern hair loss, autoimmune disease, hormonal changes, illness, medicines, nutritional problems, infection, repeated pulling, or other injury. The amount, pattern, and speed of hair loss vary with the cause.

\medterm{Alopecia Areata} An autoimmune, usually nonscarring hair-loss disorder in which immune cells attack anagen hair follicles. It often causes sharply demarcated smooth patches with exclamation-mark hairs and may involve the scalp, beard, eyebrows, or nails; regrowth is possible but relapses and extensive forms can occur.

\synonyms
AA; Patchy Hair Loss; Autoimmune Alopecia

\medterm{Alopecia Mucinosa (Follicular Mucinosis)} A skin disorder in which mucin accumulates within hair follicles and sebaceous glands, causing itchy follicular papules or plaques with hair loss. It may be a primary self-limited condition or a cutaneous sign of lymphoma, especially mycosis fungoides, so persistent or widespread disease needs evaluation.

\medterm{Alopecia Totalis} Complete loss of hair from the scalp, usually caused by alopecia areata. Hair may regrow, remain absent, or fall out again, and other autoimmune conditions may occur.

\medterm{Alopecia Universalis} Complete or near-complete loss of hair from the scalp and body, usually caused by extensive alopecia areata. Hair may regrow or recur unpredictably.

\medterm{ALP Blood Test} Alternate index label; see \textit{Alkaline Phosphatase Test}.

\medterm{ALP Isoenzyme Test} A blood test that helps identify whether a high alkaline phosphatase level comes mainly from the liver, bones, intestine, or placenta. It can help distinguish liver or bile-duct disease from increased bone activity.

\medterm{Alpha Blockers} Medicines that relax certain blood vessels and smooth muscles by blocking alpha receptors. They may lower blood pressure or improve urine flow in an enlarged prostate, but can cause dizziness or fainting.

\medterm{Alpha Cells} Cells in the pancreas that make glucagon. Glucagon raises blood glucose when it falls too low or when the body needs energy between meals.

\synonyms
Pancreatic Alpha Cells; Glucagon-Secreting Cells

\medterm{Alpha Hydroxy Acids} A group of acids used in skin products to loosen the outer layer of skin and improve roughness, fine lines, or uneven colour. Common examples are glycolic acid and lactic acid; they can irritate skin and increase sun sensitivity.

\synonyms
AHAs; Fruit Acids

\medterm{Alpha Waves} A regular pattern of electrical activity in the brain, usually eight to thirteen cycles per second. It is strongest when a relaxed, awake person has closed eyes and decreases with attention or eye opening.

\synonyms
Alpha Rhythm; Posterior Dominant Rhythm

\medterm{Alpha-1 Adrenergic Antagonist} A medicine or other substance that blocks alpha-1 adrenergic receptors. This reduces selected effects of noradrenaline on vascular smooth muscle and other tissues; the precise effects depend on the receptor distribution and the substance used.

\synonyms
Alpha-1 Adrenergic Blocker

\medterm{Alpha-1 Antitrypsin Blood Test} A blood test that measures alpha-1 antitrypsin, a protein made by the liver that helps protect lung tissue. A low level may suggest inherited alpha-1 antitrypsin deficiency, which can affect the lungs and liver.

\medterm{Alpha-1 Antitrypsin Deficiency} An inherited condition in which low or abnormal alpha-1 antitrypsin increases the risk of emphysema, especially at a young age, and liver disease. Some people also develop inflammation of fatty tissue under the skin.

\synonyms
AATD; AAT Deficiency; Alpha-1 Deficiency

\medterm{Alpha-Agonists} Medicines that activate alpha receptors on cells. Depending on the medicine and tissue, they can narrow blood vessels, raise blood pressure, reduce eye fluid, or affect the nervous system.

\medterm{Alpha-Fetoprotein} A protein made mainly by a developing fetus. Measuring it in maternal blood can help screen for some fetal conditions, while measuring it in adults can help assess selected liver or germ-cell tumours. An abnormal result does not establish a diagnosis.

\medterm{Alpha-Gal Syndrome} An allergy to a sugar found in most mammals but not humans or other primates. It can cause delayed hives, stomach symptoms, or anaphylaxis after eating mammalian meat or using some animal-derived products.

\synonyms
Alpha-Gal Allergy

\medterm{Alpha-Globulin} One of several groups of proteins in the liquid part of blood. These proteins carry substances and help with inflammation, clotting, and immune defence; their levels may change with liver disease, inflammation, or other conditions.

\medterm{Alpha-Glucosidase Inhibitors} Diabetes medicines that slow the digestion and absorption of carbohydrates in the small intestine. They mainly reduce the rise in blood glucose after meals and may cause gas, bloating, or diarrhoea.

\medterm{Alpha-Ketoglutarate} A substance made during the breakdown of nutrients for energy. It also helps the body process amino acids and supports several chemical reactions inside cells.

\medterm{Alpha-Synuclein} A protein found mainly in nerve cells that helps regulate nerve-cell functions. Abnormal forms can collect into clumps in Parkinson disease, Lewy-body dementia, and related disorders, although the clumps are not the only cause of these diseases.

\medterm{Alpha-Thalassemia} An inherited blood disorder, usually autosomal recessive, in which the body makes too little alpha-globin, a part of hemoglobin that carries oxygen. The condition may cause no symptoms, mild anemia, or severe anemia before or shortly after birth, depending on the number of affected alpha-globin genes.

\synonyms
Alpha-Globin Deficiency; HbH Disease Carrier

\medterm{Alpha1-Antitrypsin (AAT) Deficiency} An inherited deficiency or dysfunction of alpha-1 antitrypsin that predisposes to early emphysema, especially with smoking, and sometimes liver disease from abnormal protein accumulation. Testing is considered in early or unexplained COPD, basilar emphysema, or unexplained liver disease.

\medterm{Alport Syndrome} An inherited disorder of collagen that can cause progressive kidney disease, hearing loss, and changes in the eyes. It may be caused by changes in the \textit{COL4A3}, \textit{COL4A4}, or \textit{COL4A5} genes and may follow X-linked, autosomal dominant, or autosomal recessive inheritance.

\medterm{Alprazolam} A benzodiazepine medicine used for anxiety and panic disorder. It calms brain activity but can cause sleepiness, impaired coordination, dependence, and dangerous slowing of breathing, especially with opioids or alcohol.

\medterm{Alstrom Syndrome} A rare inherited disorder caused by changes in the ALMS1 gene. It may cause progressive vision and hearing loss, obesity, diabetes or insulin resistance, heart disease, and kidney or liver problems.

\synonyms
Alström Syndrome

\medterm{Alternating Hemiplegia Of Childhood} A rare neurological disorder beginning in infancy or early childhood, causing recurrent episodes of temporary paralysis that may affect either side of the body. Episodes may also involve abnormal movements, eye-movement problems, breathing difficulty, seizures, and developmental impairment; symptoms typically disappear during sleep.

\medterm{Altitude Illness} Illness caused by rapid ascent to high altitude without enough time to adjust. It includes acute mountain sickness, high-altitude brain swelling, and high-altitude lung swelling; severe forms can be fatal.

\synonyms
Altitude Sickness

\medterm{Aluminum Phosphide} A highly poisonous pesticide that releases phosphine gas when it contacts moisture or stomach acid. Poisoning can rapidly damage the heart, lungs, digestive system, and other organs and requires emergency treatment.

\medterm{Aluminum Toxicity} Toxicity from excessive aluminum exposure or accumulation, especially in advanced kidney disease or with contaminated dialysis products. Chronic exposure can cause encephalopathy, speech or movement changes, microcytic anemia, and bone pain or osteomalacia; evaluation includes exposure review, kidney function, and specialist-directed aluminum testing.

\medterm{Alvarado Score} A clinical score used to estimate the likelihood of acute appendicitis from symptoms, examination findings, and blood-test results. It supports assessment but does not by itself confirm or exclude appendicitis.

\medterm{Alveolar} Relating to an alveolus, especially one of the tiny air sacs in the lungs where oxygen enters the blood and carbon dioxide leaves it.

\medterm{Alveolar Abnormalities} Changes in the lung air sacs that can interfere with oxygen and carbon-dioxide exchange. They may occur with infection, fluid buildup, emphysema, scarring, or a developmental condition.

\medterm{Alveolar Bone} The part of the jawbone that surrounds and supports the roots of the teeth. It can shrink after tooth loss and may need rebuilding before some dental implants.

\medterm{Alveolar Bone Loss} Loss of the bone that supports the teeth, most often caused by periodontitis. It can make teeth loose and may lead to tooth loss; dental treatment can slow further loss.

\medterm{Alveolar Dead Space} The volume of inhaled air that reaches alveoli that are adequately ventilated but lack sufficient pulmonary capillary blood perfusion, preventing gas exchange from taking place.

\medterm{Alveolar Echinococcosis} A rare, serious infection caused by the young form of the Echinococcus multilocularis tapeworm, usually after swallowing contaminated food or water. It creates slowly growing, tumor-like areas in the liver that can invade nearby tissue or spread.

\synonyms
Alveolar Hydatid Disease

\medterm{Alveolar Gas Equation} A formula used to estimate the oxygen level in the lung air sacs from the inhaled oxygen level and carbon-dioxide level. Comparing this estimate with arterial blood oxygen helps assess how well the lungs exchange gases.

\medterm{Alveolar Hypoventilation} Deficient ventilation of the pulmonary alveoli resulting in reduced carbon dioxide elimination, leading to arterial hypercapnia and respiratory acidosis.

\medterm{Alveolar Macrophage} An immune cell that lives in the lung air sacs and removes inhaled dust, microbes, and damaged cells. It also helps regulate inflammation in the lungs.
\synonyms
Dust Cell

\medterm{Alveolar Osteitis} Painful inflammation of a tooth socket after extraction, usually caused by loss or breakdown of the blood clot that normally protects the socket. It is also called dry socket.

\synonyms
Dry Socket; Fibrinolytic Alveolitis

\medterm{Alveolar Proteinosis} Alternate form; see \textit{Pulmonary Alveolar Proteinosis}.

\medterm{Alveolar Rhabdomyosarcoma} A cancer of immature muscle-forming cells, usually affecting children or adolescents. It may arise in the limbs, trunk, head, or neck and can spread early to lymph nodes, lungs, bones, or other tissues.

\medterm{Alveolar Ventilation} The amount of fresh air reaching the lung air sacs each minute and available for oxygen and carbon-dioxide exchange. It depends on breathing rate, breath size, and the amount of air remaining in the breathing passages.

\medterm{Alveolar Ventilation Equation} A relationship showing how breathing affects blood carbon-dioxide levels. When the amount of fresh air reaching the lung air sacs falls, carbon dioxide rises; when effective breathing increases, carbon dioxide falls, if production stays constant.

\medterm{Alveolar-Arterial Oxygen Gradient} The difference between oxygen pressure in the lung air sacs and oxygen pressure in arterial blood. It is calculated as PAO2 minus PaO2. A high value may indicate a problem with oxygen transfer, blood flow through the lungs, or lung tissue. A normal value with low blood oxygen may suggest slow breathing or low oxygen in the surrounding air.

\synonyms
A-A Gradient; A-a Gradient; A-a Oxygen Difference

\medterm{Alveolar-Capillary Barrier} The delicate biological membrane formed by the alveolar type I pneumocyte, capillary endothelial cell, and their intervening fused basement membranes, through which oxygen and carbon dioxide diffuse between alveolar air and capillary blood.

\synonyms
Alveolar-Capillary Membrane; Blood-Air Barrier

\medterm{Alveolitis} Inflammation of the lung air sacs, usually causing cough and shortness of breath, or inflammation of a tooth socket after extraction, called dry socket. The body site identifies the meaning.

\medterm{Alveolus} One of the microscopic air sacs at the ends of the lung airways. Oxygen passes from an alveolus into the blood and carbon dioxide passes from the blood into it.

\synonyms
Alveoli

\medterm{Alzheimer Disease} A progressive disease of the brain and the most common cause of dementia. It slowly damages nerve cells and the connections between them. Early symptoms often include problems with recent memory, finding words, or thinking through tasks. As the disease advances, it can affect language, judgment, orientation, behavior, and the ability to carry out everyday activities. Abnormal proteins called amyloid plaques and tau tangles build up in the brain and are linked to this damage. Alzheimer disease is not a normal part of aging, although the risk increases with age, and its symptoms and rate of progression vary.

\synonyms
Alzheimer's Disease; Alzheimer's
Alzheimer'S
Alzheimer'S Disease

\medterm{Amalgam} A dental filling material made by combining mercury with metals such as silver, tin, and copper. It is strong and durable but has a metallic appearance and may not be suitable for every person or tooth.

\medterm{Amastia} A birth condition in which one or both breasts do not develop. The nipple, chest muscle, or other nearby structures may also be absent or underdeveloped.

\medterm{Amaurosis Fugax} Sudden, temporary loss of vision in one eye, often described as a curtain or shade, caused by briefly reduced blood flow to the retina or optic nerve. It can warn of vascular disease or stroke risk.

\medterm{Ambidextrous} Able to use the left and right hands with similar skill. It describes a pattern of handedness, not a disease.

\medterm{Ambient} Surrounding or present in the immediate environment, such as ambient temperature, light, noise, or air.

\medterm{Ambiguous Genitalia} External genital anatomy that does not have typical male or female features. It usually reflects a difference in sex development, which may involve chromosomes, hormones, or the formation of reproductive organs.

\medterm{Amblyopia} Reduced vision in one eye (rarely both) resulting from abnormal brain visual processing during early childhood, occurring despite a structurally normal eye. When the brain receives conflicting, blurry, or misaligned visual signals, it suppresses input from the weaker eye and favors the dominant eye. Common causes include crossed or misaligned eyes (strabismic amblyopia), unequal refractive focusing power between the eyes (anisometropic amblyopia), or physical blockage of the visual axis (deprivation amblyopia from congenital cataract or severe ptosis). Because neural visual pathways develop rapidly in early life, treatment is time-sensitive and most effective during the critical period before age 7 to 8. Management involves correcting refractive errors with glasses, clearing anatomical obstructions, and patching the stronger eye (or using blurring atropine drops) to force the brain to strengthen neural pathways in the amblyopic eye.

\synonyms
Lazy Eye; Functional Amblyopia; Visual Suppression Amblyopia

\medterm{Ambu Bag} A self-inflating manual resuscitation bag used with a mask or airway device to push air into the lungs when a person is not breathing adequately.

\medterm{Ambulatory Blood Pressure Monitoring} Measurement of blood pressure at regular intervals, usually for twenty-four hours, with a portable cuff worn during normal activities and sleep. It can identify blood-pressure patterns that may not appear during a clinic visit.

\medterm{Amebiasis} An infection caused by the parasite Entamoeba histolytica, usually spread by swallowing contaminated food, water, or material from contaminated hands. It may cause no symptoms, diarrhea, abdominal pain, or bloody diarrhea. In some cases, the infection can spread from the intestine to the liver.

\synonyms
Entamoeba Histolytica Infection; Amoebic Dysentery

\medterm{Amebic Colitis} Inflammation of the large intestine caused by Entamoeba histolytica. It may cause abdominal pain, diarrhoea, bloody stools, or fever.

\synonyms
Amoebic Colitis

\medterm{Amebic Dysentery} An acute, invasive intestinal infection caused by the protozoan parasite Entamoeba histolytica, characterized by extensive ulceration of the colonic mucosa ("flask-shaped" ulcers). Patients typically present with severe lower abdominal cramping, tenesmus, and frequent passage of bloody, mucoid stools, with potential complications including fulminant necrotizing colitis, toxic megacolon, bowel perforation, and secondary amebic liver abscess.

\synonyms
Amoebic Dysentery; Intestinal Amebiasis; Amebic Colitis

\medterm{Amebic Liver Abscess} A collection of infected or damaged tissue in the liver caused by the parasite \textit{Entamoeba histolytica}. It often affects the right side of the liver and may cause fever, pain under the right ribs, and an enlarged or tender liver. Ultrasound or another scan and blood tests can support the diagnosis, but results must be interpreted with the person’s symptoms and history. Treatment usually includes a medicine that kills the parasite followed by a medicine that clears it from the intestine; drainage is needed in selected cases.

\synonyms
Amoebic Liver Abscess; Hepatic Amebiasis; Entamoeba Liver Abscess

\medterm{Ameboma} A chronic, localized inflammatory granulomatous pseudotumor of the bowel wall that develops in a minority of patients with untreated or inadequately treated Entamoeba histolytica infection, most commonly affecting the cecum or ascending colon. It presents as a palpable, tender right-lower-quadrant abdominal mass associated with intermittent bloody diarrhea, obstruction, or weight loss, frequently mimicking colonic adenocarcinoma on colonoscopy and cross-sectional imaging.

\synonyms
Amoeboma; Amebic Granuloma

\medterm{Amelanotic Melanoma} A melanoma that produces little or no dark pigment and may appear pink, red, or skin-coloured instead of black or brown. Its appearance can delay recognition; diagnosis requires examination of a tissue sample.

\medterm{Amelia} A rare birth condition in which one or more arms or legs are completely absent. It may occur alone or with other developmental differences.

\medterm{Amelioration} Improvement or reduction in the severity, frequency, or effect of a symptom, disease, impairment, or other harmful condition.

\medterm{Ameloblastoma} A usually slow-growing tumour that begins in cells involved in tooth development, most often in the lower jaw. It can grow into nearby bone, deform the jaw, and return after treatment.

\medterm{Amelogenesis} The process by which developing teeth form and harden enamel, the tough outer covering of each tooth. Problems with this process can make enamel thin, weak, discoloured, or easily damaged.

\medterm{Amelogenesis Imperfecta} A group of inherited conditions in which tooth enamel does not form normally. The enamel may be thin, soft, rough, discoloured, or easily worn on most or all teeth.

\medterm{Amenorrhea} Absence of menstrual periods. Primary amenorrhea means menstruation has not started by the expected age; secondary amenorrhea means periods stop after they have begun. It may be normal during pregnancy, breastfeeding, or menopause, or may result from hormonal, reproductive, medical, nutritional, or exercise-related conditions.

\medterm{American Trypanosomiasis} Alternate name; see \textit{Chagas disease}.

\medterm{Ametropia} A focusing problem in which the eye does not bring incoming light to a sharp focus on the retina. Nearsightedness, farsightedness, and astigmatism are common forms and can cause blurred vision.

\medterm{Amino Acid} A small molecule used to build proteins. Amino acids also support energy production and other chemical processes in the body.

\synonyms
Protein Building Block; Amino Monocarboxylic Acid

\medterm{Amino Acid Catabolism} The body’s breakdown of amino acids after proteins are digested. Their nitrogen part is removed and safely processed, while the remaining part is used to make energy, glucose, ketones, or other substances.

\medterm{Amino Acid Symbols} The standard one-letter and three-letter abbreviations used to identify the 20 amino acids most commonly used to build proteins. They make genetic, laboratory, and biochemical information shorter and easier to compare.

\medterm{Aminoaciduria} Abnormally large amounts of amino acids in the urine, usually because the kidney tubules cannot reabsorb them properly or because of an inherited metabolic disorder.

\medterm{Aminoglycoside Nephrotoxicity} Kidney injury caused by aminoglycoside antibiotics, which can build up in kidney cells. Risk increases with high doses, prolonged treatment, existing kidney disease, or other kidney-toxic medicines.

\medterm{Aminoglycoside Ototoxicity} Inner-ear damage caused by aminoglycoside antibiotics, which may lead to hearing loss, ringing in the ears, dizziness, or balance problems. Risk increases with high or prolonged exposure.

\medterm{Aminoglycosides} A group of antibiotics that kill certain bacteria by stopping them from making essential proteins. Examples include gentamicin, tobramycin, amikacin, and streptomycin; they can damage the kidneys and inner ear.

\medterm{Aminophylline Overdose} Poisoning caused by too much aminophylline, a medicine that opens narrowed airways. It can cause vomiting, tremor, rapid or irregular heartbeat, low blood pressure, seizures, or dangerous changes in blood chemistry.

\medterm{Aminotransferase} An enzyme that helps cells process amino acids, the building blocks of proteins. Blood levels of ALT and AST may rise when liver, muscle, or other cells are injured, but the result does not identify the cause alone.

\medterm{Amiodarone Pulmonary Toxicity} A serious, potentially life-threatening adverse reaction to chronic amiodarone therapy, manifesting as subacute chronic interstitial pneumonitis, organizing pneumonia, or acute respiratory distress syndrome. Pathologically marked by intra-alveolar foamy lipid-laden macrophages (phospholipidosis) and diffuse alveolar damage.

\synonyms
Amiodarone Lung; Amiodarone-Induced Pneumonitis

\medterm{Amiodarone-Induced Skin Discoloration} A blue-grey or slate-grey darkening of sun-exposed skin caused by long-term amiodarone treatment. It is more likely with higher doses and may fade partly after the medicine is stopped.

\medterm{Amiodarone-Induced Thyroid Dysfunction} An overactive or underactive thyroid caused by amiodarone. The medicine contains iodine and can either increase thyroid-hormone production or inflame and damage the thyroid.

\medterm{Amiodarone-Induced Vortex Keratopathy} Whorl-shaped deposits in the surface layer of the cornea caused by amiodarone. They usually cause no symptoms but may produce halos around lights and generally improve after the medicine is stopped.

\medterm{Amitriptyline} A tricyclic antidepressant that changes serotonin and noradrenaline signalling in the brain. It is also used for some nerve pain and migraine prevention; common effects include sleepiness, dry mouth, and constipation.

\medterm{Amitriptyline And Perphenazine Overdose} Poisoning caused by taking too much amitriptyline and perphenazine. It can cause severe sleepiness, confusion, seizures, low blood pressure, abnormal heart rhythms, coma, or breathing problems.

\medterm{Amitriptyline Hydrochloride Overdose} Poisoning caused by taking too much amitriptyline. It can cause severe sleepiness, confusion, seizures, low blood pressure, abnormal heart rhythms, coma, or breathing problems.

\medterm{Ammonia} A nitrogen-containing waste product made when the body processes proteins. The liver normally changes it into urea for removal; high blood levels can affect brain function, especially in severe liver disease.

\medterm{Ammonia Blood Test} A blood test that measures ammonia, a waste product normally cleared by the liver. A high result may occur with severe liver disease, inherited metabolic conditions, some medicines, or other illnesses and must be interpreted with symptoms.

\medterm{Ammonia Poisoning} Injury caused by breathing, swallowing, or getting concentrated ammonia on the skin or eyes. It can burn tissues and cause coughing, breathing difficulty, chest pain, or serious lung damage.

\medterm{Ammonia Toxicity} Irritant or corrosive injury after inhalation, ingestion, or skin or eye contact with ammonia. Concentrated exposure can cause burning, cough, bronchospasm, laryngeal injury, pulmonary edema, or chemical burns; treatment emphasizes immediate decontamination, airway and respiratory support, and specialist toxicology guidance.

\medterm{Ammonium Hydroxide Poisoning} Injury caused by exposure to ammonium hydroxide, a corrosive chemical found in some cleaners and industrial products. It can burn the skin, eyes, mouth, throat, lungs, and digestive tract.

\medterm{Amnesia} Loss or impairment of memory beyond ordinary forgetfulness. It may prevent a person from forming new memories, recalling past events, or doing both, and may result from brain injury, illness, medicines, or emotional trauma.

\medterm{Amnestic Syndrome} A pattern of memory impairment that significantly affects learning, storing, or recalling information. It may occur with brain injury, stroke, dementia, epilepsy, substance use, or other neurological conditions.

\medterm{Amniocentesis} A procedure in which a thin needle removes a small sample of fluid from around the developing baby. The sample can be tested for genetic, chromosomal, or infection-related conditions.

\medterm{Amnioinfusion} A procedure that adds sterile fluid to the sac around the developing baby, usually during labour, to reduce pressure on the umbilical cord or dilute thick meconium.

\medterm{Amnion} The thin inner membrane forming the fluid-filled sac around the developing baby. It holds the amniotic fluid and helps protect the fetus during pregnancy.

\medterm{Amnion Nodosum} Small deposits of fetal skin material or other debris attached to the inner membrane of the pregnancy sac. They are usually linked with too little amniotic fluid, often because the developing baby is making less urine.

\medterm{Amnioscope} An instrument formerly inserted through the cervix to examine the amniotic fluid and membranes. It is rarely used now because ultrasound is safer and more practical.

\medterm{Amniotic Band} A strand of amniotic membrane that can wrap around a developing body part, restricting growth or blood flow and causing a deformity or loss of tissue.

\synonyms
Amniotic Band Sequence

\medterm{Amniotic Band Sequence} A group of birth differences caused when strands of the amniotic membrane wrap around or attach to developing fetal parts. Effects may include constriction rings, missing digits, limb defects, or facial and body-wall abnormalities.

\medterm{Amniotic Band Syndrome} Another name for a group of birth differences caused by strands of amniotic membrane constricting or entangling developing fetal parts. The effects vary widely and may involve the limbs, face, or body wall.

\synonyms
Amniotic Constriction Band Sequence; ADAM Complex; Streeter Dysplasia

\medterm{Amniotic Fluid} The liquid surrounding the developing baby inside the amniotic sac. It cushions the baby, allows movement and growth, and helps maintain a stable environment during pregnancy. It is continually produced, swallowed, and replaced.

\medterm{Amniotic Fluid Embolism} A rare, life-threatening complication of pregnancy or childbirth in which amniotic material enters the mother’s bloodstream and triggers sudden breathing or heart problems, often with severe bleeding.

\medterm{Amniotic Fluid Index} An ultrasound estimate of the amount of fluid around the developing baby. It adds the deepest fluid measurements from four areas of the uterus and helps identify unusually low or high fluid levels.

\synonyms
AFI; Amniotic Fluid Volume Assessment

\medterm{Amniotic Sac} The thin, fluid-filled membrane around the developing baby in the uterus. It contains amniotic fluid, cushions the fetus, and usually breaks near or during labour when the waters break.

\medterm{Amok} An old, culture-related term for a sudden episode of uncontrolled, often violent behaviour followed by exhaustion or memory loss. Modern assessment considers medical, neurological, psychiatric, substance-related, and cultural factors.

\medterm{AMP-Activated Protein Kinase} A master cellular energy-sensing serine/threonine kinase activated during metabolic stress when intracellular AMP-to-ATP or ADP-to-ATP ratios rise (such as during exercise, hypoxia, or fasting). Once activated, it restores energy homeostasis by promoting ATP-generating catabolic processes (such as glucose uptake via GLUT4, fatty acid oxidation, glycolysis, and autophagy) while suppressing ATP-consuming anabolic pathways (such as lipid synthesis, protein synthesis via mTOR inhibition, and gluconeogenesis). It is a primary molecular target of the antidiabetic drug metformin.

\synonyms
AMPK; 5'-AMP-Activated Protein Kinase

\medterm{Amphetamine} A stimulant medicine that increases certain chemical signals in the brain. It is used for selected attention, activity, or sleep disorders; misuse can cause dependence, high blood pressure, abnormal heart rhythms, or psychosis.

\medterm{Amphetamine Substance Use} Use of amphetamine-type stimulants, which can be associated with insomnia, anxiety, psychosis, high blood pressure, arrhythmias, overheating, dependence, and overdose.

\medterm{Amphetamine Toxicity} A sympathomimetic toxidrome caused by amphetamine or related stimulants, with agitation, mydriasis, sweating, tachycardia, hypertension, hyperthermia, chest pain, seizures, or psychosis. Severe cases may cause stroke, rhabdomyolysis, or cardiac ischemia; benzodiazepines, cooling, and treatment of complications are central.

\medterm{Amphetamine-Related Psychiatric Disorders} Psychiatric syndromes associated with amphetamine exposure, including intoxication, withdrawal, substance-induced psychosis, mood symptoms, anxiety, or a stimulant use disorder.

\medterm{Amplified Musculoskeletal Pain Syndrome} A condition, often affecting children or teenagers, in which the nervous system increases normal pain signals. Pain may be severe or widespread and can occur with sensitivity, fatigue, swelling, or difficulty using the affected area.

\medterm{Ampulla} A widened part of a tube-shaped body structure. Examples include the ampulla where the bile and pancreatic ducts join the intestine and the ampulla of a fallopian tube.

\medterm{Ampulla of Vater} An anatomical dilation located within the medial wall of the descending duodenum, formed by the union of the common bile duct and the main pancreatic duct, draining into the duodenal lumen via the major duodenal papilla.

\synonyms
Hepatopancreatic Ampulla

\medterm{Ampullary Carcinoma} A malignant epithelial tumor arising in or around the ampulla of Vater, where the bile duct and pancreatic duct enter the duodenum.
It can block bile flow and cause jaundice, abnormal liver tests, pancreatitis, or gastrointestinal bleeding.

\synonyms
Ampullary Cancer

\medterm{Amputation} Surgical removal or traumatic loss of a limb or part of a limb, such as a finger, toe, arm, or leg. Common reasons include severe injury, serious infection, poor blood flow causing tissue death or severe pain, and cancer. The level of amputation depends on healthy tissue, blood supply, healing, expected function, and rehabilitation needs.

\synonyms
Surgical Amputation; Limb Removal

\medterm{Amputation Stump} The part of a limb left after amputation. It may develop wound problems, infection, poor circulation, nerve pain, phantom-limb pain, or skin irritation from a prosthesis.

\medterm{Amsler Grid} A small square chart used to check central vision. Wavy, missing, or distorted lines may indicate a problem in the macula and are evaluated with a complete eye examination.

\medterm{Amusia} Difficulty recognizing, understanding, or producing music despite adequate hearing. It may be present from birth or develop after brain injury, stroke, or another neurological condition.

\medterm{Amygdala} A group of nerve cells deep in each side of the brain that helps process emotions, especially fear, and contributes to learning, memory, and responses to important or threatening events.

\synonyms
Amygdaloid Nucleus

\medterm{Amylase} An enzyme that begins breaking down starch. It is made mainly by the salivary glands and pancreas and can be measured in blood or urine when pancreatic or salivary-gland disease is suspected.

\medterm{Amylase And Lipase} Blood tests for digestive enzymes used when pancreatic disease is suspected. Lipase is generally more specific for pancreatic injury, but neither test alone confirms acute pancreatitis.

\medterm{Amylase Test} A blood or urine test that measures the amount of amylase, an enzyme made mainly by the salivary glands and pancreas. It may help assess pancreatitis and other problems affecting the pancreas or salivary glands.

\synonyms
Serum Amylase; Blood Amylase Test; Amylase Blood Test

\medterm{Amylin Analogues} Medicines that imitate amylin, a hormone released by the pancreas with insulin. They slow stomach emptying, reduce the rise in blood glucose after meals, and reduce glucagon release.

\medterm{Amyloid} An abnormal protein material that forms insoluble deposits in tissues. These deposits can interfere with the function of organs such as the kidneys, heart, nerves, liver, or brain.
\synonyms
Amyloid Protein

\medterm{Amyloid Nephropathy} Renal involvement occurring in systemic amyloidosis, characterized by extracellular accumulation of insoluble amyloid fibrils in glomerular, vascular, and interstitial compartments, typically leading to heavy proteinuria, nephrotic syndrome, and progressive kidney failure.

\medterm{Amyloid Neuropathy} Nerve damage caused by amyloid deposits in or around peripheral nerves. It may cause numbness, tingling, burning pain, weakness, or problems with automatic body functions such as blood-pressure control.

\medterm{Amyloid PET} A brain scan using a small radioactive tracer that attaches to amyloid deposits. It can help assess selected cognitive disorders, but the result must be interpreted with symptoms and other tests.

\medterm{Amyloid Plaque} A deposit of amyloid-beta protein outside nerve cells in the brain. Plaques are associated with Alzheimer disease, but their presence alone does not explain every person’s memory or thinking problem.

\medterm{Amyloid Precursor Protein} A normal protein found in nerve-cell membranes. Enzymes can cut it into smaller pieces, including amyloid-beta; abnormal accumulation of amyloid-beta in the brain is associated with Alzheimer disease.

\synonyms
APP; Amyloid-Beta Precursor Protein

\medterm{Amyloid Purpura} Purple or bruise-like spots, often around the eyes, caused by amyloid deposits weakening small blood vessels. It may occur with systemic amyloidosis and can be accompanied by easy bruising elsewhere.

\medterm{Amyloidosis} A group of diseases in which proteins fold abnormally, clump together, and form deposits called amyloid in tissues or organs. Amyloidosis may be localized to one site or systemic, affecting several organs such as the kidneys, heart, liver, nerves, or digestive tract. Major forms include light-chain (AL), serum amyloid A (AA), and transthyretin (ATTR) amyloidosis; effects vary with the amyloid protein and organs involved.

\synonyms
Amyloid Disease; Amyloid Infiltration

\medterm{Amyotrophic Lateral Sclerosis} A progressive disease that damages and destroys motor neurons, the nerve cells in the brain and spinal cord that control voluntary muscle movement and breathing. As these cells are lost, muscles become weak, stiff, wasted, or twitchy, eventually causing paralysis. Symptoms may begin in an arm, leg, or the muscles used for speaking and swallowing, and may later cause difficulty walking, speaking, chewing, swallowing, or breathing. Sensation is usually preserved, but some people develop changes in thinking or behavior, sometimes with frontotemporal dementia. Most cases have no known cause; about one in ten are inherited. ALS is also called Lou Gehrig disease.

\synonyms
Lou Gehrig's Disease
ALS

\medterm{Anabolic Steroid Induced Hypogonadism} Reduced testicular hormone production and fertility caused by misuse of anabolic steroids. Extra steroid hormones suppress the body’s normal hormone signals and may lead to testicular shrinkage, low sex drive, erectile problems, infertility, or depression.

\synonyms
ASIH; Anabolic-Androgenic Steroid-Induced Hypogonadism

\medterm{Anabolic Steroid Use and Abuse} Use of anabolic-androgenic steroids without a medical indication or at excessive doses. It can suppress natural hormone production and affect fertility, mood, blood pressure, cholesterol, liver function, heart structure, and growth in adolescents.

\medterm{Anabolic Steroids} Synthetic medicines related to testosterone that promote muscle and bone growth. They have limited medical uses but are sometimes misused for appearance or performance; misuse can harm the heart, liver, fertility, mood, and hormone system.

\medterm{Anabolism} The set of body processes that use energy to build larger molecules from smaller ones, such as making proteins from amino acids or storing glucose as glycogen.

\medterm{Anaemia} A condition in which the blood has too few red blood cells or too little haemoglobin. This reduces the blood’s ability to carry oxygen and may cause tiredness, weakness, breathlessness, or dizziness.

\medterm{Anaerobe} A microorganism that can grow without oxygen. Some anaerobes are harmed or killed by oxygen, while others tolerate oxygen but do not use it to produce energy.

\medterm{Anaerobic} Able to live, grow, or function without oxygen. Anaerobic bacteria can cause infections in deep wounds, abscesses, the mouth, abdomen, or pelvis.

\medterm{Anaerobic Bacteria} Bacteria that can grow without oxygen. Obligate anaerobes cannot use oxygen and may be harmed by it, while facultative anaerobes can grow with or without oxygen. Anaerobic bacteria normally live in places such as the mouth, intestine, and vagina. They can cause infections in deep wounds, abscesses, the abdomen, pelvis, or other tissues with low oxygen.

\medterm{Anaerobic Glycolysis} The metabolic breakdown of glucose into lactate in the absence or insufficiency of cellular oxygen, producing a net yield of two molecules of adenosine triphosphate per glucose molecule and regenerating oxidized nicotinamide adenine dinucleotide (NAD+) to maintain energy generation.

\medterm{Anagen} The active growth phase of the hair cycle, when cells in the hair root make the hair shaft. Scalp hairs usually remain in this phase for several years before entering a resting or shedding phase.

\medterm{Anagen Effluvium} An abrupt, extensive shedding of hair during the active growth (anagen) phase, caused by acute toxic disruption of mitotic activity in the follicular matrix cells. It typically begins within days to weeks of exposure to cytotoxic chemotherapy, ionizing radiation, or acute heavy metal poisoning.

\medterm{Anal Abscess} A collection of pus near or within the anus, usually caused by an infected anal gland. It can cause throbbing pain, swelling, fever, and difficulty sitting.

\medterm{Anal Canal} The final part of the digestive tract, extending from the rectum to the anus. It contains muscles that control the passage of stool and has different nerve and blood supplies along its length.

\medterm{Anal Cancer} Cancer arising in the anal canal or perianal skin, most often squamous-cell carcinoma associated with persistent high-risk HPV infection. Bleeding, pain, a lump, itching, or altered bowel function may occur; examination and biopsy establish the diagnosis.

\medterm{Anal Carcinoma} Cancer arising in the anal canal or anal margin, most often squamous-cell cancer associated with persistent high-risk human papillomavirus infection. It may cause bleeding, pain, itching, or a lump.

\synonyms
Anal Cancer

\medterm{Anal Condyloma} A wart-like growth on or around the anus, usually caused by low-risk types of human papillomavirus. It is usually not cancerous, but multiple or recurring warts may occur with other abnormal anal cells.

\medterm{Anal Fissure} A small tear in the lining of the anus, usually caused by hard stools, constipation, or repeated diarrhea. It can cause sharp pain during bowel movements and bright-red blood on the toilet paper.

\medterm{Anal Fistula} An abnormal tunnel between the anal canal and nearby skin, often developing after an anal abscess. It may cause repeated discharge, pain, itching, or a small opening near the anus.

\medterm{Anal Intercourse} Sexual activity involving insertion into the anus. Because the anal lining can tear and does not provide natural lubrication, it can transmit infections and cause injury if precautions are not used.

\medterm{Anal Intraepithelial Neoplasia} Abnormal cells in the lining of the anus, usually associated with persistent high-risk human papillomavirus infection. Low-grade changes may resolve, while high-grade changes can progress to anal cancer.

\medterm{Anal Itching} Itching around the anus. Common causes include moisture, irritation, dermatitis, haemorrhoids, infection, and occasionally pinworms or another disease.

\synonyms
Pruritus Ani; Perianal Itching; Rectal Itching

\medterm{Anal Mucositis} Inflammation and open sores of the anal lining, often caused by radiation or chemotherapy. It can cause pain, burning, bleeding, discharge, and difficulty passing stool.

\medterm{Anal Necrosis} Death of tissue around the anus, anal canal, or nearby skin. It may result from severe infection, poor blood flow, injury, or radiation and can cause severe pain, skin discolouration, blisters, fever, or rapid worsening.

\medterm{Anal Pain} Pain around the anus or nearby skin. Common causes include a fissure, abscess, thrombosed haemorrhoid, muscle spasm, infection, or inflammation.

\synonyms
Proctalgia

\medterm{Anal Papilla} A small, usually normal projection of the anal lining near the junction between the upper and lower types of lining in the anal canal.

\medterm{Anal Sphincter} A muscular structure that controls passage of stool and gas from the anal canal. It includes the involuntary internal anal sphincter and the voluntary external anal sphincter.

\medterm{Analgesia} Relief or absence of pain without necessarily causing unconsciousness. It may result from medicines, a nerve block, or another treatment; touch, pressure, or movement may still be felt.

\medterm{Analgesics} Medicines that reduce or relieve pain. Non-opioid analgesics reduce pain and fever. Nonsteroidal anti-inflammatory drugs (NSAIDs) reduce pain, fever, and inflammation by lowering the production of substances involved in these responses. Opioid analgesics reduce the brain's and spinal cord's response to pain but may cause sleepiness, constipation, slowed breathing, tolerance, and dependence. Adjuvant analgesics are medicines mainly used for other conditions that can help with certain types of pain, such as nerve pain. The choice depends on the type and severity of pain, the person's health, and possible risks.

\synonyms
Painkillers; Pain-Relieving Medicines

\medterm{Analog} A substance, structure, or process that resembles another and has a similar function but is not identical in chemical makeup or structure.

\medterm{Analogous} Similar in function or purpose to another structure or process, even though its origin or physical structure is different.

\medterm{Analysis} A careful examination of information, samples, measurements, symptoms, or clinical problems to identify patterns, understand causes, or support a medical decision.

\medterm{Analysis Of Urine} Laboratory examination of urine for blood, cells, protein, glucose, germs, and other substances. Results can help assess urinary infection, kidney disease, diabetes, bleeding, dehydration, and other conditions.

\medterm{Analytic Sensitivity} The smallest amount of a substance that a laboratory method can reliably detect. It describes the test’s measurement ability and is different from clinical sensitivity, which describes how well a test detects disease.

\medterm{Anaphia} Complete loss of touch sensation in a body area. It may result from nerve injury, neurological disease, reduced blood flow, or anaesthetic treatment and can allow burns, wounds, or pressure injuries to go unnoticed.

\synonyms
Anaptic

\medterm{Anaphylactic Shock} A severe form of anaphylaxis in which a sudden allergic reaction causes dangerously low blood pressure and poor blood flow to organs. It may occur with airway swelling, breathing difficulty, collapse, or cardiac arrest.

\medterm{Anaphylactoid Reaction} An immediate, severe systemic hypersensitivity reaction that closely mimics IgE-mediated anaphylaxis in its clinical manifestations, but occurs via direct non-immune or non-IgE-mediated mast cell and basophil degranulation (such as that triggered by radiocontrast media or certain drugs).

\synonyms
Non-Allergic Anaphylaxis; Pseudoallergic Reaction

\medterm{Anaphylatoxin} A small peptide fragment (predominantly C3a, C4a, and C5a) generated during activation of the complement cascade. Anaphylatoxins induce mast cell degranulation, smooth muscle contraction, increased vascular permeability, and chemotaxis of phagocytic leukocytes.

\medterm{Anaphylaxis} A sudden, potentially life-threatening systemic hypersensitivity reaction that can involve the skin, airways, lungs, circulation, or gastrointestinal tract. It may cause hives or swelling, breathing difficulty, vomiting, dizziness, dangerously low blood pressure, or collapse; skin findings may be absent. Immediate intramuscular epinephrine and emergency medical care are required.

\medterm{Anaphylaxis During Anesthesia} A severe allergic reaction occurring during or soon after anaesthesia. It may cause sudden low blood pressure, airway narrowing, breathing difficulty, collapse, or skin flushing; triggers can include medicines, antiseptics, or latex.

\medterm{Anaplasia} Loss of the normal mature features and specialized function of cells. Anaplastic cells may vary greatly in size and shape, have abnormal nuclei, and divide rapidly or unusually. Anaplasia is most often seen in high-grade cancers and usually suggests more aggressive behavior. It is different from metaplasia, in which one mature cell type changes into another cell type.

\medterm{Anaplasmosis} A tick-borne infection caused by the bacterium \textit{Anaplasma phagocytophilum}. It may cause fever, headache, muscle aches, low platelet and white-cell counts, and raised liver enzymes. PCR or paired blood-antibody tests can help confirm it. Morulae seen in white blood cells can support the diagnosis but are not always present and do not confirm it alone.

\synonyms
Human Granulocytic Anaplasmosis; HGA

\medterm{Anaplastic} Describing cells or a tumor that has lost many features of its normal mature form. Anaplastic tumors often grow quickly, invade nearby tissue, or spread, although the outlook depends on the specific cancer and its treatment.

\medterm{Anaplastic Astrocytoma} An older term for a high-grade tumour arising from star-shaped supporting cells in the brain. Modern classification relies on its genetic features as well as its appearance; symptoms depend on its location.

\medterm{Anaplastic Ependymoma} A high-grade tumour arising from cells lining the brain’s fluid-filled spaces or the spinal canal. It can grow quickly and may obstruct the flow of cerebrospinal fluid.

\medterm{Anaplastic Large Cell Lymphoma} A rare type of non-Hodgkin lymphoma arising from abnormal T cells. The cells are usually large and carry the CD30 marker; the disease may affect the skin, lymph nodes, or other organs.

\synonyms
ALCL

\medterm{Anaplastic Thyroid Cancer} A rare, fast-growing thyroid cancer that may appear as a rapidly enlarging neck lump and cause difficulty swallowing, speaking, or breathing.

\medterm{Anaplastic Thyroid Carcinoma} An extremely aggressive, undifferentiated malignancy of the thyroid follicular epithelium characterized by rapid local invasion into adjacent cervical structures and early distant metastasis. Presenting symptoms include a rapidly enlarging hard neck mass, hoarseness, dysphagia, and dyspnea, carrying a very poor median survival despite multimodality therapy.

\synonyms
Undifferentiated Thyroid Carcinoma; ATC

\medterm{Anarithmetia} A primary neuropsychological impairment in basic arithmetic reasoning, computational ability, and understanding of numerical concepts, distinguishing it from secondary calculation deficits caused by spatial or linguistic impairments. It is typically associated with lesions of the dominant inferior parietal lobule (angular gyrus).

\medterm{Anasarca} Severe, widespread swelling caused by excess fluid under the skin and sometimes around organs. It can occur with serious heart, kidney, or liver disease or very low blood protein and may cause rapid weight gain or breathing difficulty.

\medterm{Anastomosis} A surgical, traumatic, or pathological connection created between two hollow organs, tubular structures, or blood vessels (such as an arterial anastomosis, arterio-venous fistula, or bowel resection re-connection), establishing communication and allowing continuous fluid or luminal flow.

\synonyms
Surgical Anastomosis; Vascular Anastomosis; Enteroanastomosis

\medterm{Anastomotic Leak} Leakage from a connection made during surgery, often where two sections of intestine or another hollow organ were joined. An incomplete seal can allow contents into nearby tissue and cause pain, fever, infection, or life-threatening illness.

\medterm{Anastrozole} A medicine that lowers the body’s production of oestrogen. It is used mainly for hormone-sensitive breast cancer, especially after menopause.

\medterm{Anatomic Orientation Terms} Standard words used to describe the position of one body structure relative to another, including above, below, front, back, near, far, inside, and outside.

\medterm{Anatomical Axis} An imaginary line around which a body part rotates or moves. Anatomical axes help describe movements such as rotation of the shoulder, hip, spine, or other joints.

\medterm{Anatomical Conduit} A natural or surgically created passage that carries air, fluid, tissue, or another structure between body regions, such as an airway, blood vessel, or surgical channel.

\medterm{Anatomical Orifice} A natural opening into or out of the body or through a body structure, such as the mouth, nostril, anus, urethra, or an opening of a blood vessel.

\synonyms
Body Opening

\medterm{Anatomical Planes} Imaginary flat surfaces used to divide the body into sections for describing anatomy and interpreting medical scans. Common planes are sagittal, coronal, and transverse.

\medterm{Anatomical Position} The standard reference posture for describing the body: standing upright, facing forward, arms at the sides, hands turned forward, and feet directed forward.

\medterm{Anatomical Snuffbox} A triangular hollow on the thumb side of the back of the wrist that appears when the thumb is extended. It overlies the scaphoid bone and radial artery; tenderness there after a fall may indicate a scaphoid fracture.

\synonyms
Anatomic Snuffbox

\medterm{Anatomist} A scientist or health professional who studies the structure of the human or animal body, including its organs, tissues, blood vessels, nerves, and their relationships.

\medterm{Anatomy} The study of body structure and how its parts relate to one another. It includes visible structures, tissues and cells seen with magnification, and how the body develops.

\medterm{ANCA-Associated Vasculitis} A group of rare diseases in which the immune system attacks small blood vessels. It can affect the kidneys, lungs, nerves, skin, eyes, and other organs. A blood test for antineutrophil cytoplasmic antibodies (ANCA) can help with diagnosis but cannot confirm the disease on its own.

\medterm{Ancillary Testing, Pathology} Supplementary diagnostic modalities—including immunohistochemistry, special histochemical stains, flow cytometry, in situ hybridization, and molecular genetic assays—performed in conjunction with standard light microscopy to achieve precise pathological classification.

\synonyms
Ancillary Diagnostic Testing; Diagnostic Immunohistochemistry; Special Pathology Stains

\medterm{Ancylostomiasis} A parasitic nematode infection caused by species of the genus \textit{Ancylostoma} (such as \textit{Ancylostoma duodenale}), acquired through transcutaneous penetration of infective larvae from contaminated soil, characteristically leading to chronic intestinal blood loss, iron deficiency anemia, and hypoproteinemia.

\synonyms
Hookworm Infection; Old World Hookworm Disease; Ancylostoma Duodenale Infection

\medterm{Andersen-Tawil Syndrome} A rare inherited disorder causing episodes of muscle weakness, abnormal heart rhythms, and distinctive physical features such as a small lower jaw, low-set ears, or widely spaced eyes.

\synonyms
ATS; Periodic Paralysis Type 3; Andersen Syndrome

\medterm{Anderson Disease} A rare inherited disorder in which abnormal glycogen builds up in cells because of changes in the \textit{GBE1} gene. It usually begins in infancy with poor growth and an enlarged liver and may progress to scarring, high blood pressure in the liver, and liver failure.

\synonyms
Glycogen Storage Disease Type IV (GSD IV); Amylopectinosis

\medterm{Androgen} A hormone such as testosterone that helps regulate sexual development, fertility, muscle, bone, and hair growth in people of all sexes.

\medterm{Androgen Deficiency} Too little production or effect of androgen hormones, especially testosterone. In males, it may cause reduced sexual desire, erectile difficulty, infertility, loss of muscle or bone, low energy, or reduced facial and body hair.

\medterm{Androgen Excess} Abnormally high levels or effects of androgen hormones. It may cause acne, increased facial or body hair, scalp hair loss, irregular periods, infertility, or a deeper voice, especially in females.
\synonyms
Hyperandrogenism

\medterm{Androgen Insensitivity Syndrome} An inherited condition in which cells cannot respond fully to testosterone and related hormones. A person with XY chromosomes may therefore have female-appearing external genitals, incomplete male development, undescended testes, and no uterus.

\medterm{Androgen Insufficiency} A level or effect of androgen hormones that is too low for normal body function. Symptoms may include reduced sexual desire, erectile difficulty, infertility, low energy, reduced muscle mass, or weaker bones.

\synonyms
Androgen Deficiency

\medterm{Androgen Receptor} A protein in cells that binds testosterone and related hormones. This allows the hormones to influence sexual development, reproduction, muscle, bones, hair, and some cancers.

\medterm{Androgen Suppression} Lowering the production of androgen hormones or blocking their action in the body. It is used mainly for some hormone-sensitive cancers, especially prostate cancer, and can cause hot flushes, sexual problems, bone loss, weight change, and metabolic effects.

\synonyms
Androgen Deprivation

\medterm{Androgenetic Alopecia} A common, inherited form of gradual hair loss caused by sensitivity of hair follicles to androgen hormones. It usually produces a receding hairline or crown thinning in males and diffuse thinning over the crown in females.

\medterm{Androgenic} Relating to androgen hormones, such as testosterone, or producing their effects. These effects may include facial or body hair, oily skin, a deeper voice, muscle growth, and development of male reproductive features.

\medterm{Android Pelvis} A relatively narrow, heart-shaped or funnel-shaped female pelvis with reduced space in the middle. During childbirth, this shape may make passage of the baby more difficult than some other pelvic shapes.

\medterm{Andropause} Alternate informal term; see \textit{Male Hypogonadism}.

\medterm{Anejaculation} Inability to release semen during orgasm or sexual activity. It may result from nerve or spinal disease, pelvic surgery, medicines, diabetes, hormonal problems, or psychological factors and may cause infertility.

\medterm{Anemia} A condition in which the blood cannot carry enough oxygen to the body’s tissues, usually because there are too few red blood cells or too little haemoglobin. Causes include blood loss, reduced production, and increased breakdown.

\synonyms
Anaemia; Low Red Blood Cell Count

\medterm{Anemia and Thrombocytopenia in Pregnancy} Low red-cell or platelet counts during pregnancy may reflect iron deficiency, dilution, nutritional disease, immune disorders, hypertensive disease, bleeding, infection, or marrow disease. Severity, cause, symptoms, and trends guide treatment and delivery planning.

\medterm{Anemia In Pregnancy} Anemia during pregnancy, commonly caused by iron deficiency, folate deficiency, vitamin B12 deficiency, blood loss, or the normal increase in the liquid part of blood. It can cause tiredness, weakness, breathlessness, or dizziness.

\synonyms
Gestational Anemia; Pregnancy-Associated Anemia

\medterm{Anemia Of Chronic Disease} Anemia associated with long-term inflammation, infection, cancer, kidney disease, or another chronic illness. Inflammation limits the use of stored iron and reduces red-cell production; mild cases may cause no symptoms.

\medterm{Anemia Of Inflammation} Anemia caused by long-lasting infection, inflammatory disease, cancer, or another condition that activates the immune system. Inflammation traps iron in storage and slows red-cell production.

\medterm{Anemia of Malignancy} A multifactorial normocytic or microcytic anemia developing in patients with neoplastic disease, resulting from cytokine-mediated blunting of erythropoiesis, impaired iron utilization, marrow infiltration, hemorrhage, or myelosuppressive antineoplastic therapies.

\synonyms
Cancer-Related Anemia; Tumor-Induced Anemia; Anemia of Cancer

\medterm{Anemic} Having anemia, meaning that the blood carries less oxygen than the body needs because it has too few red blood cells or too little haemoglobin.

\medterm{Anencephaly} A serious neural tube birth defect in which the upper part of the neural tube fails to close during the first weeks of pregnancy. As a result, much of the brain, especially the forebrain, and parts of the skull and scalp do not develop. It is usually detected before birth and is fatal; most affected babies are stillborn or die shortly after birth. There is no treatment that can correct the defect.

\medterm{Anergy} A state in which an immune cell does not respond to a substance that would normally activate it. The cell remains present but is functionally quiet; this can help prevent attacks against the body’s own tissues.

\medterm{Anesthesia} A medically controlled loss of feeling caused by medicines during a procedure. Anesthetic medicines mainly block nerve signals, especially pain signals, from reaching the brain. Local anesthesia numbs a small area while the person remains awake. Regional anesthesia blocks nerves supplying a larger part of the body and may be used while the person remains awake. General anesthesia also changes brain activity to cause controlled unconsciousness. Before modern anesthesia, surgery was extremely painful and had to be very quick; anesthesia made longer and more complex operations possible. Breathing, blood pressure, heart rate, and oxygen levels are monitored.

\synonyms
Anaesthesia

\medterm{Anesthesia Gas Analyzers} Devices that measure the amounts of oxygen, carbon dioxide, and inhaled anesthetic gases delivered to and breathed out by a patient during anesthesia. They help detect problems with breathing or gas delivery.

\medterm{Anesthesia Informatics} The use of electronic records, monitoring data, and computer systems to support anesthesia care. It can improve documentation, identify medicine interactions or dosing errors, and help evaluate outcomes.

\medterm{Anesthesia Machine} Equipment that supplies controlled mixtures of oxygen, air, and anesthetic gases, and supports breathing during anesthesia. It includes systems that remove carbon dioxide and safety devices that warn of delivery problems.

\medterm{Anesthesia Monitoring} Continuous assessment of vital functions during anesthesia, including oxygen levels, breathing, blood pressure, heart rate, temperature, and level of unconsciousness as appropriate for the procedure.

\medterm{Anesthesia Pharmacogenomics} The study of how inherited genetic differences affect responses to anesthetic and pain medicines. Examples include prolonged muscle paralysis from pseudocholinesterase deficiency and inherited susceptibility to malignant hyperthermia.

\medterm{Anesthesiologist} A doctor trained to provide anesthesia and manage pain, breathing, blood pressure, and other vital functions during procedures. Anesthesiologists also assess patients before surgery and may treat critical illness or severe pain.

\medterm{Anesthesiology} The medical specialty covering anesthesia, monitoring and support of vital functions during procedures, care before and after surgery, critical care, and treatment of acute or chronic pain.

\medterm{Anetoderma} A rare disorder caused by loss or weakening of elastic fibers in the dermis, producing soft, wrinkled, depressed, or pouch-like macules and plaques with a characteristic button-hole feel. It may be primary or follow infection, inflammation, autoimmune disease, or another skin lesion; biopsy and selected systemic evaluation may be needed.

\medterm{Aneuploidy} An abnormal number of whole chromosomes in a cell, caused by an error during cell division. A cell may have an extra or missing chromosome, as in trisomy 21, and this can affect development or pregnancy.

\medterm{Aneurysm} A weakened area of a blood-vessel wall that bulges or enlarges. Aneurysms can occur in arteries and rarely veins; they may enlarge, form clots, tear, or rupture and cause life-threatening bleeding.

\medterm{Aneurysm Clip} A small surgical device placed across the opening of a brain aneurysm to stop blood from entering it. It is usually placed during open brain surgery; rarely, it may move, block a nearby vessel, or allow bleeding to recur.

\medterm{Aneurysm Clipping} Brain surgery to place a small clip across the opening of a brain aneurysm, a weak bulge in a brain artery. The clip stops blood from entering the bulge and lowers the risk of bleeding, but the procedure has its own risks and cannot guarantee that a stroke or future bleeding will not occur.
\synonyms
Surgical Clipping; Cerebral Aneurysm Clipping; Microvascular Clipping

\medterm{Aneurysm Coils} Soft metal coils placed inside an aneurysm through a catheter. The coils promote clotting within the bulge and reduce blood flow into it, lowering the chance of rupture.

\medterm{Aneurysm Repair} A procedure that reinforces, closes, or replaces a weakened section of a blood vessel to prevent rupture or other complications. It may be performed through open surgery or through a catheter inside the vessel.

\medterm{Aneurysms} Weakened, bulging areas in blood-vessel walls. They may remain stable or enlarge, form clots, press on nearby structures, tear, or rupture and cause life-threatening bleeding.

\medterm{Angelman Syndrome} A genetic disorder affecting brain development and function, usually because the maternal copy of the \textit{UBE3A} gene is missing or inactive. It causes severe developmental delay, very limited speech, movement and balance problems, seizures, and a characteristic happy or excitable manner.

\medterm{Anger} An emotional response to frustration, unfairness, danger, or a perceived threat. It may range from irritation to intense rage; persistent or uncontrollable anger can affect relationships, safety, and health.

\medterm{Angiitis} Pathological inflammation and necrosis of blood-vessel walls, leading to lumen stenosis, ischemia, aneurysm formation, or hemorrhage in affected organ systems.

\synonyms
Vasculitis; Blood Vessel Inflammation; Arteritis

\medterm{Angina} Pressure, tightness, burning, or discomfort caused by temporary reduction of blood flow to the heart muscle, usually because of narrowed coronary arteries. It may spread to the arm, jaw, back, or neck and may occur with activity or at rest.

\medterm{Angina Bullosa Hemorrhagica} A usually benign, sudden blood-filled blister of the oral mucosa, most often the soft palate. It may follow minor trauma, inhaled or topical steroids, or occur with diabetes; blisters usually rupture and heal without treatment, but recurrent or atypical lesions need evaluation to exclude autoimmune blistering disease.

\medterm{Angina Pectoris} Chest discomfort caused by temporary reduced blood flow and oxygen to the heart muscle, usually because of narrowed coronary arteries. It may feel like pressure, squeezing, tightness, burning, or indigestion and may spread to the arm, shoulder, neck, jaw, back, or upper abdomen. Stable angina usually occurs with exercise or emotional stress and improves with rest. Unstable angina is new, worsening, prolonged, or occurs at rest and is associated with a higher risk of heart attack.

\medterm{Angina Symptoms} Symptoms of reduced blood flow to the heart may include pressure, squeezing, burning, or tightness in the chest, with discomfort spreading to the arm, jaw, back, or shoulder. Breathlessness, sweating, nausea, or unusual tiredness may also occur.

\synonyms
Ischemic Chest Pain; Anginal Chest Pain; Anginal Discomfort

\medterm{Angioblastoma} A usually slow-growing, highly vascular tumour of the central nervous system, most often found in the cerebellum or spinal cord. It may cause headache, balance problems, weakness, or other symptoms from pressure on nearby tissue.

\medterm{Angiocardiography} Imaging of the heart and major blood vessels, usually after contrast material is given. X-ray, CT, or MRI can show the heart chambers, valves, narrowed vessels, leaks, abnormal connections, or other structural problems.

\medterm{Angiodysplasia} A common acquired vascular malformation characterized by fragile, dilated, tortuous submucosal blood vessels in the gastrointestinal tract, most frequently found in the cecum and right colon of adults over age 60. It is a major cause of recurrent, painless lower gastrointestinal bleeding and unexplained iron-deficiency anemia. It is strongly associated with aortic valve stenosis (Heyde's syndrome, due to acquired von Willebrand factor deficiency) and chronic renal failure. Diagnosis is established on colonoscopy, where lesions appear as small, bright-red, spider-like mucosal patches; active or high-risk bleeding is treated endoscopically using argon plasma coagulation (APC) or electrocautery.
\synonyms{Colonic Angiodysplasia; Arteriovenous Malformation of the Colon; Vascular Ectasia of the Colon}

\medterm{Angiodysplasia Of The Colon} Alternate name; see \textit{Angiodysplasia}.

\medterm{Angioedema} Sudden, localized swelling of deeper skin or mucosal tissues, often affecting the lips, eyelids, tongue, throat, hands, feet, or bowel wall. Histamine-mediated episodes may occur with hives and itching, whereas bradykinin-mediated forms usually lack hives; tongue or throat swelling, voice change, or breathing difficulty is an emergency.

\medterm{Angioedema And Hives} A combination of raised, itchy skin patches called hives and deeper swelling called angioedema. Common triggers include foods, medicines, insect stings, infections, and physical exposures; breathing difficulty is an emergency.

\medterm{Angioedema, Hereditary} Alternate index label; see \textit{Hereditary Angioedema}.

\medterm{Angiofibroma} A usually noncancerous growth made of blood vessels and fibrous tissue. It commonly appears as a small facial papule in adolescents or, in a different form, as a vascular tumour in the back of the nose that may cause blockage or nosebleeds.

\medterm{Angiofollicular Lymph Node Hyperplasia} An uncommon disorder in which lymph nodes contain excess small blood vessels and abnormal immune-cell growth. It is also called Castleman disease and may affect one lymph-node region or many parts of the body.

\medterm{Angiogenesis} The formation of new blood vessels from existing vessels. It is normal during growth, menstruation, and wound healing, but abnormal angiogenesis can help cancers grow and can damage vision in some eye diseases.

\medterm{Angiogenesis Inhibition} Preventing or slowing the growth of new blood vessels. This can limit a tumor’s blood supply or reduce harmful vessel growth in some eye diseases. Treatment may increase the risk of high blood pressure, bleeding, or delayed wound healing.

\synonyms
Antiangiogenic Therapy

\medterm{Angiogenesis Inhibitor} A therapeutic agent that inhibits the formation of new blood vessels, predominantly by antagonizing vascular endothelial growth factor (VEGF) or its receptors, used in oncology to restrict tumor blood supply and in ophthalmology to treat neovascular macular degeneration.

\synonyms
Antiangiogenic Agent

\medterm{Angiogram} An image of blood vessels, usually made after contrast material is given. It can show narrowing, blockage, aneurysm, tearing, abnormal connections, or bleeding and may be produced by catheter X-ray, computed tomography, or magnetic resonance imaging.

\medterm{Angiography} Imaging of blood vessels to show their shape and blood flow. It may use X-rays with injected contrast material, computed tomography (CT), magnetic resonance imaging (MRI), or a catheter placed inside a blood vessel. Angiography can show narrowing, blockage, aneurysms, abnormal connections, bleeding, or other vascular abnormalities; catheter angiography may also allow treatment during the same procedure.

\medterm{Angioid Streak} Irregular dark lines caused by breaks in a tough layer beneath the retina. They may bleed or allow abnormal blood vessels to grow, reducing vision, and can be associated with some inherited disorders.

\medterm{Angioimmunoblastic T-Cell Lymphoma} An aggressive subtype of peripheral mature T-cell lymphoma characterized by generalized lymphadenopathy, constitutional symptoms (fever, night sweats, weight loss), hepatosplenomegaly, skin rashes, and polyclonal hypergammaglobulinemia.

\synonyms
AITL

\medterm{Angiokeratoma of the Scrotum} A benign vascular papule caused by superficial capillary dilation and overlying thickened skin, usually appearing as multiple dark-red, blue, or purple papules on the scrotum. Lesions may bleed after friction and can be mistaken for melanoma or other lesions; examination establishes the diagnosis.

\medterm{Angiolymphoid Hyperplasia With Eosinophilia} An uncommon benign vascular and inflammatory lesion, also called epithelioid hemangioma, that produces solitary or grouped red-brown papules or nodules, usually on the head and neck near the ears or scalp. Eosinophilia may occur; biopsy distinguishes it from Kimura disease and other vascular or neoplastic lesions.

\medterm{Angiomyofibroblastoma} A usually noncancerous, slow-growing tumour of blood-vessel and connective-tissue cells, most often occurring as a small, well-defined, painless mass in the vulva or vagina.

\medterm{Angiomyolipoma} A usually noncancerous kidney tumour made of blood vessels, muscle, and fat. It may occur alone or with tuberous sclerosis and can bleed, especially when large.

\medterm{Angioneurotic Edema, Hereditary} An older name for hereditary angioedema, an inherited disorder causing repeated swelling of the skin, bowel, or airway. Attacks usually occur without hives and throat swelling can block breathing.

\medterm{Angiopathy} Any disease or abnormal change in blood vessels. Causes include diabetes, inflammation, fatty deposits, injury, infection, and inherited disorders; effects depend on the vessels and organs involved.

\medterm{Angioplasty} A minimally invasive procedure that widens a narrowed or blocked blood vessel to improve blood flow. A thin catheter with a small balloon is guided to the narrowed area, often using X-ray imaging. The balloon is inflated briefly and then removed. A small mesh tube called a stent may be left in place to help keep the vessel open. Angioplasty may be performed on coronary arteries or other arteries and veins.

\medterm{Angioplasty Balloon} An inflatable cylindrical balloon catheter designed to dilate stenotic vascular segments, fracture atheromatous plaques, or deploy and seat endovascular stents.

\synonyms
Balloon Catheter; Dilatation Balloon; Angioplasty Catheter

\medterm{Angiosarcoma} A rare, fast-growing cancer that begins in cells lining blood or lymph vessels. It may occur in the skin, breast, liver, or bone and can spread to other organs.

\medterm{Angiosarcoma Of The Breast} A rare cancer arising from blood-vessel cells in the breast. It may cause an enlarging breast mass or bruise-like skin colour change and can develop after breast radiation or long-term lymphoedema.

\medterm{Angiosarcoma Of The Heart} A rare, aggressive cancer arising from blood-vessel cells in the heart, most often in the right upper chamber. It may obstruct blood flow, cause fluid around the heart or abnormal rhythms, and spread to other organs.

\medterm{Angiotensin} A group of hormones that help regulate blood pressure and the body’s salt and water balance. Angiotensin II narrows blood vessels and signals the kidneys and adrenal glands to retain salt and water.

\medterm{Angiotensin Converting Enzyme} An enzyme that changes angiotensin I into angiotensin II, which raises blood pressure by narrowing blood vessels. The enzyme also breaks down bradykinin, a substance that can widen vessels.

\medterm{Angiotensin I} A mostly inactive hormone made from angiotensinogen. Angiotensin-converting enzyme changes it into angiotensin II, which narrows blood vessels and promotes salt and water retention.

\medterm{Angiotensin II} A hormone that narrows blood vessels and causes the body to retain salt and water. These effects increase blood pressure and help maintain blood flow when circulating volume is low.

\medterm{Angiotensin Receptor Blockers} Medicines that prevent angiotensin II from activating its main receptor. They relax blood vessels and reduce salt and water retention, lowering blood pressure and reducing strain on the heart and kidneys.

\medterm{Angiotensin-Converting Enzyme 2} A cell-surface zinc metalloprotease that cleaves angiotensin II into the vasodilatory and cardioprotective peptide angiotensin-(1-7), acting as an essential counter-regulatory enzyme within the renin-angiotensin-aldosterone system; also serves as the host cell receptor for the SARS-CoV-2 spike protein.

\synonyms
ACE2

\medterm{Angle-Closure Glaucoma} A type of glaucoma in which the outer edge of the iris blocks the eye’s drainage angle. Fluid then builds up, causing eye pressure to rise and potentially damaging the optic nerve and vision. It may develop suddenly with severe eye pain, blurred vision, halos, headache, nausea, or vomiting, or develop slowly without symptoms.

\medterm{Angular Cheilitis} Inflammation, soreness, and cracking at one or both corners of the mouth. Moisture and saliva, yeast or bacterial infection, dentures, irritation, or iron and vitamin deficiencies can contribute.

\medterm{Anhedonia} Reduced or absent ability to feel pleasure or interest in activities that were previously enjoyable. It can occur with depression, schizophrenia, substance use, and other conditions.

\medterm{Anhidrosis} Absent or greatly reduced sweating. It may result from nerve damage, skin disease, dehydration, medicines, or an inherited condition. Without enough sweat, the body may overheat dangerously, especially during exercise or hot weather.

\medterm{Animal-Assisted Therapy} A planned healthcare intervention in which a trained animal is used with a qualified professional to support goals such as improving mood, communication, movement, or participation. Safety, hygiene, and allergies must be considered.

\medterm{Anion} A particle with a negative electrical charge. Important anions in body fluids include chloride, bicarbonate, and phosphate, which help maintain fluid balance, electrical activity, and the blood’s acid level.

\medterm{Anion Gap} A value calculated from several blood salts to estimate acids that were not directly measured. A high result can suggest excess acid in the blood, but it must be interpreted with albumin, other tests, symptoms, and the laboratory’s method.

\medterm{Anion Gap Metabolic Acidosis} A subtype of metabolic acidosis characterized by an elevated serum anion gap (exceeding 12 mEq/L), signifying the accumulation of unmeasured organic or inorganic fixed acids in plasma, such as ketoacids in diabetic ketoacidosis, lactic acid, or exogenous toxic metabolites.

\synonyms
AGMA; High Anion Gap Metabolic Acidosis

\medterm{Aniridia} A congenital, bilateral, complete or partial absence of the iris, typically caused by mutations in the PAX6 gene. It is accompanied by photophobia, nystagmus, reduced visual acuity, foveal hypoplasia, and an increased risk of early-onset glaucoma and cataracts, and may occur as part of the WAGR syndrome.

\medterm{Anisocoria} An unequal size of the pupils, the black centers of the eyes. In about 20\% of healthy people, a slight difference is natural and harmless. When abnormal, it is caused by eye injury, certain eye drops, or nerve disorders that prevent a pupil from widening or shrinking normally. Doctors test how the pupils react in bright light and darkness to identify the underlying cause.

\medterm{Anisocytosis} Greater-than-usual variation in the size of red blood cells. It is a laboratory finding rather than a disease and may occur with iron deficiency, vitamin deficiency, or other types of anemia.

\medterm{Anisometropia} A difference in the focusing power of the two eyes. It can cause unequal image size, blurred vision, eyestrain, or reduced development of vision in one eye during childhood.

\medterm{Ankle} The joint connecting the lower leg and foot. It is formed mainly by the tibia, fibula, and talus and allows the foot to move up and down while ligaments provide stability.

\medterm{Ankle Arthroplasty} Surgical replacement of the ankle joint with an artificial joint, usually for severe painful arthritis. It aims to reduce pain while preserving movement; suitability depends on bone, alignment, activity, and surrounding-joint condition.

\medterm{Ankle Arthroscopy} A procedure using a small camera and instruments inserted through small cuts around the ankle to diagnose or treat problems such as loose bone or cartilage, inflamed lining, or tissue impingement.

\medterm{Ankle Bone} Any of the bones forming the ankle joint, chiefly the lower ends of the tibia and fibula and the upper part of the talus.

\medterm{Ankle Dislocation} Complete displacement of the ankle bones from their normal joint position, usually after a severe injury. It often occurs with a fracture or ligament injury and can damage blood vessels or nerves, so it requires urgent reduction and imaging.

\synonyms
Dislocated Ankle; Talotibial Dislocation

\medterm{Ankle Fracture} A break in one or more bones forming the ankle, usually the tibia, fibula, or talus. It may cause pain, swelling, bruising, and inability to bear weight; displaced or unstable fractures need prompt assessment.

\medterm{Ankle Fractures} Breaks in one or more bones forming the ankle joint. They may also injure ligaments, blood vessels, or nerves; the seriousness depends on bone alignment, joint stability, and whether the skin is broken.

\medterm{Ankle Joint} The hinge-like joint between the lower ends of the tibia and fibula and the talus. It mainly moves the foot upward and downward and is supported by surrounding ligaments.

\medterm{Ankle Pain} Pain in or around the ankle, commonly caused by a sprain, fracture, tendon problem, arthritis, bursitis, or infection. Deformity, marked swelling, fever, numbness, or inability to bear weight requires prompt assessment.

\medterm{Ankle Replacement} Surgery that replaces the damaged surfaces of the ankle joint with an artificial joint, usually for severe arthritis causing persistent pain and disability. It aims to reduce pain and preserve ankle movement.

\medterm{Ankle Sprain} An injury in which one or more ankle ligaments are stretched or torn, most often on the outer side after the foot twists inward. Pain, swelling, bruising, and a feeling of instability may occur.

\medterm{Ankle Swelling} Enlargement around the ankle caused by injury, inflammation, infection, vein or lymph problems, or general fluid retention. Sudden one-sided swelling with pain, warmth, fever, or breathlessness needs urgent assessment.

\medterm{Ankle-Brachial Index} A comparison of blood pressure measured at the ankle with blood pressure measured at the arm. A low result may indicate narrowed leg arteries; a very high result may mean the arteries are too stiff to measure reliably.

\medterm{Ankle-Foot Orthosis} A brace that supports the ankle and foot and controls alignment or movement during standing and walking. It may improve walking or prevent foot drop caused by weakness, spasticity, instability, or deformity.

\medterm{Ankyloblepharon} Abnormal sticking or joining of the upper and lower eyelid edges. It may be present at birth or follow inflammation, burns, injury, or surgery and can partly block vision when severe.

\medterm{Ankyloglossia} Alternate name; see \textit{Tongue-Tie}.

\medterm{Ankylose} To become stiff or unable to move at a joint, sometimes because scar tissue or bone joins the joint surfaces.

\medterm{Ankylosing} Causing or involving progressive stiffness or loss of movement in a joint, sometimes because scar tissue or bone joins the joint surfaces.

\medterm{Ankylosing Spondylitis} A long-term inflammatory arthritis that mainly affects the sacroiliac joints, where the spine joins the pelvis, and the joints and ligaments of the spine. It usually causes lower-back or hip pain and stiffness that is worse after rest or inactivity and may improve with movement. Inflammation can reduce spinal movement and, in severe cases, cause the vertebrae to fuse, making the spine rigid. The hips, ribs, knees, or feet may also be affected. Some people develop inflammation of the eye (uveitis), psoriasis, or inflammatory bowel disease. The risk is higher in people with certain genetic factors, including HLA-B27, but having HLA-B27 does not by itself cause the disease.

\synonyms
Bekhterev Disease; Axial Spondyloarthritis; Marie-Strümpell Disease

\medterm{Ankylosis} Abnormal reduction or complete loss of movement in a joint. It may result from fibrous scar tissue, bone fusion, injury, infection, chronic inflammation, or surgery. Fibrous ankylosis involves scar tissue, while bony ankylosis involves fusion of neighboring bones. In dentistry, ankylosis may describe fusion of a tooth root to the surrounding jawbone.

\medterm{Ankyrin} A protein that connects selected cell-surface proteins to the cell’s internal support framework. It helps red blood cells keep their shape; inherited changes can make the cells fragile and cause anemia.

\medterm{Annihilation} The conversion of a particle and its opposite partner into energy when they collide. In PET scanning, a positron meets an electron and produces two gamma rays that travel in opposite directions to create the image.

\medterm{Annular Pancreas} A birth defect in which pancreatic tissue forms a ring around part of the duodenum, the first section of the small intestine. The ring may narrow the intestine and cause vomiting, abdominal swelling, or feeding problems.

\medterm{Annulus} A ring-shaped structure or band in the body. Examples include the tough outer ring of an intervertebral disc and the fibrous rings surrounding the heart valves.

\medterm{Annulus Fibrosus} The tough, outer fibrocartilaginous ring of an intervertebral disc composed of concentric lamellae of collagen fibers, which encircles and restrains the central gel-like nucleus pulposus and absorbs rotational and compressive mechanical loads.

\synonyms
Anulus Fibrosus

\medterm{ANOCA} Chest pain or other angina symptoms when the major heart arteries do not have a severe blockage. The symptoms may result from poor function of the heart’s small vessels or temporary spasm of a coronary artery.

\medterm{Anodontia} Complete absence of teeth from birth, caused by failure of tooth development. It is rare and may occur with inherited conditions affecting the skin, hair, nails, and teeth.

\medterm{Anogenital Tract} The connected region containing the anus, tissues around the anus, genital organs, and nearby skin and moist linings.

\medterm{Anomalous Anatomy} A difference from the usual structure, number, position, origin, or pathway of an organ, blood vessel, duct, or other body part. It may cause no symptoms or may affect health, diagnosis, or surgery.

\medterm{Anomalous Coronary Artery} A birth defect in which one of the heart’s arteries has an unusual starting point or path. Some arteries pass between the aorta and the pulmonary artery and may be squeezed during hard exercise. This can reduce blood flow to the heart and may cause chest pain, fainting, abnormal heart rhythms, or rarely sudden cardiac death.

\synonyms
AAOCA; Anomalous Aortic Origin of a Coronary Artery; Congenital Coronary Artery Anomaly; ACAOS

\medterm{Anomaly} A difference from the usual structure, development, function, or measurement of the body. It may be harmless, or it may cause symptoms or affect health.

\medterm{Anomaly Scan} A detailed ultrasound examination usually performed around 18–22 weeks of pregnancy to assess fetal growth and organs, the placenta, and amniotic fluid, and to screen for major structural birth defects.

\synonyms
Mid-Pregnancy Ultrasound; 20-Week Scan; Fetal Anomaly Scan

\medterm{Anomia} Difficulty finding or saying a familiar word, especially the name of an object or person, even though its meaning is known. It can occur after a stroke or with other brain or language disorders.

\medterm{Anophthalmia} A birth defect in which one or both eyes fail to develop. It may occur with other birth defects and can affect growth of the eye socket and face.

\medterm{Anorchia} Complete absence of testicular tissue, present from birth or caused by loss of testes that had formed. It causes very low or absent testosterone production and infertility unless appropriate specialist care is provided.

\medterm{Anorectal Abscess} A collection of pus in tissues around the anus or rectum, usually caused by an infected anal gland. It may cause severe pain, swelling, fever, or difficulty sitting and often needs prompt drainage.

\medterm{Anorectal Fistula} An abnormal tunnel between the anus or rectum and nearby skin or another organ. It often develops after an abscess or infection and may cause persistent drainage, pain, or repeated abscesses.

\synonyms
Anal Fistula

\medterm{Anorectal Infection} Infection of the anus, rectum, or nearby tissues. It may cause pain, swelling, discharge, fever, or difficulty passing stool and can lead to an abscess that requires drainage.

\medterm{Anorectal Malformation} A birth defect in which the anus or rectum does not form normally or does not connect correctly to the skin. The opening may be absent, misplaced, or abnormally connected to another organ.

\medterm{Anorectal Manometry} A test that measures pressure and sensation in the rectum and anal canal. A thin pressure-sensitive tube assesses how well the muscles relax and tighten and may help investigate constipation or loss of bowel control.

\medterm{Anorexia} Loss or reduction of appetite. It can result from illness, pain, medicines, emotional distress, or other causes and is different from anorexia nervosa, an eating disorder.

\medterm{Anorexia Nervosa} A serious eating disorder involving persistent restriction of food, intense fear of gaining weight, and a distorted view of body size or shape. It can cause severe weight loss, malnutrition, slow heart rate, low blood pressure, and organ damage.

\synonyms
Anorexia; Restrictive Eating Disorder

\medterm{Anorgasmia} Persistent or repeated difficulty reaching orgasm or inability to reach orgasm despite adequate sexual stimulation. It may result from medicines, illness, hormonal changes, pain, stress, trauma, or relationship factors.

\medterm{Anorgasmy} Another spelling of anorgasmia: persistent or repeated inability to reach orgasm despite adequate sexual stimulation. Possible causes include medicines, illness, hormonal changes, stress, trauma, or relationship factors.

\medterm{Anoscopy} Examination of the anal canal using a short, lighted tube called an anoscope. It helps investigate rectal bleeding, haemorrhoids, fissures, fistulas, pain, and other problems near the anus.

\medterm{Anoscopy And High-Resolution Anoscopy} Examination of the anal canal and lower rectum with a short viewing tube. High-resolution anoscopy adds magnification and special lighting to identify and assess abnormal tissue, including changes related to HPV.

\medterm{Anosmia} Complete loss of the sense of smell. It may follow a viral infection, blocked or inflamed nasal passages, head injury, medicines, toxic exposure, or a neurological or inherited disorder.

\medterm{Anosodiaphoria} A neurological state of apparent emotional indifference or unconcern toward severe physical disability, such as hemiplegia, despite intellectual awareness and verbal acknowledgement of the deficit. It is most commonly associated with non-dominant (right) hemisphere parietal lobe lesions.

\medterm{Anosognosia} Lack of awareness of, or inability to recognize, one’s own illness or neurological impairment. It can occur after certain brain injuries or strokes and is not simply deliberate denial.

\medterm{Anotia} A birth defect in which the outer ear is completely absent. The ear canal and middle-ear structures may also be abnormal, causing reduced hearing on the affected side.

\medterm{Anovulation} Failure to release an egg from an ovary during a menstrual cycle. It can cause irregular or absent periods, abnormal bleeding, or infertility and may result from hormonal disorders, low body weight, stress, excessive exercise, or polycystic ovary syndrome.

\medterm{Anoxia} Complete lack of oxygen reaching body tissues. It can quickly damage or kill cells, especially in the brain, heart, and kidneys, and may result from blocked blood flow, breathing failure, or severe poisoning.

\medterm{Anrep Effect} An intrinsic physiological autoregulation mechanism of the myocardium whereby an abrupt increase in ventricular afterload causes an initial stretch followed by a sustained, positive inotropic increase in ventricular contractility over the ensuing several minutes.

\medterm{Anserine Bursa} A synovial fluid-filled sac situated on the proximal medial aspect of the tibia, interposed between the tibial collateral ligament and the conjoined tendons of the sartorius, gracilis, and semitendinosus muscles; inflammation causes pes anserine bursitis.

\synonyms
Pes Anserinus Bursa

\medterm{Antacids} Medicines that neutralise stomach acid and provide short-term relief from heartburn, indigestion, and acid reflux. Different products contain different minerals and may interact with other medicines or alter bowel movements.

\medterm{Antagonist} A drug or other substance that reduces or blocks the effect of another substance. A receptor antagonist binds to a receptor without activating it and reduces the response to an agonist, often by preventing the agonist from binding or signaling.

\medterm{Antagonist Muscle} A muscle that produces the opposite movement to the main working muscle and helps control or reverse that movement.

\medterm{Antalgic Gait} An abnormal, compensatory walking pattern adopted to avoid or minimize pain in a weight-bearing lower limb or spinal structure. It is characterized by a significantly shortened stance phase on the affected, painful limb and a rapid, hurried transition to the contralateral unaffected limb, frequently accompanied by a compensatory trunk lean toward the affected side to reduce joint contact forces.
\synonyms{Pain-Relieving Gait; Limping Gait; Antalgic Walk}

\medterm{Antebrachium} The anatomical region of the upper limb situated between the elbow and the wrist, consisting of the radius and ulna surrounded by anterior and posterior muscular compartments.

\synonyms
Forearm

\medterm{Antecubital Fossa} The shallow hollow at the front of the elbow. It contains the biceps tendon, a major artery, and a major nerve and is commonly used to draw blood or measure blood pressure.

\medterm{Antegrade} Moving or occurring in the normal forward direction, such as blood flowing away from the heart or an electrical signal travelling through its usual pathway.

\medterm{Antemortem Injury} An injury that occurred before death. It may show bleeding into living tissue, clot formation, inflammation, swelling, or early healing, which can distinguish it from damage occurring after death.

\medterm{Antenatal} Relating to the period before birth or to care during pregnancy. Antenatal care monitors the pregnant person and fetus, screens for complications, and provides health advice and support.

\synonyms
Prenatal; Antepartum

\medterm{Antepartum Fetal Surveillance} Tests used before birth to check fetal well-being, especially in higher-risk pregnancies. They may include monitoring fetal movements, the fetal heart rate, ultrasound assessment, and blood-flow measurements.

\medterm{Antepartum Hemorrhage} Bleeding from the vagina during the second half of pregnancy and before birth. Causes include placenta previa, placental abruption, and problems with the cervix or vagina; any significant bleeding needs urgent assessment.

\medterm{Anterior} Relating to the front of the body or positioned nearer the front than another structure. The term is used to describe body parts, scan findings, injuries, and disease locations.

\medterm{Anterior Cerebral Artery Syndrome} A stroke syndrome resulting from occlusion or ischemia of the anterior cerebral artery, classically causing contralateral motor weakness and sensory loss predominantly affecting the lower extremity more than the upper extremity and face. Additional findings may include abulia, urinary incontinence, and frontal release signs due to medial frontal lobe involvement.

\synonyms
ACA Syndrome; ACA Infarction

\medterm{Anterior Chamber} The fluid-filled space between the cornea and iris at the front of the eye. It contains aqueous fluid, which nourishes the cornea and lens and helps maintain eye pressure; impaired drainage can contribute to glaucoma.

\medterm{Anterior Chamber Angle} The area where the cornea and iris meet at the front of the eye. It contains the main drainage tissue for aqueous fluid; blockage here can raise eye pressure and cause glaucoma.

\medterm{Anterior Commissure} A transverse bundle of myelinated nerve fibers crossing the midline anterior to the columns of the fornix and above the lamina terminalis, interconnecting the olfactory bulbs, amygdalae, and middle and inferior temporal gyri across the cerebral hemispheres.

\medterm{Anterior Cord Syndrome} A spinal-cord injury that weakens movement and reduces pain and temperature sensation below the injury, while vibration and awareness of limb position may be relatively preserved.

\medterm{Anterior Cruciate Ligament} A strong band of tissue inside the knee that connects the thigh bone to the shin bone. It helps prevent the shin bone from sliding too far forward and controls knee rotation.

\medterm{Anterior Cruciate Ligament Injury} Damage or tearing of the main stabilizing ligament inside the knee, often after sudden turning, stopping, landing, or overextension. It may cause a popping sensation, rapid swelling, pain, and a feeling that the knee gives way.

\medterm{Anterior Cruciate Ligament Tear} A partial or complete tear of the ligament inside the knee that limits forward sliding and rotation of the shin bone. It commonly causes a pop, rapid swelling, pain, and knee instability after a sudden twist or landing.

\medterm{Anterior Drawer Sign} A physical check for a torn knee ligament (anterior cruciate ligament). With the knee bent, the doctor gently pulls the shin forward; moving too far forward means the ligament is likely torn or stretched.

\medterm{Anterior Drawer Test} A clinical orthopedic physical examination maneuver used to assess the integrity of the anterior cruciate ligament (ACL) of the knee or the anterior talofibular ligament of the ankle. Excessive anterior translation of the tibia relative to the femur with a soft or mushy endpoint indicates an ACL tear.

\synonyms
Anterior Drawer Sign

\medterm{Anterior Interosseous Nerve} A pure motor branch of the median nerve arising in the proximal forearm that courses deeply along the anterior surface of the interosseous membrane to innervate the flexor pollicis longus, the lateral half of flexor digitorum profundus, and the pronator quadratus; damage produces the classic inability to make the "OK" sign.

\synonyms
AIN

\medterm{Anterior Interosseous Nerve Syndrome} A pure motor neuropathy of the forearm resulting from compression or neuritis of the anterior interosseous nerve (a motor branch of the median nerve). Causes weakness of the flexor pollicis longus, pronator quadratus, and the radial half of the flexor digitorum profundus, classically producing an inability to form the \"OK\" sign (pinching index finger and thumb pulp-to-pulp instead of tip-to-tip).

\synonyms
Kiloh-Nevin Syndrome; AIN Neuropathy

\medterm{Anterior Knee Pain} Pain at the front of or behind the kneecap, often worse with stairs, squatting, running, or prolonged sitting. Common causes include irritation around the kneecap, tendon problems, arthritis, and injury.

\medterm{Anterior Mediastinum} The front part of the lower mediastinum, located behind the sternum and in front of the heart sac. It contains fat, connective tissue, lymph nodes, and part of the thymus.

\medterm{Anterior Midline} The imaginary vertical line running down the center of the front of the body, separating the right and left sides.

\medterm{Anterior Myocardial Infarction} A heart attack affecting the front wall of the main pumping chamber, usually because the left anterior descending coronary artery is blocked. It can seriously weaken the heart and requires emergency treatment.

\synonyms
Anterior Heart Attack; Anterior MI

\medterm{Anterior Pituitary} The front part of the pituitary gland. It releases hormones that control the adrenal and thyroid glands, growth, milk production, egg or sperm production, and reproductive function.

\medterm{Anterior Talofibular Ligament} A short, flat ligament on the lateral aspect of the ankle that spans from the anterior margin of the lateral malleolus of the fibula to the neck of the talus. It is the primary stabilizer against anterior talar displacement and excessive inversion, and represents the most frequently torn ligament in lateral ankle sprains.

\synonyms
ATFL

\medterm{Anterior Tibial Artery} An artery that carries blood down the front of the leg and continues onto the top of the foot as the dorsalis pedis artery. Its pulse can usually be felt near the front of the ankle.

\medterm{Anterior Uveitis} Inflammation in the front part of the eye, usually involving the iris and sometimes nearby tissue. It can cause eye pain, redness, light sensitivity, blurred vision, and floating spots.

\synonyms
Iritis; Iridocyclitis

\medterm{Anterior Vaginal Prolapse} Bulging of the anterior vaginal wall into the vaginal vault, commonly resulting from herniation of the urinary bladder (cystocele) due to attenuation or tearing of the pubocervical fascia and pelvic floor musculature. Symptoms may include a palpable vaginal bulge, pelvic heaviness, urinary frequency, urgency, incomplete emptying, or stress incontinence.

\synonyms
Cystocele; Bladder Prolapse; Prolapsed Bladder; Anterior Wall Prolapse

\medterm{Anterior Vaginal Wall Repair} Surgery to strengthen the front wall of the vagina when the bladder bulges into it. The procedure restores support using the person’s own tissues; the approach and risks depend on the prolapse and overall health.

\medterm{Anterior/Posterior Drawer Test} A physical examination of knee stability. Pulling the shin bone forward tests the anterior cruciate ligament; pushing it backward tests the posterior cruciate ligament.

\medterm{Anterograde Amnesia} A selective impairment of memory in which an individual is unable to form, encode, or consolidate new long-term declarative (episodic and semantic) memories after the onset of a neurological insult, while remote memories formed prior to the event remain largely intact. It typically results from bilateral damage to the medial temporal lobes, hippocampi, mammillary bodies, or diencephalic structures following traumatic brain injury, stroke, herpes simplex encephalitis, or Korsakoff syndrome.
\synonyms{Post-Traumatic Amnesia; Anterograde Memory Loss; Forward-Acting Amnesia}

\medterm{Anterolisthesis} Forward slipping of one spinal bone relative to the bone below it. It may result from wear, a stress fracture, injury, or a problem present from birth and may cause back pain or nerve symptoms.

\medterm{Anteroposterior} Directed from the front of the body toward the back. In an anteroposterior X-ray view, the beam enters through the front and leaves through the back.

\medterm{Anthocyanin} A natural plant pigment that gives many fruits, vegetables, and flowers red, purple, or blue colours. It has antioxidant effects in laboratory studies and is found in foods such as berries, grapes, and red cabbage.

\medterm{Anthracosis} A benign, non-fibrotic form of pneumoconiosis characterized by the asymptomatic accumulation of inhaled carbon and soot pigment in the lung parenchyma, alveolar macrophages, and draining hilar lymph nodes, commonly seen in urban dwellers and cigarette smokers.

\synonyms
Black Lung Pigmentation; Carbon Pigmentation

\medterm{Anthrax} A serious infection caused by the bacterium \textit{Bacillus anthracis}. People can become infected through contact with infected animals or animal products, breathing in spores, or eating contaminated food; skin, lung, or intestinal disease may result.

\medterm{Anthrax Blood Test} A laboratory test that looks for the anthrax bacterium or the body’s immune response to it in blood or another sample. It helps confirm suspected anthrax, which can affect the skin, lungs, or digestive tract.

\medterm{Anthropometry} The measurement of body size and proportions, such as height, weight, waist size, limb length, and skinfold thickness. These measurements help assess growth, nutrition, and body composition.

\medterm{Anti-Anxiety Agents} Medicines used to reduce excessive anxiety or panic. Different medicines work in different ways and may cause sleepiness, dependence, or other side effects; the choice depends on the person and the condition.

\synonyms
Anxiolytics

\medterm{Anti-CCP} A blood-test antibody marker associated with rheumatoid arthritis. A positive result supports the diagnosis when symptoms fit, but it does not prove that rheumatoid arthritis is present or predict its course by itself.

\medterm{Anti-CCP Antibody} An antibody measured in blood that targets proteins changed by a process called citrullination. It can support a diagnosis of rheumatoid arthritis, but results must be interpreted with symptoms and other tests.

\medterm{Anti-Citrullinated Peptide Antibody} Alternate index label; see \textit{Anti-Cyclic Citrullinated Peptide Antibody}.

\medterm{Anti-Cyclic Citrullinated Peptide Antibody} An autoantibody directed against citrulline-containing peptide antigens, possessing high diagnostic specificity (greater than 95 percent) for rheumatoid arthritis and correlating with aggressive, erosive joint disease.

\synonyms
Anti-CCP Antibody; ACPA; Anti-Citrullinated Peptide Antibody; Anti-Cyclic Citrullinated Peptide; Cyclic Citrullinated Peptide Antibody

\medterm{Anti-D} An injection containing antibodies given to an RhD-negative pregnant person after certain pregnancy events, or around birth, to prevent the person’s immune system making antibodies against RhD-positive fetal blood. This helps protect a current or future pregnancy from fetal red-cell destruction.

\medterm{Anti-DNase B Blood Test} A blood test that measures the body’s antibodies against a protein made by group A streptococcal bacteria. It can support evidence of a recent strep infection when doctors investigate rheumatic fever, kidney inflammation, or another complication.

\medterm{Anti-Double-Stranded DNA Antibody} An antinuclear autoantibody specifically directed against the native double-helical structure of deoxyribonucleic acid, highly specific for systemic lupus erythematosus and closely correlating with lupus nephritis disease activity and flares.

\synonyms
Anti-dsDNA Antibody; dsDNA Antibody; Anti-Double-Stranded DNA; Double-Stranded DNA Antibody

\medterm{Anti-GBM Disease} An autoimmune disease in which antibodies attack supporting layers in the kidneys and sometimes the lungs. It can rapidly cause kidney failure, coughing up blood, or bleeding into the lungs and requires urgent specialist care.

\medterm{Anti-Glomerular Basement Membrane Blood Test} A blood test that detects antibodies attacking the kidney’s filtering membranes and sometimes the lungs’ supporting membranes. A positive result supports anti-GBM disease when symptoms and other test results are consistent.

\medterm{Anti-Glomerular Basement Membrane Disease} An autoimmune disease in which antibodies damage the kidney’s filtering membranes and sometimes the lungs. It may cause rapidly worsening kidney function, blood in the urine, coughing up blood, or lung bleeding.

\medterm{Anti-HMGCR Myopathy} A rare autoimmune muscle disease in which antibodies attack a muscle enzyme. It causes ongoing weakness and muscle injury, especially around the hips, thighs, shoulders, and upper arms, and may begin after treatment with a statin medicine.
\synonyms
Statin-Induced Immune-Mediated Necrotizing Myopathy; Anti-HMG-CoA Reductase Myopathy

\medterm{Anti-IgE Therapy} Treatment using an antibody that binds free immunoglobulin E, an allergy-related immune protein. This reduces activation of allergy cells and may lessen symptoms and attacks in selected allergic asthma and other allergic conditions.

\medterm{Anti-IL4 Receptor Therapy} Treatment that blocks a cell receptor used by interleukin-4 and interleukin-13, immune signals that promote allergic inflammation. It can reduce airway inflammation in selected asthma and other allergic diseases.

\medterm{Anti-IL5 Therapy} Treatment that blocks interleukin-5 or its receptor, reducing the production or survival of eosinophils, a type of white blood cell. It is used for selected eosinophilic diseases, including some severe asthma.

\medterm{Anti-Inflammatory} Pertaining to any drug, substance, or therapeutic intervention that reduces, counteracts, or suppresses inflammation, vascular permeability, tissue edema, and associated pain. Major therapeutic classes include nonsteroidal anti-inflammatory drugs and glucocorticoids.

\medterm{Anti-Insulin Antibody Test} A blood test that detects antibodies against insulin. These antibodies can bind insulin and release it unpredictably, causing low blood sugar; the test helps investigate unexplained hypoglycaemia and insulin autoimmune syndrome.

\medterm{Anti-MDA5 Dermatomyositis} A form of dermatomyositis associated with antibodies against a protein called MDA5. It often causes distinctive skin changes with little muscle weakness but carries a high risk of rapidly worsening inflammation and scarring in the lungs.

\medterm{Anti-Mullerian Hormone} A hormone made by developing ovarian follicles and by cells in the testes before birth. Its blood level can help estimate ovarian follicle number or assess some reproductive and developmental conditions, but it does not measure fertility by itself.

\medterm{Anti-NMDA Receptor Encephalitis} A serious immune attack on the brain involving a receptor important for communication between nerve cells. It may cause sudden behavior or psychiatric changes, memory loss, seizures, abnormal movements, reduced consciousness, and unstable heart rate or blood pressure. Urgent hospital treatment is needed.

\synonyms
Anti-N-Methyl-D-Aspartate Receptor Encephalitis; NMDA Receptor Encephalitis

\medterm{Anti-Reflux Surgery} An operation to reduce the backward flow of stomach contents into the oesophagus by strengthening or reshaping the valve between the oesophagus and stomach. It may be considered for selected people with severe reflux or a hiatus hernia.

\medterm{Anti-Reflux Surgery, Pediatric} Surgical reconstruction of the gastroesophageal junction—most commonly via laparoscopic Nissen fundoplication—in infants and children with medically refractory gastroesophageal reflux disease, recurrent aspiration, or severe failure to thrive.

\synonyms
Pediatric Fundoplication; Pediatric Anti-Reflux Surgery; Nissen Fundoplication in Children

\medterm{Anti-Rust Product Poisoning} Illness caused by swallowing, breathing, or skin exposure to a product used to remove or prevent rust. It may burn the mouth and digestive tract or damage the lungs and other organs; significant exposure requires urgent medical treatment.

\medterm{Anti-Smooth Muscle Antibody} An antibody that reacts with proteins in smooth muscle cells. A blood test may support autoimmune hepatitis, but low levels can also occur with other autoimmune diseases or infections and are not diagnostic alone.

\medterm{Anti-SRP Myopathy} A rare autoimmune muscle disease caused by antibodies against a protein involved in making new proteins. It causes rapidly progressive weakness, especially around the hips, shoulders, and neck, and may cause difficulty swallowing.

\synonyms
Anti-Signal Recognition Particle Myopathy; SRP-Associated Necrotizing Myopathy

\medterm{Anti-Thyroid Peroxidase Antibody} An autoantibody directed against thyroid peroxidase, the key enzyme required for thyroid hormone synthesis. Elevated titers are the hallmark laboratory finding in Hashimoto thyroiditis and are also present in Graves disease.

\synonyms
Anti-TPO Antibody; Thyroid Peroxidase Antibody

\medterm{Anti-VEGF Injections} Injections, usually given into the eye, that block vascular endothelial growth factor, a signal that promotes abnormal blood-vessel growth and leakage. They can improve or preserve vision in several retinal diseases.

\medterm{Antiandrogen} A medicine or substance that lowers the production or blocks the effects of androgen hormones such as testosterone. It is used for selected prostate conditions, androgen-related disorders, and some gender-affirming treatments.

\medterm{Antianginal} A medicine used to prevent or relieve angina, the chest discomfort caused by reduced blood flow to heart muscle. Different medicines lower the heart’s workload, widen blood vessels, or improve coronary blood flow.

\medterm{Antiarrhythmic} A medication or therapeutic intervention designed to prevent, suppress, or terminate abnormal cardiac rhythms and restore or maintain normal sinus rhythm. Pharmacological agents are conventionally categorized according to the Vaughan Williams classification into classes I through IV.

\medterm{Antibacterial} Able to kill bacteria or stop them from growing. The term may describe a medicine, disinfectant, surface treatment, or part of the body’s immune defence.

\medterm{Antibiogram} A report showing which antibiotics are likely to work against bacteria grown from patient samples in a laboratory. It helps clinicians choose treatment and monitor local patterns of antibiotic resistance.

\medterm{Antibiotic Misuse in Viral URIs} The inappropriate prescription or self-administration of antibacterial agents for acute viral upper respiratory tract infections (such as the common cold or influenza), which confers no therapeutic benefit and drives antimicrobial resistance and adverse drug events.

\synonyms
Antibiotic Stewardship in URI; Inappropriate Antibiotic Use in Colds; Antibiotic Overuse in Flu

\medterm{Antibiotic-Resistant Bacterial Infection} An infection caused by bacteria that are not adequately stopped by one or more antibiotics normally active against them. Resistance may be intrinsic or acquired and can limit effective treatment options.

\medterm{Antibiotics} Medicines that kill bacteria or stop them from multiplying. They treat susceptible bacterial infections but do not treat viral infections such as colds or influenza.

\medterm{Antibodies} Proteins made by immune cells that recognize and attach to specific foreign or abnormal substances. They can block germs or toxins, mark them for destruction, and activate other parts of the immune system.

\medterm{Antibody} A protein made by the immune system that attaches specifically to a foreign substance, germ, toxin, or abnormal cell and may block it or mark it for removal.

\medterm{Antibody Engineering} Designing or changing antibodies so they bind more accurately, last longer, cause fewer unwanted immune reactions, or work better as medicines. Researchers may alter their target, size, immune effects, or attachment to another medicine.

\medterm{Antibody Therapy} Treatment using specially prepared antibodies to block, activate, or target a particular molecule or cell. It can be used for cancer, autoimmune disease, infection, inflammation, and other conditions.

\synonyms
Immunotherapy

\medterm{Antibody Titer Blood Test} A blood test that estimates the amount of a specific antibody by repeatedly diluting the sample until the antibody can no longer be detected. Results may help show previous exposure, immunity, or a recent infection, depending on the antibody tested.

\medterm{Antibody-Dependent Cell-Mediated Cytotox.} An immune process in which antibodies coat an infected, abnormal, or cancer cell and attract killer immune cells. The killer cells then damage and destroy the marked cell.

\medterm{Antibody-Drug Conjugate} A treatment in which an antibody carrying a linked medicine binds to a target on a cell and delivers the medicine near or inside that cell. It is used for selected cancers and can still damage healthy tissues or cause medicine-specific toxicities.

\medterm{Anticholinergic} Blocking the action of acetylcholine, a chemical messenger used by nerves. These medicines can reduce secretions and muscle spasms but may cause dry mouth, blurred vision, constipation, difficulty urinating, fast heartbeat, or confusion.

\medterm{Anticholinergic Burden} The combined anticholinergic effect of one or more medicines. A high burden can cause dry mouth, constipation, blurred vision, difficulty urinating, confusion, or delirium, especially in older adults.

\medterm{Anticholinergic Toxicity} Excess blockade of muscarinic acetylcholine receptors, producing a hot, flushed, dry, tachycardic, mydriatic, delirious patient with urinary retention and reduced bowel sounds. Treatment is supportive with cooling and benzodiazepines; physostigmine is reserved for selected severe cases after toxicology assessment.

\medterm{Anticholinergic Toxidrome} A pattern of poisoning caused by excessive blockage of acetylcholine signals. It may cause dry skin and mouth, large pupils, fast heartbeat, inability to urinate, reduced bowel activity, overheating, agitation, or delirium.

\medterm{Anticoagulants and Thrombolytics in Pregnancy} Medicines that prevent clot formation or dissolve selected clots during pregnancy or after delivery. Choice and timing depend on maternal indication, placental transfer, bleeding risk, delivery plans, and whether neuraxial anesthesia may be used; specialist management is required.

\medterm{Anticoagulation} Treatment that reduces the formation or growth of blood clots by slowing the clotting process. It is used for conditions such as clots in veins, irregular heart rhythm, and selected heart or blood-vessel disorders.

\medterm{Anticodon} A group of three RNA building blocks on a transfer-RNA molecule. It pairs with a matching sequence on messenger RNA so the cell adds the correct amino acid while making a protein.

\medterm{Anticonvulsants} Medicines used to prevent or control seizures. Some are also used for nerve pain, migraine, or mood disorders; their effects, interactions, and side effects differ between medicines.

\medterm{Antideformity Position} A planned position for a limb or joint that helps prevent or limit shortening, stiffness, or deformity after injury, surgery, burns, nerve damage, or prolonged immobility.

\medterm{Antidepressants} Medicines used mainly to treat depression. Different types change the activity of chemical messengers in the brain and may also treat anxiety, nerve pain, obsessive thoughts, or other conditions.

\medterm{Antidiabetic} Relating to a medicine or treatment that lowers high blood sugar or helps control diabetes. Different treatments increase insulin action, reduce glucose production, or remove excess glucose from the body.

\medterm{Antidiuretic Hormone} A hormone made in the brain and released by the pituitary gland that helps the kidneys conserve water. It also narrows blood vessels when blood volume or blood pressure is low. It is also called vasopressin.

\medterm{Antidiuretic Hormone Blood Test} A blood test that measures vasopressin, the hormone that helps the kidneys conserve water. It may help investigate excessive thirst, very dilute urine, inappropriate water retention, or other disorders of water balance.

\medterm{Antidiuretic Hormone Physiology} Vasopressin is made in the hypothalamus and released from the posterior pituitary when the blood becomes too concentrated or blood volume falls. It makes kidney tubules retain water and can narrow blood vessels.

\medterm{Antidiuretic Hormone Resistance} Failure of the kidneys to respond properly to vasopressin. It causes production of large amounts of dilute urine and intense thirst and may result from medicines, mineral imbalance, kidney disease, or inherited changes.

\medterm{Antidote} A substance that counteracts or limits the harmful effects of a poison. The correct antidote depends on the substance involved; some antidotes block the poison, remove it, or restore a chemical process disrupted by it.

\medterm{Antiemetics} Medicines used to prevent or relieve nausea and vomiting. Different types act on chemical signals in the brain or digestive tract, and the choice depends on the cause, age, pregnancy, other medicines, and risk of side effects.

\medterm{Antiepileptic Drugs} Medicines used to prevent or control seizures by reducing excessive electrical activity in the brain. The best choice depends on the seizure type, age, other conditions, medicines, and possible side effects.

\medterm{Antifreeze Poisoning} Poisoning after swallowing or otherwise being exposed to antifreeze, especially products containing ethylene glycol. It may first cause intoxication and vomiting, then severe acid buildup, kidney failure, brain injury, or death; severe poisoning is an emergency requiring prompt medical treatment.

\medterm{Antifungal Prophylaxis} Use of an antifungal medicine to prevent a serious fungal infection in someone at high risk, such as a person with prolonged low white-cell counts, a stem-cell transplant, or severe immune suppression.

\medterm{Antifungal Resistance} The ability of a fungus to survive a medicine that would normally stop its growth or kill it. Resistance may be present naturally or develop after genetic changes or repeated exposure to antifungal medicines.

\medterm{Antigen} A substance or part of a substance that the immune system can recognize. It may come from a germ, toxin, transplanted tissue, or abnormal cell and can trigger antibodies or other immune responses.

\medterm{Antigen Detection Test} A test that looks for a protein or other marker from a specific germ in a patient sample. It is often rapid, but it may miss an infection when the amount of germ material is small.

\medterm{Antigen Presentation} The process by which an immune or other cell displays small pieces of proteins on its surface for T cells to inspect. This helps the immune system recognize infections, abnormal cells, transplanted tissue, and sometimes the body’s own tissues.

\medterm{Antigen Receptor} A protein on an immune cell that recognizes a particular substance or protein fragment. B-cell receptors can bind whole substances, while T-cell receptors recognize fragments displayed by other cells.

\medterm{Antigen-Presenting Cell} An immune cell that breaks proteins into small pieces and displays them on its surface so T cells can recognize them. Important antigen-presenting cells include dendritic cells, macrophages, and B cells.

\synonyms
APC

\medterm{Antigenic Drift} The gradual change in a virus’s surface proteins as it multiplies. These small changes can help viruses, especially influenza viruses, partly evade existing immunity and contribute to repeated seasonal outbreaks.

\medterm{Antigenic Shift} A sudden major change in a virus’s surface proteins, usually when two different influenza A viruses exchange genetic material inside the same cell. It can produce a virus to which many people have little immunity.

\medterm{Antigenic Variation} Changes in a germ’s surface markers that help it avoid recognition by the immune system. Variation may happen gradually through mutations or suddenly when genetic material is exchanged between different strains.

\medterm{Antigout} A medicine used to prevent or treat gout. Some lower uric-acid levels over time, while others reduce the inflammation and pain of an acute gout attack.

\medterm{Antihistamines} Medicines that block the effects of histamine. They can relieve itching, sneezing, runny nose, hives, or other allergy symptoms; some cause drowsiness, while others act mainly on stomach acid.

\medterm{Antihistamines For Allergies} Medicines that block histamine and relieve allergy symptoms such as itching, sneezing, runny nose, watery eyes, or hives. Some cause drowsiness and interactions or side effects vary between medicines.

\medterm{Antihypertensives In Pregnancy} Medicines used to lower high blood pressure during pregnancy. The choice depends on blood-pressure severity and pregnancy stage; some commonly used blood-pressure medicines are avoided because they can harm the fetus.

\medterm{Antileukotriene} A medicine that blocks leukotrienes, chemical signals that can cause airway narrowing and inflammation. It is used in selected asthma and allergy-related conditions, often as part of a broader treatment plan.

\synonyms
Leukotriene Modifier; Leukotriene Pathway Inhibitor

\medterm{Antimetabolite} A substance that resembles a chemical needed by cells and interferes with the cell’s normal chemical reactions or DNA production. Some are used to treat cancer or suppress an overactive immune system.

\medterm{Antimicrobial} A substance that kills or limits the growth of microscopic organisms, including bacteria, viruses, fungi, or parasites. Antibiotics are antimicrobials, but they act only against bacteria.

\medterm{Antimicrobial Adverse Effects} Unwanted effects caused by medicines that treat infections. Depending on the medicine, they may include allergy, diarrhoea, liver or kidney injury, hearing problems, tendon injury, or harmful effects on heart rhythm.

\medterm{Antimicrobial Agent} A substance that kills microorganisms or stops them from growing. Antimicrobial agents include antibacterial, antiviral, antifungal, and antiparasitic medicines.

\medterm{Antimicrobial Desensitization} A carefully supervised process of giving gradually increasing doses of an infection-fighting medicine to temporarily reduce an immediate allergy response. It is considered only when the medicine is essential and suitable alternatives are unavailable.

\medterm{Antimicrobial Dosing Optimization} Choosing the infection medicine, dose, timing, route, and duration so enough medicine reaches the infection without unnecessary harm. Kidney function, liver function, body size, infection severity, and resistance are considered.

\medterm{Antimicrobial Hypersensitivity Reactions} Immune reactions to an infection-fighting medicine, ranging from an itchy rash to life-threatening anaphylaxis. The timing and symptoms vary with the medicine and type of immune response.

\medterm{Antimicrobial Pharmacodynamics} How an infection-fighting medicine affects a microorganism, including how much medicine and how much exposure are needed to stop or kill it. This information helps select an effective dose and schedule.

\medterm{Antimicrobial Pharmacokinetics} How the body absorbs, spreads, changes, and removes an infection-fighting medicine. Kidney or liver disease, age, body size, and other medicines can change how much medicine is needed and how often it is given.

\medterm{Antimicrobial Resistance} The ability of a microorganism to survive or grow despite a medicine that would normally kill it or stop its growth. Resistance can affect bacteria, viruses, fungi, or parasites and can spread between organisms.

\synonyms
AMR; Drug Resistance

\medterm{Antimicrobial Resistance Mechanisms} Ways microorganisms avoid being killed by infection medicines. They may destroy or change the medicine, alter its target, prevent it entering, or actively pump it out of the cell.

\medterm{Antimicrobial Resistance Surveillance} Ongoing collection and analysis of resistance results from patient samples. It shows how resistance changes over time and helps hospitals and public-health services choose effective treatment and limit spread.

\medterm{Antimicrobial Sensitivity Test} A laboratory test that shows which infection-fighting medicines can stop or kill a microorganism grown from a patient sample. It helps guide treatment for a confirmed infection.

\medterm{Antimicrobial Stewardship} Coordinated efforts to use infection-fighting medicines only when needed and to choose the right medicine, dose, route, and duration. This improves outcomes and limits resistance and side effects.

\medterm{Antimicrobial Susceptibility Testing} Laboratory testing that estimates whether a microorganism is likely to respond to particular infection-fighting medicines. Methods include measuring growth around medicine-containing discs or in liquid; results help select treatment.

\medterm{Antimicrobial Therapy} Use of antibacterial, antiviral, antifungal, or antiparasitic medicines to prevent or treat an infection. The choice depends on the microorganism, infection site, severity, resistance, and the person’s health.

\medterm{Antimitochondrial Antibody} An antibody that reacts with structures inside mitochondria, the parts of cells that produce energy. It is strongly associated with primary biliary cholangitis, an autoimmune bile-duct disease, but the result must be interpreted with symptoms and liver tests.

\medterm{Antimony} A metallic chemical element whose compounds can be poisonous. Some antimony compounds have been used to treat leishmaniasis; exposure can affect the heart, liver, kidneys, or nervous system.

\medterm{Antineoplastic} Inhibiting or halting the proliferation, maturation, or spread of malignant neoplastic cells. The term designates chemotherapy agents, targeted molecular therapies, immunotherapies, and hormonal agents used in cancer therapy.

\synonyms
Anticancer Agent; Antitumor Agent

\medterm{Antinuclear Antibody} A diverse family of autoantibodies directed against nucleic acids, histones, ribonucleoproteins, and other nuclear and cytoplasmic components, evaluated via indirect immunofluorescence as a primary serological screen for systemic lupus erythematosus, systemic sclerosis, mixed connective tissue disease, and Sjogren syndrome.

\synonyms
ANA; Antinuclear Antibodies; Antinuclear Autoantibodies; FANA

\medterm{Antinuclear Antibody Panel} A comprehensive laboratory panel of specific autoantibody assays (such as anti-dsDNA, anti-Smith, anti-Ro/SSA, anti-La/SSB, and anti-RNP) utilized to characterize specific autoimmune rheumatic diseases following a positive antinuclear antibody screen.

\medterm{Antioxidant} A substance that helps neutralize or limit damage from reactive molecules, including free radicals and some reactive oxygen or nitrogen species. The body makes antioxidant enzymes and obtains other antioxidants from food; high-dose supplements are not automatically beneficial and may be harmful or interact with treatment.

\medterm{Antiparietal Cell Antibody Test} A blood test that detects antibodies against cells in the stomach lining that make acid and help absorb vitamin B12. A positive result may support autoimmune gastritis or pernicious anemia but does not prove either condition alone.

\medterm{Antiparkinson} Relating to medicines used to reduce symptoms of Parkinson disease. These medicines may increase or imitate dopamine activity or change other brain signals that influence movement.

\medterm{Antiphospholipid Antibodies} Autoantibodies that attach to proteins linked with cell membranes. Persistent antibodies may increase the risk of blood clots and pregnancy complications, but they can also appear temporarily after an infection or with some medicines.

\medterm{Antiphospholipid Antibody} An antibody directed against proteins that bind to cell-membrane fats. Persistent antibodies may be associated with blood clots or pregnancy complications, especially when clinical features of antiphospholipid syndrome are present.

\synonyms
APA; aPL Antibody

\medterm{Antiphospholipid Antibody Syndrome} Alternate name; see \textit{Antiphospholipid Syndrome}.

\medterm{Antiphospholipid Syndrome} A systemic autoimmune disorder characterized by recurrent arterial, venous, or microvascular thrombosis and obstetric complications (such as recurrent miscarriages or late fetal loss), occurring in the presence of persistent antiphospholipid antibodies.

\synonyms
APS; Antiphospholipid Antibody Syndrome; Hughes Syndrome; aPL Syndrome

\medterm{Antiphospholipid Syndrome and Pregnancy} An autoimmune clotting disorder associated with recurrent pregnancy loss, placental insufficiency, fetal growth restriction, preeclampsia, thrombosis, or pregnancy morbidity. Diagnosis requires clinical events plus persistent laboratory antiphospholipid antibodies; pregnancy care is specialist-led.

\medterm{Antiplatelet} A medicine that reduces platelet activation or aggregation and lowers the formation of arterial blood clots; see \textit{Antiplatelet Agents}.

\medterm{Antiplatelet Agents} Medicines that reduce the ability of platelets to stick together and form clots. They are used in selected people to lower the risk of clots in arteries, including after some heart or blood-vessel events.

\medterm{Antiplatelet Agents, P2Y12 Inhibitors} A subclass of oral antiplatelet medications—including clopidogrel, ticagrelor, and prasugrel—that selectively block the platelet P2Y12 adenosine diphosphate receptor, preventing ADP-mediated platelet activation and arterial thrombosis.

\synonyms
P2Y12 Inhibitors; P2Y12 Receptor Antagonists; ADP Receptor Antagonists

\medterm{Antiplatelet Therapy} Treatment with medicines that make platelets less likely to stick together. It is used in selected people to reduce clots in arteries, including after some heart attacks, strokes, or stent procedures.

\medterm{Antiprotozoal} A medicine or substance that kills protozoa or stops them growing. Protozoa are microscopic parasites that can cause infections such as malaria, giardiasis, or amoebiasis.

\medterm{Antipsychotic Medicine} A medicine used to reduce psychosis, including hallucinations, delusions, severe disorganization, or agitation. Antipsychotics differ in their effects and can cause movement problems, sleepiness, metabolic changes, or other adverse effects.

\medterm{Antipyretic} A substance or drug that lowers an abnormally elevated body temperature in fever by acting upon the hypothalamic thermoregulatory center to reset the temperature set point back toward normal. Common clinical examples include acetaminophen, ibuprofen, and aspirin.

\medterm{Antiretroviral} A medicine that acts against retroviruses, including HIV. Different antiretroviral medicines block different steps by which HIV enters cells, copies its genetic material, or makes new virus particles.

\medterm{Antiretroviral Therapy} Treatment of HIV infection with a combination of medicines that block several steps in HIV reproduction. It lowers the amount of virus, protects immune function, reduces illness, and greatly lowers the chance of sexual transmission when the virus remains undetectable.

\medterm{Antisepsis} Use of a chemical product on living skin or tissue to reduce harmful microorganisms and lower infection risk. It does not make tissue completely sterile and must be used in a way that avoids injury.

\medterm{Antiseptic} A substance used on skin or other living tissue to kill or slow harmful microorganisms. Antiseptics can help reduce infection risk but may damage tissue if used incorrectly.

\medterm{Antiserum} Blood serum containing antibodies against a specific germ, toxin, or other substance. It can provide temporary immune protection or help neutralize a toxin.

\medterm{Antisocial Personality} A long-term pattern of disregarding other people’s rights, safety, rules, and responsibilities, often including deceit, impulsive behaviour, aggression, or lack of remorse.

\synonyms
Antisocial Personality Disorder

\medterm{Antisocial Personality Disorder} A mental-health disorder involving a persistent pattern of violating other people’s rights and social rules, deceitfulness, impulsivity, aggression, irresponsibility, or lack of remorse. The pattern begins by adolescence and continues into adult life.

\synonyms
ASPD; Dissocial Personality Disorder; Sociopathy

\medterm{Antispasmodic Drugs} Medicines that relax involuntary muscle and reduce cramping or spasms, most often in the digestive tract. Some are also used for bladder or bile-duct spasms; side effects depend on the medicine.

\synonyms
Antispasmodics; Spasmolytics

\medterm{Antistreptolysin O Titer} A blood test that measures antibodies against a toxin made by group A streptococcal bacteria. A raised or rising result may show a recent strep infection when investigating rheumatic fever or kidney inflammation, but does not prove an active infection.

\medterm{Antisynthetase Syndrome} An autoimmune disease in which the immune system attacks proteins needed to make other proteins. It may cause muscle inflammation and weakness, arthritis, fever, colour changes in the fingers, rough cracked skin on the hands, or lung inflammation and scarring.

\medterm{Antithrombin} A natural blood protein that slows clotting by blocking thrombin and other clotting proteins, especially factor Xa. Too little antithrombin, whether inherited or acquired, increases the risk of clots in veins.

\medterm{Antithrombin Deficiency} A condition in which the blood has too little or too little active antithrombin, a natural anti-clotting protein. It may be inherited or acquired and increases the risk of clots, especially in the veins.

\medterm{Antithrombin III Blood Test} A blood test that measures the amount or activity of antithrombin, a protein that limits clotting. A low result may indicate inherited deficiency or may occur with liver disease, kidney protein loss, severe illness, or some treatments.

\medterm{Antithrombin III Deficiency} Reduced amount or activity of antithrombin, causing increased risk of blood clots in veins. It may be inherited or result from liver disease, kidney disease, severe illness, or other causes.

\medterm{Antithymocyte Globulin} A preparation of antibodies that removes or suppresses T cells, immune cells involved in transplant rejection and some autoimmune reactions. It is used in selected transplants and blood-cell disorders and can cause allergic reactions, low blood counts, and infections.

\synonyms
Rabbit Antithymocyte Globulin
rATG
Thymoglobulin

\medterm{Antithyroglobulin Antibody Test} A blood test that detects antibodies against thyroglobulin, a protein used to make thyroid hormones. Raised levels may support Hashimoto thyroiditis or other autoimmune thyroid disease but do not determine thyroid function by themselves.

\medterm{Antitoxin} An antibody or antibody preparation that binds to and neutralizes a toxin. It can provide temporary protection or treatment for selected poisonings and infections that produce toxins.

\medterm{Antitubercular Agents} Medicines used to treat tuberculosis, an infection caused by \textit{Mycobacterium tuberculosis}. Several medicines are used together for months to prevent resistance, with monitoring for liver injury, nerve problems, vision changes, and other side effects.

\medterm{Antitussive} An agent that suppresses or relieves coughing, either centrally by dampening cough centers within the medulla oblongata (such as dextromethorphan or codeine) or peripherally by desensitizing sensory stretch receptors in the respiratory tract.

\synonyms
Cough Suppressant

\medterm{Anton Syndrome} A rare neurological condition characterized by cortical blindness accompanied by visual anosognosia (complete denial of blindness) and visual confabulation. Patients genuinely behave and report as if they can see, despite profound vision loss, classically caused by bilateral occipital lobe infarctions.

\synonyms
Anton-Babinski Syndrome; Visual Anosognosia

\medterm{Antral Follicle Count} The number of small developing egg-containing sacs seen in both ovaries, usually measured by vaginal ultrasound early in the menstrual cycle. It helps estimate ovarian reserve and response to fertility medicines but does not by itself predict natural fertility.

\synonyms
AFC; Basal Antral Follicle Count

\medterm{Antral Web} A thin, circumferential mucosal diaphragm in the gastric antrum, usually located one to two centimetres proximal to the pylorus. It can be congenital or acquired and may cause gastric outlet obstruction, nausea, postprandial vomiting, and early satiety.

\synonyms
Gastric Antral Diaphragm; Antral Septum

\medterm{Antrochoanal Polyp} A usually noncancerous growth that starts in the lining of a cheek sinus and extends into the back of the nose. It is more common in children and young adults and may cause blockage or drainage on one side.

\medterm{Antrum} A cavity or chamber within an organ. The term commonly refers to the lower part of the stomach near its outlet or the air-filled sinus in the cheek.

\medterm{Anuria} Almost complete absence of urine production, often defined in an adult as less than about 100 millilitres in 24 hours. It can signal severe kidney failure, blockage, shock, or another emergency.

\medterm{Anus} The opening at the end of the digestive tract through which stool leaves the body. Two ring-shaped muscles help control when stool and gas are released.

\medterm{Anxiety} A feeling of worry, fear, or unease that can be a normal response to danger or stress. It becomes a disorder when it is excessive, persistent, difficult to control, or interferes with daily life.

\medterm{Anxiety Disorder} Alternate index label; see \textit{Anxiety Disorders}.

\medterm{Anxiety Disorders} A group of mental-health conditions involving excessive fear, worry, or anxiety that persists, is difficult to control, or causes significant distress or impairment. Major types include generalized anxiety disorder, panic disorder, social anxiety disorder, specific phobia, agoraphobia, separation anxiety disorder, and selective mutism. Physical symptoms may include a racing heart, sweating, muscle tension, poor sleep, or stomach upset.

\medterm{Aorta} The body’s largest artery, carrying oxygen-rich blood from the left side of the heart to the rest of the body. It passes through the chest and abdomen and divides into the arteries supplying the pelvis and legs.

\medterm{Aortal} Relating to the aorta, the large artery that carries blood from the heart to the body.

\medterm{Aortic} Relating to the aorta, the main artery carrying blood from the heart to the body.

\medterm{Aortic Aneurysm} A permanent pathological dilation of the aortic lumen exceeding normal diameter by at least 50 percent, most commonly involving the infrarenal abdominal aorta or ascending thoracic aorta. Risk factors include atherosclerosis, hypertension, smoking, and connective tissue disorders.

\medterm{Aortic Aneurysm Rupture} A break through the wall of an aortic aneurysm that causes severe internal bleeding. It may produce sudden chest, back, or abdominal pain, collapse, and shock and is a life-threatening emergency.

\medterm{Aortic Aneurysm Screening} Testing people at increased risk for an aortic aneurysm before symptoms appear. Abdominal ultrasound is commonly used; eligibility depends on age, sex, smoking history, and family history.

\medterm{Aortic Angiography} Imaging of the aorta after contrast material is given, usually with X-ray, computed tomography, or magnetic resonance imaging. It can show narrowing, blockage, aneurysm, tearing, or abnormal blood flow.

\medterm{Aortic Arch} The curved part of the aorta between its upward and downward sections. Three major arteries branch from it to supply the head, neck, and arms; disease here can affect blood flow or press on nearby structures.

\medterm{Aortic Arch Syndrome} Reduced blood flow caused by narrowing or blockage of the aortic arch or its major branches. It may cause unequal pulses or blood pressure, arm fatigue, dizziness, or neurological symptoms.

\medterm{Aortic Balloon Pump} A temporary heart-support device. A catheter with a balloon is placed in the main artery in the chest; the balloon inflates and deflates with each heartbeat to improve blood flow to the heart and reduce the heart's workload.

\synonyms
Intra-Aortic Balloon Pump; IABP

\medterm{Aortic Body} A small group of sensory cells near the aortic arch that detects changes in blood oxygen, carbon dioxide, and acidity. It sends signals to the brain to help adjust breathing and circulation.

\medterm{Aortic Coarctation} A narrowing of the aorta present from birth, usually near the part that connects to the fetal circulation. It can cause high blood pressure in the arms, weak leg pulses, poor growth, or heart failure in a newborn.

\medterm{Aortic Disease} Any disorder affecting the aorta, the body’s main artery, including aneurysm, dissection, narrowing, inflammation, or fatty deposits. Some cases cause no symptoms; sudden severe chest, back, or abdominal pain is an emergency.

\medterm{Aortic Dissection} A tear in the inner lining of the aorta that allows blood to split the layers of its wall. It can suddenly cause severe chest or back pain, stroke, organ damage, or death and needs emergency assessment.

\synonyms
Dissecting Aortic Aneurysm; Dissecting Aneurysm; Aortic Tear

\medterm{Aortic Embolism} Blockage of an artery by clot, fatty material, or other debris that breaks away from the aorta and travels downstream. It can suddenly reduce blood flow to a limb or organ.

\medterm{Aortic Graft Infection} Infection involving an artificial graft used to repair or replace part of the aorta. It can cause fever, bloodstream infection, bleeding, or a new abnormal connection with the intestine and requires urgent specialist care.

\medterm{Aortic Injury} Damage to the aorta, often caused by severe chest trauma or a medical procedure. A tear can cause internal bleeding, reduced blood flow, shock, and a life-threatening emergency.

\medterm{Aortic Insufficiency} Leakage through the aortic valve that allows blood to flow backward into the heart’s main pumping chamber. It can make the heart work harder and may eventually cause heart failure.

\medterm{Aortic Regurgitation} Leakage through an aortic valve that does not close completely, allowing blood to flow from the aorta back into the left ventricle during diastole. Acute regurgitation can cause sudden pulmonary edema or shock; chronic regurgitation may enlarge and weaken the left ventricle and cause breathlessness, chest discomfort, palpitations, or heart failure. Echocardiography assesses the cause and severity.

\medterm{Aortic Regurgitation Murmur} A blowing sound heard during the part of the heartbeat when the heart relaxes, caused by blood leaking backward through the aortic valve. It is often loudest along the left side of the breastbone.

\medterm{Aortic Rupture} A life-threatening break through the wall of the aorta, often following an aneurysm or dissection. It can cause sudden severe chest, back, or abdominal pain, internal bleeding, collapse, and shock.

\medterm{Aortic Sclerosis} Calcification, fibrosis, and thickening of the aortic valve leaflets without significant hemodynamic obstruction or restriction of valve opening. It is a common echocardiographic finding in older adults and is associated with increased cardiovascular risk.

\medterm{Aortic Stenosis} Narrowing of the aortic valve that obstructs blood leaving the left ventricle. Common causes include age-related calcification, a bicuspid valve, and rheumatic disease; severe disease can cause exertional breathlessness, chest pain, fainting, left-ventricular thickening, or heart failure. Echocardiography assesses severity.

\medterm{Aortic Stenosis Grading} Assessment of how severely the aortic valve is narrowed using ultrasound measurements of valve area and blood-flow speed or pressure difference. Results are generally described as mild, moderate, or severe and help guide monitoring and treatment.

\medterm{Aortic Valve} A valve between the heart’s left pumping chamber and the aorta. It opens to let blood leave the heart and closes to stop blood flowing backward.

\medterm{Aortic Valve Disease} Narrowing, leakage, or another abnormality of the valve between the heart’s left pumping chamber and the aorta. It may cause breathlessness, chest pain, tiredness, fainting, or no symptoms at first.

\medterm{Aortic Valve Insufficiency} Leakage of the aortic valve that allows blood to flow backward into the left pumping chamber between heartbeats. The extra workload can enlarge or weaken the heart over time.

\medterm{Aortic Valve Regurgitation} Backward leakage of blood through an aortic valve that does not close completely. The heart must pump harder, and severe or long-standing leakage can cause enlargement or heart failure.

\synonyms
Aortic Regurgitation; Aortic Insufficiency; Aortic Incompetence

\medterm{Aortic Valve Replacement} A cardiac surgical or catheter-directed interventional procedure in which a severely diseased or calcified aortic valve is excised or displaced and replaced with a mechanical prosthesis or a bioprosthetic tissue valve.

\synonyms
AVR; TAVR; TAVI

\medterm{Aortic Valve Stenosis} Alternate index label; see \textit{Aortic Stenosis}.

\medterm{Aortic Valve Surgery, Minimally Invasive} Surgical or transcatheter repair or replacement of the aortic valve performed without conventional full median sternotomy, utilizing mini-thoracotomy, upper hemisternotomy, or transcatheter arterial access.

\synonyms
Minimally Invasive Aortic Valve Replacement; MIAVR; TAVR; TAVI

\medterm{Aortic Valve Surgery, Open} Conventional surgical repair or replacement of a diseased aortic valve performed through a full median sternotomy with cardiopulmonary bypass and cardioplegic arrest.

\synonyms
Surgical Aortic Valve Replacement; SAVR; Open Aortic Valve Replacement

\medterm{Aortitis} Inflammation of the aorta, the body’s largest artery. It may result from an autoimmune disease or infection and can weaken the artery or cause narrowing, aneurysm, dissection, or leakage of the aortic valve.

\medterm{Aortocaval Compression Syndrome} Hemodynamic compromise occurring in late pregnancy when the gravid uterus compresses the inferior vena cava and abdominal aorta while the mother is in the supine position. Leads to decreased venous return, reduced cardiac output, hypotension, dizziness, and fetal distress, rapidly relieved by placing the patient in the left lateral decubitus position.

\synonyms
Supine Hypotensive Syndrome of Pregnancy

\medterm{Aortoenteric Fistula} A life-threatening direct communication between the aorta and the gastrointestinal tract, most frequently involving the third or fourth portion of the duodenum. It may occur primarily from an eroding abdominal aortic aneurysm or secondarily following prosthetic aortic reconstruction, classically presenting with minor herald gastrointestinal bleeding followed by catastrophic exsanguination.

\medterm{Aortography} Imaging of the aorta after contrast material is injected, often through a catheter. It can show aneurysms, tears, narrowing, blockages, or abnormal blood-vessel connections; computed tomography or magnetic resonance imaging is often used instead.

\medterm{Aortoiliac Occlusive Disease} Severe atherosclerotic stenosis or complete occlusion of the distal abdominal aorta and common iliac arteries. It is clinically characterized by the triad of buttock and thigh claudication, diminished or absent femoral pulses, and erectile dysfunction in men.

\synonyms
Leriche Syndrome

\medterm{Aortopulmonary Window} A rare birth defect in which an opening connects the aorta and pulmonary artery. Too much blood may flow to the lungs, causing breathing problems, poor growth, or heart failure in infancy.

\medterm{AP X-Ray} An X-ray view in which the beam enters through the front of the body and exits through the back. It is called an anteroposterior view.

\medterm{APACHE II Score} A hospital scoring system that combines vital signs, laboratory results, age, and serious long-term illnesses to estimate the severity of critical illness. A higher score generally indicates greater risk, but it does not predict an individual outcome with certainty.

\medterm{Apathy} Reduced motivation, interest, or emotional response. It can occur with depression, dementia, Parkinson disease, brain injury, other neurological illness, or medicines and is different from ordinary tiredness or sadness.

\medterm{APC Tumor Suppressor} A protein made from the APC gene that helps control cell growth and the removal of damaged cells. Harmful changes in this gene can allow excess intestinal polyps to form and increase colorectal-cancer risk.

\medterm{Apert Syndrome} A rare inherited condition in which some skull bones join too early. It causes an unusually shaped head and face and often joined fingers or toes; breathing, vision, hearing, or development may also be affected.

\medterm{Apex} The tip or pointed end of a body structure. Examples include the lower pointed end of the heart, formed mainly by the left ventricle, and the upper tip of a lung.

\medterm{Apex Beat} The outermost and lowest point on the precordium at which the cardiac impulse can be palpated or inspected during ventricular systole, typically located in the left fifth intercostal space along the midclavicular line.

\synonyms
Apical Impulse

\medterm{Apgar Score} A quick assessment of a newborn’s colour, heart rate, reflex response, muscle tone, and breathing, usually recorded one and five minutes after birth. It describes immediate adaptation and helps guide care but does not predict long-term health by itself.

\synonyms
Apgar Assessment; Newborn Apgar Score

\medterm{Aphagia} Complete inability to swallow. It may result from a blockage, neurological disease, severe pain, or loss of normal swallowing-muscle function and can lead to choking, dehydration, or aspiration.

\medterm{Aphakia} Absence of the eye’s natural lens, usually after cataract surgery, injury, or a birth difference. It causes major focusing problems and is usually corrected with an intraocular lens, contact lens, or strong glasses.

\medterm{Aphasia} A brain-related disorder that impairs speaking, understanding, reading, or writing language. It most often follows a stroke but can also result from head injury, brain tumour, infection, or degeneration.

\medterm{Apheresis} A procedure in which blood passes through a machine that removes a selected component, such as plasma, platelets, red cells, or abnormal cells, and returns the remaining blood to the person.

\medterm{Apheresis Donation} A blood donation in which a machine collects one component, such as platelets or plasma, and returns the other blood components to the donor. The process may allow more of the needed component to be collected than whole-blood donation.

\medterm{Aphonia} Complete loss of the ability to produce a normal voice, although whispering may remain possible. Causes include inflammation, vocal-cord paralysis, injury, growths, nerve disease, or psychological factors.

\medterm{Aphrasia} Loss or impairment of the ability to speak. The term is older and less precise than aphasia, which can affect speaking, understanding, reading, or writing.

\medterm{Aphthous Stomatitis} Repeated painful ulcers on the moist lining of the mouth, commonly called canker sores. They are usually shallow, heal without scarring, and are not contagious.

\medterm{Aphthous Ulcer} A small, painful ulcer on the moist lining of the mouth, commonly called a canker sore. It usually heals within one to two weeks and is not contagious.

\medterm{Apical} Relating to or located at the tip of a body structure, such as the tip of a lung, heart, or tooth.

\medterm{Apical Ballooning Syndrome} A temporary weakening and ballooning of the main pumping chamber of the heart, often after severe emotional or physical stress. It can resemble a heart attack but usually occurs without a blocked coronary artery.

\medterm{Apical Periodontitis} Inflammation or infection around the tip of a tooth root, usually caused by infected or dead pulp. It may cause pain when biting, swelling, or no symptoms and is assessed with dental examination and imaging.

\medterm{Apical Pulse} The heartbeat counted by listening over the tip of the heart, usually on the left side of the chest. It may be used when the pulse at the wrist is difficult to feel or when an accurate heart rate is needed.

\medterm{Apicoectomy} Dental surgery that removes the tip of a tooth root and nearby infected tissue, usually when infection persists after root-canal treatment. The end of the root is then sealed.

\medterm{Aplasia} Failure of an organ, tissue, or group of cells to develop, causing it to be absent or severely underdeveloped. It may be present from birth or result from inherited disease, injury, infection, or severe damage to the cells that produce it.

\medterm{Aplasia Cutis Congenita} A birth defect in which a localized area of skin, most often on the scalp, is absent or incompletely formed. It may heal as a scar or hairless patch, sometimes with an underlying skull or vascular abnormality; large, ulcerated, or atypical lesions need specialist evaluation.

\medterm{Aplastic} Describing a tissue that has failed to develop or produce its normal cells. The term is often used for bone-marrow failure, in which too few red cells, white cells, and platelets are made.

\medterm{Aplastic Anemia} A serious condition in which the bone marrow does not make enough red cells, white cells, or platelets. It can cause tiredness, infections, and bleeding and may result from immune attack, medicines, chemicals, infection, or inherited disease.

\medterm{Aplastic Crisis} A sudden, temporary stop in red-cell production that causes a rapid worsening of anemia. It most often follows parvovirus B19 infection in people whose red cells already have a shortened lifespan, such as those with sickle-cell disease.

\medterm{Apley Compression Test} A physical examination manoeuvre used to evaluate for meniscal tears of the knee joint. With the patient prone and the knee flexed to ninety degrees, the examiner applies downward axial force to the heel while rotating the tibia internally and externally; pain or joint line clicking indicates meniscal pathology.

\synonyms
Apley Grinding Test; Apley's Compression Test

\medterm{Apley Distraction Test} A clinical examination technique used to differentiate collateral ligament injuries from meniscal tears. With the patient prone and the knee flexed to ninety degrees, the examiner stabilizes the posterior thigh and applies upward traction to the lower leg with internal and external rotation; localized pain suggests ligamentous sprain rather than meniscal damage.

\synonyms
Apley's Distraction Test

\medterm{Apley Grind Test} A physical examination maneuver used in orthopedics to evaluate meniscus tears of the knee. With the patient prone and the knee flexed to 90 degrees, the examiner stabilizes the thigh, applies downward axial compressive force onto the heel while externally and internally rotating the tibia; pain or mechanical clicking indicates a torn meniscus.

\synonyms
Apley Compression Test; Apley Meniscal Test

\medterm{Apley Scratch Test} A rapid physical examination maneuver used in musculoskeletal medicine to assess the active range of motion of the shoulder joint complex. The patient is asked to reach behind the head and touch the superior medial angle of the opposite scapula (evaluating glenohumeral abduction and external rotation) and reach behind the back to touch the inferior angle of the opposite scapula (evaluating adduction and internal rotation). Asymmetry or restricted reach indicates rotator cuff pathology, adhesive capsulitis, or glenohumeral arthritis.
\synonyms{Apley Scratch Maneuver; Apley Shoulder Test}

\medterm{Apnea} A temporary pause in breathing. Prolonged or repeated episodes can lower oxygen, raise carbon dioxide, slow the heartbeat, cause loss of consciousness, or lead to death.

\medterm{Apnea Of Prematurity} Repeated pauses in breathing in a premature infant because breathing control is not fully developed. Episodes may lower oxygen or slow the heartbeat and require monitoring and, when necessary, breathing support.

\medterm{Apnea-Hypopnea Index} The average number of complete or partial breathing interruptions during each hour of sleep, measured in a sleep study. It helps describe sleep-apnea severity, but symptoms, oxygen levels, and other findings also matter.

\medterm{Apneic} Relating to a period when breathing has stopped temporarily. Apneic episodes may occur during sleep, illness, poisoning, seizures, or severe airway blockage and can quickly lower oxygen.

\medterm{Apneustic Breathing} An abnormal pattern with an unusually long pause after breathing in. It usually indicates serious injury to the brainstem and may occur with other signs of severe neurological damage.

\medterm{Apnoea} The temporary stopping of breathing; the spelling is commonly used in British English. It may result from airway blockage, reduced breathing signals from the brain, illness, medicines, or sleep-related breathing disorders.

\medterm{Apocrine Hidrocystoma} A benign cyst of an apocrine sweat gland, commonly presenting as a translucent or bluish, dome-shaped papule near the eyelid or elsewhere on the skin. It may enlarge with heat or sweating and is usually diagnosed clinically.

\medterm{Apoenzyme} The protein part of an enzyme before its required helper substance has attached. It is usually inactive or less active until the helper, called a cofactor, binds to it.

\medterm{Apolipoprotein B100} A protein on cholesterol-carrying particles such as low-density and very-low-density lipoproteins. Its blood level estimates the number of particles that can contribute to fatty buildup in artery walls.

\medterm{Apolipoprotein CII} A protein on fat-carrying particles in the blood that activates an enzyme needed to break down triglycerides. Too little or faulty apolipoprotein CII can cause triglycerides to rise to very high levels.

\medterm{Apolipoproteins} Proteins attached to fats in the blood that help transport and process cholesterol and triglycerides. Different apolipoproteins are found on different fat-carrying particles and have different effects.

\medterm{Aponeurosis} A broad, flat sheet of tough connective tissue that attaches a muscle to bone, another muscle, or other tissues and helps transmit the muscle’s force.

\medterm{Apophenia} Seeing meaningful patterns or connections in unrelated or random events. Mild experiences can occur normally, while marked or distressing experiences may occur with some psychiatric or neurological conditions.

\medterm{Apophyses} Bony projections that grow from or arise near a bone and usually provide attachment for muscles or ligaments. Unlike the main growth plate of a long bone, an apophysis does not normally increase the bone’s length.

\medterm{Apoptosis} A controlled process of cell death in which a cell dismantles itself and is removed with little inflammation. It helps shape developing organs and remove old, damaged, or abnormal cells; failure of this process can contribute to cancer.

\medterm{Apoptosis Extrinsic Pathway} A controlled cell-death process that begins when an outside signal attaches to a death receptor on a cell’s surface. Internal enzymes then take the cell apart in an orderly way, usually without causing the swelling and inflammation of an injury.

\medterm{Apoptosis Intrinsic Pathway} A controlled cell-death process that begins with severe stress inside a cell, such as DNA damage, lack of growth signals, or lack of oxygen. The cell’s energy-producing structures then release signals that activate enzymes to dismantle it.

\medterm{Apoptotic Body} A small membrane-bound piece of a cell that has begun controlled cell death. Immune or neighbouring cells can recognize and remove these pieces without causing much inflammation.

\medterm{Apparent Mineralocorticoid Excess} A rare inherited disorder in which cortisol acts too strongly on the kidneys because it is not changed into an inactive form normally. It can cause severe high blood pressure, low potassium, and unusually low levels of renin and aldosterone.

\medterm{Appendageal} Relating to structures that develop from the skin, including hair follicles, sweat glands, and oil-producing glands.

\medterm{Appendectomy} The surgical removal of the vermiform appendix, a small finger-shaped pouch attached to the cecum at the beginning of the large intestine. It is one of the most common emergency surgical operations worldwide, primarily performed to treat acute appendicitis (sudden inflammation and infection of the appendix) before it perforates or ruptures and causes widespread peritonitis (infection of the abdominal cavity) or an intra-abdominal abscess. An appendectomy can be performed using minimally invasive laparoscopic surgery (keyhole surgery through small puncture incisions using a miniature camera and instruments, enabling faster recovery and less scarring) or open surgery (through a single incision in the lower right abdomen, typically used if the appendix has ruptured or severe adhesions are present). Most patients recover rapidly and fully without long-term changes to digestion.

\synonyms
Appendicectomy; Laparoscopic Appendectomy; Appendix Removal; Open Appendectomy

\medterm{Appendiceal} Relating to the appendix, a narrow, finger-shaped pouch attached to the beginning of the large intestine.

\medterm{Appendiceal Adenocarcinoma} A cancer arising from gland-forming cells in the appendix. It may be found after appendicitis or unexpectedly and can spread through the abdominal lining; outlook depends on its type, size, and spread.

\medterm{Appendiceal Cancer} Cancer arising in the appendix. Types include neuroendocrine tumours, mucin-producing tumours, and adenocarcinoma; some can rupture and spread tumour cells or mucus through the abdomen.

\medterm{Appendiceal Rupture} A tear or hole in an inflamed appendix that allows infected material into the abdomen. It can cause peritonitis, an abscess, sepsis, or severe abdominal pain and requires urgent care.

\medterm{Appendiceal Tumors} Tumours arising in the appendix, including neuroendocrine tumours, mucin-producing tumours, and adenocarcinoma. Some cause no symptoms, while others cause appendicitis, pain, blockage, or abdominal swelling.

\medterm{Appendicitis} Inflammation of the appendix, often after its opening becomes blocked by hardened stool, mucus, or swollen tissue. It commonly causes pain that moves to the lower right abdomen, nausea, loss of appetite, and fever and can lead to rupture.

\synonyms
Epityphlitis; Acute Appendicitis

\medterm{Appendicular Skeleton} The bones of the arms and legs and the shoulder and pelvic girdles that attach them to the central skeleton.

\medterm{Appendix} A small, finger-shaped pouch attached to the beginning of the large intestine. It contains immune tissue and may help support healthy gut bacteria, although its exact function is not fully understood and it is not essential for digestion. Inflammation of the appendix is called appendicitis.

\medterm{Appendix Cancer} A malignant neoplasm arising within the appendix, including neuroendocrine tumors, goblet cell carcinoids, mucinous cystadenocarcinomas, and colonic-type adenocarcinomas.

\medterm{Appendix Testis} A small remnant of tissue near the upper part of the testis, left over from early development. It is usually harmless but can sometimes twist and cause pain in boys.

\medterm{Appetite} The psychological desire to eat food, distinct from physiological hunger, modulated by central nervous system pathways, neuroendocrine signaling, sensory stimuli, emotional states, and metabolic demands.

\medterm{Appetite Regulation} The control of hunger and fullness by the brain, digestive tract, hormones, emotions, habits, sleep, and surroundings. Signals such as ghrelin promote hunger, while leptin and other signals help promote fullness.

\medterm{Applied Behavior Analysis} A structured approach that observes behaviour and uses learning principles, practice, and reinforcement to build useful skills or reduce behaviour that interferes with daily life. Goals should be individualized and respectful of the person.

\medterm{Apposition} Placement of two structures next to or against each other. In wound care, good apposition means that the edges of a cut are accurately aligned so they can heal.

\medterm{Apprehension Test} A provocative orthopedic examination maneuver designed to detect joint instability, most commonly applied to the shoulder (glenohumeral joint) or knee (patellofemoral joint). In the anterior shoulder apprehension test, the examiner places the arm in 90 degrees of abduction and maximal external rotation; a positive test occurs when the patient exhibits visible fear, alarm, or resistance due to a feeling of impending anterior subluxation or dislocation. In the patellar apprehension test (Fairbank's sign), lateral displacement of the patella elicits anxiety and quadriceps contraction.
\synonyms{Anterior Apprehension Test; Crank Test; Fairbank's Apprehension Test}

\medterm{Appropriate For Gestational Age} Describing a fetus or newborn whose estimated size or weight is within the expected range for the number of weeks of pregnancy, commonly between the 10th and 90th percentiles.

\medterm{Apraxia} Difficulty performing a learned, purposeful movement even though strength, sensation, coordination, and understanding are sufficient. It results from a problem in the brain’s planning of movement rather than simple weakness.

\medterm{Apraxia of Speech} A neurogenic motor speech disorder characterized by an impaired ability to plan, sequence, and program the sensorimotor commands necessary for phonetically normal speech production in the absence of significant neuromuscular weakness or paralysis. Marked by inconsistent articulatory errors, effortful groping, and disrupted prosody.

\synonyms
Verbal Apraxia; Speech Apraxia

\medterm{Apraxic Gait} A higher-level neurological gait disorder characterized by the loss of the ability to coordinate voluntary stepping and leg movements in the absence of primary sensory loss, weakness, or cerebellar ataxia. Patients exhibit hesitation in initiating walking, a wide-based stance, short shuffling steps, and a tendency for the feet to appear glued to the floor (\"magnetic gait\"), commonly seen in normal pressure hydrocephalus, bilateral frontal lobe lesions, and advanced subcortical vascular encephalopathy.
\synonyms{Gait Apraxia; Frontal Gait Disorder; Magnetic Gait; Bruns Apraxia}

\medterm{Aprosodia} A neurological speech and language disorder characterized by an inability to express or comprehend the emotional tone, rhythm, pitch, and melody (prosody) of verbal communication. It is classically associated with lesions of the non-dominant (right) cerebral hemisphere in areas homologous to Broca and Wernicke expressive and receptive language regions.

\synonyms
Dysprosody; Affective Agnosia of Speech

\medterm{Aquagenic Keratoderma} Temporary painful whitening, thickening, and wrinkling of the palms or soles after contact with water. It may cause prickling or burning, is often associated with cystic fibrosis, and usually settles after the skin dries.

\medterm{Aquaporin} A channel in a cell membrane that lets water pass through. Aquaporin-2 in kidney tubules responds to vasopressin and helps the kidneys return water to the blood instead of losing it in urine.

\medterm{Aquatic Therapy} Exercise or rehabilitation performed in water. Water buoyancy reduces load on joints while water resistance strengthens muscles and supports balance, movement, and conditioning.

\medterm{Aqueous Flare} A hazy beam seen in the clear fluid at the front of the eye during slit-lamp examination. It means inflammation has allowed proteins to leak into the eye and may occur with uveitis, infection, injury, or surgery.

\medterm{Aqueous Humor} The clear fluid in the front part of the eye. It nourishes the cornea and lens, removes waste, and helps maintain pressure and the shape of the eye.

\medterm{Arachidonic Acid Pathway} A major cellular lipid metabolic cascade in which arachidonic acid, released from cell membrane phospholipids by phospholipase A2, is converted into bioactive eicosanoids (prostaglandins, prostacyclin, thromboxane A2, and leukotrienes) via cyclooxygenase and lipoxygenase enzymes.

\medterm{Arachnodactyly} Abnormally long, thin fingers or toes. It may occur alone or as a feature of inherited connective-tissue disorders such as Marfan syndrome.

\medterm{Arachnoid Cyst} A usually noncancerous, fluid-filled sac between layers of the membrane covering the brain or spinal cord. Many cause no symptoms, but a large cyst may cause headache, seizures, weakness, or pressure on nearby tissue.

\medterm{Arachnoid Cysts} Fluid-filled sacs between layers of the membrane covering the brain or spinal cord. Many are discovered by chance and cause no symptoms; larger cysts can press on nearby tissue and cause neurological problems.

\medterm{Arachnoid Granulations} Macroscopic tufted outgrowths of the arachnoid membrane that project through the dura mater into the dural venous sinuses, predominantly the superior sagittal sinus, serving as the primary structural pathway for the passive reabsorption of cerebrospinal fluid into venous blood.

\synonyms
Pacchionian Bodies; Arachnoid Villi

\medterm{Arachnoid Mater} The middle of three thin membranes covering the brain and spinal cord. The space beneath it contains the fluid that cushions the nervous system and carries important blood vessels.

\medterm{Arachnoiditis} Inflammation of the arachnoid, one of the membranes covering the brain and spinal cord, most often in the spine. Inflammation may cause thickening, scarring, and bands that make nerve roots stick together or interfere with the normal flow of cerebrospinal fluid. Possible causes include infection, bleeding around the spinal cord, spinal injury, chronic nerve compression, surgery or another spinal procedure, and chemical irritation. Symptoms may include severe or burning back and limb pain, numbness, tingling, weakness, abnormal reflexes, or problems with bladder or bowel function. The course may remain stable or worsen over time.

\medterm{Arbovirus} A virus spread by a biting insect or tick, such as a mosquito, tick, or sandfly. Arboviruses can cause fever, rash, joint pain, bleeding, or inflammation of the brain, depending on the virus.

\medterm{ARBs} Medicines that block the main receptor for angiotensin II, relaxing blood vessels and reducing salt and water retention. They lower blood pressure and can protect the heart and kidneys in selected conditions, but are usually not combined with ACE inhibitors.

\medterm{Arcanobacterium haemolyticum Infection} A bacterial infection that commonly causes pharyngitis in adolescents and young adults, with fever, sore throat, tonsillar exudate, cervical lymph-node enlargement, and sometimes a scarlatiniform or morbilliform rash. It can resemble group A streptococcal infection, so throat testing and clinical assessment guide treatment.

\medterm{Arciform} Having a curved or arch-like shape. The term may describe the appearance of a skin lesion, rash, blood vessel, or other body structure.

\medterm{Arcuate Uterus} A mild variation in the shape of the inside top of the uterus, present from birth, that creates a shallow curved indentation. It is usually considered a normal variant and generally has little or no effect on fertility or pregnancy.

\medterm{Arcus Senilis} A greyish-white, circular lipid deposition ring appearing in the peripheral corneal stroma, commonly seen in elderly individuals as a benign finding, or in younger patients (arcus juvenilis) associated with familial hyperlipidemia.

\synonyms
Corneal Arcus; Arcus Lipoides; Gerontoxon

\medterm{Area Under The Concentration-Time Curve} A measure of the total amount of a medicine in the blood over a period of time. It helps compare how medicines are absorbed and removed and can guide dosing for selected treatments.

\medterm{Arenaviruses} A family of viruses carried mainly by rodents. Some cause serious human diseases, including viral haemorrhagic fevers, after exposure to infected rodent urine, droppings, saliva, or contaminated dust.

\medterm{Areola} The darker circular area surrounding the nipple. It contains small glands that help lubricate and protect the skin, especially during breastfeeding.

\medterm{Argentine Hemorrhagic Fever} A serious viral illness caused by Junín virus, an arenavirus carried by rodents in parts of Argentina. It can cause fever, weakness, bleeding, and nervous-system problems and requires urgent medical care.

\medterm{Argininosuccinate Lyase Deficiency} An inherited urea-cycle disorder caused by deficient argininosuccinate lyase. Argininosuccinic acid and ammonia can accumulate, causing poor feeding, vomiting, liver problems, developmental impairment, or life-threatening hyperammonemia.

\medterm{Argininosuccinic Aciduria} A rare inherited disorder in which the body cannot safely remove waste nitrogen produced when protein is broken down. Ammonia and argininosuccinic acid build up and may cause vomiting, sleepiness, seizures, liver problems, or developmental impairment.

\medterm{Argon Plasma Coagulation} A procedure that uses electrically charged argon gas to deliver heat to the surface of tissue. It can stop bleeding or destroy selected abnormal tissue, often during an endoscopic examination.

\synonyms
APC; Argon Plasma Hemostasis

\medterm{Argyll Robertson Pupil} A small, irregular pupil that becomes smaller when focusing on a nearby object but reacts poorly or not at all to light. It is classically associated with late syphilis affecting the nervous system, but other causes are possible.

\medterm{Arias-Stella Reaction} A benign change in cells lining the uterus or cervix caused by pregnancy hormones or hormone medicines. Under a microscope it can resemble cancer, so the clinical setting and other findings are important.

\medterm{Arm} The part of the upper limb between the shoulder and elbow. In everyday language, “arm” may also mean the entire upper limb, including the forearm and hand.

\medterm{Arm CT Scan} A scan that uses X-rays and computer processing to create cross-sectional images of the arm. It can show fractures, bleeding, inflammation, masses, and other structural problems.

\medterm{Arm Fracture} A break in one or more bones of the upper limb, usually the upper arm bone or one of the forearm bones. It may cause pain, swelling, bruising, deformity, or inability to move the arm normally.

\medterm{Arm MRI Scan} A scan that uses a magnetic field and radio waves to create detailed images of the arm. It can show muscles, tendons, ligaments, nerves, bone marrow, blood vessels, and masses without X-ray radiation.

\medterm{Arm Pain} Pain in the shoulder, upper arm, elbow, forearm, or nearby tissues. Causes include muscle or tendon injury, nerve compression, arthritis, and blood-vessel disease; sudden severe pain with chest discomfort or breathlessness is an emergency.

\synonyms
Brachialgia; Upper-Limb Pain

\medterm{Armpit Lump} A lump or swelling in the armpit, commonly caused by an enlarged lymph node, skin infection, cyst, or noncancerous growth. A hard, fixed, enlarging, persistent, or unexplained lump needs medical assessment.

\medterm{Arnold-Chiari Malformation} A birth defect in which part of the lower brain extends through the opening at the base of the skull into the spinal canal. It may block the normal flow of brain fluid and cause headache, neck pain, balance problems, swallowing difficulty, or weakness.

\medterm{Aromatase} An enzyme that changes androgen hormones into oestrogen. It is found in tissues including the ovaries, fat, and breast and helps regulate hormone levels.

\medterm{Aromatase Excess Syndrome} A rare inherited condition in which too much aromatase causes excess oestrogen production. It may cause breast development in males and early puberty, rapid growth followed by short adult height, or other hormone-related changes.

\medterm{Arousal Index} The average number of brief awakenings or shifts to lighter sleep during each hour of a sleep study. A high result suggests repeatedly interrupted sleep, but its importance depends on symptoms and other sleep-study findings.

\medterm{Arrector Pili} A small involuntary muscle attached to a hair follicle. When it contracts, it raises the hair and produces goosebumps.

\medterm{Arrector Pili Muscle} The small involuntary muscle attached to each hair follicle. Nerve signals can make it contract, raising the hair and producing temporary goosebumps.

\synonyms
Piloerector Muscle; Arrectores Pilorum

\medterm{Arrested Caries} Tooth decay that has stopped progressing, often because the affected surface has become hard, smooth, and easier to clean. It may remain stable but still requires dental assessment and prevention of new decay.

\medterm{Arrhythmia} An abnormal heart rate or rhythm. The heartbeat may be too fast, too slow, irregular, or start from the wrong place; some arrhythmias are harmless, while others reduce blood flow or cause cardiac arrest.

\medterm{Arrhythmias} Abnormal heart rates or rhythms caused by changes in the heart’s electrical signals. They range from harmless extra beats to rhythms that cause fainting, poor circulation, or cardiac arrest.

\medterm{Arrhythmogenic Right Ventricular Cardiomyopathy} An inherited disease in which heart muscle, especially in the right pumping chamber, is replaced by scar tissue and sometimes fat. This can cause dangerous irregular rhythms, fainting, or sudden cardiac arrest, particularly during exercise.

\medterm{Arsenic} A naturally occurring element that is poisonous in many forms. Exposure can cause vomiting, diarrhoea, nerve problems, heart problems, skin changes, and increased cancer risk; some compounds have limited medical uses.

\medterm{Arsenic Poisoning} Illness caused by sudden or long-term exposure to arsenic. It may cause vomiting, diarrhoea, abdominal pain, heart problems, nerve damage, anemia, skin changes, or increased cancer risk and is an emergency requiring urgent medical treatment.

\synonyms
Arsenic Toxicity; Arsenicosis

\medterm{Arsenic Toxicity} Harm caused by arsenic exposure. It may affect the digestive system, heart, nerves, blood, skin, and nails; severe exposure can be fatal.

\medterm{Artemisinin} A compound originally obtained from the sweet-wormwood plant and used, with other medicines, to treat malaria. Combination treatment is important because malaria parasites can become resistant when artemisinin is used alone.

\medterm{Arteria Femoralis} The femoral artery, the main artery supplying the thigh and lower limb. It continues from the external iliac artery below the groin and can be felt or accessed near the groin.

\medterm{Arterial Blood Gas} A blood test, usually from an artery at the wrist, that measures acidity, oxygen, carbon dioxide, and bicarbonate. It helps assess breathing, oxygen delivery, and acid–base disorders.

\medterm{Arterial Blood Gas Analyzer} A machine that measures acidity, oxygen, carbon dioxide, and related values in arterial blood. The results help assess how well a person is breathing and whether the blood is too acidic or alkaline.

\medterm{Arterial Catheter} A thin tube placed in an artery, often at the wrist, to measure blood pressure continuously and take repeated blood samples. Possible complications include bleeding, infection, clotting, or reduced blood flow.

\medterm{Arterial Catheterization} Placement of a thin tube in an artery, usually at the wrist, for continuous blood-pressure monitoring or repeated arterial blood tests. It is used when close monitoring is needed and can cause bleeding, infection, clotting, or reduced circulation.

\medterm{Arterial Dissection} A pathological process in which a disruption or tear in the tunica intima allows circulating blood under pressure to enter and cleave the media, creating a false lumen that tracks longitudinally along the arterial wall and may lead to vessel occlusion or rupture.

\medterm{Arterial Duplex} An ultrasound test that shows the structure of arteries and measures blood flow through them. It can detect narrowing, blockage, aneurysm, or whether a vascular graft remains open.

\medterm{Arterial Embolism} Sudden blockage of an artery by material carried in the bloodstream, usually a blood clot. It can abruptly cut off blood flow to a limb, brain, kidney, intestine, or other organ and is an emergency.

\medterm{Arterial Embolization} A procedure that deliberately blocks a selected artery, usually to stop bleeding, reduce blood supply to a tumour, or treat an abnormal blood-vessel connection.

\medterm{Arterial Injury} Damage to an artery that may cause severe bleeding, a clot, narrowing, reduced blood flow, tissue loss, or life-threatening complications. It can result from trauma, surgery, or a medical procedure.

\medterm{Arterial Insufficiency} Reduced blood flow through an artery to a tissue or limb, usually because the artery is narrowed or blocked. It may cause pain with walking, pain at rest, cool skin, slow-healing wounds, or weak pulses.

\medterm{Arterial Line} A thin tube placed in an artery, usually at the wrist, for continuous blood-pressure monitoring and repeated arterial blood sampling. It is used in critical illness or major surgery and can cause bleeding, infection, clotting, or reduced circulation.

\medterm{Arterial Resistance} The opposition blood encounters as it flows through arteries, especially the small arteries that can narrow or widen. It helps control blood pressure and how much blood reaches each tissue.

\synonyms
Vascular Resistance; Systemic Vascular Resistance

\medterm{Arterial Spin Labeling} An MRI method that uses magnetically labeled blood as a natural tracer to measure blood flow, especially in the brain, without injecting a tracer.

\medterm{Arterial Stenosis} Narrowing of an artery that reduces blood flow to the tissues beyond it. Common causes include fatty plaque, inflammation, injury, or an inherited abnormality of the artery wall.

\medterm{Arterial Stick} Puncture of an artery with a needle, usually to obtain an arterial blood sample. It may cause temporary pain or bruising and, rarely, bleeding, arterial injury, or reduced blood flow.

\medterm{Arterial Stiffness} Reduced ability of arteries to expand and recoil with each heartbeat. It is associated with ageing, high blood pressure, diabetes, kidney disease, and fatty artery disease and is linked with higher cardiovascular risk.

\medterm{Arterial Switch Operation} Open-heart surgery usually performed in newborns whose two main arteries leave the wrong pumping chambers. The arteries and coronary arteries are reconnected so blood flows normally; lifelong heart follow-up is needed.

\synonyms
Jatene Procedure; ASO

\medterm{Arterial Tension} A measurement taken from arterial blood. Depending on the context, it may mean blood pressure inside an artery or the amount of oxygen or carbon dioxide dissolved in the blood.

\medterm{Arterial Thromboembolism} Blockage of an artery by a blood clot that formed elsewhere and travelled through the bloodstream. It can suddenly reduce blood flow to the brain, heart, limb, kidney, intestine, or another organ.

\medterm{Arterial Thrombosis} Formation of a blood clot inside an artery, often where a fatty plaque has damaged the vessel lining. The clot can block blood flow to the heart, brain, limb, intestine, or another organ.

\medterm{Arterial Ulcer} A slow-healing wound caused by poor blood flow through an artery, usually on a toe, heel, or other pressure point. It may be deep and sharply outlined, with a pale or dead-looking base and cool, shiny, hairless surrounding skin.

\medterm{Arteries} Blood vessels that carry blood away from the heart. Most carry oxygen-rich blood, but the pulmonary and umbilical arteries carry oxygen-poor blood; strong elastic and muscular walls help withstand and regulate pressure.

\synonyms
Arterial Vessels

\medterm{Arteriogram} An image of an artery, usually made after contrast material is injected. It can show narrowing, blockage, aneurysm, abnormal connections, or other changes in blood flow.

\medterm{Arteriole} A small branch of an artery that delivers blood to a network of tiny vessels called capillaries. Its muscular wall can narrow or widen to control blood flow and blood pressure.

\medterm{Arterioles} Small branches of arteries that control blood flow into capillaries by narrowing or widening their muscular walls. They are important in regulating blood pressure and the supply of blood to tissues.

\synonyms
Resistance Vessels; Small Arteries

\medterm{Arteriolosclerosis} A form of vascular disease affecting small arteries and arterioles, characterized by mural thickening, luminal narrowing, and loss of elasticity, classically manifesting as hyaline arteriolosclerosis (associated with aging and diabetes) or hyperplastic arteriolosclerosis (associated with severe hypertension).

\medterm{Arteriosclerosis} Thickening, hardening, and loss of elasticity in artery walls. It includes several processes, including age-related stiffening and plaque buildup, and can reduce blood flow to the heart, brain, legs, or other organs.

\medterm{Arteriosclerosis / Atherosclerosis} Arteriosclerosis means hardening or thickening of artery walls. Atherosclerosis is a type in which fatty, inflammatory plaque builds up inside the walls and may restrict blood flow or rupture.

\medterm{Arteriosclerotic Retinopathy} Damage or changes in the retina caused by thickened, narrowed, or hardened small blood vessels, most often after long-standing high blood pressure. Severe disease can reduce vision.

\medterm{Arteriovenous Fistula} An abnormal connection between an artery and a vein that bypasses the normal capillary network. It may be present from birth, follow injury or a procedure, or be created to provide access for haemodialysis.

\medterm{Arteriovenous Fistula Creation} Surgery to connect an artery to a nearby vein to create blood access for hemodialysis. Over time, the vein may become wider and stronger so it can be used for dialysis needles. It often has advantages over a graft or catheter, but its success, risks, and long-term usefulness vary from person to person.

\synonyms
AV Fistula Creation; Arteriovenous Fistula Surgery; Hemodialysis Fistula Placement; AVF Construction

\medterm{Arteriovenous Graft} An artificial tube connecting an artery and vein to provide access for haemodialysis, usually when a person’s veins cannot form a suitable fistula. It can clot or become infected and needs regular monitoring.

\medterm{Arteriovenous Malformation} An abnormal tangle of arteries and veins joined without the usual capillaries between them. It can bleed, cause seizures or other neurological symptoms, or put extra strain on the heart.

\medterm{Arteriovenous Malformation Resection} A complex neurosurgical operation performed to remove an arteriovenous malformation (AVM)—an abnormal, high-pressure tangle of blood vessels connecting arteries directly to veins in the brain without intervening capillaries. Because the thin-walled veins cannot tolerate direct arterial pressure, AVMs carry high risks of catastrophic brain bleeding (hemorrhagic stroke), seizures, and chronic headaches. Under high-magnification operating microscopes and continuous image guidance, a neurosurgeon enters through a craniotomy, painstakingly identifies and seals off every feeding artery supplying the malformation, coagulates and detaches the draining veins, and completely excises the vascular nidus from surrounding healthy brain tissue to eliminate the danger of hemorrhage permanently.
\synonyms
AVM Resection; Surgical Excision of Brain AVM; Cerebral AVM Removal

\medterm{Arteriovenous Nicking} A clinical fundoscopic finding in chronic hypertensive retinopathy characterized by the apparent tapering, deflection, or constriction of a retinal venule where it is crossed by a thickened, sclerotic retinal arteriole (Gunn's sign). Because the arteriole and venule share a common adventitial sheath at crossing points, progressive arteriolosclerosis compresses the thinner-walled venule, indicating sustained moderate-to-severe systemic vascular hypertension and microvascular end-organ damage.
\synonyms{AV Nicking; Gunn's Sign; Arteriovenous Crossing Changes}

\medterm{Arteritis} Inflammation of an artery wall. The inflammation can narrow, weaken, or block the artery, reducing blood flow and causing pain or organ damage in the affected area.

\medterm{Arteritis, Giant Cell} See \textit{Giant cell arteritis}.

\medterm{Arteritis, Takayasu's} See \textit{Takayasu's arteritis}.

\medterm{Artery} A blood vessel that carries blood away from the heart. Most arteries carry oxygen-rich blood, except the pulmonary arteries, which carry oxygen-poor blood to the lungs. Their elastic, muscular walls withstand and regulate pressure.

\medterm{Artery Occlusion} A blockage inside an artery, often caused by a clot, fatty plaque, or material that travelled from elsewhere. It can cut off blood flow; sudden pain, paleness, coldness, numbness, weakness, or absent pulse is an emergency.

\medterm{Arthralgia} Pain in one or more joints without clear signs of inflammation, such as swelling, warmth, or redness. Causes include overuse, injury, infection, autoimmune disease, and early joint wear.

\medterm{Arthritis} A group of conditions that affect joints and the tissues around them; it is not one single disease. Depending on the type, arthritis may involve inflammation, gradual breakdown of joint tissues, infection, or crystal deposits. It can cause joint pain, swelling, warmth, stiffness, and reduced movement. Some types can also affect the skin, eyes, heart, lungs, or other organs. Examples include osteoarthritis, rheumatoid arthritis, gout, and ankylosing spondylitis.

\synonyms
Joint Inflammation; Articular Inflammation

\medterm{Arthritis of the Wrist} Inflammation or degenerative disease affecting one or more wrist joints. Pain, swelling, stiffness, reduced grip, and limited motion may result from osteoarthritis, rheumatoid arthritis, crystal disease, injury, infection, or other inflammatory disorders.

\medterm{Arthritis, Ankylosing Spondylitis} See \textit{Ankylosing Spondylitis}.

\medterm{Arthritis, Basal Joint} See \textit{Thumb Arthritis}.

\medterm{Arthritis, Degenerative} See \textit{Osteoarthritis}.

\medterm{Arthritis, Gouty} See \textit{Gout}.

\medterm{Arthritis, Infectious} See \textit{Septic Arthritis}.

\medterm{Arthritis, Juvenile Idiopathic} See \textit{Juvenile Idiopathic Arthritis}.

\medterm{Arthritis, Lyme} See \textit{Lyme Disease}.

\medterm{Arthritis, Mixed Connective Tissue Disease} See \textit{Mixed Connective Tissue Disease}.

\medterm{Arthritis, Osteoarthritis} See \textit{Osteoarthritis}.

\medterm{Arthritis, Plant Thorn} See \textit{Plant Thorn Synovitis}.

\medterm{Arthritis, Pseudogout} See \textit{Pseudogout}.

\medterm{Arthritis, Psoriatic} See \textit{Psoriatic Arthritis}.

\medterm{Arthritis, Reactive} See \textit{Reactive Arthritis}.

\medterm{Arthritis, Reiter} See \textit{Reactive Arthritis}.

\medterm{Arthritis, Rheumatoid} See \textit{Rheumatoid Arthritis}.

\medterm{Arthritis, Sarcoid} See \textit{Sarcoidosis}.

\medterm{Arthritis, Scleroderma} See \textit{Systemic Sclerosis}.

\medterm{Arthritis, Septic} See \textit{Septic Arthritis}.

\medterm{Arthritis, Sjogren} See \textit{Sjogren Syndrome}.

\medterm{Arthritis, Still} See \textit{Still Disease}.

\medterm{Arthritis, Systemic Lupus Erythematosus} See \textit{Systemic Lupus Erythematosus}.

\medterm{Arthritis, Thumb} See \textit{Thumb Arthritis}.

\medterm{Arthrocentesis} A procedure in which a needle removes fluid from a joint. The fluid can be tested for infection, blood, or crystals and may be removed to reduce pain and pressure.

\medterm{Arthrodesis} Surgery that permanently joins the bones of a damaged joint so they heal as one solid unit. It can relieve severe pain and improve stability, but the fused joint no longer moves.

\medterm{Arthrogram} An image of a joint taken after contrast material is placed inside or around it. It can show tears or problems in cartilage, ligaments, the joint capsule, or other joint structures.

\medterm{Arthrography} Imaging of a joint after contrast material is injected into it. It can help examine the shoulder, hip, wrist, knee, or another joint when ordinary imaging does not show enough detail.

\medterm{Arthrogryposis} A group of conditions present from birth in which two or more joints are permanently stiff or limited in movement. It usually results from reduced movement before birth and may involve muscles, nerves, or connective tissue.

\medterm{Arthrogryposis Multiplex Congenita} A group of birth conditions causing stiffness and limited movement in several joints. The affected joints may be in the arms, legs, jaw, or spine, and muscle weakness or abnormal nerves and connective tissue may also be present.

\medterm{Arthrolysis} A procedure that releases scar tissue or other restrictions around a stiff joint to improve movement. It may be performed through small cuts with a camera or through open surgery.

\medterm{Arthropathy} Any disease or abnormal condition affecting a joint. It may result from inflammation, wear, injury, infection, bleeding, nerve damage, or crystals such as those that cause gout.

\synonyms
Joint Disease

\medterm{Arthroplasty} Surgery to repair, reshape, or replace a damaged joint to reduce pain and improve movement or alignment. It may replace part or all of the joint and is used for severe arthritis, fractures, or other joint damage.

\medterm{Arthropod Vector} An insect, tick, or related animal that carries and transmits an infectious organism between hosts. Examples include mosquitoes spreading malaria or ticks spreading Lyme disease.
\synonyms
Disease Vector

\medterm{Arthroscope} A thin instrument with a light and camera used to look inside a joint through a small cut. It can help diagnose damage and guide keyhole surgery to repair or remove abnormal tissue.

\medterm{Arthroscopy} Keyhole surgery in which a camera and small instruments are inserted through small cuts to examine or treat a joint. It may be used to repair cartilage or ligaments, remove loose tissue, or treat an inflamed joint lining.

\medterm{Arthrosis} Degenerative damage or wear of a joint, involving loss of cartilage and changes in the underlying bone. Osteoarthritis is the most common example and can cause pain, stiffness, and reduced movement.

\medterm{Arthus Reaction} A local immune reaction that causes redness, swelling, pain, and sometimes tissue damage several hours after exposure to a substance. It occurs when antibodies and the substance form deposits that activate inflammation in the affected tissue.

\medterm{Articular Cartilage} Smooth, tough tissue covering the ends of bones inside a joint. It allows bones to move with little friction and helps spread forces across the joint.

\medterm{Articular Cartilage Injuries} Damage to the smooth cartilage covering the ends of bones in a joint, sometimes extending into the underlying bone. It may cause pain, swelling, catching, or reduced movement and can follow injury, repeated stress, or joint disease.

\medterm{Articular Ligament} A strong band of connective tissue joining bones around a joint. It stabilizes the joint and limits excessive bending, sliding, or twisting.

\synonyms
Joint Ligament

\medterm{Articulation} A joint or connection between bones. In speech, articulation also means forming clear speech sounds using the tongue, lips, teeth, and other speech organs.

\medterm{Artifacts} Distortions or false features in a medical image that do not represent real body structures. They may result from movement, metal, air, the imaging beam, or limits of the scanner and image-processing method.

\medterm{Artificial Abortion} An older term for an induced abortion: ending a pregnancy with medicines, a procedure, or both, rather than through spontaneous miscarriage.

\medterm{Artificial Heart} A mechanical device that supports or replaces the heart’s pumping action. It may support one or both lower chambers temporarily or long term and requires management of power, infection, bleeding, and blood clots.

\medterm{Artificial Insemination} A fertility procedure in which prepared sperm is placed into the cervix or uterus to increase the chance of pregnancy. The sperm may come from a partner or a donor.

\medterm{Artificial Insemination By Donor} Fertility treatment in which sperm from a screened donor is placed into the uterus or cervix to attempt pregnancy. It may be used when partner sperm is unavailable or unlikely to result in pregnancy.

\medterm{Artificial Insemination By Husband} Fertility treatment in which prepared sperm from a woman’s husband or male partner is placed into the uterus or cervix to increase the chance of fertilization and pregnancy.

\medterm{Artificial Intelligence} Computer systems that learn from data to recognize patterns, make predictions, or support healthcare decisions. Medical AI can assist diagnosis, treatment, and research, but it must be tested, monitored for bias, and supervised by qualified professionals.

\synonyms
AI

\medterm{Artificial Kidney} A dialysis device that removes metabolic waste and excess fluid from blood and helps correct electrolyte disturbances when renal function is inadequate.

\medterm{Artificial Pacemaker} An implanted device that sends electrical signals to the heart when its natural rhythm is too slow or irregular. It helps maintain an adequate heart rate and may have one or more leads.

\medterm{Artificial Rupture Of Membranes} Deliberate breaking of the amniotic sac during labour or an examination. It may be used to start or strengthen labour, assess the amniotic fluid, or attach internal monitors; it carries risks such as infection and umbilical-cord prolapse.

\medterm{Artificial Saliva} A product that moistens the mouth and temporarily relieves discomfort from dry mouth. It may improve speaking, swallowing, and comfort but does not replace fluids, dental care, or treatment of the cause.

\synonyms
Saliva Substitute

\medterm{Artificial Skin} A biological or manufactured covering placed over a wound to protect it and support growth of new skin. It may reduce fluid loss and contamination, but some wounds later need a permanent skin graft.

\synonyms
Skin Substitute

\medterm{Artificial Sweeteners} Sweet-tasting substances that provide few or no calories and are used instead of sugar. They can reduce added sugar for some people but do not by themselves make a diet healthy; effects and safety depend on the product and amount.

\medterm{Artificial Urinary Sphincter} An implanted device that controls urine leakage by gently closing and releasing the urethra. It is used for selected severe stress incontinence and requires surgery and follow-up to operate and maintain the device.

\medterm{Artificial Ventilation} Support or replacement of breathing with a machine that moves air and oxygen into and out of the lungs when a person cannot breathe well enough alone. It may use a mask or a tube and requires close monitoring.

\medterm{Artificial Ventricle} A mechanical device that helps or replaces the pumping action of one or both lower heart chambers in severe heart failure. The exact device and treatment goal vary; it is not the same as every ventricular assist device or a total artificial heart.

\synonyms
Total Artificial Heart; TAH; Ventricular Assist Device

\medterm{Aryepiglottic Fold} A triangular mucosal fold enclosing ligamentous and muscular fibers, extending bilaterally from the lateral margins of the epiglottis to the arytenoid cartilages, forming the lateral anatomical boundaries of the laryngeal inlet.

\medterm{Arytenoid Cartilage} One of a pair of small cartilages at the back of the voice box. They anchor the vocal cords and help the laryngeal muscles open, close, and adjust the cords for breathing and speech.

\medterm{ASA Physical Status Classification System} A scale used before anesthesia to summarize a patient’s overall health and major medical conditions, from healthy to severely ill. It helps the care team communicate risk but does not by itself predict complications or outcome.

\medterm{Asbestos-Related Disorders} Diseases that may develop years after breathing asbestos fibres, including pleural plaques, lung scarring, lung cancer, and mesothelioma. Risk increases with greater exposure and smoking; symptoms may include breathlessness, cough, or chest pain.

\medterm{Asbestosis} Permanent scarring of the lungs caused by long-term inhalation of asbestos fibres. It can gradually cause breathlessness, dry cough, and reduced lung function and increases the risk of lung cancer and mesothelioma.

\medterm{Ascariasis} Infection with a roundworm acquired by swallowing eggs from contaminated soil, food, or water. It may cause no symptoms, abdominal pain, poor nutrition, intestinal blockage, or cough while larvae pass through the lungs.

\medterm{Ascending Aorta} The first section of the aorta, rising from the heart’s left pumping chamber. It gives rise to the coronary arteries and continues into the aortic arch; an aneurysm or dissection here can be life-threatening.

\medterm{Ascending Aortic Aneurysm} An enlarged, weakened area in the part of the aorta that rises from the heart. It may be associated with a bicuspid aortic valve, inherited connective-tissue disease, high blood pressure, or artery-wall disease and may rupture or dissect.

\medterm{Ascending Cholangitis} A life-threatening bacterial infection of the biliary tree caused by a combination of bile duct obstruction (most commonly a gallstone in the common bile duct or a tumor) and bacterial ascent from the duodenum into stagnant bile under pressure. It characteristically presents with Charcot's triad (fever with chills, right-upper-quadrant pain, and jaundice), which can rapidly progress to Reynolds' pentad (adding confusion and septic shock/hypotension in acute suppurative cases). Blood tests show leukocytosis and hyperbilirubinemia, and imaging (ultrasound, CT, or MRCP) confirms biliary duct dilation. Emergency management consists of resuscitation with intravenous fluids, broad-spectrum IV antibiotics, and urgent biliary decompression via endoscopic retrograde cholangiopancreatography (ERCP).
\synonyms{Acute Cholangitis; Acute Suppurative Cholangitis; Biliary Sepsis}

\medterm{Ascending Colon} The first part of the large intestine, running upward on the right side of the abdomen from the appendix region to the bend beneath the liver. It absorbs water and helps form stool.

\medterm{Ascites} Abnormal buildup of fluid inside the abdomen. It commonly results from advanced liver disease but may also be caused by cancer, heart failure, kidney disease, infection, or pancreatitis and can cause swelling, discomfort, or breathlessness.

\medterm{Ascitic Fluid} Fluid collected from the abdomen in a person with ascites. Laboratory testing can help identify liver disease, infection, cancer, pancreatitis, or other causes of the fluid buildup.

\synonyms
Peritoneal Fluid; Ascitic Effusion

\medterm{Ascorbic Acid} Vitamin C, a water-soluble vitamin needed to make collagen, heal wounds, absorb iron, protect cells, and support normal immune and connective-tissue function.

\medterm{Ascus} A Pap-test result meaning that some cells from the cervix look slightly unusual, but it is unclear whether the change is important. Follow-up depends on age, HPV testing, previous results, and clinical guidance; it is not a cancer diagnosis.

\medterm{Asepsis} The condition or practice of preventing contamination by disease-causing microorganisms, especially during medical or surgical procedures. It includes sterile technique, environmental controls, hand hygiene, and appropriate handling of instruments and materials.

\medterm{Aseptic} Free from disease-causing contamination, or performed using methods that prevent microorganisms from entering a wound, specimen, medicine, or body site. Aseptic does not necessarily mean completely sterile.

\medterm{Aseptic Loosening} Loss of firm attachment of an artificial joint or other implant without an active infection. Repeated stress or a reaction to tiny wear particles can cause pain, instability, and eventual implant failure.

\medterm{Aseptic Meningitis} Inflammation of the membranes around the brain and spinal cord without the usual evidence of bacterial meningitis. Viruses are common causes, but medicines, autoimmune disease, and other infections can also be responsible.

\medterm{Aseptic Necrosis} See \textit{Avascular necrosis (osteonecrosis)}.

\synonyms
Avascular Necrosis; Osteonecrosis; Ischemic Bone Necrosis

\medterm{Aseptic Technique} Practices used to prevent harmful microorganisms contaminating a wound, body site, medicine, or specimen. They include hand hygiene, clean or sterile equipment, surface cleaning, and limiting contact with exposed sterile items.

\medterm{Asherman Syndrome} Scar tissue inside the uterus that partly or completely joins the uterine walls. It most often follows a procedure inside the uterus or complications of pregnancy and may cause light or absent periods, infertility, or repeated miscarriage.

\medterm{Ashman Phenomenon} A temporary wide-looking heartbeat that occurs when a short beat follows a long pause, often during atrial fibrillation. It can resemble a dangerous rhythm on an ECG, so clinicians interpret it with the full tracing and the person’s symptoms.

\synonyms
Ashman's Phenomenon; Aberrant Intraventricular Conduction

\medterm{ASL} American Sign Language, a complete visual language used by many Deaf and hard-of-hearing people. Healthcare communication should respect each person’s language preference and provide qualified interpreters when needed.

\medterm{ASOT} An abbreviation for the antistreptolysin O test, a blood test measuring antibodies against a toxin made by group A streptococcal bacteria. It can help show a recent strep infection when investigating later complications.

\medterm{Asp} The standard three-letter abbreviation for aspartic acid, an amino acid used to build proteins and involved in energy production and other cell reactions.

\medterm{Aspartame} A low-calorie sweetener made from two amino-acid components. It is not suitable for people with phenylketonuria because it releases phenylalanine when digested.

\medterm{Aspartate Aminotransferase} An enzyme found mainly in the liver, heart, muscles, and red blood cells. When these tissues are injured, the enzyme can enter the blood and its level may rise.

\medterm{Aspartate Aminotransferase Blood Test} A blood test measuring the level of the enzyme aspartate aminotransferase. A high result can indicate injury to the liver, heart, muscles, or red blood cells and must be interpreted with other tests.

\medterm{Aspartic Acid} An amino acid that the body can make and uses to build proteins and carry out cell reactions. It is one of the components of aspartame but is not the same substance as the sweetener.

\medterm{Asperger Syndrome} A former diagnostic classification for a neurodevelopmental profile within the autism spectrum characterized by qualitative differences in social communication and social reciprocity, accompanied by restricted, repetitive patterns of behavior and intense focal interests, without significant delays in general language acquisition or cognitive development.

\synonyms
Asperger Disorder; Asperger's Syndrome; Asperger's Disorder; Autistic Psychopathy

\medterm{Aspergillosis} A group of infections and allergic conditions caused by Aspergillus molds. It may affect the airways, sinuses, or lungs and can include allergic disease, a fungal growth in an existing lung cavity, or invasive infection.

\synonyms
Aspergillus Infection; Aspergillus Disease

\medterm{Aspergillosis Precipitin} A qualitative or semi-quantitative serological test detecting circulating precipitating IgG antibodies against \textit{Aspergillus} antigens, useful in diagnosing chronic pulmonary aspergillosis and allergic bronchopulmonary aspergillosis.

\synonyms
Aspergillus Antibodies; Aspergillus IgG Precipitins

\medterm{Aspergillus} A genus of saprophytic molds whose airborne conidia are ubiquitous in the environment, capable of causing allergic bronchopulmonary aspergillosis, aspergilloma formation in pre-existing cavitary lung lesions, or severe invasive tissue infections in immunocompromised hosts.

\medterm{Aspergillus Infection} Alternate index label; see \textit{Aspergillosis}.

\medterm{Asphalt Cement Poisoning} Illness caused by swallowing, inhaling, or having significant skin contact with hot or solvent-containing asphalt cement. Burns, breathing problems, dizziness, or chemical injury can occur; significant exposure requires urgent medical treatment.

\medterm{Asphyxia} A life-threatening lack of oxygen caused by inability to breathe or by blocked blood oxygen delivery, such as from choking, drowning, strangulation, or severe airway blockage. It can quickly cause unconsciousness, brain injury, or death.

\medterm{Asphyxia Neonatorum} Failure of a newborn to start or maintain effective breathing at or soon after birth. It can reduce oxygen to the brain and other organs and requires immediate assessment and breathing support.

\medterm{Asphyxial Death Pathology} Autopsy findings associated with impaired oxygenation or ventilation, which may include congestion, petechiae, pulmonary edema, and visceral hypoxia. These findings are usually nonspecific and must be interpreted with scene, history, and toxicology.

\medterm{Asphyxiating Thoracic Dysplasia} A rare inherited disorder in which short ribs create a narrow, stiff chest that restricts breathing, especially in infancy. Kidney, liver, and limb abnormalities may also occur.

\synonyms
Jeune Syndrome

\medterm{Aspirate} Fluid, cellular material, tissue fragments, or particulate matter evacuated from a cavity, cyst, organ, or lesion via negative-pressure suction; or foreign liquid and solid material accidentally inhaled into the tracheobronchial tree.

\synonyms
Aspirated Material; Biopsy Aspirate

\medterm{Aspiration} Breathing food, liquid, saliva, vomit, blood, or another substance into the voice box or lungs instead of swallowing it into the oesophagus. It can cause choking, airway blockage, inflammation, or pneumonia.

\medterm{Aspiration Pneumonia} A lung infection that develops after inhaling saliva, food, liquid, or stomach contents. It is more likely with swallowing difficulty, reduced consciousness, stroke, or tube feeding and may cause cough, fever, breathlessness, or chest pain.

\medterm{Aspiration Pneumonitis} Chemical inflammation of the lungs after inhaling stomach contents, especially acidic fluid. It can cause coughing, wheezing, low oxygen, breathing difficulty, and fluid or collapse in parts of the lungs.

\medterm{Aspiration Pneumonitis and Pneumonia} Aspiration pneumonitis is an acute chemical lung injury after inhaling gastric contents; aspiration pneumonia is infection after inhaling contaminated material. Fever, cough, hypoxia, or infiltrates after an aspiration event require assessment, with treatment guided by severity and evidence of infection.

\medterm{Aspiration Risk} The likelihood that food, liquid, saliva, or stomach contents will enter the airway or lungs. Risk increases with swallowing difficulty, stroke, reduced alertness, nerve or muscle disease, reflux, and some dental or feeding problems.

\medterm{Aspirin} A medicine that can relieve pain, lower fever, reduce inflammation, and make platelets less likely to form blood clots. Low-dose aspirin is often used after a heart attack, stroke, or some heart procedures, but it is not suitable for everyone because it can cause serious bleeding. It should not be taken daily for prevention without medical advice.

\synonyms
Acetylsalicylic Acid; ASA

\medterm{Aspirin Overdose} Acute or chronic poisoning resulting from excessive salicylate ingestion. Characteristic features include tinnitus, nausea, vomiting, hyperventilation causing early respiratory alkalosis, followed by severe uncoupled oxidative phosphorylation, high anion gap metabolic acidosis, hyperthermia, confusion, seizures, and noncardiogenic pulmonary edema.

\synonyms
Salicylate Toxicity; Aspirin Poisoning; Salicylism

\medterm{Aspirin Poisoning} Alternate name; see \textit{Aspirin Overdose}.

\medterm{Aspirin-Exacerbated Respiratory Disease} A condition involving asthma, chronic nasal or sinus inflammation with polyps, and respiratory reactions to aspirin or some similar pain medicines. Exposure may trigger wheezing, nasal blockage, cough, or shortness of breath.

\medterm{Assay} A laboratory method used to detect or measure a substance, biological activity, genetic feature, or response in a sample.

\medterm{Assist-Control Ventilation} A ventilator setting that gives a full preset breath whenever the patient starts one and also gives breaths at a minimum rate if the patient does not breathe enough. Each breath is set by volume or pressure.

\medterm{Assisted Breech Delivery} A vaginal birth in which the baby’s buttocks or feet come first and a trained clinician helps deliver the body and head. It is considered only in carefully selected circumstances because breech birth can be dangerous for the baby or pregnant person.

\medterm{Assisted Circulation} Mechanical or procedural support that maintains blood flow when the heart cannot pump enough on its own. It may use a pump, oxygenator, or other device for a short time or longer-term support.
\synonyms
Mechanical Circulatory Support

\medterm{Assisted Delivery With Forceps} A vaginal birth in which a trained clinician places specially shaped instruments around the baby’s head to help complete delivery. It may be considered when labour is prolonged or the baby or pregnant person needs delivery sooner.

\medterm{Assisted Living} A residential setting where people receive help with daily activities such as bathing, dressing, meals, medicines, or transport while retaining some independence. It generally does not provide the continuous skilled nursing care of a hospital or nursing home.

\medterm{Assisted Reproductive Technology} Clinical, surgical, and laboratory procedures that involve the in vitro manipulation of human oocytes, spermatozoa, or embryos to achieve pregnancy, including in vitro fertilization (IVF), intracytoplasmic sperm injection (ICSI), gamete intrafallopian transfer (GIFT), and frozen embryo transfer.

\synonyms
ART; Assisted Reproduction; Assisted Conception

\medterm{Assistive Device} Equipment that helps a person move, communicate, perform daily activities, or remain independent. Examples include a cane, wheelchair, hearing aid, communication tool, brace, walker, or adapted household equipment.

\medterm{Assistive Device Training} Instruction in choosing, fitting, using, and maintaining equipment such as walkers, canes, crutches, braces, or wheelchairs, including safe walking, transfers, and stair use.

\medterm{Association} A relationship in which two findings, events, or conditions occur together more often than expected. An association alone does not prove that one causes the other.

\medterm{Association Cortex} Areas of the brain’s outer layer that combine information from different senses. They support memory, language, planning, judgment, learning, and other complex behavior rather than controlling one simple movement or sensation.

\medterm{Association Study} A research study examining whether a genetic variant, exposure, or other factor occurs more often with a particular disease or finding. Such a relationship does not prove that the factor causes the disease.

\medterm{Asteatotic Eczema} A dry-skin eczema pattern, also called eczema craquelé, in which fine fissures create a cracked or porcelain-like appearance, usually on the shins, arms, or trunk of older adults. Cold, dry weather, frequent washing, and diuretics can worsen it; emollients and treatment of inflammation help.

\medterm{Astereognosis} Inability to recognize a familiar object by touch alone even though basic touch sensation works. It usually results from damage to a brain area that combines sensory information, often after a stroke or other brain injury.

\medterm{Asterixis} Brief, repeated lapses in holding a posture that cause a flapping movement, often when the arms are stretched out and the wrists are bent back. It can occur with liver or kidney failure, breathing disorders, or some medicines.

\medterm{Asteroid Hyalosis} A usually harmless eye condition in which tiny calcium-and-fat particles float in the gel inside the eye and appear as bright spots. It often affects one eye and causes few symptoms, but dense particles can obscure the retina.

\medterm{Asthenia} Abnormal lack of physical strength or energy. It is a symptom rather than a disease and may occur with infection, chronic illness, malnutrition, depression, or medicine effects.

\medterm{Asthenic} Describing a person or condition marked by weakness, or sometimes a slender body build. In a clinical note, the surrounding context shows which meaning is intended; asthenia specifically means reduced strength or energy.

\medterm{Asthenopia} Eye strain or visual fatigue that may cause aching, burning, blurred vision, headache, or difficulty focusing. Refractive errors, dry eyes, poor lighting, and prolonged screen use can contribute.

\medterm{Asthenozoospermia} Abnormally reduced movement of sperm, which can make fertilization more difficult. It is assessed by semen analysis and may result from infection, a varicocele, toxins, illness, or an unknown cause.

\synonyms
Asthenospermia

\medterm{Asthma} A long-term inflammatory disease of the airways in the lungs. The airways are overly sensitive and can narrow when triggered because their lining swells, the surrounding muscles tighten, and extra mucus forms. This variable narrowing makes it harder for air to move, especially when breathing out. Symptoms may include wheezing, cough, chest tightness, and shortness of breath; they can come and go and may be triggered by allergens, infections, smoke, air pollution, exercise, certain medicines, or cold air.

\medterm{Asthma Attack} A sudden worsening of asthma in which inflamed, narrowed airways make breathing difficult. It may cause wheezing, cough, chest tightness, or breathlessness; trouble speaking, blue lips, severe exhaustion, or poor response to a reliever is an emergency.

\medterm{Asthma Controller Medications} Medicines used regularly to reduce airway inflammation and prevent asthma symptoms and attacks. Inhaled corticosteroids are the main controller medicine. Other medicines may be added when needed; long-acting bronchodilators should not be used alone for asthma.

\synonyms
Asthma Controllers; Long-Term Control Medications; Inhaled Corticosteroids for Asthma

\medterm{Asthma Exacerbation} A sudden or gradual worsening of asthma symptoms and airflow through the lungs. It may cause increasing wheeze, cough, chest tightness, or breathlessness and can become life-threatening.

\medterm{Asthma in Older Adults} Asthma in later life may be underdiagnosed because symptoms overlap with COPD, heart disease, deconditioning, and medication effects. Diagnosis uses history and objective airflow testing when possible; inhaler technique, comorbidities, cognition, and treatment adverse effects need special attention.

\medterm{Asthma in Pregnancy} Asthma may improve, worsen, or remain unchanged during pregnancy. Maintaining control with appropriate inhaled treatment is generally safer than allowing maternal hypoxia; worsening wheeze, breathlessness, or poor response to rescue medicine requires prompt care.

\medterm{Asthma Medications} Medicines for asthma include regular treatments that reduce airway inflammation and reliever treatments that quickly ease narrowed airways. The regimen depends on symptoms, lung function, attack risk, age, pregnancy, inhaler technique, and response.

\medterm{Asthma Quick-Relief Medications} Medicines used to ease asthma symptoms quickly by opening narrowed airways. They may include a short-acting bronchodilator or, for some people, an inhaled corticosteroid--formoterol inhaler. The correct reliever depends on the person’s treatment plan.

\synonyms
Asthma Relievers; Rescue Inhalers; Short-Acting Beta-Agonists; SABA

\medterm{Asthma, Adult-Onset} Alternate index label; see \textit{Adult-Onset Asthma}.

\medterm{Asthma, Exercise-Induced} See \textit{Exercise-induced asthma}.

\medterm{Asthma, Occupational} See \textit{Occupational asthma}.

\medterm{Asthma-COPD Overlap} A pattern in which a person has features of both asthma and chronic obstructive pulmonary disease, such as variable symptoms together with persistent airflow limitation. Diagnosis and treatment require individual assessment because the conditions and risks differ.

\medterm{Astigmatism} A focusing problem caused by an unevenly curved cornea or lens. Light is not focused at one clear point on the retina, causing blurred or distorted vision at near or far distances.

\medterm{Astringent} A substance that makes surface tissues contract and may reduce minor secretions or bleeding. It is used in some skin, mouth, or wound products but can irritate tissue and does not replace assessment of persistent bleeding or serious injury.

\medterm{Astrocyte} A star-shaped support cell in the brain and spinal cord. It helps nourish nerve cells, maintain the blood–brain barrier, regulate the chemical environment, and respond to injury.

\medterm{Astrocytoma} A tumour arising from astrocytes, the star-shaped support cells of the brain or spinal cord. Astrocytomas vary from slow-growing to aggressive, and symptoms depend on their location and size.

\medterm{Astrovirus Infection} An acute viral infection of the gastrointestinal tract caused by human astroviruses, a leading cause of non-bacterial gastroenteritis primarily in infants, young children, the elderly, and immunocompromised patients. It presents with watery diarrhea, mild abdominal pain, low-grade fever, and vomiting, typically resolving spontaneously within 2 to 4 days with oral rehydration therapy.

\synonyms
Astrovirus Gastroenteritis; Human Astrovirus Infection

\medterm{Asymmetrical} Unequal in shape, size, distribution, or appearance between corresponding sides. Asymmetry in a pigmented skin lesion can be a warning sign that needs assessment, especially with other changes.

\medterm{Asymptomatic} Having a disease, infection, or abnormal finding without noticeable symptoms. A person may still have detectable disease, develop complications, or transmit an infection; this differs from presymptomatic illness, in which symptoms appear later.

\medterm{Asymptomatic Adults} Adults who have no noticeable symptoms of the condition being assessed. Absence of symptoms does not exclude early or hidden disease, risk factors, or the ability to transmit an infection.

\medterm{Asymptomatic Bacteriuria} Bacteria detected in a properly collected urine sample when the person has no urinary symptoms. It does not always require treatment, but management is important in selected situations such as pregnancy or before some urinary-tract procedures.

\medterm{Asymptomatic Bacteriuria In Pregnancy} Bacteria growing in a pregnant person’s urine without burning, urgency, or other urinary symptoms. It is identified by urine culture and is important because untreated infection increases the risk of kidney infection and pregnancy complications.

\medterm{Asymptomatic Carrier} A person who carries an infectious organism without symptoms but may still transmit it to others. Whether transmission occurs depends on the organism, body site, and contact.

\medterm{Asymptomatic HIV Infection} HIV infection without noticeable symptoms at present. HIV can still multiply and gradually weaken the immune system, and the person may transmit it without effective treatment.

\medterm{Asymptomatic Infections} Infections in which the person has no noticeable symptoms. The organism may still be detectable, cause later complications, or spread to others; this differs from presymptomatic infection, in which symptoms develop later.

\synonyms
Subclinical Infections

\medterm{Asystole} Absence of effective electrical activity in the heart’s lower chambers and absence of useful pumping. It causes cardiac arrest and requires immediate resuscitation; a flat-looking ECG line must first be checked to ensure it is not an equipment or lead problem.

\medterm{Ataxia} Loss of smooth control and coordination of movement. It is a sign rather than one specific disease and may cause an unsteady walk, poor balance, clumsy movements, slurred speech, shaking, or abnormal eye movements. It can result from a problem in the cerebellum, spinal cord, sensory nerves, or balance organs and is different from muscle weakness alone. Causes include inherited disorders, stroke, multiple sclerosis, infection, vitamin deficiency, medicines, alcohol, or other toxins. The onset may be sudden or gradual, and the course depends on the cause.

\medterm{Ataxia Telangiectasia} A rare autosomal-recessive DNA-repair disorder causing progressive cerebellar ataxia, telangiectasias of the eyes or skin, immune deficiency, radiosensitivity, and increased risk of leukemia, lymphoma, or other cancers. Respiratory infections and elevated alpha-fetoprotein may occur.

\synonyms
A-T; Ataxia-Telangiectasia

\medterm{Ataxic Breathing} An irregular, unpredictable pattern in which the rate and depth of breaths vary randomly. It usually indicates severe injury to the lower brainstem and may precede breathing failure.

\medterm{Ataxic Gait} An abnormal, unsteady, and poorly coordinated walking pattern resulting from dysfunction in the motor coordination or sensory balance pathways. In cerebellar ataxia, the gait is characteristically wide-based, irregular, lurching, and clumsy, with difficulty maintaining a straight line and inability to perform tandem walking. In sensory ataxia (caused by dorsal column or peripheral proprioceptive loss), the gait is wide-based with high-stepping stamping of the feet, markedly worsening with eye closure (positive Romberg sign).
\synonyms{Cerebellar Gait; Sensory Ataxic Gait; Uncoordinated Gait}

\medterm{Atelectasis} Partial or complete collapse of a region of lung, reducing the amount of air it contains. It may result from blocked airways, pressure from outside the lung, shallow breathing, or scarring and can reduce oxygen exchange.

\synonyms
Lung Collapse

\medterm{Athelia} Absence of one or both nipples, present from birth or acquired after surgery, injury, or disease. It may occur alone or with other differences in breast or chest-wall development.

\medterm{Atherectomy} A minimally invasive endovascular catheter technique utilizing directional, rotational, orbital, or laser devices to physically excise, shave, vaporize, or pulverize calcified or atheromatous plaque from the intraluminal wall of diseased arteries.

\synonyms
Endovascular Atherectomy; Directional Atherectomy; Rotational Atherectomy; Orbital Atherectomy

\medterm{Atheroembolic Renal Disease} Kidney injury caused when cholesterol crystals break from an artery plaque and block small kidney vessels, often after a vascular procedure. It may cause worsening kidney function, mottled skin, painful toes, or raised eosinophils.

\medterm{Atheroma} A fatty, inflamed buildup within an artery wall that forms part of an atherosclerotic plaque. It may narrow the artery or rupture and trigger a blood clot.

\synonyms
Atheromatous Plaque; Arterial Plaque

\medterm{Atheromatous Plaque} An asymmetrical, fibrofatty lesion developing within the tunica intima of medium and large arteries, comprising a necrotic lipid-rich core of oxidized cholesterol crystals and apoptotic debris covered by a fibrous cap of collagen and vascular smooth muscle cells.

\synonyms
Atheroma; Atherosclerotic Plaque

\medterm{Atherosclerosis} A disease in which fatty, cholesterol-rich plaques build up inside artery walls. Plaques can narrow arteries or rupture and form clots, reducing blood flow to the heart, brain, legs, or other organs.

\medterm{Atherosclerosis Pathology} A chronic arterial intimal disease characterized by lipid-rich plaques with a fibrous cap, inflammation, smooth-muscle proliferation, calcification, and possible ulceration or thrombosis.

\medterm{Atherosclerotic} Relating to fatty plaque buildup and hardening inside artery walls. Atherosclerotic changes can narrow an artery or encourage a clot to form.

\medterm{Atherosclerotic Heart Disease} A disease in which cholesterol-containing plaque builds up inside the heart’s arteries, narrowing them and reducing blood flow. It may cause chest pain, heart attack, or other heart problems.

\medterm{Atherosclerotic Plaque} A buildup of cholesterol, fat, inflammatory cells, and scar-like tissue inside an artery wall. It can gradually narrow the artery or rupture and cause a blood clot.

\medterm{Atherothrombotic Stroke} An ischemic stroke caused by a clot forming on fatty plaque in a large artery supplying the brain or travelling from that plaque into a brain artery. It deprives brain tissue of oxygen.

\synonyms
Large-Artery Atherosclerotic Stroke; Atherothrombotic Ischemic Stroke

\medterm{Athetosis} Slow, involuntary, writhing movements that may affect the hands, fingers, feet, face, neck, or trunk. They result from abnormal function in brain areas that help control movement and may occur with inherited or acquired brain injury.

\medterm{Athlete Foot} Alternate name; see \textit{Athlete's Foot}.

\medterm{Athlete's Foot} A common superficial fungal infection (dermatophytosis) of the skin of the feet, most commonly caused by \textit{Trichophyton rubrum} or \textit{Trichophyton interdigitale}. Classically presents with pruritus, erythema, scaling, maceration, and fissuring, especially in the interdigital web spaces.

\synonyms
Tinea Pedis; Foot Ringworm; Foot Dermatophytosis

\medterm{Athlete's Heart} Normal physiological cardiac remodeling induced by intensive, prolonged physical endurance or strength training, manifested as physiological left ventricular hypertrophy, chamber dilation, and resting sinus bradycardia, without diastolic dysfunction.

\synonyms
Athletic Heart Syndrome; Athlete Heart Remodeling

\medterm{Atlanto-Occipital Dislocation} A severe injury in which the first neck bone separates from the base of the skull. It is extremely unstable and can damage the spinal cord or brainstem, so the neck must be immobilized and emergency care provided immediately.

\medterm{Atlanto-Occipital Joints} The paired joints between the upper surface of the atlas, the first cervical vertebra, and the occipital condyles at the base of the skull. They mainly allow nodding and small side-to-side movements while helping stabilize the head and upper neck.

\medterm{Atlantoaxial Instability} Excessive movement between the first and second cervical vertebrae, sometimes caused by trauma, rheumatoid arthritis, congenital conditions, or ligament weakness. It may cause neck pain, occipital headache, neurologic symptoms, or spinal-cord compression.

\medterm{Atlantoaxial Instability in Down Syndrome} Excessive movement between the first and second cervical vertebrae in a person with Down syndrome, often related to ligament laxity or altered bone development. Many people have no symptoms, but neck pain, weakness, gait change, or other neurologic signs require prompt assessment before activities that stress the neck.

\medterm{Atlantoaxial Joint} The complex synovial articulation between the first cervical vertebra (atlas, C1) and the second cervical vertebra (axis, C2). Comprising one median pivot joint (dens with anterior arch) and two lateral facet joints, it provides approximately 50\% of total cervical rotational range of motion. Instability may arise from trauma (odontoid fractures), rheumatoid arthritis, or Down syndrome.

\synonyms
Atlas-Axis Joint; C1-C2 Articulation

\medterm{Atlas} The first bone of the neck, supporting the skull and helping the head nod. It is ring-shaped and has openings through which important blood vessels pass.

\medterm{Atom} The smallest unit of a chemical element that retains the element’s properties. It contains a central nucleus surrounded by electrons.

\medterm{Atonic} Lacking normal muscle tone or ability to contract. For example, an atonic bladder does not squeeze effectively to empty urine.

\medterm{Atonic Seizure} A seizure causing sudden loss of muscle tone, often for only a few seconds. The person may drop their head, knees, or whole body and fall; injuries can occur during the fall.

\synonyms
Drop Attack; Akinetic Seizure

\medterm{Atopic} Having a tendency to develop allergic conditions such as eczema, allergic rhinitis, or asthma. This tendency may run in families and reflects increased sensitivity of the immune system.

\medterm{Atopic Dermatitis} A chronic, relapsing, intensely itchy inflammatory skin disease involving epidermal-barrier dysfunction and immune dysregulation. It causes dry, inflamed, scaly, or lichenified lesions whose distribution varies with age; scratching can produce infection, excoriations, and sleep disturbance.

\synonyms
Atopic Eczema; Eczema; Infantile Eczema; Allergic Dermatitis

\medterm{Atopic Eruption Of Pregnancy} An itchy eczema-like rash that begins or worsens during pregnancy, often on the face, neck, chest, or skin folds. It usually improves after delivery and is not harmful to the fetus.

\medterm{Atopic Rhinitis} Allergy-related inflammation inside the nose. Exposure to pollen, dust mites, mould, animals, or other triggers can cause sneezing, itching, a blocked or runny nose, and reduced smell.

\synonyms
Allergic Rhinitis

\medterm{Atopy} A genetic tendency to develop allergic conditions such as eczema, allergic rhinitis, asthma, or food allergy after exposure to common triggers.

\medterm{Atorvastatin} A medicine that lowers low-density lipoprotein cholesterol by reducing cholesterol production in the liver. It lowers the risk of heart attack and stroke in appropriate people; muscle symptoms and liver problems are uncommon but important side effects.

\medterm{Atovaquone–Proguanil} A combination of medicines used to prevent or treat some forms of malaria. The medicines interfere with energy production and folate use in malaria parasites; suitability depends on the type of malaria and the person’s health.

\medterm{ATP} Adenosine triphosphate, the main short-term energy-carrying molecule in cells. Cells release usable energy when ATP is broken down and remake it during metabolism.

\medterm{ATP Synthase} A protein machine in mitochondria and some other cell membranes that uses the movement of hydrogen ions to make ATP, the main molecule cells use for immediate energy.

\medterm{Atransferrinemia} A very rare inherited disorder in which the blood has little or no transferrin, the protein that carries iron. It can cause severe anemia, poor growth, and harmful iron buildup in organs.

\medterm{Atresia} A birth defect in which a normal opening or passage is absent, closed, or severely narrowed. Examples include intestinal, bile-duct, heart-valve, and ear-canal atresia; effects depend on the structure involved.

\medterm{Atresia Of Bile Ducts} See \textit{Biliary atresia}.

\medterm{Atria} The two upper chambers of the heart. The right atrium receives blood from the body and the left atrium receives blood from the lungs; each passes blood to the lower pumping chamber on its side.

\medterm{Atrial} Relating to either of the heart’s two upper chambers, which receive blood and pass it to the lower pumping chambers.

\medterm{Atrial Ablation} A procedure that uses heat, cold, or another form of energy to create small scars in heart tissue responsible for an abnormal atrial rhythm. The scars block unwanted electrical signals.

\synonyms
Cardiac Ablation

\medterm{Atrial Arrhythmia} An abnormal heart rhythm that begins in one of the heart’s upper chambers. It may cause palpitations, breathlessness, dizziness, chest discomfort, or fainting; examples include atrial fibrillation and atrial flutter.

\medterm{Atrial Fib} Alternate name; see \textit{Atrial Fibrillation}.

\medterm{Atrial Fibrillation} A supraventricular tachyarrhythmia characterized by rapid, disorganized, and chaotic electrical activation of the atria without effective mechanical atrial contraction. Electrocardiographically manifested by absent P waves and irregularly irregular QRS complexes, it predisposes to left atrial appendage thrombus formation, thromboembolic stroke, and heart failure. Management encompasses rate control, rhythm control, and systemic anticoagulation based on CHA2DS2-VASc risk scoring.

\synonyms
AF; AFib; Auricular Fibrillation

\medterm{Atrial Fibrillation Ablation} A catheter procedure that creates small scars, usually around the pulmonary veins, to block electrical signals that trigger atrial fibrillation. It may reduce symptoms, but recurrence and procedure-related risks remain possible.

\medterm{Atrial Fibrillation Medicines} Medicines used for atrial fibrillation may slow the heart rate, restore or maintain a regular rhythm, or reduce stroke risk by preventing clots. Choice depends on symptoms, heart structure, kidney function, bleeding risk, and other conditions.

\medterm{Atrial Flutter} A supraventricular arrhythmia characterized by a rapid, organized, macro-reentrant electrical circuit within the atria, typically circulating around the cavotricuspid isthmus of the right atrium. Manifested on ECG as rapid, regular sawtooth flutter waves (typically 250–350 bpm) with variable AV nodal conduction block (often 2:1), it presents with palpitations, fatigue, and elevated thromboembolic risk.

\synonyms
AFlutter; Cavotricuspid Isthmus-Dependent Flutter

\medterm{Atrial Myxoma} A usually noncancerous tumour arising inside the heart, most often in the left atrium. It may obstruct blood flow, cause breathlessness or fainting, or release fragments that block blood vessels.

\medterm{Atrial Natriuretic Peptide} A hormone released by the heart’s upper chambers when they stretch from extra blood. It tells the kidneys to remove more salt and water and helps relax blood vessels, lowering blood volume and pressure.

\medterm{Atrial Septal Defect} A hole present from birth in the wall between the heart’s upper chambers. A small hole may cause no problems, while a larger one can overwork the right side of the heart, raise lung pressure, or increase the risk of abnormal rhythms and stroke.

\medterm{Atrial Septum} The wall separating the heart’s right and left upper chambers. A hole in this wall is called an atrial septal defect.

\medterm{Atrial Standstill} Failure of the heart’s upper chambers to produce normal electrical signals or contractions. It can cause a very slow heartbeat, dizziness, fainting, poor circulation, or heart failure.

\medterm{Atrial Switch Operation} An older operation for transposition of the major heart arteries that redirects blood inside the atria using a tunnel-like passage. It restores circulation but leaves the right pumping chamber supporting the body and may cause later rhythm or heart-muscle problems.

\synonyms
Mustard Procedure; Senning Procedure

\medterm{Atrial Tachycardia} A fast heart rhythm that starts in an area of an upper chamber outside the normal pacemaker. It may cause palpitations, breathlessness, dizziness, chest discomfort, or fainting.

\medterm{Atrial Thrombosis} Formation of a blood clot inside one of the heart’s upper chambers, especially during atrial fibrillation. The clot can break loose and block blood flow to the brain or another organ.

\medterm{Atrioventricular} Relating to the connection between a heart’s upper chamber and the lower chamber on the same side. The term includes the mitral and tricuspid valves and the electrical pathway that coordinates the timing of their contractions.

\medterm{Atrioventricular Block} Delayed or interrupted electrical signals between the heart’s upper and lower chambers. It may cause a slow heartbeat, tiredness, dizziness, fainting, or no symptoms; severe block can require urgent treatment and pacing.

\medterm{Atrioventricular Bundle} A bundle of specialized electrical fibers that carries signals from the heart’s central relay point into the lower chambers. Its branches help both pumping chambers contract in a coordinated way.

\synonyms
Bundle of His

\medterm{Atrioventricular Canal Defect} A birth defect in which the wall between the heart chambers and the valves between upper and lower chambers do not form normally. It can allow abnormal blood flow and often causes heart failure or poor growth in infancy.

\synonyms
Atrioventricular Septal Defect (AVSD); Endocardial Cushion Defect

\medterm{Atrioventricular Dissociation} A cardiac rhythm abnormality in which the atria and ventricles depolarize and contract independently under the control of separate pacemakers, occurring when ventricular pacing is faster than atrial rate or during complete atrioventricular block.

\synonyms
AV Dissociation

\medterm{Atrioventricular Nodal Reentrant Tachycardia} A sudden, regular, fast heartbeat caused by an electrical loop near the heart’s central relay point. Episodes may cause pounding heartbeat, chest discomfort, breathlessness, dizziness, or fainting and may stop as suddenly as they begin.

\synonyms
AV Nodal Reentrant Tachycardia; AVNRT

\medterm{Atrioventricular Nodal Reentry Tachycardia} A rapid, regular heartbeat caused by a looping electrical signal near the atrioventricular node. It commonly begins and ends suddenly and may cause palpitations, dizziness, breathlessness, or fainting.

\medterm{Atrioventricular Node} A small area of specialized tissue near the junction of the heart’s upper and lower chambers that receives electrical signals and briefly delays them before passing them to the lower chambers. This delay allows the upper chambers to empty first.

\medterm{Atrioventricular Septal Defect} Alternate name; see \textit{Atrioventricular canal defect}.

\medterm{Atrioventricular Valves} The two heart valves between the upper and lower chambers: the tricuspid valve on the right and the mitral valve on the left. They open to let blood pass forward and close to prevent backflow.

\medterm{Atrium} One of the heart’s two upper chambers, which receive returning blood and pass it to the lower pumping chambers. The right atrium receives blood from the body; the left receives oxygen-rich blood from the lungs.

\medterm{Atrophic Glossitis} An inflammatory condition of the tongue characterized by the diffuse atrophy and loss of lingual filiform and fungiform papillae, resulting in a smooth, shiny, erythematous, and \"beefy red\" appearance often accompanied by lingual burning, pain, or altered taste. It is a key physical finding in nutritional deficiencies including vitamin B12 (Hunter's glossitis), folate, iron (Plummer-Vinson syndrome), niacin, and riboflavin, as well as systemic conditions such as celiac disease and severe protein-calorie malnutrition.
\synonyms{Hunter's Glossitis; Bald Tongue; Smooth Tongue; Moeller-Hunter Glossitis}

\medterm{Atrophic Vaginitis} Alternate name; see \textit{Vaginal Atrophy}.

\medterm{Atrophoderma Maculatum} A rare skin condition causing flat, slightly sunken patches in which the skin becomes thinner and changes colour. The patches usually have no open sore or active infection.

\medterm{Atrophy} Shrinkage and loss of function of a tissue, organ, or muscle because its cells become smaller or fewer. Causes include disuse, nerve damage, poor blood supply, malnutrition, hormone deficiency, chronic disease, and ageing; it may be partly reversible.

\medterm{Atrophy, Vaginal} Alternate index label; see \textit{Vaginal Atrophy}.

\medterm{Atropine} A medicine that blocks acetylcholine signals at certain nerve endings. It can increase heart rate, reduce secretions, widen the pupils, and counteract selected poisonings or dangerously slow heartbeats.

\medterm{Atropine In Organophosphate Poisoning} A medicine used in organophosphate poisoning to block excessive acetylcholine effects, such as wet secretions, narrowed airways, slow heartbeat, vomiting, and diarrhoea. Severe poisoning requires emergency treatment and monitoring.

\medterm{Attachment Disorder} See \textit{Reactive Attachment Disorder}.

\medterm{Attachment, Dental} A precision mechanical or magnetic interlocking device used to secure, stabilize, and retain a removable partial denture or overdenture to retained tooth roots, natural abutments, or dental implants.

\synonyms
Denture Attachment; Precision Attachment; Overdenture Attachment; Dental Attachment

\medterm{Attack Rate} The proportion of people at risk who develop a particular illness during a defined outbreak or exposure period. It is commonly used to describe foodborne outbreaks and other limited events.

\medterm{Attention} The ability to select information, focus on it, and maintain that focus while limiting distraction. Attention can be affected by sleep loss, stress, illness, medicines, or brain disorders.

\medterm{Attention Deficit Hyperactivity Disorder} Alternate index label; see \textit{Attention-Deficit/Hyperactivity Disorder}.

\medterm{Attention-Deficit/Hyperactivity Disorder} A neurodevelopmental disorder involving persistent inattention, hyperactivity, and impulsive behaviour that affects functioning or development. Presentations may be mainly inattentive, mainly hyperactive-impulsive, or combined.

\synonyms
ADHD; Hyperkinetic Disorder

\medterm{Attenuated} Weakened or reduced in force, severity, concentration, or biological activity. An attenuated vaccine contains a weakened form of a microorganism that is intended to stimulate immunity without usually causing disease.

\medterm{Attenuation Correction} A computer adjustment for the way body tissues weaken signals during medical imaging. It helps make PET, SPECT, and combined scan images more accurate by compensating for signal loss.

\medterm{Attributable Risk} The difference in disease risk between people exposed and not exposed to a possible cause. It estimates the extra risk associated with the exposure, but does not by itself prove causation.

\medterm{Atypia} An unusual appearance of cells or tissue seen under a microscope. It may result from irritation, inflammation, ageing, a precancerous change, or cancer and must be interpreted with the tissue pattern and clinical context.

\medterm{Atypia Morphology} An unusual appearance or arrangement of cells or tissue under a microscope. It may result from irritation or inflammation, or represent a precancerous or cancerous change; its meaning depends on the whole sample and clinical setting.

\synonyms
Cytologic Atypia

\medterm{Atypical} Different from the usual appearance, behaviour, structure, or course of a condition. Atypical does not by itself mean that a finding is harmless, precancerous, or cancerous.

\medterm{Atypical (Dysplastic) Nevus} A melanocytic mole with unusual clinical or microscopic features such as asymmetry, irregular border, varied color, or larger size. Most are benign, but multiple atypical nevi increase melanoma risk; a changing, bleeding, or distinctly different lesion should be examined and may require biopsy.

\medterm{Atypical Antipsychotics} A group of medicines used mainly for schizophrenia, bipolar disorder, and selected other conditions. They change several brain signals and can cause sleepiness, weight gain, movement problems, high blood sugar, or other metabolic changes.

\medterm{Atypical Endometrial Hyperplasia} An overgrowth of abnormal cells in the lining of the uterus. It is not cancer, but it can develop into uterine-lining cancer or occur beside cancer, so medical assessment and follow-up are important.

\medterm{Atypical Facial Pain} Persistent or recurring pain in the face that does not fit a clear nerve, dental, sinus, or other recognized cause. Evaluation looks for treatable causes before considering persistent facial-pain disorders.

\medterm{Atypical Features} A pattern sometimes used to describe depression in which mood can improve temporarily in response to positive events, with increased sleep, increased appetite or weight, heavy-feeling limbs, and marked sensitivity to rejection.

\medterm{Atypical Fibroxanthoma} A low-grade cutaneous sarcoma that usually arises as a rapidly growing red, pink, or flesh-coloured nodule on sun-damaged head or neck skin of an older adult. Diagnosis requires biopsy because it can mimic squamous-cell carcinoma or melanoma; complete excision or Mohs surgery is typical treatment.

\medterm{Atypical Genitalia} External genital anatomy that is not typically classified as male or female. It may result from differences in chromosomes, hormones, gonads, or genital development and requires respectful specialist assessment.

\medterm{Atypical Hemolytic Uremic Syndrome} A rare disease in which inherited or acquired abnormalities of the complement system cause uncontrolled immune activity and damage to small blood vessels, especially in the kidneys. It can destroy red blood cells, lower platelets, and cause sudden kidney injury, high blood pressure, or neurological problems. Unlike typical hemolytic uremic syndrome, it is not primarily caused by a toxin-producing diarrheal infection.

\synonyms
aHUS; Complement-Mediated HUS

\medterm{Atypical Hyperplasia Of The Breast} An overgrowth of unusual cells in a breast duct or lobule that is not cancer. It can increase the future risk of breast cancer and may require closer screening or preventive discussion.

\medterm{Atypical Lymphocyte} A lymphocyte, a type of white blood cell, with an unusual size or appearance under a microscope. It often reflects immune activity during a viral infection, but the finding alone does not identify a specific disease.

\medterm{Atypical Mole} A mole with unusual clinical or microscopic features, also called a dysplastic nevus. Most are noncancerous, but multiple atypical moles can increase melanoma risk; a changing, bleeding, or distinctly different mole needs examination.

\medterm{Atypical Nevi} Alternate plural name; see \textit{Dysplastic Nevi}.

\medterm{Atypical Pneumonia} A lung infection that often develops gradually and may cause more general symptoms than chest symptoms, such as headache, tiredness, sore throat, or muscle aches. Causes include certain bacteria and viruses.

\medterm{Atypical Spitz Nevus} An unusual mole made of pigment-producing cells that resembles a Spitz nevus but has concerning features. It can resemble melanoma under a microscope, so expert skin examination and tissue review may be needed.

\medterm{Atypical Teratoid Rhabdoid Tumor} A rare, fast-growing cancer of the brain or spinal cord that occurs mainly in infants and young children. It may cause vomiting, headache, sleepiness, balance problems, seizures, or developmental changes.

\synonyms
ATRT; Rhabdoid Brain Tumor

\medterm{Audiogram} A chart that shows a person’s ability to hear at different pitches or frequencies.

\medterm{Audiologist} A health professional who assesses hearing and fits hearing aids.

\medterm{Audiology} The healthcare field concerned with hearing, balance, and related communication problems. Audiologists test hearing and balance, provide hearing aids and other devices, offer training and support, and identify when medical or surgical care may be needed.

\medterm{Audiometry} Measurement of hearing sensitivity at different sound frequencies and intensities, usually using headphones or insert earphones.

\medterm{Auditory Acuity} The sensitivity or sharpness of hearing, commonly assessed by determining the softest sounds a person can detect at different frequencies.

\medterm{Auditory Area} A region of the cerebral cortex that processes sound information.

\synonyms
Auditory Cortex

\medterm{Auditory Fatigue} A temporary decrease in hearing sensitivity or listening ability after prolonged loud sound exposure; rest usually helps, but repeated exposure can cause lasting damage.

\medterm{Auditory Hallucinations} False sensory perceptions of sound occurring in the absence of any external acoustic stimulus. They may be non-verbal (such as clicks, ringing, or music) or verbal (such as single words, conversations, running commentary, or command voices). Verbal auditory hallucinations are a hallmark feature of psychotic disorders such as schizophrenia, but can also occur in severe mood disorders with psychotic features, delirium, temporal lobe epilepsy, sensory deprivation, and substance withdrawal.
\synonyms{Paracusia; Phonemes; Auditory Sensory Hallucinations}

\medterm{Auditory Information Processing Disorder} See \textit{Auditory processing disorder}.

\medterm{Auditory Nerve} The portion of the vestibulocochlear nerve that carries hearing signals from the inner ear to the brain.

\medterm{Auditory Neuropathy Spectrum Disorder} A form of sensorineural hearing impairment characterized by preserved cochlear outer hair cell function (demonstrated by normal otoacoustic emissions or cochlear microphonics) but abnormal, dyssynchronous neural transmission along the eighth cranial nerve.

\synonyms
ANSD; Auditory Neuropathy

\medterm{Auditory Ossicle} One of three small middle-ear bones—the malleus, incus, and stapes—that transmit and amplify vibrations from the eardrum to the inner ear. Disease, injury, or fixation of these bones can cause conductive hearing loss.

\synonyms
Ear Ossicle

\medterm{Auditory Pathway} A chain of nerve structures carrying sound signals from the inner ear through the brainstem to the auditory cortex for recognition and interpretation.

\medterm{Auditory Perception} The brain's process of detecting, recognizing, and interpreting sounds received through the auditory system.

\medterm{Auditory Processing Disorder} Difficulty recognizing or interpreting spoken sounds even though basic hearing sensitivity may be normal. Affected people may struggle to understand speech in noise, follow rapid instructions, or distinguish similar sounds; assessment requires specialized testing and management may include communication strategies and support.

\synonyms
APD

\medterm{Auditory Threshold} The quietest sound a person can hear at a particular pitch or frequency. It is measured during a hearing test and helps describe the degree and type of hearing loss in each ear.

\synonyms
Hearing Threshold

\medterm{Auer Rods} Needle-shaped structures sometimes seen inside immature blood cells under a microscope. They are made from fused cell granules and strongly suggest a cancer of immature myeloid blood cells, especially acute myeloid leukemia.

\medterm{Auerbach Plexus} An extensive intrinsic neural network of unmyelinated nerve fibers and postganglionic parasympathetic ganglia situated in the muscularis externa between the circular and longitudinal smooth muscle layers throughout the gastrointestinal tract, coordinating peristalsis and motor activity.

\synonyms
Myenteric Plexus

\medterm{Auger Electrons} Very low-energy electrons released by some radioactive atoms after a change inside the atom. They travel only a tiny distance and can damage nearby DNA, a property being studied for precisely targeted radiation treatment.

\medterm{Augmentation Mammaplasty} An operation to increase breast size or change breast shape, usually with implants or transferred fat. Possible risks include infection, bleeding, changes in sensation, scarring, and the need for further surgery.

\medterm{Augmentation Mammoplasty} An operation to increase breast size or change breast shape, usually with implants or transferred fat. Possible risks include infection, bleeding, changes in sensation, scarring, and the need for further surgery.
\synonyms
Breast Augmentation

\medterm{Augmentation Of Labor} Treatment used when labour has started but contractions are not strong or frequent enough for the cervix to continue opening. It may involve oxytocin or deliberately breaking the waters, with monitoring of parent and baby.

\medterm{Aura} A temporary change in vision, sensation, speech, movement, or awareness that may occur before or during a migraine or seizure. Symptoms usually develop gradually and then resolve, but a new or prolonged aura needs urgent assessment.

\medterm{Aural Fullness} A feeling of pressure, blockage, or fullness in the ear. It can result from poor pressure equalization, fluid, earwax, inner-ear disease, or jaw-joint problems and may occur with hearing loss, ringing, or dizziness.

\synonyms
Ear Pressure; Plugged-Ear Sensation; Blocked-Ear Sensation

\medterm{Aural Polyps} Soft, abnormal tissue growths in the ear canal or middle ear. They often result from long-lasting infection or inflammation, but persistent discharge, pain, or bleeding needs specialist examination to exclude a more serious cause.

\medterm{Auricle} The visible outer part of the ear attached to the side of the head. Its ridges collect and direct sound into the ear canal, helping the brain judge where sounds come from.

\medterm{Auricular Hematoma} A collection of blood between the cartilage and skin covering of the outer ear, usually after a blow or sports injury. Prompt drainage and compression help prevent permanent thickening and deformity of the ear.

\synonyms
Pinna Hematoma; Subperichondrial Auricular Hematoma

\medterm{Auscultation} Listening to sounds produced inside the body, usually with a stethoscope. Clinicians use it to assess the heart, lungs, abdomen, and blood vessels for findings such as murmurs, wheezing, crackles, or abnormal blood-flow sounds.

\medterm{Auscultatory Gap} A temporary disappearance of the sounds heard while a blood-pressure cuff is slowly deflated, followed by their return. If it is missed, the top blood-pressure number may be recorded falsely low.

\medterm{Auspitz Sign} A skin check for psoriasis where tiny pinpoint bleeding spots appear when the silvery scales of a rash are gently scraped away.

\medterm{Austin Flint Murmur} A low-pitched sound heard during the middle part of the heartbeat over the heart's apex. It is associated with severe leakage of the aortic valve and altered filling of the mitral valve.

\medterm{Autism} Alternate name; see \textit{Autism Spectrum Disorder}.

\medterm{Autism Spectrum Disorder} A lifelong neurodevelopmental condition involving persistent differences in social communication and interaction, along with restricted or repetitive behaviors, interests, or sensory patterns. Features begin during the developmental period, vary widely among individuals, and may affect communication, learning, relationships, daily activities, or the type and amount of support needed.

\synonyms
ASD; Autistic Spectrum Disorder; Pervasive Developmental Disorder

\medterm{Autism Spectrum Disorder, Childhood Disintegrative Disorder} Alternate index label; see \textit{Childhood Disintegrative Disorder}.

\medterm{Autistic Savant} An autistic person with an unusually strong, specific ability—such as music, art, calculation, or memory—compared with their other abilities. The term describes a pattern of skills and is not itself a medical diagnosis.

\medterm{Autoantibodies} Immune proteins that mistakenly recognize and attach to the body's own cells, tissues, or proteins. They may cause disease or serve as test clues; their presence alone does not prove that a person has an autoimmune disease.

\medterm{Autoantibody} An immune protein that mistakenly recognizes and binds a substance from the person's own body. It may damage tissues or interfere with their function, or it may be present only as a laboratory sign without causing disease.

\medterm{Autoantigen} A substance normally found in the body that is mistakenly recognized as foreign by the immune system. This mistaken response can lead to inflammation or damage to the person's own cells or tissues.

\medterm{Autoclave} A machine that sterilizes equipment by exposing it to pressurized steam at a controlled temperature and time. It destroys bacteria, viruses, fungi, and spores when the correct cycle is used.

\medterm{Autocrine Signaling} A mode of intercellular and intracellular communication in which a cell synthesizes and secretes a messenger molecule (such as a cytokine, hormone, or growth factor) that binds to surface receptors on the very same cell that produced it, initiating a feedback response.

\medterm{Autoerythrocyte Sensitivity} A rare condition causing painful bruises or raised purple patches, thought to involve an abnormal reaction to a person's own red-cell components. Other bleeding disorders, medicines, and causes of bruising must be excluded first.

\medterm{Autogenic Drainage} A controlled-breathing method that moves mucus from the smaller airways toward the throat so it can be coughed out. It uses different breathing depths and requires instruction and practice.

\medterm{Autograft} Tissue moved from one part of a person's body to another part of the same body. Because the tissue is genetically the person's own, rejection is unlikely, but the original site may have pain, scarring, or other complications.

\medterm{Autoimmune} Describing an immune response that mistakenly targets the body's own cells, tissues, or proteins. It can cause inflammation, tissue damage, or impaired organ function.

\medterm{Autoimmune Associated Epilepsy} See \textit{Autoimmune epilepsy}.

\medterm{Autoimmune Autonomic Ganglionopathy} A rare, severe autoimmune disorder of the autonomic nervous system characterized by widespread sympathetic and parasympathetic failure. Manifests clinically with severe orthostatic hypotension, anhidrosis, dry mouth and eyes, tonic pupils, gastrointestinal dysmotility (gastroparesis, intestinal pseudo-obstruction), and urinary retention; strongly associated with autoantibodies directed against the ganglionic nicotinic acetylcholine receptor ($\alpha_3\beta_4$ subtype).

\synonyms
AAG; Subacute Autonomic Neuropathy

\medterm{Autoimmune Disease} A disease in which the immune system mistakenly attacks the body's own cells, tissues, or organs. It may cause inflammation, pain, tissue damage, or loss of organ function.

\medterm{Autoimmune Diseases} Diseases caused by immune responses against the body's own tissues. They may affect one organ, such as the thyroid, or many parts of the body and can produce inflammation and lasting damage.

\medterm{Autoimmune Disorder} A condition in which the immune system mistakenly attacks the body's own cells, tissues, or organs. It may be limited to one organ or affect many organs, causing inflammation and impaired function.

\medterm{Autoimmune Encephalitis} Inflammation of the brain caused by an immune attack on nerve cells or their proteins. It may cause seizures, confusion, memory problems, unusual movements, behaviour changes, sleep problems, or reduced consciousness.

\medterm{Autoimmune Encephalopathy} See \textit{Autoimmune Encephalitis}.

\medterm{Autoimmune Endocrine Disease} An autoimmune disorder that damages a hormone-producing gland or hormone-producing cells. Examples include autoimmune thyroid disease, Addison disease, and type 1 diabetes; it may cause too little or, less often, abnormal hormone production.

\medterm{Autoimmune Enteropathy} A rare disease in which the immune system damages the lining of the small intestine, causing ongoing diarrhoea, poor nutrient absorption, weight loss, or growth problems. It may occur with other immune disorders.

\medterm{Autoimmune Epilepsy} Recurrent seizures caused by an immune attack on the brain. Seizures may begin suddenly or be difficult to control and can occur with memory problems, confusion, movement changes, or other signs of brain inflammation.

\medterm{Autoimmune Gastritis} Long-term inflammation that damages acid- and vitamin-B12-producing cells in the stomach. It can cause low stomach acid, iron deficiency, vitamin B12 deficiency, anaemia, digestive symptoms, and an increased risk of certain stomach growths.

\medterm{Autoimmune Hemolytic Anemia} Anaemia caused when immune proteins attach to and destroy red blood cells faster than the body can replace them. It may cause tiredness, pale or yellow skin, dark urine, rapid heartbeat, or shortness of breath.

\medterm{Autoimmune Hepatitis} Long-term liver inflammation caused by an immune attack on liver cells. It may cause tiredness, nausea, joint pain, or jaundice, while some people have no symptoms; untreated disease can lead to scarring and liver failure.

\medterm{Autoimmune Keratitis} Inflammation of the clear front surface of the eye caused by or associated with immune disease. It may cause eye pain, redness, light sensitivity, blurred vision, or a corneal scar and needs prompt eye assessment.

\medterm{Autoimmune Liver Disease Panel} A group of blood tests that looks for immune proteins linked with autoimmune liver diseases. Results are interpreted with liver tests, symptoms, examination, and sometimes a liver biopsy; the panel alone cannot establish a diagnosis.

\medterm{Autoimmune Lymphoproliferative Syndrome} An inherited immune disorder in which lymphocytes do not die off normally after an immune response, causing noncancerous enlargement of lymph nodes and spleen from early childhood. Autoimmune destruction of blood cells may cause anemia, low platelets, or low white-cell counts. It is often due to a defect in the Fas cell-death pathway; diagnosis uses blood counts, specialized immune-cell testing, and gene testing, and there is an increased risk of lymphoma.

\medterm{Autoimmune Pancreatitis} A rare form of pancreas inflammation caused by an abnormal immune response. It may cause painless jaundice, abdominal pain, weight loss, or a mass that resembles cancer and often improves with immune-suppressing treatment.

\medterm{Autoimmune Polyendocrine Syndromes} A group of inherited or acquired disorders in which immune attacks affect several hormone-producing glands. Depending on the type, features may include adrenal failure, low parathyroid hormone, thyroid disease, diabetes, chronic fungal infection, or skin changes.

\medterm{Autoimmune Polyendocrinopathy} A disorder in which immune attacks affect several hormone-producing glands. One inherited form can cause chronic fungal infections of skin or mouth, low parathyroid hormone, adrenal failure, and other autoimmune conditions.

\medterm{Autoimmune Polyglandular Syndrome} A group of disorders in which the immune system damages two or more hormone-producing glands. Possible effects include adrenal failure, thyroid disease, low parathyroid hormone, diabetes, or problems with reproductive hormones.

\medterm{Autoimmune Thyroid Disease and Pregnancy} Thyroid autoimmunity can cause hypothyroidism, hyperthyroidism, pregnancy loss, preterm birth, or fetal thyroid effects. Thyroid function and medicine doses are monitored closely before and during pregnancy and after delivery.

\medterm{Autoimmune Thyroiditis} Immune-mediated inflammation and damage of the thyroid gland. See \textit{Hashimoto Thyroiditis}.

\medterm{Autoimmunity} An immune response directed against the body's own cells, tissues, or proteins. It may cause no harm, or it may produce inflammation and disease when it damages tissues or interferes with organ function.

\medterm{Autoimmunity: Mechanisms} Autoimmunity can develop when the immune system fails to remove or control self-reactive cells, or when an infection triggers responses that also recognize body tissues. Genetic factors, immune regulation, and environmental triggers may all contribute.

\medterm{Autoinflammatory Diseases} Conditions in which the body's early, nonspecific immune response becomes overactive and causes repeated or ongoing inflammation. They may cause fever, rash, joint pain, abdominal pain, or organ damage, often without disease-specific autoantibodies.

\medterm{Autoinoculation} The spread of an infection from one part of a person's body to another, often by touching or scratching. This can transfer germs to broken skin, the eyes, the mouth, or another vulnerable site.

\medterm{Autologous} Taken from and returned to the same person. Examples include storing a person's own blood for later transfusion or collecting and returning their own stem cells during treatment.

\medterm{Autologous Blood Donation} Collection and storage of a person's own blood before an operation for possible transfusion during or after surgery. It avoids incompatibility with donor blood but may not be suitable or necessary for everyone.

\medterm{Autologous Breast Reconstruction} Surgical reconstruction of the breast contour utilizing the patient's own vascularized tissue flaps harvested from donor sites such as the abdomen (e.g., deep inferior epigastric perforator [DIEP] or transverse rectus abdominis myocutaneous [TRAM] flap), back (latissimus dorsi flap), or gluteal/thigh regions.

\synonyms
Autologous Breast Reconstruction; Flap Breast Reconstruction; Autogenous Tissue Reconstruction; DIEP Flap Reconstruction

\medterm{Autologous Fat Transplant} A procedure that removes fat from one area of a person's body and injects it into another area to restore volume or change shape. Some of the transferred fat may be reabsorbed, and repeat procedures may be needed.

\synonyms
Autologous Fat Grafting; Fat Transfer; Lipofilling

\medterm{Autologous Stem Cell Transplantation} A clinical procedure in which a patient's own hematopoietic stem cells are mobilized, harvested from peripheral blood or bone marrow, cryopreserved, and reinfused following high-dose myeloablative chemotherapy or radiation to rescue bone marrow hematopoiesis.

\synonyms
Autologous Bone Marrow Transplant; Auto-SCT

\medterm{Autolysis} The breakdown of cells or tissues by their own enzymes, usually after the cells die. This natural process contributes to tissue decomposition and can affect specimens examined after death.

\medterm{Automated External Defibrillator} A portable device that checks the heart rhythm during cardiac arrest and gives an electrical shock when an abnormal rhythm may respond to one. It provides spoken or visual instructions for using it safely.

\medterm{Automaticity} The physiological capacity of specialized cardiac pacemaker myocytes (principally in the sinoatrial node, atrioventricular junction, and His-Purkinje system) to undergo spontaneous diastolic phase 4 depolarization and generate action potentials without external neurohumoral stimulation.

\medterm{Autonomic Disorder} A problem affecting nerves that control automatic body functions such as blood pressure, heart rate, digestion, sweating, bladder emptying, sexual function, and pupil size. Causes include diabetes, nerve disease, infection, immune disease, and medicines; symptoms depend on the affected functions.

\medterm{Autonomic Dysreflexia} A sudden, potentially life-threatening rise in blood pressure in a person with a spinal-cord injury, usually at or above the upper chest. A problem such as a full bladder or bowel below the injury can trigger pounding headache, sweating, flushing, or heart-rate changes.

\medterm{Autonomic Ganglion} A small group of nerve-cell bodies outside the brain and spinal cord that passes signals between the autonomic nervous system and organs such as the heart, intestines, glands, and bladder.

\medterm{Autonomic Nervous System} The part of the nervous system that controls body functions that usually happen without conscious effort, including heart rate, blood pressure, digestion, sweating, pupil size, bladder activity, and sexual responses.

\medterm{Autonomic Neuropathy} Damage to nerves that control automatic functions such as blood pressure, heart rate, digestion, sweating, bladder emptying, and sexual function. Diabetes is common, but other diseases, toxins, and medicines can also cause it.

\synonyms
Dysautonomia; Autonomic Nerve Dysfunction

\medterm{Autonomic Neuropathy, Diabetic} Alternate index label; see \textit{Diabetic Autonomic Neuropathy}.

\medterm{Autonomy, Patient} A person's right and ability to make informed, voluntary decisions about medical care, including accepting or refusing tests or treatment. It requires understandable information and freedom from pressure.

\medterm{Autophagy} A normal process in which cells break down and recycle worn-out proteins and cell parts. It helps cells survive stress and remove damaged material, but abnormal autophagy may contribute to some cancers and other diseases.

\medterm{Autopsy} A medical examination of a body after death. It may include external examination, examination of internal organs, tissue testing, and sometimes toxicology to help determine the cause of death or study disease.

\synonyms
Post-Mortem Examination; Necropsy; Obduction

\medterm{Autopsy, Medicolegal} A specialized post-mortem examination performed under statutory or judicial authority by a forensic pathologist or medical examiner to determine the cause and manner of violent, sudden, suspicious, unnatural, or unattended deaths and preserve legal evidence.

\synonyms
Forensic Autopsy; Medico-Legal Autopsy; Coroner Post-Mortem

\medterm{Autoradiography} A laboratory method that uses a radioactive label and a photographic or digital detector to show where a substance is located in a tissue, cell, or biological sample.

\medterm{Autosomal Dominant} A pattern of inheritance in which one altered copy of a gene on a non-sex chromosome can cause a condition. A person with one altered copy may pass it to about half of their children.

\medterm{Autosomal Dominant Inheritance} Inheritance in which one altered copy of a gene on a non-sex chromosome can cause a condition. Each child of a parent with one altered copy commonly has a one-in-two chance of inheriting it.

\medterm{Autosomal Dominant PKD} An inherited disease in which many fluid-filled cysts gradually develop in the kidneys and may also occur in the liver. Cysts can cause high blood pressure, pain, blood in the urine, infections, and progressive kidney failure.

\medterm{Autosomal Dominant Polycystic Kidney Disease} An autosomal dominant inherited disease, usually caused by a change in the \textit{PKD1} or \textit{PKD2} gene, in which many cysts gradually enlarge the kidneys and can reduce kidney function. It may cause high blood pressure, pain, blood in the urine, kidney infections, liver cysts, or rarely brain-artery aneurysms.

\synonyms
ADPKD

\medterm{Autosomal Dominant Tubulointerstitial Kidney Disease} An inherited disease that gradually damages the kidney's small tubules and surrounding tissue. Kidney function declines, often without much protein or blood in the urine; some forms cause early gout, low blood pressure, or anaemia.

\medterm{Autosomal Inheritance} The passing of a genetic trait through one of the 22 pairs of chromosomes that are not sex chromosomes. Such traits can affect people of any sex and may follow dominant or recessive patterns.

\medterm{Autosomal Recessive} A pattern of inheritance in which a condition usually develops only when both copies of a gene are altered. Parents with one altered copy are often healthy carriers but can pass the condition to a child.

\medterm{Autosomal Recessive Disorders} Conditions that usually occur when a person inherits an altered copy of the same gene from each parent. People with one altered copy are often healthy carriers.

\medterm{Autosomal Recessive Inheritance} Inheritance in which a condition usually requires altered copies of a gene from both parents. When both parents are carriers, each pregnancy commonly has a one-in-four chance of producing an affected child.

\medterm{Autosomal Recessive PKD} An inherited disease that usually causes enlarged, poorly functioning kidneys in babies or children and scarring of the liver. Severe disease can interfere with lung development before birth; milder forms may appear later.

\medterm{Autosomal Recessive Polycystic Kidney Disease} An inherited disease that usually affects babies or children, causing enlarged kidneys with small cysts and progressive kidney problems. It can also cause liver scarring, high blood pressure, and breathing problems in severe early disease.

\synonyms
ARPKD

\medterm{Autosome} Any chromosome that is not a sex chromosome. Humans usually have 22 pairs of autosomes. Genes on these chromosomes can cause inherited conditions in people of any sex.

\synonyms
Non-Sex Chromosome

\medterm{Avascular Necrosis} Death of bone cellular components resulting from compromised subchondral microvascular blood supply, leading to bone trabecular death, structural collapse, and progressive secondary osteoarthritis. It most commonly affects the femoral head, scaphoid, talus, and humeral head; major etiologies include trauma/fractures, chronic corticosteroid therapy, excessive alcohol use, and sickle cell disease.

\synonyms
Osteonecrosis; Aseptic Necrosis; Ischemic Bone Necrosis; Femoral Head Avascular Necrosis; AVN

\medterm{Avastin} The brand name for bevacizumab, a monoclonal antibody that blocks vascular endothelial growth factor (VEGF). It is used for several cancers and, in ophthalmology, by injection into the eye for selected retinal conditions; risks include high blood pressure, bleeding, clotting, impaired wound healing, and rare gastrointestinal perforation.

\medterm{Avian Flu} Alternate name; see \textit{Avian Influenza}.

\medterm{Avian Influenza} An infection caused by avian-adapted influenza A virus strains (such as H5N1, H5N6, or H7N9) transmitted to humans primarily via direct contact with infected poultry or contaminated environments. Clinical manifestations range from conjunctivitis and upper respiratory symptoms to fulminant viral pneumonia, acute respiratory distress syndrome (ARDS), multi-organ failure, and death.

\synonyms
Bird Flu; Avian Flu; Highly Pathogenic Avian Influenza; HPAI

\medterm{AVN} Alternate name; see \textit{Avascular Necrosis}.

\medterm{Avoidance Behavior} Actions taken to stay away from a feared or distressing situation, thought, place, person, or activity. Avoidance may briefly reduce anxiety but can keep the fear going and interfere with treatment, work, relationships, or daily life.

\medterm{Avoidant Personality Disorder} A long-term pattern of avoiding social contact because of strong feelings of inadequacy and fear of criticism, embarrassment, or rejection. It can limit relationships and work; psychological treatment can help.

\medterm{Avoidant/Restrictive Food Intake Disorder} An eating disorder characterized by persistent failure to meet appropriate nutritional or energy needs, resulting in significant weight loss, nutritional deficiency, dependence on enteral feeding/supplements, or marked interference with psychosocial functioning. Unlike anorexia nervosa, the restriction is driven by sensory sensitivity, lack of interest in eating, or fear of aversive consequences (e.g., choking or vomiting), rather than body-image disturbance.

\synonyms
ARFID; Selective Eating Disorder

\medterm{Avolition} A marked reduction in motivation or ability to begin and complete purposeful activities. A person may neglect hygiene, household tasks, work, or social contact. It can occur in schizophrenia, severe depression, and other conditions.

\medterm{Avulsed Tooth} A tooth completely knocked out of its socket by an injury. It is a dental emergency; prompt treatment may allow the tooth to be replaced in its socket if it is handled by the crown and kept moist.

\medterm{Avulsion} An injury in which tissue is partly or completely torn away by force. It may involve skin, muscle, a nail, a fingertip, or a tooth and can cause severe bleeding, tissue loss, infection, and loss of function.

\medterm{Avulsion Fracture} A break in which a tendon or ligament pulls off a small piece of attached bone, usually during a sudden movement or injury. Treatment depends on the bone, fragment position, joint involvement, age, and associated damage.

\medterm{Avuncular} Relating to an uncle. In medical family-history or genetic writing, it describes the relationship between a person and their parent's brother or sister's husband.

\medterm{Awake Craniotomy} A specialized neurosurgical procedure in which the patient is deliberately kept awake, alert, and responsive during portions of open brain surgery. Performed primarily to remove tumors (such as gliomas) or epileptic seizure foci located within or adjacent to critical functional brain areas (the "eloquent cortex," which controls speech, language, motor function, and sensory perception), the operation begins under general anesthesia or conscious sedation while the skull bone flap is temporarily opened. The neurosurgical team then awakens the patient and applies gentle, low-voltage electrical stimulation to the exposed brain surface (cortical mapping) while a speech-language pathologist or neuropsychologist asks the patient to name pictures, read words, or move fingers and toes. If stimulating a particular brain spot temporarily stops speech or movement, surgeons know to spare that exact pathway. This real-time functional mapping maximizes tumor resection while safeguarding essential neurological abilities, reducing post-operative paralysis or loss of speech.
\synonyms
Awake Brain Surgery; Intraoperative Brain Mapping; Functional Craniotomy

\medterm{Awake Fiberoptic Intubation} A method of placing a breathing tube while the person remains awake and breathes independently. Numbing medicine and a flexible camera are used to guide the tube when a difficult airway is expected.

\medterm{Awareness} Conscious recognition of oneself, surroundings, thoughts, or sensations. Clinicians assess it along with alertness, attention, memory, and orientation when examining brain function.
\synonyms
Consciousness

\medterm{Awareness Under Anesthesia} Unintended awareness or memory of events during general anaesthesia. A person may hear sounds, feel pressure or pain, or recall part of the operation; although uncommon, it can cause significant distress and needs support.

\medterm{Axial Skeleton} The bones forming the central framework of the body: the skull, spine, ribs, breastbone, and associated small bones of the neck. It supports the body and protects the brain, spinal cord, heart, and lungs.

\medterm{Axial Spondyloarthritis} Long-term inflammation affecting the joints between the spine and pelvis and other parts of the spine. It commonly causes back pain and stiffness that improve with movement and may eventually limit spinal flexibility.

\medterm{Axial Tomography, Computerized} An imaging test that uses X-rays and computer processing to create detailed cross-sectional pictures of the body. It can show bones, organs, blood vessels, bleeding, tumours, and other abnormalities.

\medterm{Axilla} The hollow area beneath the shoulder, commonly called the armpit. It contains major blood vessels, nerves, lymph nodes, and tissues connecting the chest and neck with the arm.

\medterm{Axillary} Relating to the armpit or to structures within it, including the axillary blood vessels, nerves, and lymph nodes.

\synonyms
Armpit-Related; Pertaining to the Axilla

\medterm{Axillary Artery} The main artery passing through the armpit and supplying blood to the shoulder, chest wall, and arm. It continues as the brachial artery near the lower border of the armpit.

\medterm{Axillary Dissection} Surgery that removes selected lymph nodes and nearby tissue from the armpit, often for cancer assessment or treatment. It may cause numbness, shoulder stiffness, fluid collection, nerve injury, or swelling of the arm.

\synonyms
Axillary Node Dissection

\medterm{Axillary Lymph Node} A lymph node in the armpit that drains lymph from the arm, chest wall, and parts of the breast. Its size or involvement by cancer, infection, or inflammation may be assessed by examination or imaging.

\medterm{Axillary Lymphadenectomy} An operation to remove lymph nodes from the armpit, usually to assess or treat breast cancer or melanoma. Possible complications include pain, numbness, shoulder stiffness, a fluid collection, nerve injury, and arm swelling.

\medterm{Axillary Lymphadenopathy} An abnormal enlargement, induration, or matted consistency of the lymph nodes within the axilla. It receives lymphatic drainage from the ipsilateral breast, upper extremity, and chest wall. Major causes include metastatic breast carcinoma, lymphomas, melanoma, acute pyogenic infections of the upper limb, cat-scratch disease (Bartonella henselae), tuberculosis, and systemic autoimmune diseases.

\synonyms
Axillary Adenopathy; Enlarged Axillary Nodes

\medterm{Axillary Nerve} A nerve from the neck that passes around the shoulder and supplies muscles that lift the arm. It also provides feeling over the outer shoulder and may be injured by shoulder dislocation or upper-arm fracture.

\medterm{Axillary Nerve Dysfunction} Impairment of the axillary nerve, often from shoulder trauma, dislocation, compression, or surgery. It can cause weakness of shoulder abduction, deltoid wasting, and numbness over the lateral shoulder.

\medterm{Axillary Tail Of The Breast} The extension of breast tissue toward the armpit. It is also called the tail of Spence.

\medterm{Axillary Vein} A large vein in the armpit that drains blood from the arm and shoulder. It continues beneath the collarbone as the subclavian vein and returns blood toward the heart.

\medterm{Axis} The second bone in the neck. Its tooth-like projection acts as a pivot for the first neck bone, allowing the head and upper neck to rotate.

\medterm{Axis Vertebra} The second neck vertebra, whose tooth-like projection allows the first vertebra and head to rotate while protecting the spinal cord.

\synonyms
C2; Axis

\medterm{Axolemma} The thin outer membrane surrounding a nerve cell's axon. It maintains the electrical difference needed for nerve impulses and contains channels that allow those impulses to travel along the axon.

\medterm{Axon} The long extension of a nerve cell that carries electrical signals away from the cell body toward another nerve cell, muscle, or gland.

\medterm{Axon Fasciculation} The process during development in which growing nerve fibres group together to form organized nerve pathways. Chemical guidance signals and cell connections help them reach their correct targets.

\medterm{Axon Guidance} The process by which growing nerve fibres find and connect with their correct targets during development. Chemical signals and surrounding cells help direct their path.

\medterm{Axon Hillock} The cone-shaped part where a nerve cell's main body joins its axon, the long fiber that carries signals away. It combines incoming electrical signals and usually starts the nerve impulse when those signals are strong enough.

\medterm{Axon Terminal} The end of a nerve-cell axon where chemical messengers are released to communicate with another nerve cell, muscle, or gland.

\medterm{Axonal Transport} The movement of proteins, cell parts, and other materials along a nerve cell's axon between the cell body and the axon terminal. It supports nerve-cell maintenance and communication.

\medterm{Axoneme} The internal framework of tiny tubes inside a cilium or sperm tail. Together with motor proteins, it produces the bending movements that help move fluid over cells or propel sperm.

\medterm{Axonotmesis} A moderate degree of peripheral nerve injury in Seddon's classification characterized by disruption of the internal axon and its myelin sheath with subsequent Wallerian degeneration distal to the lesion site, but with anatomical preservation of the surrounding endoneurium, perineurium, and epineurium, allowing spontaneous axonal regeneration.

\medterm{Ayurveda} A traditional system of health care from India that uses ideas about bodily balance and may include diet, lifestyle practices, massage, and herbal products. Some products can interact with medicines or contain harmful contaminants.

\medterm{Ayurvedic Medicine} Health care based on traditional Indian theories of bodily balance, using individualized diet, lifestyle practices, physical treatments, and herbal products. Evidence and safety vary, and some products may interact with medicines or contain toxic ingredients.

\medterm{Azole Antifungals} A group of medicines that stop fungi making ergosterol, a substance needed for strong cell membranes. This weakens or kills the fungus. They can interact with other medicines and may affect the liver or heart rhythm.

\medterm{Azoospermia} The absence of sperm in semen on laboratory examination. It may result from a blockage that prevents sperm from leaving the reproductive tract or from reduced or absent sperm production in the testes.

\medterm{Azotaemia} Alternate spelling; see \textit{Azotemia}.

\medterm{Azotemia} A build-up of nitrogen-containing waste products in the blood, usually because the kidneys are filtering less effectively. It may occur with dehydration, reduced kidney blood flow, or kidney disease.

\medterm{Azygos Vein} A large vein running along the right side of the spine that drains blood from the chest wall into the superior vena cava. It provides an alternate route for blood when major veins are blocked.

\medterm{Azygous Vein} Alternate spelling; see \textit{Azygos Vein}.

\medterm{A1AT Deficiency} Abbreviation for alpha-1 antitrypsin deficiency, an inherited disorder that can cause early emphysema, liver disease, or both because protective alpha-1 antitrypsin levels are too low or abnormal.

\medterm{AAT} Abbreviation for alpha-1 antitrypsin, a protein made mainly in the liver that helps protect lung tissue from enzyme-related damage. Blood testing can measure its amount and type.

\medterm{AAT Deficiency} Alternate name for alpha-1 antitrypsin deficiency, an inherited condition associated with emphysema, liver disease, panniculitis, and rarely vasculitis.

\medterm{Abercrombie Syndrome} An older name for systemic amyloidosis, in which abnormal amyloid proteins accumulate in organs and impair their function.

\medterm{Abnormal Pap Test} A cervical screening result showing abnormal cells or evidence of human papillomavirus. It does not by itself mean cancer; follow-up depends on the cell changes, HPV result, age, and prior screening history.

\medterm{Abruption, Placental} Alternate name for placental abruption, premature separation of the placenta from the uterine wall before delivery. It can cause vaginal bleeding, abdominal pain, uterine tenderness, fetal distress, and maternal shock.

\medterm{Absence of Menstruation, Primary} Failure to have the first menstrual period by the expected age, called primary amenorrhea. Causes include normal variation, pregnancy, hormonal disorders, anatomic differences, chronic illness, undernutrition, and excessive exercise.

\medterm{Abuse, Child} Physical, sexual, or emotional harm to a child, or neglect of a child’s basic needs. Suspected abuse requires prompt attention to safety and reporting according to local law and professional responsibilities.

\medterm{ACC -- Adenoid Cystic Carcinoma} Abbreviation for adenoid cystic carcinoma, a malignant tumor of glandular tissue that may grow slowly but can invade nerves and recur or spread years after treatment.

\medterm{ACC -- Agenesis of Corpus Callosum} Abbreviation for agenesis of the corpus callosum, a congenital absence or incomplete development of the major nerve-fiber bridge connecting the brain’s two hemispheres.

\medterm{ACC -- Aplasia Cutis Congenita} Abbreviation for aplasia cutis congenita, a birth defect in which an area of skin, most often on the scalp, is absent or incompletely formed.

\medterm{ACH} Abbreviation commonly used for achondroplasia, a genetic skeletal disorder causing disproportionate short stature and characteristic changes in the skull, spine, and long bones.

\medterm{Achilles Tendon Problems} Disorders affecting the Achilles tendon, including tendinitis, tendinopathy, bursitis, partial tear, and rupture. Pain, stiffness, swelling, or a sudden loss of push-off strength warrants assessment.

\medterm{ADD/ADHD} Attention-deficit/hyperactivity disorder, a neurodevelopmental condition involving persistent inattention, hyperactivity, impulsivity, or combinations of these symptoms that impair functioning.

\medterm{Adenosine Deaminase Severe Combined Immunodeficiency (ADA-SCID)} A rare inherited immune deficiency caused by adenosine deaminase deficiency. Toxic metabolites prevent normal development of lymphocytes, leaving infants highly vulnerable to severe infections.

\medterm{ADHD in Children} Attention-deficit/hyperactivity disorder in childhood, characterized by developmentally inappropriate inattention, hyperactivity, or impulsivity across settings. Diagnosis requires persistent impairment and information from caregivers, teachers, and clinical assessment.

\medterm{Aganglionic Megacolon} Congenital absence of nerve cells in part of the large intestine, causing Hirschsprung disease. The affected segment cannot relax normally, leading to intestinal obstruction, abdominal swelling, and delayed passage of stool.

\medterm{Agenesis of Commissura Magna Cerebri} An older anatomic name for agenesis of the corpus callosum, in which the major commissural fiber tract connecting the cerebral hemispheres is absent or underdeveloped.

\medterm{Agnogenic Myeloid Metaplasia (AMM)} An older name for primary myelofibrosis, a chronic bone-marrow disorder in which abnormal scar tissue and blood-cell production outside the marrow cause anemia, enlarged spleen, constitutional symptoms, and abnormal blood counts.

\medterm{Agnosia, Primary Visual} An inability to recognize or identify objects by sight despite adequate vision. It results from dysfunction in visual association areas of the brain and may follow stroke, trauma, dementia, or other neurologic disease.

\medterm{Agnosis, Primary} An inability to recognize familiar objects, people, sounds, or other stimuli despite adequate sensation and understanding. The specific form depends on the affected sensory association system.

\medterm{Airsickness} Motion sickness caused by travel through the air. Conflicting signals from the eyes and inner ear can produce nausea, dizziness, sweating, headache, and vomiting.

\medterm{Albinismus} A genetic condition involving reduced or absent melanin production, causing lighter skin, hair, or eyes and often visual problems such as reduced acuity, nystagmus, and light sensitivity.

\medterm{ALD} Abbreviation commonly referring to adrenoleukodystrophy, an inherited disorder of fatty-acid metabolism that can damage the nervous system and adrenal glands.

\medterm{Aldosteronism With Normal Blood Pressure} A form of primary aldosteronism in which excess aldosterone is present without obvious hypertension. It may still cause low potassium, metabolic alkalosis, or cardiovascular and kidney risk and requires evaluation of the cause.

\medterm{Algodystrophy} An older term for complex regional pain syndrome, a persistent pain condition after injury or surgery that may cause swelling, color or temperature change, sweating abnormalities, and stiffness.

\medterm{Algoneurodystrophy} An older term for complex regional pain syndrome, involving disproportionate regional pain with sensory, vasomotor, sudomotor, and motor or trophic changes.

\medterm{Allergies to Insect Stings} Allergic reactions to venom from bees, wasps, hornets, or fire ants. Severe reactions can cause hives, throat swelling, wheezing, low blood pressure, or collapse and require immediate epinephrine and emergency care.

\medterm{Allergy to Natural Rubber (Latex)} An immune reaction to natural rubber latex that may cause contact dermatitis, hives, wheezing, or anaphylaxis. Avoidance and clearly documented latex precautions are important for people with significant reactions.

\medterm{AloBar Holoprosencephaly} A severe form of holoprosencephaly in which the developing forebrain fails to divide and the cerebral hemispheres remain largely fused. It is usually associated with profound developmental impairment and major facial or other congenital abnormalities.

\medterm{Alopecia Celsi} An older name for alopecia areata, an autoimmune disorder causing sharply defined patches of nonscarring hair loss.

\medterm{Alopecia Cicatrisata} An older term for scarring alopecia, in which inflammation or injury destroys hair follicles and replaces them with scar tissue, causing permanent hair loss in affected areas.

\medterm{Alopecia Circumscripta} An older name for alopecia areata, characterized by one or more circumscribed patches of nonscarring hair loss.

\medterm{Alpha-1} A shortened name commonly used for alpha-1 antitrypsin deficiency or alpha-1 antitrypsin testing; the intended meaning should be confirmed from context.

\medterm{Altitude Headache} Headache related to ascent to high altitude, often as part of acute mountain sickness. It may occur with nausea, dizziness, fatigue, or sleep disturbance and can signal more serious altitude illness when severe or accompanied by neurologic or breathing symptoms.

\medterm{Alzheimer’s} A common possessive form of Alzheimer disease, a progressive neurodegenerative disorder causing worsening memory, thinking, behavior, and ability to perform daily activities.

\medterm{Anaplastic Large-cell Lymphoma} A usually aggressive T-cell non-Hodgkin lymphoma characterized by large anaplastic cells and CD30 expression. It may involve lymph nodes, skin, bone, or other organs and is classified as systemic or cutaneous.

\medterm{Anxiety & Panic} Anxiety disorders involve excessive fear, worry, or physiological arousal; panic attacks are sudden episodes of intense fear with symptoms such as palpitations, breathlessness, chest discomfort, trembling, or fear of dying.

\medterm{APLS} Abbreviation for antiphospholipid syndrome, an autoimmune disorder associated with venous or arterial thrombosis and pregnancy complications, with persistent antiphospholipid antibodies.

\medterm{Apnea, Infantile} Repeated pauses in breathing in an infant. Causes vary with age and context and may include prematurity, infection, airway obstruction, reflux, seizures, metabolic disease, or unexplained infant events; color change, poor responsiveness, or prolonged pauses are emergencies.

\medterm{Appalachian Type Amyloidosis} A rare hereditary transthyretin amyloidosis variant associated with abnormal amyloid deposition, often causing neuropathy, autonomic symptoms, or heart involvement. Diagnosis requires genetic and tissue or specialized laboratory evaluation.

\medterm{APS} Abbreviation for antiphospholipid syndrome; it is also used for other conditions in different contexts, so the intended meaning should be confirmed.

\medterm{Arachnitis} An older spelling of arachnoiditis, chronic inflammation and scarring of the arachnoid membrane around spinal or cranial nerves that can cause persistent pain, sensory symptoms, weakness, or bladder and bowel dysfunction.

\medterm{Arachnoidal Fibroblastoma} A rare benign or low-grade tumor arising from fibroblastic cells of the arachnoid coverings. Symptoms depend on location and may result from compression of the brain, spinal cord, or nerves.

\medterm{Arbovirus A Chikungunya Type} An older classification term for chikungunya virus, a mosquito-borne alphavirus that causes fever and often severe inflammatory joint pain, rash, headache, and muscle pain.

\medterm{ARDS} Abbreviation for acute respiratory distress syndrome, a life-threatening inflammatory lung injury causing fluid leakage into the air sacs, severe oxygenation failure, and often the need for ventilatory support.

\medterm{Aregenerative Anemia} Anemia with an inappropriately low reticulocyte response, indicating inadequate red-cell production by the bone marrow. Causes include iron, vitamin B12, or folate deficiency, kidney disease, inflammation, endocrine disease, marrow failure, and some medicines.

\medterm{ARG Deficiency} Abbreviation for arginase deficiency, an inherited urea-cycle disorder in which arginine accumulates and can cause progressive spasticity, developmental problems, seizures, and elevated ammonia.

\medterm{Arteritis, Takayasu} Alternate name for Takayasu arteritis, a large-vessel vasculitis that can narrow or block the aorta and its major branches, causing pulse differences, limb claudication, blood-pressure differences, and organ ischemia.

\medterm{Arthritis (Osteoarthritis)} Osteoarthritis, a degenerative joint disorder involving cartilage loss, bone remodeling, pain, stiffness, and reduced function. It commonly affects the knees, hips, hands, spine, and feet.

\medterm{Arthritis Urethritica} An older name for reactive arthritis, inflammatory arthritis occurring after certain gastrointestinal or genitourinary infections and sometimes accompanied by urethritis, conjunctivitis, or skin findings.

\medterm{Arthritis, Juvenile Rheumatoid} An older name for juvenile idiopathic arthritis, chronic inflammatory arthritis beginning in childhood. It may affect one or many joints and can involve the eyes, skin, or other organs.

\medterm{AVM} Abbreviation for arteriovenous malformation, an abnormal direct connection between arteries and veins that bypasses capillaries. AVMs can bleed, cause seizures or neurologic deficits, or produce symptoms from high-flow effects.

\medterm{Axonal Neuropathy, Giant} A rare inherited neurodegenerative disorder in which giant axons develop in peripheral nerves and the central nervous system. It usually begins in childhood with progressive ataxia, weakness, sensory loss, and developmental or cognitive changes.
